\documentclass[12pt,a4paper,twoside,openright]{book}
% ---------------------------------------------------------------------------
% Preambolo condiviso della monografia NanoFaaS.
% Compilazione: latexmk -pdf nanofaas.tex   (pdfLaTeX, TeX Live 2026+)
% ---------------------------------------------------------------------------

\usepackage[T1]{fontenc}
\usepackage[utf8]{inputenc}
\usepackage[italian]{babel}
\usepackage{lmodern}
\usepackage{microtype}

\usepackage{geometry}
\geometry{a4paper, top=3cm, bottom=3cm, inner=3cm, outer=2.5cm}

\usepackage{amsmath, amssymb}
\usepackage{booktabs}
\usepackage{longtable}
\usepackage{array}
\usepackage{multirow}
\usepackage{tabularx}
\usepackage{caption}
\usepackage{subcaption}
\usepackage{enumitem}
\usepackage{float}
\usepackage{adjustbox}

% Figure TikZ: le sorgenti in figures/ sono documenti `standalone`.
% Caricando il pacchetto `standalone` qui, \input{figures/fig-x} le include
% ignorandone il \documentclass e riusando questo preambolo.
\usepackage{figures/nanofaas-figs}
\usepackage{standalone}
\usepackage{pgfplots}
\pgfplotsset{compat=1.18}

\usepackage{listings}
\usepackage{xcolor}

\usepackage[hidelinks]{hyperref}
\usepackage[nameinlink,italian]{cleveref}

% ---- stile dei listati -----------------------------------------------------
\lstdefinestyle{nf}{
  basicstyle=\ttfamily\footnotesize,
  keywordstyle=\color{nfblue}\bfseries,
  commentstyle=\color{nfslate}\itshape,
  stringstyle=\color{nfgreen},
  numbers=left, numberstyle=\tiny\color{nfslate}, numbersep=7pt,
  frame=leftline, framerule=1pt, rulecolor=\color{nfslate!40},
  xleftmargin=10pt, breaklines=true, breakatwhitespace=false,
  postbreak=\mbox{\textcolor{nfslate}{$\hookrightarrow$}\space}, columns=flexible,
  showstringspaces=false, tabsize=2, captionpos=b,
  literate={à}{{\`a}}1 {è}{{\`e}}1 {é}{{\'e}}1 {ì}{{\`i}}1 {ò}{{\`o}}1 {ù}{{\`u}}1
           {—}{{---}}1 {–}{{--}}1 {’}{{'}}1 {‘}{{'}}1 {“}{{"}}1 {”}{{"}}1 {≥}{{$\geq$}}1 {≤}{{$\leq$}}1 {×}{{$\times$}}1 {λ}{{$\lambda$}}1 {μ}{{$\mu$}}1,
}
\lstset{style=nf}
\lstdefinelanguage{yamlnf}{
  keywords={true,false,null,enabled,mode,strategy,metric},
  sensitive=false, comment=[l]{\#}, morestring=[b]',morestring=[b]"
}

% ---- elementi ricorrenti ---------------------------------------------------
\newcommand{\code}[1]{\texttt{#1}}
\newcommand{\nf}{NanoFaaS}

% Riquadro "da misurare": marca un dato non ancora rilevato.
\usepackage[most]{tcolorbox}
\newtcolorbox{damisurare}[1][]{
  colback=nfamberbg, colframe=nfamber, boxrule=0.6pt, arc=2pt,
  fonttitle=\bfseries\sffamily\small, title={Da misurare}, #1}
\newcommand{\todomis}[1]{\textcolor{nfamber}{\textbf{[da misurare: #1]}}}

% Numerazione profonda ma indice compatto.
% Tolleranza extra: il testo italiano con molti \code{...} non spezzabili
% produrrebbe altrimenti righe che escono dal margine.
\setlength{\emergencystretch}{6em}
\hyphenpenalty=1000
\tolerance=3000

\setcounter{secnumdepth}{3}
\setcounter{tocdepth}{1}


\title{\textbf{\nf}\\[0.4em]\large Un control plane FaaS minimale e modulare:\\ architettura, controllo della concorrenza e valutazione sperimentale}
\author{Universit\`a degli Studi di Milano--Bicocca\\ Dipartimento di Informatica, Sistemistica e Comunicazione}
\date{\today}

\begin{document}

\frontmatter
\maketitle
\chapter*{Sommario}
\addcontentsline{toc}{chapter}{Sommario}

\nf{} \`e un control plane Function-as-a-Service progettato attorno a un vincolo
insolito: essere deliberatamente pi\`u piccolo delle piattaforme con cui
condivide il dominio. Non offre alta disponibilit\`a, non persiste il proprio
stato, non autentica i chiamanti e gira in un solo processo. Questa monografia
sostiene che tali rinunce non sono un debito da colmare ma il metodo che rende
la piattaforma utilizzabile come \emph{strumento di misura}: un sistema in cui
il percorso critico di un'invocazione \`e abbastanza corto da essere letto per
intero, e abbastanza governabile da permettere di attribuire una variazione di
latenza alla politica che l'ha causata anzich\'e all'infrastruttura che la
circonda.

Il documento descrive l'architettura del control plane --- un nucleo minimale
reattivo attorno a cui si innestano moduli opzionali selezionati a tempo di
build attraverso un SPI --- e dedica un capitolo a ciascun modulo: le code
asincrone e sincrone, l'autoscaler di repliche, il governo della concorrenza per
funzione, l'offload trasparente verso un'istanza remota, la riconfigurazione a
caldo e i provider di deployment su Kubernetes e su runtime container locale.

Il contributo pi\`u sostanziale riguarda il controllo della concorrenza. Il
modulo \code{concurrency-control} implementa cinque politiche, di cui tre
adattive, e la loro valutazione sperimentale ha prodotto due risultati negativi
utili quanto quelli positivi: un controllore che minimizza il \emph{tempo di
servizio} non ha ottimo interno e cammina verso il proprio limite inferiore, e
un carico a ciclo chiuso non \`e in grado di porre al sistema la domanda a cui
il controllore dovrebbe rispondere. Da qui la formulazione di una politica che
decide dal \emph{tempo di sojourn} --- attesa pi\`u servizio --- che l'ottimo
interno ce l'ha, e di una politica a budget di piattaforma con equit\`a max-min
pesata che separa quanto una funzione \emph{chiede} da quanto la piattaforma
\emph{pu\`o dare}.

La valutazione sperimentale copre prestazioni di base (throughput, latenza,
cold start) attraverso una serie di release su profilo hardware fissato,
l'impronta di memoria del control plane su JVM e in immagine nativa GraalVM, il
comportamento del garbage collector sotto limiti di memoria, e il confronto fra
i controllori di concorrenza. Un capitolo finale enumera esplicitamente ci\`o
che i dati raccolti \emph{non} stabiliscono.

\tableofcontents

\mainmatter

\part{Contesto, motivazione, metodo}
\chapter{Introduzione}
\label{cap:01-introduzione}

\section{Il contesto: che cosa promette il modello serverless}

Il modello \emph{Function-as-a-Service} (FaaS) propone al programmatore un
contratto tanto semplice quanto ambizioso: si consegna alla piattaforma una
funzione, ossia un frammento di codice senza stato con una firma nota, e la
piattaforma si assume ogni responsabilit\`a residua --- collocarla su una
macchina, replicarla quando gli arrivi crescono, ritirarla quando cessano,
instradare le richieste, riprovare i fallimenti, contabilizzare l'uso. Il
programmatore non dimensiona nulla e non paga nulla quando nessuno invoca.

\`E un contratto attraente perch\'e sposta un intero registro di decisioni
--- capacit\`a, elasticit\`a, collocazione --- dall'utente al fornitore. Ed \`e
anche un contratto costoso, perch\'e quelle decisioni non spariscono: si
concentrano dentro la piattaforma, dove diventano politiche automatiche che
l'utente non vede e, quando il comportamento osservato non lo soddisfa, non sa
come correggere. Il \emph{cold start} \`e l'esempio pi\`u noto di questa
opacit\`a: un ritardo che il chiamante misura ma che dipende da una decisione
di collocazione presa altrove, in un istante che il chiamante non conosce, sulla
base di uno stato che non gli viene esposto.

La conseguenza per chi fa ricerca su queste piattaforme \`e concreta. Studiare
una politica di scaling, di ammissione o di controllo della concorrenza richiede
di poterla sostituire e di poter attribuire una differenza misurata a quella
sostituzione e non ad altro. Sulle piattaforme FaaS mature questo \`e difficile
per una ragione strutturale, non accidentale: sono sistemi distribuiti composti
da molti processi cooperanti --- un ingress, un autoscaler, un attivatore, un
controller, una coda, un data plane --- ciascuno con la propria dinamica, il
proprio periodo di campionamento e le proprie isteresi. Il percorso di una
richiesta attraversa diversi confini di processo prima di raggiungere il codice
dell'utente, e ogni confine contribuisce alla latenza in un modo che
l'osservatore non pu\`o disaggregare senza strumentare l'intero sistema.

\section{Il problema}

Da qui la domanda che motiva questo lavoro: \emph{qual \`e la piattaforma pi\`u
piccola che \`e ancora una piattaforma FaaS, e che cosa diventa misurabile una
volta che la si \`e ridotta a quella dimensione?}

La domanda non \`e retorica. Una piattaforma pu\`o essere ridotta lungo molte
dimensioni, e non tutte le riduzioni sono equivalenti. Togliere
l'autenticazione riduce il codice ma non accorcia il percorso critico di
un'invocazione; togliere la persistenza dello stato lo accorcia ma rinuncia alla
sopravvivenza a un riavvio; togliere l'alta disponibilit\`a elimina un intero
registro di problemi di consenso ma vincola la piattaforma a un singolo punto di
fallimento. La tesi di questo lavoro \`e che esiste un insieme di rinunce
--- pod singolo, stato in memoria, nessuna autenticazione --- che, prese
insieme, producono un sistema il cui percorso critico \`e leggibile per intero
da una sola persona e in cui una modifica di politica \`e osservabile senza
rumore di fondo, e che tale sistema resta abbastanza completo da eseguire
carichi reali su Kubernetes.

\nf{} \`e la realizzazione di questa tesi. \`E un control plane scritto in Java
su Spring WebFlux, il cui nucleo conta circa 5\,000 righe di codice di
produzione, esteso da nove moduli opzionali che ne aggiungono complessivamente
circa 6\,800. Il rapporto fra le due cifre \`e esso stesso una scelta di
progetto: i moduli, presi insieme, sono pi\`u grandi del nucleo che estendono,
il che significa che la maggior parte di ci\`o che rende \nf{} interessante \`e
anche ci\`o che si pu\`o togliere senza che smetta di funzionare.

\section{Domande di ricerca}

Il lavoro \`e organizzato attorno a quattro domande.

\begin{description}[style=nextline, leftmargin=0pt]
  \item[\textbf{RQ1 --- Modularit\`a a costo zero sul percorso critico.}]
  \`E possibile costruire un control plane in cui funzionalit\`a come le code,
  l'autoscaling e il controllo della concorrenza siano \emph{assenti} anzich\'e
  disattivate, senza che la loro assenza richieda rami condizionali sul percorso
  di invocazione? Il \cref{cap:08-sistema-moduli} descrive la soluzione adottata
  --- selezione a tempo di build attraverso un SPI, con implementazioni no-op di
  default per ogni punto di estensione --- e ne discute i limiti.

  \item[\textbf{RQ2 --- Da quale segnale deve decidere un controllore di
  concorrenza.}] Il limite di concorrenza per funzione \`e la grandezza che
  determina simultaneamente l'utilizzazione delle repliche e il tempo che un
  chiamante passa in attesa. La domanda \`e da quale segnale misurato tale limite
  debba essere derivato. Il \cref{cap:12-concurrency-control} mostra
  analiticamente e sperimentalmente che il tempo di servizio --- la grandezza
  che il control plane misura pi\`u naturalmente, fra dispatch e completamento
  --- \`e la scelta sbagliata, e che il tempo di sojourn \`e quella corretta.

  \item[\textbf{RQ3 --- Concorrenza come risorsa condivisa.}] Quando pi\`u
  funzioni insistono sulle stesse repliche o sullo stesso nodo, un controllore
  che chiede a ciascuna funzione di trovare da sola il proprio limite le sta
  chiedendo di decidere su una risorsa condivisa senza informazioni sui vicini.
  Il \cref{cap:12-concurrency-control} formula un'alternativa (modo
  \code{BUDGETED}) che separa la stima del fabbisogno dall'allocazione, e il
  \cref{cap:25-valutazione-concorrenza} riporta la misura del cross-talk fra
  funzioni co-residenti.

  \item[\textbf{RQ4 --- Qual \`e il costo reale di un control plane.}] Un
  processo che non esegue codice utente ma lo orchestra ha un'impronta di
  memoria e un profilo di avvio propri. Il \cref{cap:24-footprint} misura dove
  va la memoria di una JVM Spring, che cosa cambia compilando lo stesso codice
  in immagine nativa GraalVM, e come la scelta del garbage collector interagisce
  con i limiti di memoria imposti dal container.
\end{description}

\section{Contributi}

Il documento presenta i seguenti contributi.

\begin{enumerate}[leftmargin=*]
  \item \textbf{Un'architettura di control plane FaaS a nucleo minimale e moduli
  opzionali}, con un meccanismo di composizione a tempo di build che consente di
  produrre binari privi delle funzionalit\`a non selezionate, e un insieme di
  punti di estensione i cui default no-op degradano il sistema in modo esplicito
  anzich\'e farlo fallire (\cref{cap:04-architettura-control-plane,cap:06-nucleo,cap:08-sistema-moduli}).

  \item \textbf{Una tassonomia operativa di cinque politiche di controllo della
  concorrenza per funzione} --- fissa, statica per replica, adattiva per replica,
  a budget di piattaforma, e a minimizzazione del tempo di sojourn --- con la
  discussione analitica del perch\'e le prime quattro non possono essere valutate
  con lo stesso esperimento (\cref{cap:12-concurrency-control}).

  \item \textbf{Due risultati sperimentali negativi}, riportati come tali: un
  controllore che decide dal tempo di servizio non ammette ottimo interno e
  converge al proprio limite inferiore; e un generatore di carico a ciclo chiuso
  non pu\`o distinguere fra le politiche, perch\'e retroaziona sulla stessa
  grandezza che il controllore sta manipolando
  (\cref{cap:25-valutazione-concorrenza}).

  \item \textbf{Un modulo di offload trasparente edge$\to$cloud} attivabile in
  modo eager per funzione o reattivo sulla pressione della coda sincrona, con una
  semantica a singolo salto e senza fallback locale, e la giustificazione di
  entrambe le rinunce (\cref{cap:13-offload}).

  \item \textbf{Una caratterizzazione dell'impronta di runtime} del control plane
  su JVM e in immagine nativa, inclusa l'interazione fra scelta del collector e
  limite di memoria del container (\cref{cap:24-footprint}).

  \item \textbf{Un impianto sperimentale riproducibile}, esterno al codice della
  piattaforma, che provvisiona l'ambiente, esegue i carichi e confronta i
  risultati contro record versionati su profilo hardware fissato
  (\cref{cap:22-metodo-sperimentale}).
\end{enumerate}

\section{Che cosa questo lavoro non \`e}

\`E utile dichiarare subito i confini, per non farli scoprire al lettore a
met\`a documento.

\nf{} non compete con Knative, OpenWhisk o OpenFaaS su completezza funzionale, e
non \`e proposto come loro sostituto in produzione. Non implementa
autenticazione n\'e autorizzazione, quindi non \`e esponibile su una rete non
fidata. Non replica il proprio stato: un riavvio del pod perde il registro delle
funzioni e lo storico delle esecuzioni. Non implementa alta disponibilit\`a, e
l'esistenza di un singolo processo di controllo \`e assunta ovunque nel codice,
non incapsulata dietro un'astrazione che un giorno potrebbe diventare
distribuita.

Sul piano sperimentale, i numeri riportati nel \cref{cap:23-prestazioni-base} e
seguenti sono ottenuti su un profilo hardware specifico e dichiarato. Non sono
proposti come caratterizzazione assoluta della piattaforma n\'e come base per un
confronto con altre piattaforme misurate altrove. Il
\cref{cap:27-discussione} enumera in modo esplicito le affermazioni che i dati
raccolti non sostengono.

\section{Struttura del documento}

Il documento \`e diviso in sei parti.

La \textbf{Parte I} (\crefrange{cap:01-introduzione}{cap:03-requisiti-principi})
inquadra il problema, rivede lo stato dell'arte e formalizza i vincoli di
progetto, mostrando per ciascuno quale esperimento abilita.

La \textbf{Parte II} (\crefrange{cap:04-architettura-control-plane}{cap:08-sistema-moduli})
descrive l'architettura: la vista d'insieme, i contratti di dati e l'API REST,
il nucleo e i suoi punti di estensione, le modalit\`a di esecuzione e il
meccanismo di composizione modulare.

La \textbf{Parte III} (\crefrange{cap:09-async-queue}{cap:16-moduli-diagnostici})
dedica un capitolo a ciascun modulo opzionale. Il
\cref{cap:12-concurrency-control} \`e il pi\`u lungo della parte ed \`e quello a
contenuto scientifico maggiore.

La \textbf{Parte IV} (\crefrange{cap:17-runtime-funzione}{cap:20-build-packaging})
attraversa il lato funzione: il contratto di runtime, gli SDK nei vari
linguaggi, gli strumenti a riga di comando e la catena di build.

La \textbf{Parte V} (\crefrange{cap:21-osservabilita}{cap:25-valutazione-concorrenza})
presenta il modello di osservabilit\`a, il metodo sperimentale e le tre
campagne di misura.

La \textbf{Parte VI} (\crefrange{cap:26-qualita-verifica}{cap:29-conclusioni})
discute la verifica del software, i limiti dei risultati e le direzioni
successive.

\section{Convenzioni}

Il codice sorgente \`e citato con il percorso relativo alla radice del
repository, per esempio \code{platform/control-plane/src/main/java/\ldots}. Il
package radice del codice Java \`e \code{it.unimib.datai.nanofaas}; nel testo i
nomi di classe sono riportati senza package quando non ambigui. I nomi dei
moduli opzionali compaiono nella forma con cui sono selezionabili dalla riga di
comando (\code{async-queue}, \code{concurrency-control}, \ldots), che coincide
con il nome della directory sotto \code{platform/modules/}.

Le grandezze temporali sono espresse in millisecondi salvo indicazione contraria.
Con \emph{tempo di servizio} si intende l'intervallo fra l'istante in cui il
control plane consegna l'invocazione al runtime e l'istante in cui ne riceve la
risposta; con \emph{tempo di attesa} l'intervallo fra l'ammissione della
richiesta e la consegna; con \emph{tempo di sojourn} la loro somma, che \`e ci\`o
che il chiamante osserva al netto della rete. La distinzione \`e centrale nel
\cref{cap:12-concurrency-control} e viene mantenuta rigorosamente in tutto il
documento.

Infine, i riquadri marcati \textbf{Da misurare} segnalano affermazioni per cui
il progetto dispone di uno strumento di misura ma non ancora di un dato
raccolto. Sono lasciati visibili di proposito: un documento che presenta una
piattaforma come strumento sperimentale deve distinguere ci\`o che ha misurato
da ci\`o che potrebbe misurare.

\chapter{Background e stato dell'arte}
\label{cap:02-stato-dell-arte}

Questo capitolo fissa il vocabolario e colloca \nf{} rispetto al lavoro
esistente. Non \`e una rassegna esaustiva del serverless --- ne esistono di
migliori~\cite{jonas2019berkeley, baldini2017trends} --- ma una ricostruzione
mirata delle quattro linee di ricerca su cui il progetto interviene: la
struttura delle piattaforme FaaS, il costo di avvio a freddo, le politiche di
elasticit\`a e di ammissione, e l'offloading fra dispositivi periferici e
cloud.

\section{Una definizione operativa di FaaS}

Nella letteratura il termine \emph{serverless} copre sia un modello economico
(pagamento a consumo, granularit\`a fine, nessun costo a riposo) sia un modello
di programmazione. Qui interessa il secondo, e conviene definirlo per ci\`o che
richiede alla piattaforma anzich\'e per ci\`o che promette all'utente. Una
piattaforma FaaS \`e un sistema che soddisfa quattro obblighi:

\begin{enumerate}[leftmargin=*]
  \item \textbf{Registrazione.} Accetta la definizione di un'unit\`a di calcolo
  --- un'immagine, un comando, un insieme di variabili d'ambiente, dei limiti di
  risorsa --- e le associa un nome logico stabile.
  \item \textbf{Attivazione.} Su richiesta, garantisce l'esistenza di almeno
  un'istanza in esecuzione capace di servire quel nome, creandola se assente.
  \item \textbf{Instradamento.} Consegna una richiesta a una delle istanze e ne
  restituisce l'esito al chiamante, rispettando un contratto di
  serializzazione noto.
  \item \textbf{Ritiro.} Riduce il numero di istanze quando la domanda cala,
  eventualmente fino a zero.
\end{enumerate}

Sono obblighi minimi: tutto ci\`o che le piattaforme reali offrono in pi\`u ---
composizione di funzioni, trigger da sorgenti di eventi, versioning, traffic
splitting, isolamento multi-tenant --- \`e sovrastruttura rispetto a questo
nucleo. \nf{} implementa i quattro obblighi e quasi nient'altro; il
\cref{cap:03-requisiti-principi} argomenta perch\'e.

Va osservato che gli obblighi 2 e 4 sono in tensione fra loro, ed \`e da questa
tensione che discende la maggior parte dei problemi interessanti. Ritirare
istanze riduce il costo a riposo ma allunga l'attivazione successiva; mantenerle
accorcia l'attivazione ma paga capacit\`a inutilizzata. Il punto in cui una
piattaforma colloca il proprio compromesso fra i due \`e, di fatto, la sua
identit\`a.

\section{Anatomia delle piattaforme di riferimento}

Le piattaforme FaaS open source per Kubernetes condividono una struttura
ricorrente, che vale la pena esplicitare perch\'e \`e precisamente la struttura
da cui \nf{} si discosta.

\paragraph{Knative Serving~\cite{knative}.} \`E la realizzazione pi\`u
articolata del modello. Una richiesta entra da un ingress di rete, raggiunge un
componente chiamato \emph{activator} quando la revisione ha zero repliche,
attende che l'\emph{autoscaler} (KPA) osservi la domanda e comandi allo
scheduler di Kubernetes la creazione di pod, e viene infine instradata a un
sidecar (\emph{queue-proxy}) che, all'interno del pod applicativo, applica il
limite di concorrenza e riporta le metriche all'autoscaler. Il modello di
concorrenza \`e esplicito e centrale: la grandezza \code{containerConcurrency}
governa contemporaneamente l'ammissione locale e il segnale di scaling. La
qualit\`a del progetto \`e alta e il numero di componenti che una singola
invocazione attraversa \`e, di conseguenza, elevato.

\paragraph{Apache OpenWhisk~\cite{openwhisk}.} Adotta un'architettura a
\emph{controller} e \emph{invoker}: il controller riceve l'attivazione, consulta
un database (CouchDB) per la definizione dell'azione, e la accoda su Kafka verso
un invoker, che mantiene un pool di container caldi e ne riusa uno se
disponibile. La distinzione fra container \emph{cold}, \emph{prewarm} e
\emph{warm} \`e resa esplicita nel progetto, e la letteratura sul cold start ha
spesso usato OpenWhisk come piattaforma di riferimento~\cite{wang2018peeking}.
Il costo \`e la dipendenza da due sistemi di stato esterni sul percorso di ogni
invocazione.

\paragraph{OpenFaaS~\cite{openfaas}.} Semplifica considerevolmente: un gateway
riceve la richiesta e la inoltra a un Service Kubernetes che bilancia sulle
repliche del Deployment della funzione; lo scaling \`e delegato a componenti
esterni (HPA o KEDA). Il percorso sincrono \`e corto, e l'asincrono passa per
una coda (NATS). \`E la piattaforma pi\`u vicina a \nf{} per filosofia, ma
delega all'esterno esattamente le politiche che \nf{} vuole rendere
sostituibili.

\paragraph{Fission~\cite{fission}.} Introduce il concetto di \emph{pool} di pod
generici pre-avviati in cui il codice della funzione viene iniettato al momento
dell'uso, riducendo drasticamente il tempo di attivazione a scapito
dell'isolamento e della generalit\`a del runtime.

\Cref{tab:piattaforme} riassume il confronto. La colonna che conta ai fini di
questo lavoro \`e l'ultima: il numero di processi distinti che una singola
invocazione sincrona attraversa nel caso caldo.

\begin{table}[htbp]
\centering
\caption{Confronto strutturale fra piattaforme FaaS per Kubernetes.
Il conteggio dei salti \`e indicativo e riferito al percorso sincrono con
repliche gi\`a calde, escluso l'ingress di rete.}
\label{tab:piattaforme}
\small
\begin{tabularx}{\textwidth}{@{}lXccc@{}}
\toprule
\textbf{Piattaforma} & \textbf{Componenti sul percorso} & \textbf{Stato esterno} &
\textbf{Scale-to-zero} & \textbf{Salti} \\
\midrule
Knative     & ingress, activator, queue-proxy, utente & no & s\`i & 3--4 \\
OpenWhisk   & controller, Kafka, invoker, utente      & CouchDB, Kafka & s\`i & 3--4 \\
OpenFaaS    & gateway, Service, utente                & NATS (async) & con add-on & 2 \\
Fission     & router, executor, pod generico, utente  & etcd & s\`i & 2--3 \\
\nf{}       & control plane, utente                   & nessuno & no & 1 \\
\bottomrule
\end{tabularx}
\end{table}

La riga di \nf{} non \`e un vanto prestazionale: un singolo salto \`e
raggiungibile perch\'e la piattaforma rinuncia a tutto ci\`o che gli altri salti
fornivano. \`E per\`o la condizione che rende possibile l'esperimento descritto
nel \cref{cap:25-valutazione-concorrenza}, in cui una variazione di alcuni
millisecondi nel tempo di attesa dev'essere attribuita a una politica di
ammissione e non a una catena di proxy interposti.

\section{Il cold start e le sue mitigazioni}

L'attivazione di una funzione che non ha istanze pronte impone al chiamante un
ritardo composto da pi\`u contributi: la decisione di scheduling, il pull
dell'immagine, la creazione del sandbox, l'inizializzazione del runtime (che per
una JVM o un interprete Python con dipendenze pesanti domina spesso tutti gli
altri) e infine l'esecuzione dell'handler.

La letteratura ha attaccato ciascun contributo separatamente.

\begin{itemize}[leftmargin=*]
  \item \textbf{Sandbox pi\`u leggeri.} Firecracker~\cite{agache2020firecracker}
  dimostra che una microVM con superficie hardware ridotta pu\`o avviarsi in
  poche decine di millisecondi mantenendo isolamento a livello di hypervisor.
  SOCK~\cite{oakes2018sock} riprogetta il container per il caso serverless,
  eliminando dal percorso di avvio le operazioni del runtime container generico
  che il caso d'uso non richiede.

  \item \textbf{Riuso dello stato inizializzato.}
  Catalyzer~\cite{du2020catalyzer} e i lavori su snapshot e
  restore~\cite{ustiugov2021snapshots} sostituiscono l'inizializzazione con il
  ripristino di un'immagine di memoria gi\`a inizializzata, spostando il costo
  dal percorso critico a una fase preparatoria.

  \item \textbf{Politiche di ritenzione.} FaasCache~\cite{fuerst2021faascache}
  riformula il mantenimento dei container caldi come un problema di caching e
  applica politiche di sostituzione greedy-dual; \cite{shahrad2020wild} mostra su
  traccia di produzione che le distribuzioni di inter-arrivo reali sono
  sufficientemente strutturate da rendere efficace la predizione.
\end{itemize}

\nf{} non contribuisce a questa linea. Misura il cold start
(\code{function\_cold\_start\_ms}, \cref{cap:21-osservabilita}) e lo tratta come
una grandezza osservabile del proprio ambiente, ma non introduce meccanismi per
ridurlo. La scelta \`e coerente con il posizionamento del progetto: il cold
start su \nf{} \`e dominato dal tempo di creazione del pod da parte di
Kubernetes e dall'avvio del runtime della funzione, entrambi fuori dal
controllo del control plane.

\section{Elasticit\`a: quante repliche}

Il problema dell'autoscaling di applicazioni elastiche \`e antecedente al
serverless e ne esiste una rassegna
consolidata~\cite{loridobotran2014autoscaling}, che classifica gli approcci in
reattivi a soglia, basati su teoria del controllo, su modelli di code, su
apprendimento per rinforzo e su serie temporali.

Nel contesto Kubernetes le due realizzazioni dominanti sono reattive a soglia.
L'Horizontal Pod Autoscaler~\cite{k8shpa} confronta periodicamente un valore
osservato (tipicamente l'utilizzazione di CPU, oppure una metrica custom) con
un valore obiettivo e calcola il numero di repliche come rapporto fra i due,
applicando bande morte e finestre di stabilizzazione per evitare oscillazioni.
KEDA~\cite{keda} generalizza la sorgente della metrica a un catalogo esteso di
\emph{scaler} orientati agli eventi, e --- soprattutto --- consente lo scaling
a zero, che l'HPA di base non permette.

Knative adotta un autoscaler dedicato la cui metrica primaria non \`e
l'utilizzazione di CPU ma la \emph{concorrenza media per replica}, ossia il
numero medio di richieste contemporaneamente in carico su ciascun pod. \`E una
scelta significativa, perch\'e lega la decisione di scaling alla stessa
grandezza che governa l'ammissione locale.

\nf{} riproduce questa impostazione nel modulo \code{autoscaler}
(\cref{cap:11-autoscaler}), che pu\`o decidere da richieste al secondo,
profondit\`a di coda o richieste in volo. La differenza rilevante \`e che \nf{}
tratta il numero di repliche e il limite di concorrenza per funzione come due
decisioni \emph{separabili e separatamente sostituibili}, mentre nelle
piattaforme citate sono accoppiate nello stesso componente.

\section{Ammissione, backpressure e controllo della concorrenza}

Quante richieste ammettere contemporaneamente \`e un problema pi\`u antico del
serverless e con una letteratura pi\`u matura.

\paragraph{Il risultato strutturale.} La legge di Little~\cite{little1961proof}
lega il numero medio di richieste presenti nel sistema $L$, il tasso di arrivo
$\lambda$ e il tempo medio di permanenza $W$ dalla relazione $L = \lambda W$.
Applicata a un sistema con limite di concorrenza $C$ saturo, fornisce
immediatamente $W = C / \lambda$: fissato il tasso di arrivo, il tempo di
permanenza \`e proporzionale al limite. Questa relazione, banale in s\'e, ha una
conseguenza che il \cref{cap:12-concurrency-control} sviluppa in dettaglio: un
controllore che aumenta $C$ per aumentare il throughput sta necessariamente
allungando la permanenza, e i due obiettivi non possono essere ottimizzati
indipendentemente. La teoria delle code~\cite{kleinrock1975queueing} fornisce
l'apparato per rendere quantitativa l'osservazione.

\paragraph{Architetture a stadi.} SEDA~\cite{welsh2001seda} propone di
strutturare un servizio come una pipeline di stadi separati da code esplicite,
ciascuno con il proprio pool di thread dimensionato dinamicamente. Il contributo
concettuale rilevante qui \`e l'idea che le code \emph{debbano essere esplicite e
osservabili}: un servizio ben condizionato \`e uno in cui il punto in cui le
richieste si accumulano \`e noto e strumentato, anzich\'e disperso nei buffer
del sistema operativo. Il modulo \code{sync-queue} di \nf{}
(\cref{cap:10-sync-queue}) segue direttamente questa impostazione.

\paragraph{Controllo dal ritardo.} La linea di lavoro sul controllo di
congestione fornisce due meccanismi che \nf{} adatta.
TCP Vegas~\cite{brakmo1995vegas} stima il ritardo di coda confrontando il
round-trip time corrente con il minimo osservato --- assunto come stima del
ritardo di propagazione puro --- e regola la finestra di trasmissione in base
alla differenza, anzich\'e attendere la perdita di pacchetti come segnale.
CoDel~\cite{nichols2012codel} osserva che il segnale utile non \`e la lunghezza
della coda ma il \emph{ritardo minimo sperimentato in una finestra temporale},
e scarta pacchetti quando tale minimo resta sopra soglia per un intervallo
prolungato. Entrambi condividono la stessa intuizione: il ritardo di coda \`e
distinguibile dal costo di servizio irriducibile solo confrontando l'osservazione
corrente con un riferimento storico.

La libreria \code{concurrency-limits} di Netflix~\cite{netflixlimits} porta
esplicitamente questi algoritmi nel dominio delle chiamate RPC, implementando
limitatori in stile Vegas e Gradient che aggiustano un limite di concorrenza
osservando la latenza. \nf{} riprende il segnale in stile Vegas nel modo
\code{ADAPTIVE\_PER\_POD} (\cref{cap:12-concurrency-control}), e il
\cref{cap:25-valutazione-concorrenza} riporta il motivo per cui, nel contesto
FaaS, questo segnale si \`e rivelato insufficiente.

\paragraph{Equit\`a nella ripartizione.} Quando pi\`u flussi competono per una
risorsa condivisa e la somma delle richieste eccede la disponibilit\`a, la
nozione di allocazione max-min~\cite{bertsekas1992datanetworks} fornisce un
criterio: si massimizza l'allocazione minima, e chi chiede meno della propria
quota equa la ottiene per intero, mentre l'eccedenza viene ridistribuita fra i
restanti. Nella variante pesata ciascun flusso porta un peso che scala la
propria quota. Il modo \code{BUDGETED} di \nf{} applica questo criterio ai limiti
di concorrenza per funzione, trattando la capacit\`a della piattaforma come la
risorsa condivisa e il fabbisogno stimato di ciascuna funzione come la sua
domanda.

\section{Offloading fra periferia e cloud}

L'idea di spostare il calcolo dal dispositivo o dal nodo periferico verso
un'infrastruttura pi\`u capiente quando la periferia non basta risale a prima
del serverless~\cite{bonomi2012fog, shi2016edge, satya2017edge}. La formulazione
classica \`e un problema di decisione: dato un compito, valutare se eseguirlo
localmente o trasferirlo, sulla base di un modello del costo di trasferimento e
del guadagno di esecuzione.

La declinazione che \nf{} adotta nel modulo \code{offload}
(\cref{cap:13-offload}) \`e volutamente pi\`u modesta e pi\`u operativa. La
decisione non \`e presa da un modello ma da due segnali osservati: una politica
per funzione dichiarata dall'operatore (\emph{eager}), oppure il rifiuto della
coda di ammissione locale per profondit\`a o per attesa stimata eccessiva
(\emph{pressure}). Il trasferimento \`e a singolo salto e senza fallback locale.
Questa combinazione --- decidere dal rifiuto anzich\'e da un modello di costo ---
non risolve il problema di ottimizzazione classico, ma ha la propriet\`a di non
richiedere alcuna calibrazione preventiva, e di degradare in modo prevedibile.

\section{Benchmarking delle piattaforme serverless}

Confrontare piattaforme FaaS \`e notoriamente difficile: il risultato dipende dal
carico, dalla dimensione del payload, dal profilo di arrivo, dall'hardware, e
dallo stato caldo o freddo del sistema al momento della misura.
SeBS~\cite{copik2021sebs} e ServerlessBench~\cite{yu2020serverlessbench}
propongono suite di carichi standardizzati e metriche comparabili, e
\cite{wang2018peeking} mostra quanto informazione si possa estrarre sulle
politiche interne di una piattaforma commerciale per sola osservazione esterna.

Il metodo sperimentale di \nf{} (\cref{cap:22-metodo-sperimentale}) non aspira
alla comparabilit\`a inter-piattaforma. Persegue invece la comparabilit\`a
\emph{intra}-piattaforma nel tempo: fissato un profilo hardware e un carico, i
risultati di ogni release sono confrontati con quelli della release precedente
sullo stesso profilo, e la deviazione \`e trattata come un segnale di
regressione. \`E un obiettivo pi\`u ristretto e pi\`u difendibile.

\section{Posizionamento}

Riassumendo, \nf{} non contende alle piattaforme mature il terreno della
completezza funzionale, e non contribuisce alla riduzione del cold start. Si
colloca invece nello spazio, meno affollato, di una piattaforma FaaS
sufficientemente piccola da essere usata come banco di prova per politiche di
ammissione, scaling e concorrenza, con tre elementi di novit\`a rispetto a
quanto sopra:

\begin{enumerate}[leftmargin=*]
  \item la \textbf{separazione esplicita} fra la decisione sul numero di repliche
  e la decisione sul limite di concorrenza per funzione, mantenute in due moduli
  indipendenti e sostituibili separatamente;
  \item la formulazione di un controllore che decide dal \textbf{tempo di
  sojourn} anzich\'e dal tempo di servizio, con la dimostrazione --- analitica e
  sperimentale --- che la seconda scelta non ammette ottimo interno;
  \item un modo di allocazione \textbf{a budget di piattaforma} che tratta la
  concorrenza come risorsa condivisa soggetta a equit\`a max-min pesata, anzich\'e
  come parametro che ogni funzione negozia per conto proprio.
\end{enumerate}

\chapter{Requisiti e principi di progetto}
\label{cap:03-requisiti-principi}

Il \cref{cap:01-introduzione} ha enunciato la tesi del lavoro: un insieme
selezionato di rinunce produce un sistema pi\`u utile come strumento di misura
di quanto lo sarebbe la sua versione completa. Questo capitolo rende operativa
quella tesi. Per ciascun vincolo dichiara che cosa esclude, che cosa
\emph{abilita}, e in quale punto del documento se ne trova la conseguenza
sperimentale.

\section{I vincoli, e che cosa comprano}

\subsection{Processo di controllo singolo}

\nf{} gira in un solo processo. Non esiste modalit\`a distribuita, non esiste
elezione di leader, non esiste replicazione dello stato di controllo. Il codice
non incapsula questa assunzione dietro un'astrazione: le strutture dati di
coordinamento sono mappe concorrenti in memoria, non registri replicati.

\paragraph{Che cosa esclude.} L'alta disponibilit\`a. Un riavvio del pod
interrompe tutte le invocazioni in volo e azzera il registro delle funzioni. La
piattaforma non \`e adatta a un servizio la cui indisponibilit\`a abbia costo.

\paragraph{Che cosa abilita.} Tre cose distinte, tutte necessarie nei capitoli
sperimentali.

Primo, l'\textbf{osservabilit\`a totale dello stato di controllo}. Con un solo
processo, la profondit\`a di una coda, il numero di richieste in volo per
funzione e il limite di concorrenza corrente sono variabili di un unico spazio
di indirizzamento, leggibili in modo consistente nello stesso istante. In un
control plane distribuito la stessa lettura richiederebbe un'aggregazione, che
introduce ritardo e incertezza proprio nelle grandezze che il
\cref{cap:25-valutazione-concorrenza} deve misurare con risoluzione di
millisecondi.

Secondo, l'\textbf{assenza di dinamiche di coordinamento} nel segnale misurato.
Un controllore di concorrenza replicato deve riconciliare decisioni prese da
istanze diverse su osservazioni parzialmente sovrapposte; le oscillazioni che ne
derivano sono un fenomeno interessante ma \emph{diverso} da quello studiato qui,
e sarebbero indistinguibili da esso.

Terzo, la \textbf{brevit\`a del percorso critico}, quantificata nella
\cref{tab:piattaforme}: una singola invocazione sincrona attraversa un solo
processo di controllo prima di raggiungere il codice utente.

\subsection{Stato in memoria}

Il registro delle funzioni, lo storico delle esecuzioni, le chiavi di
idempotenza e le code sono tutte strutture in memoria. Nulla \`e persistito su
disco o su un sistema di stato esterno.

\paragraph{Che cosa esclude.} La sopravvivenza a un riavvio, e con essa
qualunque garanzia di durabilit\`a sulle invocazioni asincrone accettate ma non
ancora eseguite.

\paragraph{Che cosa abilita.} L'eliminazione dal percorso critico di ogni
operazione di I/O che non sia la chiamata alla funzione stessa. Nelle
piattaforme che consultano un database per risolvere la definizione di una
funzione, quella lettura \`e sul percorso di ogni invocazione; in \nf{} \`e
l'accesso a una \code{ConcurrentHashMap}. La conseguenza \`e che la latenza
misurata dal control plane, al netto dell'esecuzione, \`e dominata da
elaborazione in memoria e non da una coda di I/O il cui comportamento sotto
carico avrebbe dinamiche proprie.

Va notato che questa scelta impone il proprio prezzo anche a regime, non solo al
riavvio: uno store in memoria che non dimentica cresce senza limite. Il
\code{ExecutionStore} (\cref{cap:06-nucleo}) implementa quindi tre orizzonti di
ritenzione distinti --- TTL di conservazione, TTL di alleggerimento del record e
durata massima assoluta --- perch\'e un'esecuzione che non arriva mai a uno stato
terminale non deve poter restare in memoria per sempre.

\subsection{Nessuna autenticazione n\'e autorizzazione}

L'API non autentica i chiamanti e non applica alcun controllo di accesso.

\paragraph{Che cosa esclude.} L'esposizione su una rete non fidata. \nf{} \`e
pensato per girare dentro un perimetro gi\`a protetto --- un cluster di
laboratorio, una VM di test, una rete privata.

\paragraph{Che cosa abilita.} Meno di quanto si potrebbe pensare sul piano delle
prestazioni: il costo di una verifica di token \`e reale ma modesto. Il guadagno
\`e principalmente sulla superficie di codice e sul numero di modi in cui un
esperimento pu\`o fallire per ragioni non attinenti all'esperimento. Un carico di
prova che riceve \code{401} per un token scaduto \`e tempo perso in un contesto
in cui l'unica cosa che si vuole osservare \`e la distribuzione della latenza.

Questa \`e la rinuncia meno difendibile delle tre, e il \cref{cap:27-discussione}
la riprende: \`e la sola che impedisce alla piattaforma di essere usata al di
fuori di un laboratorio, ed \`e anche la sola per cui l'aggiunta non
comprometterebbe nessuno degli obiettivi sperimentali.

\subsection{Prestazioni e latenza prima delle funzionalit\`a}

Dove una funzionalit\`a aggiuntiva comporterebbe lavoro sul percorso di
invocazione, la funzionalit\`a viene rifiutata o spostata fuori dal percorso.

Questo principio ha conseguenze visibili nel codice. Il \code{ExecutionStore}
espone due varianti di lettura, una che restituisce \code{Optional} e una che
restituisce \code{null}, con il commento che la seconda \`e destinata al percorso
caldo per evitare l'allocazione dell'\code{Optional}. Il metodo che risolve
l'idempotenza pu\`o attendere brevemente su una rivendicazione contesa, e per
questo \`e esplicitamente confinato su uno scheduler elastico anzich\'e sul
thread dell'event loop Netty. Sono micro-decisioni; ci\`o che conta \`e che
esistano e siano commentate come tali, perch\'e definiscono il criterio con cui
il progetto arbitra i propri compromessi.

\section{Il principio architetturale: nucleo minimale e moduli}

Il vincolo pi\`u caratterizzante non \`e una rinuncia ma una struttura: il
control plane \`e diviso in un \textbf{nucleo} che implementa i quattro obblighi
del \cref{cap:02-stato-dell-arte} e in \textbf{moduli opzionali} che ne
estendono il comportamento.

La regola che governa la divisione \`e la seguente:

\begin{quote}
Il nucleo deve poter essere compilato ed eseguito senza alcun modulo, e in
quella configurazione deve restare una piattaforma FaaS funzionante --- ridotta,
ma non rotta.
\end{quote}

\`E una regola pi\`u forte di quanto sembri. Non dice che i moduli sono
disattivabili con un flag: dice che il binario prodotto senza moduli non li
contiene affatto. La selezione avviene a tempo di compilazione
(\cref{cap:08-sistema-moduli}), e quindi il codice del nucleo non pu\`o
contenere rami condizionali che testano la presenza di un modulo. Deve invece
dipendere da \emph{interfacce} di cui esiste sempre almeno un'implementazione
neutra.

\subsection{I punti di estensione e i loro default neutri}

Il nucleo dichiara quattro interfacce e fornisce per ciascuna
un'implementazione no-op, attiva quando nessun modulo ne registra una:

\begin{table}[htbp]
\centering
\caption{Punti di estensione del nucleo e comportamento senza moduli.}
\label{tab:estensioni}
\small
\begin{tabularx}{\textwidth}{@{}lXX@{}}
\toprule
\textbf{Interfaccia} & \textbf{Fornita da} & \textbf{Comportamento del default} \\
\midrule
\code{InvocationEnqueuer} & \code{async-queue} &
Dichiara \code{enabled() == false}; il percorso asincrono risponde
\code{501 Not Implemented}, il sincrono dispaccia in linea. \\
\code{SyncQueueGateway} & \code{sync-queue} &
Nessuna ammissione: ogni richiesta sincrona procede senza controllo di
backpressure. \\
\code{ScalingMetricsSource} & \code{async-queue} &
Restituisce zero per profondit\`a e richieste in volo, e dichiara di non essere
una sorgente reale. \\
\code{ImageValidator} & provider di deployment &
Accetta ogni immagine senza verificarne l'esistenza. \\
\bottomrule
\end{tabularx}
\end{table}

Il quarto caso merita una nota, perch\'e non segue lo schema degli altri tre. La
validazione delle immagini non \`e un modulo autonomo: ciascun provider di
deployment possiede il proprio validatore
(\code{KubernetesImageValidator}, \code{DockerImageValidator}) e lo attiva
quando viene selezionato come backend. La ragione \`e che verificare l'esistenza
di un'immagine \`e un'operazione specifica del backend --- interrogare un
registry attraverso l'API Kubernetes non \`e la stessa cosa che interrogare un
demone Docker locale --- e un modulo generico avrebbe dovuto reintrodurre la
distinzione al proprio interno.

\subsection{Quando un default neutro non \`e accettabile}

Il principio ha un'eccezione codificata, ed \`e istruttiva. Il modulo
\code{concurrency-control} legge profondit\`a di coda e richieste in volo da
\code{ScalingMetricsSource}, e --- soprattutto --- \emph{riscrive} attraverso la
stessa interfaccia i limiti che calcola. \`E quella scrittura a rendere i limiti
effettivi. Selezionato senza \code{async-queue}, il modulo leggerebbe zeri e
scriverebbe verso un'implementazione che scarta: girerebbe, produrrebbe metriche
plausibili, e non governerebbe nulla.

Il modulo quindi \textbf{rifiuta di avviarsi} e nomina il modulo mancante. La
distinzione che il progetto traccia \`e fra degradazione \emph{osservabile} e
degradazione \emph{silenziosa}: la prima \`e accettabile, la seconda no. Lo
stesso criterio spiega perch\'e il modulo \code{autoscaler}, che legge la stessa
sorgente ma la cui metrica principale (\code{rps}) proviene da un contatore
indipendente, si limita invece a registrare un avviso quando una delle sue
metriche derivate diventa cieca.

\section{Requisiti funzionali}

\begin{description}[style=nextline, leftmargin=0pt]
  \item[\textbf{RF1 --- Registrazione e ciclo di vita.}] La piattaforma deve
  consentire di registrare, aggiornare, elencare, ispezionare e cancellare una
  funzione tramite API REST, con validazione della specifica al momento della
  registrazione (\cref{cap:05-contratti-dati}).

  \item[\textbf{RF2 --- Invocazione sincrona e asincrona.}] Deve offrire
  un'invocazione sincrona che restituisce l'esito nella stessa risposta HTTP e
  un'invocazione asincrona che restituisce un identificatore di esecuzione
  consultabile in seguito (\cref{cap:06-nucleo}).

  \item[\textbf{RF3 --- Idempotenza.}] Due invocazioni che portano la stessa
  chiave di idempotenza sulla stessa funzione devono corrispondere a una sola
  esecuzione, e la seconda deve poter attendere o riprodurre l'esito della prima
  (\cref{cap:06-nucleo}).

  \item[\textbf{RF4 --- Modalit\`a di esecuzione multiple.}] Deve supportare
  l'esecuzione in-process (per il collaudo), l'inoltro a un endpoint esterno gi\`a
  esistente, e la gestione del ciclo di vita di un deployment attraverso un
  backend sostituibile (\cref{cap:07-modalita-esecuzione}).

  \item[\textbf{RF5 --- Riprova e timeout.}] Deve applicare un timeout per
  invocazione e un numero massimo di tentativi, entrambi configurabili per
  funzione con default di piattaforma.

  \item[\textbf{RF6 --- Osservabilit\`a.}] Deve esporre metriche in formato
  Prometheus su una porta separata da quella dell'API, con granularit\`a per
  funzione su latenza, esito, cold start, profondit\`a di coda e concorrenza
  effettiva (\cref{cap:21-osservabilita}).
\end{description}

\section{Requisiti non funzionali}

\begin{description}[style=nextline, leftmargin=0pt]
  \item[\textbf{RNF1 --- Componibilit\`a verificata.}] Ogni combinazione di moduli
  dichiarata valida deve essere costruibile e la suite di test deve poter girare
  su di essa. La configurazione priva di moduli deve avere una copertura di test
  propria (\cref{cap:16-moduli-diagnostici}).

  \item[\textbf{RNF2 --- Isolamento fra moduli.}] Un modulo non deve dipendere da
  un altro modulo se non attraverso interfacce dichiarate nel nucleo o nel
  modulo condiviso dei contratti. La regola \`e verificata automaticamente
  (\cref{cap:26-qualita-verifica}).

  \item[\textbf{RNF3 --- Non bloccante sul percorso di invocazione.}] Il percorso
  sincrono non deve bloccare thread dell'event loop. Le operazioni
  potenzialmente bloccanti devono essere esplicitamente confinate su scheduler
  dedicati (\cref{cap:04-architettura-control-plane}).

  \item[\textbf{RNF4 --- Compilabilit\`a in immagine nativa.}] L'intero control
  plane, con tutti i moduli, deve essere compilabile con GraalVM Native Image
  senza rinunciare a funzionalit\`a, il che vincola l'uso della riflessione e
  richiede la dichiarazione esplicita degli \emph{hint} di runtime
  (\cref{cap:20-build-packaging}).

  \item[\textbf{RNF5 --- Riproducibilit\`a delle misure.}] Ogni numero
  prestazionale pubblicato deve essere associato a un profilo hardware
  identificato e a un record versionato, e confrontabile con quello della
  release precedente sullo stesso profilo (\cref{cap:22-metodo-sperimentale}).
\end{description}

\section{I default di piattaforma}

I valori di default definiscono il comportamento della piattaforma in assenza di
configurazione esplicita, e vale la pena esaminarli perch\'e ne rivelano le
priorit\`a. Sono dichiarati in \code{application.yml}:

\begin{center}
\small
\begin{adjustbox}{max width=\textwidth}
\begin{tabular}{@{}llp{7.2cm}@{}}
\toprule
\textbf{Chiave} & \textbf{Default} & \textbf{Interpretazione} \\
\midrule
\code{defaults.timeoutMs} & 30\,000 & Generoso: pensato per non far fallire un
esperimento a causa del timeout. \\
\code{defaults.concurrency} & 4 & Conservativo: quattro invocazioni contemporanee
per funzione. \\
\code{defaults.queueSize} & 100 & Coda per funzione sul percorso asincrono. \\
\code{defaults.maxRetries} & 3 & \\
\code{rate.maxPerSecond} & 1\,000\,000 & Di fatto disattivato: il limitatore
esiste ma il default non vincola. \\
\code{sync-queue.max-depth} & 200 & \\
\code{sync-queue.max-estimated-wait} & 2\,s & Soglia di rifiuto per attesa
stimata. \\
\code{scaling.poll-interval-ms} & 5\,000 & Periodo di campionamento
dell'autoscaler. \\
\code{metrics.profile} & \code{basic} & Il profilo esteso \`e opt-in. \\
\code{admin.runtime-config.enabled} & \code{false} & L'API di riconfigurazione a
caldo \`e disattivata per default. \\
\bottomrule
\end{tabular}
\end{adjustbox}
\end{center}

Due default meritano commento. Il limite di frequenza a un milione di richieste
al secondo \`e un limitatore \emph{presente ma neutro}: il meccanismo esiste ed
\`e sul percorso, cos\`i che il suo costo sia incluso in ogni misura, ma non
respinge nulla finch\'e un operatore non lo configura. \`E l'opposto della scelta
usuale, e discende dal principio che un esperimento non deve fallire per una
protezione che non stava studiando. Simmetricamente, l'API di riconfigurazione a
caldo \`e disattivata per default perch\'e espone operazioni di scrittura su un
sistema che non autentica: attivarla \`e una decisione che l'operatore deve
prendere consapevolmente.


\part{Architettura}
\chapter{Architettura del control plane}
\label{cap:04-architettura-control-plane}

\section{Vista d'insieme}

Il control plane di \nf{} \`e un'applicazione Spring Boot che espone due porte:
la \code{8080} per l'API di gestione e invocazione, e la \code{8081} per gli
endpoint di management --- salute e metriche Prometheus. La separazione non \`e
cosmetica: consente di esporre l'API al traffico applicativo senza esporre gli
endpoint diagnostici, e soprattutto fa s\`i che lo scraping delle metriche non
condivida n\'e le code n\'e i limiti di ammissione del traffico che sta
misurando. Uno strumento di misura che si mette in coda dietro il carico che
osserva riporta la propria attesa insieme a quella del sistema.

La \cref{fig:architettura} d\`a la vista d'insieme. Il riquadro esterno \`e il
processo del control plane; al suo interno la fascia superiore contiene i quattro
componenti del nucleo, la fascia inferiore tratteggiata i moduli opzionali. Fuori
dal riquadro, in basso, i tre tipi di destinazione verso cui un'invocazione pu\`o
essere dispacciata.

\begin{figure}[htbp]
  \centering
  \resizebox{0.95\textwidth}{!}{% nanofaas system architecture (icon / "card" style).
% Grouping containers live on the background layer and are declared AFTER every
% blur-shadow node (shadows.blur manipulates layers; a shadowed node emitted
% after background content suppresses it). Keep that ordering when editing.
\documentclass[border=8pt]{standalone}
\usepackage{nanofaas-figs}
\begin{document}
\begin{tikzpicture}[nf]

  % ================= core components (shadowed cards) =================
  \node[iconcore, minimum width=44mm, minimum height=27mm] (registry)
    {{\color{nfnavy}\large\ic{book}}\ \ \textbf{Function Registry}\\[5pt]
     {\footnotesize\color{nfslate}in-memory}};
  \node[iconcore, right=6mm of registry, minimum width=48mm, minimum height=27mm] (invoc)
    {{\color{nfnavy}\large\ic{server}}\ \ \textbf{Invocation Service}\\[5pt]
     {\footnotesize\color{nfslate}sync/async $\cdot$ retries}\\
     {\footnotesize\color{nfslate}idempotency $\cdot$ rate limit}};
  \node[iconcore, right=6mm of invoc, minimum width=46mm, minimum height=27mm] (dispatch)
    {{\color{nfnavy}\large\ic{sitemap}}\ \ \textbf{Pool Dispatcher}\\[5pt]
     {\footnotesize\color{nfslate}LOCAL / POOL /}\\
     {\footnotesize\color{nfslate}DEPLOYMENT}};
  \node[iconcore, right=6mm of dispatch, minimum width=44mm, minimum height=27mm] (store)
    {{\color{nfnavy}\large\ic{database}}\ \ \textbf{Execution Store}\\[5pt]
     {\footnotesize\color{nfslate}lifecycle $\cdot$ TTL eviction}};

  % ================= optional modules (shadowed cards in a grid) =================
  \matrix (mods) [matrix of nodes, anchor=north,
      nodes={iconmod, minimum width=59mm, anchor=center},
      column sep=5mm, row sep=3mm]
      at ($(registry.south)!0.5!(store.south)+(0,-17mm)$) {
    |[iconmod]| {\color{nfgreen}\ic{envelope}}~~async-queue &
    |[iconmod]| {\color{nfgreen}\ic{sync-alt}}~~sync-queue &
    |[iconmod]| {\color{nfgreen}\ic{chart-bar}}~~autoscaler \\
    |[iconmod]| {\color{nfgreen}\ic{cloud-upload-alt}}~~offload &
    |[iconmod]| {\color{nfgreen}\ic{cog}}~~runtime-config &
    |[iconmod]| {\color{nfgreen}\ic{shield-alt}}~~image-validator \\
    |[iconmod]| {\color{nfgreen}\ic{dharmachakra}}~~k8s-deployment-provider &
    |[iconmod]| {\color{nfgreen}\ic{docker}}~~container-deployment-provider &
    |[iconmod]| {\color{nfgreen}\ic{tag}}~~build-metadata \\
  };

  % ================= execution targets (shadowed cards) =================
  \node[iconback, minimum width=50mm, minimum height=21mm] (clocal)
    at ($(mods.south)+(0,-22mm)$)
    {{\color{nfamber}\large\ic{docker}}\ \ \textbf{Local containers}\\[4pt]
     {\footnotesize\color{nfslate}container-local $\cdot$ Docker CLI}};
  \node[iconback, left=9mm of clocal, minimum width=48mm, minimum height=21mm] (k8s)
    {{\color{nfamber}\large\ic{dharmachakra}}\ \ \textbf{Kubernetes}\\[4pt]
     {\footnotesize\color{nfslate}Deployment + Service}};
  \node[iconback, right=9mm of clocal, minimum width=42mm, minimum height=21mm] (pool)
    {{\color{nfamber}\large\ic{fire}}\ \ \textbf{Warm pod pool}\\[4pt]
     {\footnotesize\color{nfslate}OpenWhisk-style}};

  % ================= client =================
  \node[iconcore, minimum width=104mm, minimum height=23mm] (client)
    at ($(registry.north)!0.5!(store.north)+(0,44mm)$)
    {{\color{nfnavy}\LARGE\ic{users}}\ \ \textbf{\large Client / SDK}\quad
     \begin{tabular}[c]{@{}l@{}}
       {\footnotesize\ttfamily POST /v1/functions/\{name\}:invoke~(sync)}\\
       {\footnotesize\ttfamily POST /v1/functions/\{name\}:enqueue~(async)}
     \end{tabular}};

  % strip reserving room for the banner at the top of the control plane
  \node[inner sep=0, minimum height=10mm, anchor=south] (topstrip)
    at ($(registry.north)!0.5!(store.north)+(0,3mm)$) {};

  % ================= grouping containers (background, AFTER all shadowed nodes) ===
  \begin{scope}[on background layer]
    \node[container, draw=nfblue, line width=1.6pt, fill=nfblue!4,
      fit=(topstrip)(registry)(store)(mods)] (cp) {};
    \node[container, draw=nfgreen, line width=1.1pt, dash pattern=on 3pt off 2.4pt,
      fill=nfgreen!7, fit=(mods)] (modbox) {};
  \end{scope}

  % ================= banner + labels (no shadow) =================
  \node[banner, anchor=west] at ($(registry.north west)+(-1mm,8mm)$)
    {NanoFaaS Control Plane\, {\normalsize\mdseries(Java / Spring Boot)}};
  \node[font=\sffamily\small, text=nfnavy, anchor=east]
    at ($(store.north east)+(1mm,8mm)$) {\texttt{:8080} API \; $\cdot$ \; \texttt{:8081} metrics};
  \node[clabel, text=nfgreen, anchor=south west] at ($(mods.north west)+(-1mm,4mm)$)
    {\textbf{Optional modules}\, {\mdseries\footnotesize(loaded via \texttt{ControlPlaneModule} SPI)}};

  % ================= edges =================
  \draw[flow] (client.south -| invoc) -- (invoc |- cp.north);

  \coordinate (hub) at ($(cp.south)+(0,-9mm)$);
  \draw[flowc] (cp.south) -- (hub)
    node[tag, midway, right=1mm, fill=white, inner sep=1.5pt]
      {Dispatch to execution targets $\cdot$ LOCAL / POOL / DEPLOYMENT};
  \draw[flowc] (hub) -| (k8s.north);
  \draw[flowc] (hub) -- (clocal.north);
  \draw[flowc] (hub) -| (pool.north);

\end{tikzpicture}
\end{document}
}
  \caption{Architettura di \nf{}: control plane a processo singolo, moduli
  opzionali caricati via SPI, backend di esecuzione sostituibili.}
  \label{fig:architettura}
\end{figure}

I quattro componenti del nucleo si dividono le responsabilit\`a nel modo
seguente. Il primo e il quarto vivono in moduli diversi da quello del control
plane (\cref{cap:08-sistema-moduli}): lo store, la macchina a stati dei
tentativi e il motore di scheduling stanno in \code{execution-runtime}, mentre
i contratti con cui i moduli si innestano nel nucleo (le interfacce
\code{ControlPlaneModule}, dei provider e dei punti di estensione) stanno nella
libreria \code{control-plane-spi}.

\begin{description}[style=nextline, leftmargin=0pt]
  \item[\code{FunctionRegistry}] custodisce le specifiche delle funzioni
  registrate. \`E uno snapshot immutabile sostituito per intero a ogni scrittura
  e reso durevole su file; il \cref{cap:06-nucleo} ne descrive la
  persistenza, la risoluzione dei default e l'ordine di registrazione e
  rimozione.

  \item[\code{InvocationService}] orchestra le invocazioni: risolve la
  specifica, crea o riusa un record di esecuzione, e affida il resto al
  coordinatore reattivo. Il limite di frequenza non \`e pi\`u qui: sta in un
  filtro che precede la lettura del corpo (\cref{sec:04-percorso}).

  \item[\code{DispatcherRouter}] sceglie il dispatcher concreto in base alla
  modalit\`a di esecuzione della funzione (\cref{cap:07-modalita-esecuzione}).

  \item[\code{ExecutionStore}] mantiene le esecuzioni vive e gli esiti di quelle
  concluse, in due cache con scadenza e limite propri; la transizione terminale
  \`e di \code{ExecutionLifecycle}.
\end{description}

\section{Perch\'e reattivo}

Il control plane usa Spring WebFlux~\cite{springwebflux} su Netty, non lo stack
servlet. La scelta \`e determinata dal profilo di carico che il componente
sostiene.

Il control plane non calcola: attende. Per ogni invocazione sincrona, il lavoro
che gli compete --- validazione, risoluzione della specifica, creazione del
record --- occupa qualche decina di microsecondi; poi il processo resta in
attesa della risposta della funzione, che dura da qualche millisecondo a
decine di secondi. Con un modello a thread-per-richiesta, sostenere $N$
invocazioni contemporanee richiede $N$ thread bloccati, ciascuno con il proprio
stack. A $N$ nell'ordine delle migliaia il costo di memoria degli stack diventa
comparabile all'intera impronta del processo, e il costo di context switch
inizia a comparire nella latenza misurata --- ossia nella grandezza che gli
esperimenti di questo documento devono osservare pulita.

Con il modello reattivo la richiesta in attesa \`e un oggetto in heap, non un
thread. Il costo per richiesta in volo scende di due ordini di grandezza, e il
numero di thread resta legato al numero di core anzich\'e al numero di
richieste.

\subsection{Il prezzo: la disciplina del non bloccare}

Il modello reattivo impone un vincolo che la programmazione bloccante non ha:
nessuna operazione che possa attendere deve essere eseguita su un thread
dell'event loop. Un event loop bloccato non rallenta una richiesta, le blocca
tutte.

Nel codice di \nf{} questo vincolo compare in due punti espliciti, entrambi
commentati come tali. Il primo \`e la preparazione dell'invocazione sincrona, e
la sua forma \`e cambiata in modo istruttivo:

\begin{lstlisting}[language=Java, caption={Confinamento condizionale della fase
preparatoria (\code{InvocationService}).}]
Mono<Prepared> prepared = Mono.fromCallable(() -> {
    FunctionSpec spec = functionService.get(functionName)
            .orElseThrow(FunctionNotFoundException::new);
    refuseEarlyIfQueueFull(functionName, spec, idempotencyKey, InvocationKind.SYNC);
    return new Prepared(spec, executionFactory.createOrReuseExecution(
            functionName, spec, request, idempotencyKey, traceId, InvocationKind.SYNC));
});
// Hop only when there is something to hop for. ...
if (idempotencyKey != null && !idempotencyKey.isBlank()) {
    prepared = prepared.subscribeOn(Schedulers.boundedElastic());
}
return prepared.flatMap(p -> reactiveCoordinator.invoke(
        p.lookup(), p.spec(), timeoutOverrideMs, offloadContext));
\end{lstlisting}

La risoluzione dell'idempotenza usa un ciclo di \emph{compare-and-set} su una
cache concorrente, che sotto contesa pu\`o attendere. Una versione precedente
confinava sullo scheduler elastico \emph{ogni} richiesta; il commento attuale
spiega perch\'e non pi\`u: \code{createOrReuseExecution} sosta solo dentro quel
ciclo e, senza chiave, ritorna al primo ramo senza toccare alcun lock, e il resto
del blocco \`e un controllo di coda e una ricerca in mappa. Il passaggio a un
altro thread costa 6,1~$\mu$s misurati contro 0,085~$\mu$s di lavoro, settanta
volte ci\`o che protegge, e lo pagavano anche le richieste che sarebbero state
rifiutate, per le quali decidere un \code{429} non dovrebbe richiedere un secondo
thread. Oggi il salto avviene \emph{solo} se la richiesta porta una chiave.

Il secondo punto \`e l'intero percorso asincrono (\code{:enqueue}), che il
controller esegue sempre sullo scheduler elastico: l'accodamento pu\`o contendere
sullo stato della coda della funzione, e non deve toccare l'event loop.

\subsection{Il completamento come futuro condiviso}

Il punto architetturalmente pi\`u interessante del percorso sincrono \`e come
il control plane attende l'esito. Ogni \code{ExecutionRecord} possiede un
\code{CompletableFuture} che verr\`a completato da chiunque porti l'esito ---
il dispatcher che riceve la risposta HTTP, oppure la callback che la funzione
invia all'endpoint interno \code{/v1/internal/executions/\{id\}:complete}. Il
percorso sincrono si limita a sottoscrivere quel futuro:

\begin{lstlisting}[language=Java, caption={Attesa dell'esito nel
\code{ReactiveInvocationCoordinator}.}]
// suppressCancel=true: a single subscriber's timeout/disconnect must not cancel
// the shared completion future other idempotent waiters depend on.
return Mono.fromFuture(executionRecord.completion(), true)
        .onErrorMap(CompletionException.class,
                ex -> ex.getCause() != null ? ex.getCause() : ex)
        .timeout(Duration.ofMillis(timeoutMs))
        ...
\end{lstlisting}

Il parametro \code{suppressCancel} merita attenzione perch\'e codifica un
requisito che nasce dall'idempotenza. Se due chiamanti inviano la stessa chiave
di idempotenza, entrambi finiscono ad attendere lo \emph{stesso} futuro. Se il
primo scade o chiude la connessione, la cancellazione del suo \code{Mono} non
deve propagarsi al futuro condiviso, o il secondo chiamante osserverebbe un
fallimento causato dall'abbandono di qualcun altro. La soppressione della
cancellazione \`e ci\`o che rende l'idempotenza corretta anche in presenza di
disconnessioni.

\section{Il percorso di una richiesta}
\label{sec:04-percorso}

La \cref{fig:flusso} contrappone i due percorsi. Vale la pena percorrerli in
prosa, perch\'e l'ordine delle decisioni \`e esso stesso un elemento di
progetto.

\begin{figure}[htbp]
  \centering
  \resizebox{0.95\textwidth}{!}{% nanofaas invocation flow: synchronous (:invoke) and asynchronous (:enqueue)
% paths as a UML-style sequence diagram with an alt combined fragment.
\documentclass[border=8pt]{standalone}
\usepackage{nanofaas-figs}
\begin{document}
\begin{tikzpicture}[nf]

  % ---- lifeline x-coordinates (cm) ----
  \def\cx{0}\def\px{4.0}\def\qx{8.2}\def\dx{11.5}\def\rx{15.2}
  \def\ybot{-10.6}

  % ---- participant heads ----
  \node[client, minimum width=22mm] (cH) at (\cx,0) {Client / SDK};
  \node[core,   minimum width=34mm] (pH) at (\px,0)
    {Control Plane\\{\scriptsize validate $\cdot$ rate limit $\cdot$ ExecutionStore}};
  \node[module, minimum width=30mm] (qH) at (\qx,0)
    {Admission queue\\{\scriptsize sync-queue / async-queue}};
  \node[core,   minimum width=24mm] (dH) at (\dx,0) {DispatcherRouter};
  \node[backend,minimum width=26mm] (rH) at (\rx,0)
    {Function runtime\\{\scriptsize LOCAL $\cdot$ EXTERNAL $\cdot$ DEPLOYMENT}};

  % ---- lifelines ----
  \foreach \n in {cH,pH,qH,dH,rH}{\draw[lifeline] (\n.south) -- ($(\n.south)-(0,10.1)$);}

  % ---- message helpers ----
  \newcommand{\msg}[4]{\draw[flow] (#1,#3) -- (#2,#3)
      node[tag, midway, above=0.2mm]{#4};}
  \newcommand{\rmsg}[4]{\draw[ret] (#1,#3) -- (#2,#3)
      node[tagc, midway, above=0.2mm]{#4};}

  % ================= alt combined fragment frame =================
  \coordinate (fnw) at (-1.35,-0.7);
  \coordinate (fse) at (16.15,\ybot+0.35);
  \draw[draw=nfslate, line width=0.7pt] (fnw) rectangle (fse);
  \node[anchor=north west, fill=nfslatebg, draw=nfslate, line width=0.5pt,
        font=\sffamily\scriptsize\bfseries, inner sep=2pt, text=nfink] at (fnw) {alt};
  \node[anchor=west, tagc, font=\sffamily\scriptsize\itshape]
    at ($(fnw)+(0.95,-0.16)$) {[\,synchronous $\cdot$ \texttt{:invoke}\,]};

  % ---------- synchronous operand ----------
  \msg{\cx}{\px}{-1.35}{\texttt{POST \dots:invoke}}
  % self-activity on the control plane
  \draw[flow] ($(\px,-1.95)$) -- ++(0.9,0) -- ++(0,-0.5) -- ++(-0.9,0);
  \node[tag, anchor=west] at ($(\px+0.95,-2.2)$) {validate $\cdot$ rate limit $\cdot$ create exec state};
  \msg{\px}{\qx}{-2.9}{admit \scriptsize(else \texttt{429})}
  \msg{\qx}{\dx}{-3.5}{dispatch}
  \msg{\dx}{\rx}{-4.1}{forward \texttt{/invoke}}
  \rmsg{\rx}{\dx}{-4.6}{result}
  \rmsg{\dx}{\px}{-5.0}{update store}
  \rmsg{\px}{\cx}{-5.4}{\texttt{200} + body}

  % ---------- operand divider ----------
  \draw[dashed, draw=nfslate] (-1.35,-5.9) -- (16.15,-5.9);
  \node[anchor=west, tagc, font=\sffamily\scriptsize\itshape]
    at (-1.30,-6.06) {[\,asynchronous $\cdot$ \texttt{:enqueue}\,]};

  % ---------- asynchronous operand ----------
  \msg{\cx}{\px}{-6.6}{\texttt{POST \dots:enqueue}}
  \msg{\px}{\qx}{-7.15}{enqueue \scriptsize(needs async-queue; else \texttt{501})}
  \rmsg{\px}{\cx}{-7.6}{\texttt{202} + execution id}

  \draw[dotted, draw=nfslate] (-1.0,-8.0) -- (15.2,-8.0);
  \node[fill=white, tagc, font=\sffamily\scriptsize\itshape] at (7.1,-8.0)
    {scheduler drains the queue asynchronously};

  \msg{\qx}{\dx}{-8.55}{dispatch}
  \msg{\dx}{\rx}{-9.1}{forward \texttt{/invoke}}
  \rmsg{\rx}{\dx}{-9.55}{result}
  \rmsg{\dx}{\px}{-9.95}{update store}
  \msg{\px}{\cx}{-10.35}{callback \scriptsize POST \texttt{callbackUrl}}

\end{tikzpicture}
\end{document}
}
  \caption{Invocazione sincrona (\code{:invoke}) e asincrona (\code{:enqueue})
  a confronto.}
  \label{fig:flusso}
\end{figure}

\subsection{Percorso sincrono}

\begin{enumerate}[leftmargin=*]
  \item Due filtri precedono il controller. Il primo, con precedenza massima,
  applica il confine del corpo (\code{413}, \cref{sec:06-capacita}); il
  secondo \`e il limitatore di frequenza globale, ristretto ai suffissi
  \code{:invoke} e \code{:enqueue}: una richiesta respinta risponde \code{429} con
  \code{Retry-After: 1} senza aver decodificato il JSON n\'e cercato la funzione,
  e il suo corpo viene comunque svuotato per non impedire il riuso della
  connessione.

  \item Il controller estrae gli header dedicati (chiave di idempotenza, trace
  id, override del timeout, marcatore di salto di offload) e
  \textbf{ricostruisce} la richiesta con gli header effettivi del chiamante.

  \item Il nome viene risolto nel registro; se assente, \code{404}. Se la coda
  della funzione \`e certamente piena e nulla pu\`o deviare la richiesta altrove, il
  rifiuto avviene qui, prima di costruire un'esecuzione destinata a essere
  abbandonata --- salvo quando \`e presente una chiave, perch\'e la replica di
  un'esecuzione esistente non deve essere respinta (nel commento, 590 rifiuti al
  secondo contro 299 dispatch).

  \item Si crea o si riusa il record di esecuzione, dopo aver riservato le
  quote di ammissione (\cref{sec:06-capacita}). Se la chiave di idempotenza
  corrisponde a un'esecuzione gi\`a conclusa, l'esito viene riprodotto
  immediatamente senza toccare la funzione (\emph{replay}); se l'esito \`e stato
  espulso per capacit\`a, la risposta \`e \code{410}.

  \item \textbf{Ammissione.} \`E il punto in cui il comportamento dipende dai
  moduli presenti, e la logica \`e una catena di ripiego ordinata: se il modulo
  di offload \`e attivo e la politica della funzione impone l'offload eager, la
  richiesta parte verso l'istanza remota; altrimenti si tenta l'ammissione
  locale, che a sua volta usa la coda sincrona se presente, la coda per funzione se
  presente, e in ultima istanza dispaccia direttamente, acquisendo un lease di
  capacit\`a: senza spazio la risposta \`e \code{429}, mai una coda implicita. Se l'ammissione locale
  viene rifiutata per pressione (profondit\`a o attesa stimata) e l'offload
  reattivo \`e abilitato, la richiesta viene deviata anzich\'e respinta.

  \item Il chiamante attende il futuro di completamento, con un proprio budget
  (\code{X-Timeout-Ms} o \code{timeoutMs} della funzione). Allo scadere risponde
  \code{408} soltanto \emph{quel} chiamante: l'esecuzione condivisa prosegue
  (\cref{cap:06-nucleo}).

  \item L'esito viene tradotto in risposta HTTP. Se la funzione ha dichiarato un
  proprio codice di stato, questo diventa il codice della risposta e viene
  segnalato dall'header \code{X-NanoFaaS-Function-Status}.
\end{enumerate}

L'ordine del punto dell'ammissione \`e la parte non ovvia. La sequenza \emph{coda sincrona
$\to$ coda asincrona $\to$ dispatch diretto} significa che il percorso sincrono
degrada in modo continuo al ridursi dei moduli presenti, senza che il codice
contenga un test sul nome dei moduli: ciascun ripiego \`e il default no-op del
punto di estensione precedente. Con un modulo di coda la scelta fra le due code
e l'ordine di servizio sono quelle del motore di scheduling composto
(\cref{cap:09-async-queue,cap:10-sync-queue}).

\subsection{Percorso asincrono}

Il percorso \code{:enqueue} \`e pi\`u corto e ha una differenza qualitativa: non
degrada. Se nessun modulo di coda fornisce l'accodamento asincrono, l'accodatore dichiara
di non supportarlo e l'API risponde \code{501 Not Implemented}. La scelta
\`e deliberata: un'invocazione asincrona eseguita in linea non sarebbe
un'invocazione asincrona degradata, sarebbe una risposta scorretta al contratto,
perch\'e il chiamante ha ricevuto un identificatore assumendo di poter
disconnettersi.

\section{Superficie dell'API}

L'API \`e piccola, e le sue tre famiglie di endpoint corrispondono alle tre
responsabilit\`a della piattaforma.

\begin{center}
\small
\begin{adjustbox}{max width=\textwidth}
\begin{tabular}{@{}llp{6.6cm}@{}}
\toprule
\textbf{Metodo} & \textbf{Percorso} & \textbf{Ruolo} \\
\midrule
\code{POST}   & \code{/v1/functions} & Registra una funzione. \\
\code{GET}    & \code{/v1/functions} & Elenca le funzioni. \\
\code{GET}    & \code{/v1/functions/\{name\}} & Ispeziona una funzione. \\
\code{PATCH}  & \code{/v1/functions/\{name\}} & Aggiorna parzialmente la specifica. \\
\code{DELETE} & \code{/v1/functions/\{name\}} & Rimuove funzione e risorse gestite. \\
\code{PUT/GET} & \code{/v1/functions/\{name\}/replicas} & Imposta / legge le repliche (\code{DEPLOYMENT}). \\
\midrule
\code{POST} & \code{/v1/functions/\{name\}:invoke}  & Invocazione sincrona. \\
\code{POST} & \code{/v1/functions/\{name\}:enqueue} & Invocazione asincrona. \\
\code{GET}  & \code{/v1/executions/\{id\}}          & Stato di un'esecuzione. \\
\midrule
\code{POST} & \code{/v1/internal/executions/\{id\}:complete} &
Callback di completamento dal runtime. \\
\midrule
\code{GET} & \code{/actuator/health} \emph{(:8081)}     & Sonde di liveness/readiness. \\
\code{GET} & \code{/actuator/prometheus} \emph{(:8081)} & Metriche. \\
\midrule
\code{GET/PATCH} & \code{/v1/admin/runtime-config} & Configurazione a caldo (\cref{cap:14-runtime-config}). \\
\bottomrule
\end{tabular}
\end{adjustbox}
\end{center}

La convenzione \code{:verbo} sul nome della risorsa (\code{:invoke},
\code{:enqueue}) segue lo stile delle API di Google Cloud per le azioni che non
sono operazioni CRUD sulla risorsa stessa. \`E preferibile a un sottopercorso
(\code{/functions/\{name\}/invoke}) perch\'e non suggerisce l'esistenza di una
sotto-risorsa che si possa elencare o cancellare.

\subsection{Trattamento degli header}

Il controller applica una regola che vale la pena isolare perch\'e riguarda la
sicurezza del contratto pi\`u che l'ergonomia:

\begin{lstlisting}[language=Java]
/**
 * Rebuild the request with the caller's real headers, lower-cased. The body's own
 * {@code headers} field is overwritten, never merged: a caller must not be able to
 * forge a header it did not send.
 */
\end{lstlisting}

Il corpo della richiesta contiene un campo \code{headers} che l'handler della
funzione vedr\`a. Se il control plane fondesse quel campo con gli header HTTP
reali, un chiamante potrebbe iniettare nel campo un header che non ha inviato,
e una funzione che decide in base agli header sarebbe ingannabile. Il campo
viene quindi \emph{sovrascritto}, mai unito. Dagli header propagati sono esclusi
quelli hop-by-hop e quelli che il control plane lega gi\`a a un parametro
dedicato.

\subsection{Mappatura degli errori}

I codici di stato che l'API produce sono pochi e ciascuno corrisponde a una
condizione precisa; il \cref{cap:A-api} li elenca per endpoint. Vale la pena
anticipare i meno ovvi. Ogni corpo di errore generato dal gestore globale ha i
campi in ordine fisso (\code{error}, \code{message}, poi \code{details} se
presente): una mappa non ordinata cambia l'ordine di iterazione da una JVM
all'altra, e un chiamante che confronta il corpo come testo lo vedrebbe cambiare.

\code{429 Too Many Requests} ha quattro origini distinguibili. Il limitatore di
frequenza risponde senza corpo, con \code{Retry-After: 1}. Il rifiuto della coda
sincrona porta \code{Retry-After} e \code{X-Queue-Reject-Reason} con valore
\code{depth}, \code{est\_wait} o \code{timeout}. Il superamento di una quota di
ammissione porta \code{Retry-After: 1} e un corpo
\code{\{"error":"invocation\_quota\_exceeded","resource":...\}}. La coda per funzione
piena, o l'assenza di capacit\`a nel dispatch diretto, risponde \code{429} senza
altro. Rendere distinguibile la causa \`e necessario perch\'e le condizioni
richiedono reazioni diverse dal chiamante: la piattaforma \`e globalmente satura,
lo \`e la singola funzione, o \`e una risorsa di memoria a esserlo.

\code{410 Gone} segnala la replica di una chiave il cui esito \`e stato espulso
per capacit\`a, con \code{X-Execution-Id}; \code{413} un corpo oltre il confine di
ingresso; \code{409} con codice \code{FUNCTION\_REMOVAL\_PENDING} una funzione la cui
rimozione ha lasciato risorse (\cref{cap:06-nucleo}); \code{408} lo scadere del
budget di un singolo chiamante sincrono.

\code{502 Bad Gateway} e \code{504 Gateway Timeout} sono emessi solo dal
percorso di offload, e segnalano che la richiesta \`e stata deviata verso
un'istanza remota che ha fallito o non ha risposto in tempo. Non esiste
fallback locale, e il \cref{cap:13-offload} argomenta perch\'e.

\section{Modello a thread}

Riassumendo i thread che il processo utilizza a regime. La colonna
\emph{Presenza} distingue ci\`o che esiste sempre da ci\`o che dipende dai moduli o
dalla configurazione.

\begin{center}
\small
\begin{adjustbox}{max width=\textwidth}
\begin{tabular}{@{}lll@{}}
\toprule
\textbf{Gruppo} & \textbf{Presenza} & \textbf{Ruolo} \\
\midrule
Event loop Netty & sempre; pari ai core & I/O HTTP in ingresso e in uscita \\
\code{boundedElastic} & sempre; elastico, limitato & Percorso \code{:enqueue}; invocazioni con chiave \\
\code{execution-expiry} & sempre; 1 (daemon) & Scadenza amministrativa dei record vivi \\
\code{nanofaas-core-retry-timer} & sempre; 1 (daemon) & Riprove ritardate del profilo diretto \\
Timer del cancello di prontezza & sempre; 1 (daemon) & Timeout dell'attesa dei deployment \\
\code{nanofaas-scheduler-engine} & con un modulo di coda; 1 & Motore di scheduling composto (\cref{cap:09-async-queue}) \\
Manutenzione dello stimatore & con \code{sync-queue}; 1 & Stima dell'attesa (\cref{cap:10-sync-queue}) \\
\code{nanofaas-internal-scaler} & con \code{autoscaler}; 1 & Campionamento e scalatura \\
\code{nanofaas-concurrency-governor} & con \code{concurrency-control}; 1 & Controllo della concorrenza \\
\code{nanofaas-admin-config-} & con \code{runtime-config}; 1 & Cambi di configurazione \\
\bottomrule
\end{tabular}
\end{adjustbox}
\end{center}

Il janitor periodico dello store, che la prima versione elencava, non c'\`e pi\`u:
la scadenza \`e affidata alle cache Caffeine (\cref{cap:06-nucleo}) e a un unico
timer di sola pianificazione.
Nella configurazione priva di moduli restano l'event loop, lo scheduler elastico
e i tre timer daemon: \`e ancora un numero ridotto di gruppi per un'intera
piattaforma FaaS, e la misura pi\`u compatta di che cosa significhi il vincolo di
minimalit\`a enunciato nel \cref{cap:03-requisiti-principi}. Va detto per\`o che
non tutti i gruppi sono stati verificati sul processo in esecuzione: l'elenco
proviene dal codice dei componenti, non da un dump dei thread.

\chapter{Contratti e modello dei dati}
\label{cap:05-contratti-dati}

Il modulo \code{platform/common} contiene i tipi che attraversano ogni confine
del sistema: quelli che il control plane scambia con i chiamanti, con i runtime
delle funzioni e con gli SDK nei vari linguaggi. Sono meno di cinquecento righe
di codice, e la loro stabilit\`a \`e ci\`o che consente ai moduli e agli SDK di
evolvere indipendentemente. Questo capitolo li descrive e ne discute le scelte
non ovvie.

\section{\code{FunctionSpec}: la definizione di una funzione}

\code{FunctionSpec} \`e il tipo che il chiamante invia per registrare una
funzione e che il control plane conserva. \`E un \code{record} Java con sedici
componenti, raggruppabili in cinque famiglie.

\begin{center}
\small
\begin{adjustbox}{max width=\textwidth}
\begin{tabular}{@{}llp{6.2cm}@{}}
\toprule
\textbf{Famiglia} & \textbf{Campi} & \textbf{Ruolo} \\
\midrule
Identit\`a & \code{name}, \code{image} &
Obbligatori, non vuoti. Il nome \`e la chiave logica. \\
Esecuzione & \code{command}, \code{env}, \code{runtimeMode},
\code{runtimeCommand} &
Come avviare e come parlare con il processo funzione. \\
Risorse & \code{resources} &
Richieste e limiti di CPU e memoria, validati come coerenti fra loro. \\
Politiche & \code{timeoutMs}, \code{concurrency}, \code{queueSize},
\code{maxRetries} &
Sovrascrivono i default di piattaforma. \\
Collocazione & \code{executionMode}, \code{endpointUrl},
\code{imagePullSecrets}, \code{scalingConfig}, \code{offload} &
Dove la funzione vive e come viene governata. \\
\bottomrule
\end{tabular}
\end{adjustbox}
\end{center}

Tutti i campi delle ultime tre famiglie sono tipi \emph{boxed}
(\code{Integer}, non \code{int}). La scelta \`e deliberata e ricorre in tutto il
modello: il valore \code{null} significa ``non specificato'', ed \`e
distinguibile da qualunque valore esplicito. Un campo \code{int} con default
zero renderebbe indistinguibili ``l'utente non ha detto nulla'' e ``l'utente ha
chiesto zero'', e la risoluzione dei default (\cref{cap:06-nucleo}) non
potrebbe sapere quando intervenire.

\subsection{Validazione dichiarativa}

I vincoli sono espressi con annotazioni Jakarta Validation applicate direttamente
ai componenti del record:

\begin{lstlisting}[language=Java]
public record FunctionSpec(
        @NotBlank String name,
        @NotBlank String image,
        ...
        @Valid ResourceSpec resources,
        @Min(1) Integer timeoutMs,
        @Min(1) Integer concurrency,
        @Min(1) Integer queueSize,
        @Min(0) Integer maxRetries,
        ...
) { ... }
\end{lstlisting}

Il vincolo pi\`u interessante \`e quello che le annotazioni sui singoli campi non
possono esprimere, ossia la coerenza fra richiesta e limite di risorsa. \`E
codificato come metodo di validazione sul record che lo riguarda:

\begin{lstlisting}[language=Java, caption={Vincolo relazionale in
\code{ResourceSpec}.}]
@AssertTrue(message = "resource request must not exceed limit")
public boolean isRequestWithinLimit() {
    if (requests == null || limits == null) {
        return true;
    }
    return notGreater(requests.cpu(), limits.cpu())
            && notGreater(requests.memoryMiB(), limits.memoryMiB());
}
\end{lstlisting}

La conseguenza pratica \`e che una specifica incoerente viene respinta con
\code{400} al momento della registrazione, e non produce un deployment che
Kubernetes rifiuterebbe pi\`u tardi con un messaggio meno leggibile.

\subsection{Costruttori di compatibilit\`a}

\code{FunctionSpec} espone tre costruttori: quello canonico a sedici componenti
e due pi\`u corti che delegano al canonico passando \code{null} per i campi
aggiunti nelle evoluzioni successive (\code{offload} e, prima, \code{scalingConfig}).

\`E una tecnica esplicita di gestione della compatibilit\`a in un progetto che
usa i \code{record}: un record aggiunge un componente e il costruttore canonico
cambia firma, il che rompe ogni chiamante --- comprese, in questo progetto,
diverse centinaia di istanziazioni nei test. I costruttori di compatibilit\`a
assorbono la rottura. Il prezzo \`e che il tipo porta con s\'e la propria storia
di versioni, e il beneficio \`e che l'aggiunta di un campo resta un'operazione a
diff contenuto.

\section{\code{InvocationRequest} e \code{InvocationResponse}}

Il contratto di invocazione \`e volutamente povero.

\begin{lstlisting}[language=Java]
public record InvocationRequest(
        @NotNull(message = "Input payload is required") Object input,
        Map<String, String> metadata,
        Map<String, String> headers
) { ... }
\end{lstlisting}

Il campo \code{input} \`e di tipo \code{Object}: la piattaforma non impone uno
schema al payload e si limita a trasportarlo. \`E \code{@NotNull} ma pu\`o essere
qualunque struttura JSON. La distinzione fra \code{metadata} e \code{headers} \`e
di provenienza: \code{metadata} \`e ci\`o che il chiamante dichiara nel corpo,
\code{headers} \`e ci\`o che il control plane ha effettivamente osservato sulla
richiesta HTTP e ricostruisce prima di consegnarla
(\cref{cap:04-architettura-control-plane}).

La risposta trasporta l'esito e, opzionalmente, una \emph{busta} che la funzione
pu\`o usare per governare la risposta HTTP:

\begin{lstlisting}[language=Java]
public record InvocationResponse(
        String executionId, String status, Object output, ErrorInfo error,
        Integer statusCode, Map<String, String> headers, String encoding
) { ... }
\end{lstlisting}

I tre campi finali sono il meccanismo con cui una funzione pu\`o restituire, per
esempio, un \code{404} applicativo o un contenuto binario codificato in base64.
Il control plane li applica alla risposta soltanto se il codice di stato ricade
nell'intervallo valido, e in tal caso marca la risposta con
\code{X-NanoFaaS-Function-Status: true} cos\`i che il chiamante possa distinguere
un errore deciso dalla funzione da un errore deciso dalla piattaforma. \`E una
distinzione che le piattaforme FaaS spesso non rendono osservabile e la cui
assenza rende ambigua l'interpretazione di un tasso di errore.

\section{\code{ExecutionStatus} e il ciclo di vita}

\code{ExecutionStatus} \`e la proiezione esterna di un'esecuzione, restituita da
\code{GET /v1/executions/\{id\}}. Oltre all'esito porta due campi diagnostici
--- \code{coldStart} e \code{initDurationMs} --- che rendono osservabile per
singola invocazione se questa ha pagato un'inizializzazione e quanto. Sono la
sorgente della metrica \code{function\_cold\_start\_ms}
(\cref{cap:21-osservabilita}).

Lo stato \`e rappresentato come stringa e non come enumerazione, il che \`e una
scelta discutibile su cui vale la pena essere espliciti: agevola l'evoluzione del
vocabolario di stato senza rompere i client, ma sposta a runtime un controllo che
il compilatore avrebbe potuto fare. I valori usati sono \code{queued},
\code{running}, \code{succeeded}, \code{failed} e \code{timeout}.

\section{\code{InvocationResult}: il contratto verso il runtime}

Mentre \code{InvocationResponse} \`e ci\`o che il control plane restituisce al
chiamante, \code{InvocationResult} \`e ci\`o che il \emph{runtime della funzione}
restituisce al control plane --- sia come corpo della risposta HTTP sincrona,
sia come corpo della callback asincrona verso
\code{/v1/internal/executions/\{id\}:complete}.

\begin{lstlisting}[language=Java]
public record InvocationResult(
        boolean success, Object output, ErrorInfo error,
        Integer statusCode, Map<String, String> headers, String encoding
) {
    public static InvocationResult success(Object output) { ... }
    public static InvocationResult successWithEnvelope(Object output, Integer statusCode,
                                                       Map<String, String> headers, String encoding) { ... }
    public static InvocationResult error(String code, String message) { ... }
}
\end{lstlisting}

I metodi factory statici sono l'API che gli SDK usano
(\cref{cap:18-sdk}); l'esistenza di \code{successWithEnvelope} accanto a
\code{success} rende esplicito nel codice della funzione quando si sta
esercitando il controllo sulla risposta HTTP e quando no.

\section{Le enumerazioni di politica}

Cinque enumerazioni definiscono lo spazio delle politiche configurabili. Sono
elencate qui per completezza; ciascuna \`e trattata nel capitolo del modulo che
la implementa.

\begin{center}
\small
\begin{adjustbox}{max width=\textwidth}
\begin{tabular}{@{}lll@{}}
\toprule
\textbf{Enumerazione} & \textbf{Valori} & \textbf{Trattata in} \\
\midrule
\code{ExecutionMode} & \code{LOCAL}, \code{EXTERNAL}, \code{DEPLOYMENT} &
\cref{cap:07-modalita-esecuzione} \\
\code{RuntimeMode} & \code{HTTP}, \code{STDIO}, \code{FILE} &
\cref{cap:17-runtime-funzione} \\
\code{ScalingStrategy} & \code{HPA}, \code{INTERNAL}, \code{NONE} &
\cref{cap:11-autoscaler} \\
\code{ConcurrencyControlMode} & \code{FIXED}, \code{STATIC\_PER\_POD},
\code{ADAPTIVE\_PER\_POD}, \code{BUDGETED}, \code{SOJOURN} &
\cref{cap:12-concurrency-control} \\
\code{SyncQueueRejectReason} & \code{DEPTH}, \code{EST\_WAIT}, \code{TIMEOUT} &
\cref{cap:10-sync-queue} \\
\bottomrule
\end{tabular}
\end{adjustbox}
\end{center}

\section{Il blocco di scaling e la sua composizione}

\code{ScalingConfig} contiene la strategia di scaling, i limiti di repliche, una
lista di metriche e --- annidato --- il blocco di controllo della concorrenza:

\begin{lstlisting}[language=Java]
public record ScalingConfig(
        ScalingStrategy strategy, Integer minReplicas, Integer maxReplicas,
        List<ScalingMetric> metrics, ConcurrencyControlConfig concurrencyControl
) { ... }
\end{lstlisting}

L'annidamento merita una difesa, perch\'e sembra contraddire l'affermazione del
\cref{cap:02-stato-dell-arte} secondo cui \nf{} separa la decisione sul numero di
repliche da quella sul limite di concorrenza. La separazione \`e sul piano dei
\emph{meccanismi} --- due moduli distinti, sostituibili indipendentemente, che
non si conoscono --- non su quello della configurazione, dove entrambe le
politiche descrivono come una funzione consuma capacit\`a e conviene che siano
dichiarate insieme. Il modo \code{STATIC\_PER\_POD} rende evidente perch\'e i due
piani non possano essere del tutto disgiunti: calcola il limite come
\emph{repliche pronte $\times$ obiettivo per replica}, e quindi legge una
grandezza che l'autoscaler determina.

\section{Il punto di estensione dei moduli}

L'interfaccia che definisce un modulo del control plane \`e minima:

\begin{lstlisting}[language=Java]
public interface ControlPlaneModule {
    default String name() { return getClass().getSimpleName(); }
    Set<Class<?>> configurationClasses();
}
\end{lstlisting}

Un modulo dichiara soltanto quali classi di configurazione Spring devono essere
aggiunte al contesto. Non dichiara dipendenze, non dichiara un ordine, non
espone un ciclo di vita. Tutto ci\`o che gli serve lo ottiene dal contenitore
attraverso i punti di estensione del nucleo. Il \cref{cap:08-sistema-moduli}
descrive il meccanismo di caricamento e discute perch\'e questa povert\`a
dell'interfaccia sia un pregio e non una lacuna.

\section{Il contratto dell'handler}

Infine, il contratto che il codice utente implementa:

\begin{lstlisting}[language=Java]
public interface FunctionHandler {
    Object handle(InvocationRequest request);
}
\end{lstlisting}

Un solo metodo, un solo argomento, nessuna eccezione dichiarata, nessun contesto
di esecuzione iniettato. La povert\`a \`e voluta: qualunque cosa l'handler
riceva come contesto diventerebbe parte del contratto e vincolerebbe gli SDK
degli altri linguaggi a riprodurla. Ci\`o che l'handler pu\`o restituire \`e
governato da \code{HandlerResponse} e \code{ResponseHeaderPolicy}, discussi nel
\cref{cap:17-runtime-funzione}.

\section{L'API come artefatto versionato}

La descrizione formale dell'API vive in \code{openapi.yaml} nella radice del
repository. \`E mantenuta a mano, e questa scelta ha un costo che il progetto ha
sperimentato: essendo esterna al codice, pu\`o divergere in silenzio. La verifica
automatica esistente accerta che il file sia presente, non che sia allineato.

\begin{damisurare}
Il disallineamento fra \code{openapi.yaml} e il codice \`e stato rilevato e
corretto manualmente in almeno un'occasione. Una verifica che confronti gli
endpoint dichiarati nel file con quelli effettivamente esposti dal contesto
Spring \`e realizzabile con gli strumenti gi\`a presenti nel progetto, ma non \`e
stata implementata. \todomis{copertura effettiva di \code{openapi.yaml} rispetto
alle rotte registrate}
\end{damisurare}

\chapter{Il nucleo del control plane}
\label{cap:06-nucleo}

Il nucleo \`e ci\`o che resta quando si compila \nf{} senza alcun modulo di coda:
i quattro obblighi di una piattaforma FaaS enunciati nel
\cref{cap:02-stato-dell-arte} vi sono tutti implementati, e il modulo pu\`o solo
aggiungere politiche. Il codice non sta pi\`u in un solo progetto: la macchina a
stati dei tentativi, lo store delle esecuzioni e il registro della capacit\`a
vivono in \code{platform/libs/execution-runtime}, i contratti verso i moduli in
\code{platform/libs/control-plane-spi}, e il resto (registro delle funzioni, API,
dispatcher, metriche) in \code{platform/control-plane}
(\cref{cap:08-sistema-moduli}). Insieme alle tre librerie obbligatorie il modulo
\code{control-plane} conta 17\,334 righe non vuote di produzione, contro le 6\,299 dei
dieci moduli opzionali (misura del 28~settembre~2026, la stessa del
\cref{cap:01-introduzione}). Il capitolo attraversa i componenti
nell'ordine in cui li incontra una richiesta, soffermandosi sulle decisioni che il
codice motiva esplicitamente. Le garanzie sul ciclo di vita, che attraversano pi\`u
componenti, sono raccolte in un contratto scritto
(\code{docs/architecture/0001-execution-lifecycle-contract.md}, nel seguito ADR~0001)
che il capitolo cita per nome.

\section{Registro e risoluzione della specifica}

\subsection{\code{FunctionRegistry}, \code{FunctionService} e il catalogo persistente}

Il registro non \`e pi\`u una mappa mutabile: \code{FunctionRegistry} conserva uno
\emph{snapshot} immutabile, sostituito per intero a ogni scrittura, con due viste
--- quella pubblica, che le API leggono, e quella di recupero, che include anche
le funzioni indisponibili (in rimozione parziale o non ancora ripristinate).
La lettura non prende lock; le scritture sono \code{synchronized} e pubblicano
lo snapshot solo dopo averlo reso durevole.

La durabilit\`a \`e la novit\`a pi\`u rilevante rispetto alla prima versione del
progetto. \code{FunctionCatalog} scrive l'intero catalogo in un file
(\code{nanofaas.registry.path}): file temporaneo nella stessa directory, permessi
\code{0600} sul file e \code{0700} sulla directory (in modo best-effort, perch\'e un
volume montato da root non lo consente), \code{fsync} del file, spostamento atomico
e \code{fsync} della directory. Il commento nel sorgente nomina il costo senza
nasconderlo: ogni mutazione riscrive e sincronizza tutto il catalogo, serializzata
sul monitor del registro, e la scrittura differita con flush coalescente \`e la
via di miglioramento se la frequenza diventa un collo di bottiglia. Un catalogo
illeggibile o con schema sconosciuto fa fallire l'avvio, anzich\'e ripartire vuoto.
Nel chart Helm il percorso sta su un volume persistente (\code{/var/lib/nanofaas}).

Sopra il registro, \code{FunctionService} orchestra le operazioni che la
registrazione comporta oltre alla memorizzazione, tutte sotto un lock per nome
(\code{FunctionOperationLocks}) e tutte rifiutate finch\'e il catalogo non \`e
stato ripristinato (\code{FunctionRestoreGate}). L'ordine della registrazione
\`e esplicito:

\begin{enumerate}
  \item risoluzione dei default (\cref{sec:06-resolver});
  \item rifiuto se il nome \`e in rimozione parziale o gi\`a presente;
  \item validazione dell'immagine attraverso il validatore del backend attivo;
  \item provisioning del deployment gestito, se la modalit\`a lo prevede;
  \item notifica ai listener di registrazione (capacit\`a, code, metriche);
  \item commit durevole nel registro --- l'ultimo passo.
\end{enumerate}

Se un passo fallisce, i listener gi\`a notificati vengono ritirati e le risorse
gi\`a create vengono rilasciate: la funzione non compare a met\`a. La cancellazione
percorre l'ordine inverso e mette il commit durevole \emph{per ultimo}.

La registrazione non \`e idempotente: registrare un nome gi\`a presente
restituisce \code{409 Conflict}. La modifica avviene invece per
\code{PATCH}, con una regola dichiarata nel tipo della richiesta:

\begin{lstlisting}[language=Java, caption={\code{FunctionUpdateRequest}: il
commento nel sorgente enuncia la regola.}]
/**
 * Partial update of a registered function. Only fields that the control plane can apply without
 * touching the deployment are accepted; a null field leaves the current value alone.
 *
 * <p>Unknown properties are rejected rather than ignored, so patching an immutable field such as
 * {@code image} or {@code executionMode} fails loudly instead of appearing to succeed.</p>
 */
@JsonIgnoreProperties(ignoreUnknown = false)
public record FunctionUpdateRequest(
        @Min(1) Integer concurrency, @Min(1) Integer timeoutMs,
        @Min(0) Integer maxRetries, ConcurrencyControlConfig concurrencyControl
) { ... }
\end{lstlisting}

La superficie modificabile a caldo \`e quindi esattamente quattro campi, e sono
i quattro che governano il comportamento del control plane senza toccare il
deployment sottostante. Il rifiuto delle propriet\`a sconosciute \`e la parte
importante: senza di esso, un \code{PATCH} che tentasse di cambiare
\code{image} riceverebbe \code{200} e non cambierebbe nulla, che \`e il modo
peggiore di fallire. Per una funzione gestita, l'aggiornamento raggiunge anche il
provider: il proxy del backend container deriva dalla specifica il timeout del
singolo salto e il limite di ammissione, e senza questo passo il deployment
continuerebbe a imporre i valori originali mentre il chiamante crede in quelli
modificati. Un errore di un listener dopo l'aggiornamento \emph{non} produce un
rollback (il commento lo giustifica: il percorso di rollback della registrazione
cancella la coda della funzione, che sarebbe peggio di un listener che manca un
aggiornamento; i listener sono idempotenti, quindi ripetere lo stesso \code{PATCH}
converge).

Una nota sulla semantica del blocco annidato: \code{concurrencyControl} viene
sostituito \emph{in blocco}, non fuso campo per campo. La ragione \`e che i campi
di quel blocco hanno significati diversi a seconda del modo
(\cref{cap:12-concurrency-control}), e una fusione parziale produrrebbe
configurazioni ibride prive di senso --- per esempio una soglia di carico
ereditata da \code{ADAPTIVE\_PER\_POD} sopravvissuta al passaggio a
\code{BUDGETED}.

\paragraph{Rimozione parziale.} Una cancellazione pu\`o finire in tre stati, e solo
due possono essere presentati come esito affidabile. La rimozione completata
cancella la voce dal catalogo per ultima. Se il backend segnala una
\code{PartialDeprovisionException} --- risorse che non ha potuto liberare --- la
funzione entra in \emph{rimozione in sospeso}: la voce resta nel catalogo (\`e ci\`o
che tiene tracciabili i residui e fa s\`i che un \code{DELETE} ripetuto riprenda la
pulizia invece di rispondere \code{404}), i listener \emph{non} vengono
riattivati, e ogni via d'ingresso alla funzione --- invocazione, enqueue, patch,
repliche, ri-registrazione --- risponde \code{409} con codice
\code{FUNCTION\_REMOVAL\_PENDING}. Il rollback operativo \`e ammesso solo quando
nulla di necessario \`e andato perso. La rimozione in sospeso \`e uno stato di
processo, non persistito: dopo un riavvio la funzione viene ripristinata come
qualunque altra, e se la pulizia \`e ancora voluta il \code{DELETE} va ripetuto
(ADR~0001, \S 8.3). \`E una conseguenza del vincolo di un singolo pod con stato in
memoria, non una svista.

\subsection{\code{FunctionSpecResolver}: dove i default diventano valori}
\label{sec:06-resolver}

La risoluzione trasforma una specifica parzialmente compilata in una completa.
\`E il punto in cui il modello ``\code{null} significa non specificato''
(\cref{cap:05-contratti-dati}) viene consumato:

\begin{lstlisting}[language=Java]
Optional.ofNullable(spec.timeoutMs()).orElse(defaults.timeoutMs()),
Optional.ofNullable(spec.concurrency()).orElse(defaults.concurrency()),
Optional.ofNullable(spec.queueSize()).orElse(defaults.queueSize()),
Optional.ofNullable(spec.maxRetries()).orElse(defaults.maxRetries()),
\end{lstlisting}

I default delle politiche di concorrenza non vivono pi\`u nel resolver: il resolver
delega a \code{ConcurrencyControlConfig.normalize}, nel progetto dei contratti
condivisi, cos\`i che la stessa normalizzazione valga ovunque il blocco venga letto.
Due default sono commentati in modo che vale la pena riportare perch\'e esprimono
principi, non valori.

\paragraph{Le soglie del controllore adattivo.}
\begin{lstlisting}[language=Java]
// Fractions of latency degradation over the function's best observed service time:
// back off above 2x, grow again below ~1.18x. See the concurrency-control module.
private static final double DEFAULT_HIGH_LOAD_THRESHOLD = 0.5;
private static final double DEFAULT_LOW_LOAD_THRESHOLD = 0.15;
\end{lstlisting}

Le soglie sono espresse come frazioni di degradazione rispetto al \emph{miglior}
tempo di servizio osservato per quella funzione, non come valori assoluti. \`E
la stessa costruzione di TCP Vegas~\cite{brakmo1995vegas}: il riferimento \`e il
minimo storico, assunto come stima del costo di servizio irriducibile, e il
segnale \`e lo scostamento da esso. Il \cref{cap:12-concurrency-control} discute
perch\'e questa costruzione, corretta nel contesto originale, si riveli
insufficiente qui.

\paragraph{L'SLO di default.}
\begin{lstlisting}[language=Java]
// A generic service-time SLO for a function that did not state one. Deliberately not
// derived from anything the platform measures: a target the controller inferred from
// observed latency would move whenever the function got slower, which is the one thing
// an SLO must not do.
private static final long DEFAULT_TARGET_LATENCY_MS = 250L;
\end{lstlisting}

L'osservazione \`e sottile e generale: un obiettivo derivato dall'osservazione
non \`e un obiettivo. Se il sistema stimasse l'SLO dalla latenza misurata, un
degrado prestazionale sposterebbe automaticamente l'asticella, e il controllore
si dichiarerebbe conforme mentre peggiora. Il default \`e quindi un numero
arbitrario e dichiarato tale.

\paragraph{Default dipendenti dalla modalit\`a.} La risoluzione del blocco di
scaling avviene solo per le funzioni in modalit\`a \code{DEPLOYMENT}: per
\code{LOCAL} e \code{EXTERNAL} il blocco resta come fornito, perch\'e il control
plane non gestisce repliche che non ha creato. Per \code{DEPLOYMENT} senza blocco
esplicito il default \`e strategia \code{INTERNAL}, da una a dieci repliche,
metrica \code{queue\_depth} con obiettivo 5. Sono validati anche i limiti
(\code{minReplicas} non negativo, \code{maxReplicas} almeno 1, minimo non
superiore al massimo).

Le metriche ammesse per la strategia interna sono validate contro un insieme
chiuso (\code{queue\_depth}, \code{in\_flight}, \code{rps}): una metrica
sconosciuta fa fallire la registrazione con \code{400} anzich\'e produrre un
autoscaler che non scala e non lo dice.

\section{Capacit\`a di ammissione}
\label{sec:06-capacita}

Prima di creare un'esecuzione il nucleo riserva le risorse che essa
tratterr\`a. Le riserve sono finite per configurazione, e i valori di default
sono calibrati sul contenitore minimo supportato dal chart (512~MiB):
\code{nanofaas.invocation-capacity.*}.

\begin{center}
\small
\begin{adjustbox}{max width=\textwidth}
\begin{tabular}{@{}lll@{}}
\toprule
\textbf{Risorsa} & \textbf{Globale / per funzione} & \textbf{Che cosa conta} \\
\midrule
esecuzioni & 4\,096 / 512 & proprietari logici di un'esecuzione viva \\
input canonico & 128~MiB / 32~MiB & rappresentazione Java trattenuta \\
copie fisiche & 64~MiB / 16~MiB & materializzazioni di runtime e trasporto \\
attese (waiter) & 8\,192 / 1\,024 & chiamanti sincroni attaccati, repliche incluse \\
corpo in ingresso & 1~MiB (unico) & byte ricevuti prima del parsing \\
\bottomrule
\end{tabular}
\end{adjustbox}
\end{center}

L'ordine sul percorso di una invocazione HTTP \`e fissato: confine del corpo
(\code{413}), limitatore di frequenza, decodifica JSON e validazione, risoluzione
della funzione, e solo poi la ricerca della chiave di idempotenza. Una \code{Content-Length}
dichiarata oltre il limite \`e rifiutata senza nemmeno sottoscrivere il corpo; per
un corpo a chunk, il primo buffer che supera il confine viene rilasciato e la
sottoscrizione cancellata. Il vincolo che ne segue \`e importante: una
\emph{replica} non pu\`o aggirare il limite del corpo, perch\'e il controllo
precede la ricerca. Una richiesta genuinamente nuova reclama la chiave,
canonicalizza l'input, riserva i proprietari e pubblica il record vivo \emph{prima}
che un waiter sincrono sia ammesso; se l'ammissione del waiter fallisce,
record, chiave e proprietari sono annullati prima del \code{429}.

Una saturazione di esecuzioni, input, copie o waiter risponde \code{429} con
\code{Retry-After: 1} e un corpo stabile,
\code{\{"error":"invocation\_quota\_exceeded","resource":"execution|input|input\_copy|waiter"\}}.
Un \code{413} non \`e sovraccarico. Le stime coprono gli oggetti Java trattenuti dal
control plane, non la memoria nativa dei buffer Netty n\'e quella remota della
funzione: sono una busta di ammissione, non una garanzia sull'RSS del processo.
Quando il modulo \code{runtime-config} \`e attivo, le quattro coppie
globale/per funzione sono modificabili a caldo (\cref{cap:14-runtime-config}); abbassare
un tetto sotto l'occupazione attuale non espelle nulla, e nuove riserve sono
rifiutate finch\'e il drenaggio non riporta l'uso entro il nuovo limite.

\section{Il ciclo di vita di un'esecuzione}

\subsection{\code{ExecutionRecord}}

\code{ExecutionRecord} \`e l'unica struttura mutabile significativa del nucleo.
Contiene lo stato, i timestamp, l'esito, i flag di cold start, la busta di
risposta, i proprietari delle risorse d'input e --- centralmente --- il
\code{CompletableFuture} su cui i chiamanti sincroni attendono.

La mutabilit\`a \`e governata da invarianti dichiarati nel codice.

Il primo \`e che ogni transizione avviene sotto il monitor dell'oggetto e che gli
stati terminali sono definitivi:

\begin{lstlisting}[language=Java]
/**
 * Terminal states (SUCCESS, ERROR, TIMEOUT) are final; every transition between
 * non-terminal states (including RUNNING -> QUEUED for retries) is allowed.
 */
\end{lstlisting}

La permissivit\`a sulle transizioni non terminali \`e ci\`o che rende possibile
la riprova: un'esecuzione fallita torna in coda e ripercorre il ciclo. La
rigidit\`a su quelle terminali \`e ci\`o che rende sicuro il completamento
concorrente: se il dispatcher e la callback arrivano entrambi --- il che accade
quando la funzione risponde \emph{e} invia la callback --- il secondo viene
scartato invece di sovrascrivere l'esito.

Il secondo invariante \`e che il futuro non viene mai completato tenendo il lock.
Oggi il record \emph{sceglie} la risposta terminale sotto il proprio monitor e la
\emph{pubblica} altrove, con un metodo dedicato:

\begin{lstlisting}[language=Java]
/** Publish only the selected terminal answer; callbacks run outside the monitor. */
public void publishTerminal() {
    ... // sotto il monitor: sceglie risultato o eccezione
    if (failure != null) completion.completeExceptionally(failure);
    else completion.complete(result);
}
\end{lstlisting}

Completare un \code{CompletableFuture} esegue in linea le continuazioni dei
sottoscrittori. Farlo sotto il monitor del record significherebbe eseguire codice
applicativo arbitrario --- la mappatura della risposta, la scrittura sulla socket
--- mentre si tiene un lock su cui il completamento concorrente contende.
Separare la scelta dalla pubblicazione ha un secondo effetto: due completamenti
concorrenti scelgono la stessa risposta, perch\'e la scelta \`e una sola.

Due valori sono catturati \emph{una volta} alla costruzione e non riletti dal
task corrente: se l'esecuzione sar\`a leggibile dopo la fine
(\code{readableAfterFinishing}: vero per le asincrone e per quelle con chiave, falso
per una sincrona semplice, la cui risposta \`e gi\`a tornata sulla connessione che il
chiamante teneva) e la chiave di idempotenza. La ragione \`e la stessa per
entrambi: la riprova sostituisce il task con uno privo di chiave, e rileggere il
task dopo di essa declasserebbe proprio le esecuzioni che hanno avuto
problemi, facendo smettere silenziosamente di funzionare l'idempotenza per esse.

\subsection{Il proprietario della transizione terminale}

Fino a un certo punto la transizione terminale era divisa fra pi\`u collaboratori:
lo store archiviava l'esito, e lo spostamento della chiave di idempotenza verso il
suo stato terminale era un listener in una catena best-effort. La conseguenza, nel
sorgente, \`e chiamata per nome (\emph{review finding R2}): un esito troppo grande,
espulso fra i due passi, lasciava una chiave ancora \code{published} che una
replica poteva rivendicare, autorizzando una seconda esecuzione di una funzione
gi\`a eseguita.

\code{ExecutionLifecycle.settle} \`e oggi l'unico proprietario della transizione, e
i passi hanno un ordine che nessun listener pu\`o saltare:

\begin{enumerate}
  \item \emph{protezione della chiave}: la chiave diventa non rivendicabile
    (\code{markTerminal}) prima che l'ultimo riferimento vivo possa sparire, sotto
    il monitor del record, lo stesso che il percorso di ammissione usa per
    pubblicare, cos\`i che un record terminale non possa mai essere pubblicato sotto
    una chiave rivendicabile;
  \item \emph{pubblicazione} della risposta terminale scelta dal record;
  \item \emph{archiviazione} dell'esito, se il suo peso rientra nel budget, e
    rimozione del record vivo;
  \item \emph{notifica} degli osservatori best-effort (metriche, conclusione
    end-to-end): un listener che lancia un'eccezione non pu\`o interrompere la
    pulizia n\'e impedire a un altro proprietario di chiudere.
\end{enumerate}

Il metodo \`e idempotente (un record non terminale \`e ignorato, uno gi\`a assestato
non rifa nulla). Non c'\`e un lock globale sul percorso caldo: il record \`e il proprio
monitor, e nessun listener, dispatcher o codice esterno gira mentre lo si tiene.
\code{ExecutionStore.settle} resta come adattatore per i chiamanti dei moduli e
fallisce esplicitamente se un record \emph{con chiave} viene assestato senza un
proprietario collegato, anzich\'e perdere in silenzio la garanzia di deduplicazione.

\subsection{Il contratto del ciclo di vita}

Le regole che attraversano pi\`u componenti sono state raccolte nell'ADR~0001, con
dieci invarianti (I1--I10). Quattro meritano di essere riportate perch\'e
cambiano cosa un lettore pu\`o assumere.

\paragraph{Il timeout di un chiamante non \`e il timeout dell'esecuzione.}
Un waiter che scade (\code{X-Timeout-Ms}, altrimenti \code{spec.timeoutMs}) ottiene
la propria risposta \code{408}, ma il record condiviso, la chiave, lo store, il
lease e i contatori non cambiano: l'esecuzione prosegue, e una replica successiva
o un polling ne osserva l'esito reale. Il codice lo rende esplicito
sottoscrivendo il futuro condiviso con \code{suppressCancel}, cos\`i che n\'e il
timeout n\'e la disconnessione di un singolo chiamante possano cancellare il lavoro
condiviso. Il \code{TIMEOUT} condiviso \`e riservato alle scadenze di livello
esecuzione, come l'attesa massima in coda (\cref{cap:10-sync-queue}).

\paragraph{Quattro orologi.}
\begin{center}
\small
\begin{adjustbox}{max width=\textwidth}
\begin{tabular}{@{}lll@{}}
\toprule
\textbf{Orologio} & \textbf{Chi lo possiede} & \textbf{Effetto alla scadenza} \\
\midrule
budget del waiter & il chiamante & solo quel waiter riceve \code{408} \\
scadenza del tentativo & il dispatch (\code{timeoutMs}) & segue la politica di riprova \\
politica di riprova & l'esecuzione & nuovo tentativo o errore terminale \\
\code{maxLifetime} & lo store & \code{EXECUTION\_EXPIRED} fabbricato \\
\bottomrule
\end{tabular}
\end{adjustbox}
\end{center}

\paragraph{Tentativo, lease e generazione.} Un tentativo \`e la singola
spedizione di un'esecuzione, e possiede un \code{DispatchLease} sotto la
\emph{generazione} corrente della funzione. Il rilascio del lease \`e idempotente
ed \`e guidato dal futuro \emph{grezzo} del trasporto, non dallo specchio della
scadenza: un handler LOCAL non interrompibile che continua a girare dopo la
scadenza continua a occupare la capacit\`a che gli \`e attribuita, e questa non
viene ceduta a un secondo tentativo. Il futuro terminale e la fine fisica del
lavoro sono due momenti diversi e solo il primo \`e una transizione di stato.
Una funzione rimossa e ri-registrata con lo stesso nome ha una nuova
generazione: un lease vecchio rilascia solo lo slot della vecchia, e una
callback tardiva non pu\`o ricreare stato n\'e toccare le metriche della nuova
(\code{ACTIVE}$\to$\code{RETIRING}$\to$\code{CLOSED}; ADR~0001 \S 8.2).

\paragraph{La cancellazione richiesta prima che esista l'handle.} Una scadenza
amministrativa o uno shutdown possono correre contro un dispatch che non ha
ancora restituito il proprio futuro. La richiesta viene ricordata sul record e
applicata quando l'handle arriva; non \`e mai scartata per essere arrivata presto.

\subsection{\code{ExecutionStore}: vivi e postumi}

La prima versione descritta in questo capitolo era una mappa concorrente con un
thread janitor che ogni minuto eliminava le voci scadute, con tre orizzonti
(\code{ttl}, \code{cleanupTtl}, \code{maxLifetime}). Il sorgente attuale racconta
perch\'e \`e stata sostituita, con due misure nel commento della classe: in una
esecuzione a 1\,093 ammissioni al secondo, \code{function\_inFlight} ha toccato un
massimo di \textbf{2} mentre lo store ne conteneva \textbf{165\,786}: il 99,999\%
di ci\`o che conservava era postumo, e si portava dietro l'apparato dei vivi (il
futuro, il task, la richiesta, l'insieme dei tentativi rilasciati): 35 oggetti
per record, cinque milioni e ottocentomila in tutto, ciascuno da percorrere a
ogni raccolta completa. Il janitor impiegava 12,4~ms su 166\,000 record. E la
ritenzione era dichiarata in tempo e illimitata in numero, quindi la memoria
cresceva con il tasso d'arrivo: il commento riporta 1,05~GB, il 50,6\% del tempo
in GC e una probe di liveness mancata tre volte di seguito.

Le due vite sono oggi due strutture, entrambe cache Caffeine.

\begin{center}
\small
\begin{adjustbox}{max width=\textwidth}
\begin{tabular}{@{}lll@{}}
\toprule
\textbf{Struttura} & \textbf{Contiene} & \textbf{Scadenza e limite} \\
\midrule
\code{inFlight} & \code{ExecutionRecord} mutabili &
  \code{maxLifetime} (30~min) dalla scrittura \\
\code{outcomes} & \code{Outcome} piatto e immutabile &
  \code{ttl} (5~min) se leggibile, \code{sync-ttl} (30~s) altrimenti; \\
 & & tetto in \emph{byte} (\code{max-outcome-bytes}) \\
\bottomrule
\end{tabular}
\end{adjustbox}
\end{center}

Un \code{Outcome} \`e quattro oggetti e 116 byte (contro 35 oggetti e 1\,330 byte di
un record terminale, misurati su 200\,000 record il 2026-08-26 secondo il commento),
ed \`e esattamente ci\`o che serve a due scopi: rispondere a
\code{GET /v1/executions/\{id\}} e servire la replica di una chiave. I tempi sono
\code{long} e non \code{Instant} per la stessa ragione di peso.

Il limite \`e in byte e non in numero, perch\'e un esito leggibile trattiene il
payload del chiamante: 20\,000 esiti da 64~KB occupano 1,28~GB (misura citata nel
sorgente da \code{docs/experiments/control-plane-tuning-2026-09/RESULTS.md}).
Il peso \`e una \emph{stima} conservativa calcolata una volta all'inserimento da
\code{OutcomeWeigher}, che congela in profondit\`a il payload in una copia
immutabile e non condivisa (cos\`i che nulla possa mutarlo, e farne crescere la
dimensione, dopo che \`e stato pesato) e non lo riserializza mai. La visita \`e
limitata in profondit\`a e ampiezza, ma \emph{mai} valuta zero la parte saltata:
se la dimensione non pu\`o essere limitata --- struttura troppo profonda,
contenitore troppo largo, ciclo, riferimento condiviso, oggetto opaco --- l'esito
non viene ritenuto. Anche un esito che da solo supera l'intero budget \`e
rifiutato, anzich\'e inserito per essere espulso un istante dopo. Il default di
\code{max-outcome-bytes} \`e derivato: \code{max-outcomes} $\times$ 116 byte, cio\`e
il budget \`e il vincolo che davvero limita lo heap.

Il rifiuto di ritenere il payload non \`e la perdita della deduplicazione: il
tombstone \`e la chiave, non il payload, ed \`e il proprietario della transizione
a proteggerlo prima di tutto. Una replica di una chiave il cui esito \`e stato
espulso per capacit\`a riceve \code{410 Gone} (\code{OutcomeGoneException}); la
funzione \emph{non} viene rieseguita.

La scadenza amministrativa (\code{maxLifetime}) \`e il secondo orizzonte che
resta, e la sua condizione \`e dichiarata: un'esecuzione la cui callback non
arriva mai non raggiunge uno stato terminale, quindi gli orizzonti ancorati
all'istante di completamento non la toccherebbero mai. La cache dei vivi ha uno
scheduler proprio, con la motivazione riportata nel commento: senza di esso
Caffeine controlla la scadenza solo quando qualcosa tocca la cache, e un record
bloccato su un dispatch perso, che nessuno rilegge pi\`u, restava scaduto ma vivo
a tempo indeterminato --- e con esso lo slot di concorrenza che teneva.
L'espulsione notifica i listener di scadenza amministrativa; il coordinatore dei
tentativi fabbrica un esito \code{EXECUTION\_EXPIRED}, marcato come tale per
distinguerlo da un errore genuino del runtime, e lo assesta.

Due metriche misurano la distanza fra ci\`o che la piattaforma ricorda e ci\`o che
sta eseguendo: \code{execution\_store\_size} (esiti archiviati) e
\code{execution\_in\_flight\_records}. Senza la seconda, la prima potrebbe essere
scambiata per lavoro in corso.

\section{Idempotenza}

L'idempotenza \`e attivata dall'header \code{Idempotency-Key} ed \`e implementata
da \code{IdempotencyStore}, una cache Caffeine. La chiave \`e composta con il nome
della funzione: due funzioni con lo stesso valore di chiave sono associazioni
indipendenti. Il problema che risolve non \`e la deduplicazione a posteriori ma la
\emph{corsa}: due richieste con la stessa chiave che arrivano contemporaneamente
devono produrre una sola esecuzione.

Il meccanismo \`e una rivendicazione a fasi, e le fasi sono quattro, ciascuna con
la propria scadenza:

\begin{center}
\small
\begin{adjustbox}{max width=\textwidth}
\begin{tabular}{@{}lll@{}}
\toprule
\textbf{Stato} & \textbf{Significato} & \textbf{Scadenza} \\
\midrule
\emph{pending} & rivendicata da un chiamante che sta creando l'esecuzione &
  \code{maxLifetime} \\
\emph{published} & associata a un'esecuzione ancora viva & \code{maxLifetime} \\
\emph{abandoned} & pubblicata e poi esplicitamente abbandonata: l'esecuzione
  non girer\`a mai, il legame pu\`o essere rivendicato di nuovo & --- \\
\emph{terminal} & l'esecuzione \`e archiviata: \`e il tombstone & \code{ttl} da
  completamento \\
\bottomrule
\end{tabular}
\end{adjustbox}
\end{center}

Solo lo stato \emph{abandoned} \`e rivendicabile: l'assenza di un record vivo o di
un esito non \`e mai la prova che un legame pubblicato sia stato abbandonato.
La scadenza per stato non \`e configurata a parte, ma derivata dalle proprieta'
dello store delle esecuzioni, e la ragione \`e nel commento: una chiave \`e utile
solo finch\'e esiste la risposta che nomina, e una chiave che scade prima del
suo esito esegue silenziosamente due volte la funzione, che \`e esattamente il
guasto che la chiave esiste per evitare. Il passaggio a \emph{terminal} riavvia
l'orologio da completamento.

\begin{lstlisting}[language=Java, caption={Il ciclo di rivendicazione (schema).}]
public AcquireResult acquireOrGet(String functionName, String key) {
    ...
    // riserva prima uno slot di budget, poi tenta putIfAbsent;
    // se perde la corsa, restituisce la riserva
    ... claimed(token) | pending() | existing(id, terminal) | budgetExhausted()
}
\end{lstlisting}

Gli esiti possibili --- \code{CLAIMED}, \code{PENDING}, \code{EXISTING},
\code{BUDGET\_EXHAUSTED}, \code{MISSING} --- corrispondono a reazioni distinte del
chiamante: procedere alla creazione, attendere, riusare l'esecuzione esistente,
rifiutare, ricominciare. La pubblicazione e l'abbandono sono condizionati al
possesso del token.

Il numero di chiavi \`e limitato da \code{max-keys}, e il budget \`e una
\emph{riserva atomica}, non la lettura di una dimensione: una nuova
rivendicazione riserva il proprio slot \emph{prima} di creare il legame, cos\`i che
chiavi distinte non possano osservare tutte spazio disponibile e inserirsi tutte
(\emph{finding R7}). Ogni voce possiede esattamente uno slot, legato alla
sua identit\`a e non al suo stato, e lo restituisce una volta sola, quando la voce
lascia la cache. Le repliche di una chiave gi\`a posseduta non toccano mai il budget,
quindi la saturazione non le blocca, e il tetto non espelle mai una chiave viva
per far posto a una nuova. Un ingresso rifiutato per budget esaurito \`e contato
(\code{idempotency\_key\_budget\_rejections}) e l'occupazione \`e esposta come
\code{idempotency\_keys\_held}, senza etichette che rivelerebbero un identificatore
di esecuzione o una chiave dell'utente.

Il ciclo di rivendicazione pu\`o attendere sotto contesa, ed \`e questo a motivare il
confinamento sullo scheduler elastico descritto nel
\cref{cap:04-architettura-control-plane}: una richiesta \emph{senza} chiave non lo
attraversa e resta sull'event loop.

\section{Dispatch}

\subsection{Il router e il trasporto}

\code{DispatcherRouter} \`e volutamente banale: due metodi che delegano a
\code{LocalDispatcher} o a \code{ExternalDispatcher}. La modalit\`a
\code{DEPLOYMENT} non ha un dispatcher proprio --- una volta che il deployment
esiste, invocarlo \`e la stessa operazione che invocare un endpoint esterno, e
la differenza sta nel \emph{chi} ha creato l'endpoint
(\cref{cap:07-modalita-esecuzione}).

La scelta della modalit\`a non vive pi\`u nel gestore di completamento:
\code{AttemptTransportAdapter}, nel control plane, \`e l'unico posto da cui il
percorso del tentativo chiama il router e la condizione di prontezza. Il
coordinatore dei tentativi, in \code{execution-runtime}, non importa n\'e l'uno n\'e
l'altra, e il commento spiega perch\'e: il cancello del motore di scheduling aveva
gi\`a dovuto essere indebolito una volta per prendere un lock per funzione, e una
chiamata al provider raggiungibile dalle vicinanze del cancello sarebbe una
classe di problema ben peggiore.

\subsection{\code{ExternalDispatcher} e il contratto sul filo}

Il dispatcher esterno effettua una \code{POST} verso l'endpoint della funzione,
allegando quattro header di contesto: identificatore di esecuzione, numero di
tentativo, trace id e chiave di idempotenza. Il timeout \`e applicato in due
punti: come \code{responseTimeout} sulla richiesta Reactor Netty sottostante,
che interrompe effettivamente la connessione anzich\'e limitarsi ad abbandonare
l'attesa, e come timeout del \code{Mono}, che produce l'errore
\code{EXTERNAL\_TIMEOUT}. Il secondo ha un effetto che il commento del
\code{HttpClientConfig} nomina: cancella anche l'acquisizione pendente dal pool.

Nella direzione di ritorno, il dispatcher interpreta tre header prodotti dal
runtime:

\begin{center}
\small
\begin{adjustbox}{max width=\textwidth}
\begin{tabular}{@{}ll@{}}
\toprule
\textbf{Header} & \textbf{Significato} \\
\midrule
\code{X-Cold-Start} & La richiesta ha pagato un'inizializzazione. \\
\code{X-Init-Duration-Ms} & Durata di tale inizializzazione. \\
\code{X-NanoFaaS-Function-Status} & Il codice di stato HTTP \`e deciso
dall'handler. \\
\bottomrule
\end{tabular}
\end{adjustbox}
\end{center}

Il terzo \`e il pi\`u delicato, e il codice lo tratta con diffidenza esplicita:

\begin{lstlisting}[language=Java]
// An out-of-range status on a marker-bearing response is not spec-legal
// ([200,599] only). Fall through to the platform-error path rather than
// trusting it: EXTERNAL/DEPLOYMENT endpoints are arbitrary unauthenticated
// URLs, and HTTP's status-line grammar is 3DIGIT, so a misbehaving upstream
// can emit one.
\end{lstlisting}

L'endpoint di una funzione \`e un URL arbitrario e non autenticato: il control
plane non pu\`o assumere che rispetti il contratto. Un marcatore che dichiara
``il mio codice di stato \`e intenzionale'' accompagnato da un codice fuori
intervallo viene trattato come errore di piattaforma, non propagato.

\subsection{Il client HTTP: pool, DNS e connessioni}

Il client dei dispatcher non usa il pool globale di Reactor Netty: il pool
(\code{DispatchConnectionPool}) appartiene al contesto e ha un bilancio
dichiarato, per \emph{destinazione}: 500 connessioni, coda di acquisizione pari
a due volte tale valore, timeout di acquisizione 45~s, chiusura delle connessioni
inattive dopo 30~s, e smaltimento dei pool di destinazioni inattive. Il commento
riporta una scelta di \emph{non} adottare un timeout di acquisizione pi\`u breve:
misurato come grande vantaggio solo in un harness che non cancellava mai, e
senza effetto misurabile quando, come fa \code{ExternalDispatcher}, si cancella.

Il client risolve il DNS a \emph{ogni} nuova connessione invece di conservarne le
risposte (\code{NoopDnsCache}; non un TTL zero, che Netty rifiuta). La
motivazione, nel commento, \`e un incidente: una funzione cancellata e
ri-registrata con lo stesso nome riceve un nuovo Service con un nuovo ClusterIP;
con la cache del record (30~s su CoreDNS di minikube) i dispatch continuavano a
chiamare il vecchio ClusterIP, dove nulla rispondeva, e ogni connessione restava
appesa fino al timeout di connessione pur essendo il nuovo pod Ready. Poich\'e le
connessioni sono in pool, la risoluzione avviene solo all'apertura di una nuova.

\section{Riprova}
\label{sec:06-riprova}

La riprova \`e gestita da \code{AttemptCoordinator}, in \code{execution-runtime};
\code{ExecutionCompletionHandler} ne \`e oggi una facciata Spring. La regola di
base non \`e cambiata --- si riprova finch\'e il numero di tentativo non supera
\code{maxRetries} --- ma i dettagli s\`i.

\paragraph{La riprova ha un ritardo.} Nella prima versione una riprova era
immediata, e il suo costo era misurabile: dopo un deploy, gli errori
\code{EXTERNAL\_ERROR} erano riprove ravvicinate contro un endpoint che non era ancora
pronto, che si esaurivano in pochi millisecondi. Oggi \code{RetryBackoff} calcola
l'istante di eleggibilit\`a del prossimo tentativo: la base cresce come
$\min(\text{max},\,\text{iniziale}\cdot 2^{n-1})$ e il ritardo effettivo \`e
campionato con \emph{jitter uguale} in $[\text{base}/2,\,\text{base}]$; il minimo
non \`e mai zero. Un \code{Retry-After} del server a monte (accettato in secondi o
come data, solo per \code{429} e \code{503}, e solo se coerente fra pi\`u valori) \`e un
\emph{minimo}, non un tetto: se \`e pi\`u lontano del ritardo locale, vince.
I default sono \code{nanofaas.retry.initial-backoff}=100~ms e
\code{max-backoff}=2~s.

Una riprova in attesa \emph{conserva la propria riserva} e non occupa capacit\`a di
dispatch: nel profilo diretto il lavoro in attesa \`e tenuto da
\code{ExecutorBackedInvocationEnqueuer}, con un timer e un numero massimo di lavori
pendenti; con un modulo di coda, il biglietto rientra nel motore con \code{notBefore}
e attende nell'insieme dei \emph{ritardati} (\cref{cap:09-async-queue}).

\paragraph{Misure di riprova.} Il confronto \`e in
\code{docs/experiments/scheduler-switching-2026-09} (base \code{a17289c2}). Nella
riproduzione dal vivo di un deploy, la base ha completato 742 esecuzioni su 900 e il
candidato 900 su 900. Con un \code{429} controllato e \code{Retry-After: 1}, la base
ha tentato 1\,195 volte contro 600, ha concluso 3 esecuzioni con successo e 297 con
errore contro 300 e 0, con un intervallo minimo di 0,000727~s contro 1,000698~s e 892
intervalli sotto il secondo contro nessuno. Va detto che un caso su dodici del
deploy vivo resta senza spiegazione: un timeout che i log non chiariscono.

\paragraph{Il cancello di prontezza per DEPLOYMENT.} Un Deployment appena creato
non \`e pronto per i secondi che il primo pod impiega ad avviarsi, e un dispatch
inviato allora incontra un Service senza endpoint. Invece di spendere riprove,
\code{DeploymentWakeUpGate} trattiene il dispatch finch\'e la generazione non \`e stata
osservata pronta, una volta sola: la funzione gi\`a vista pronta costa una ricerca
e nessuna chiamata al provider. L'attesa \`e cancellabile, ha un timeout
(\code{nanofaas.deployment.wakeup.timeout}, 30~s), un intervallo di polling
(250~ms) e un'et\`a massima dell'osservazione di pronto (5~s); la stessa attesa
serve la scale-to-zero.

\paragraph{Tre dettagli che restano.}
Primo, il tentativo di riprova viene costruito come un \code{InvocationTask}
nuovo, con la chiave di idempotenza \emph{azzerata}:

\begin{lstlisting}[language=Java]
InvocationTask retryTask = new InvocationTask(
        ..., null,  // No idempotency key for retry - retry is internal
        ..., currentTask.attempt() + 1, currentTask.kind());
\end{lstlisting}

La riprova \`e un'operazione interna sulla \emph{stessa} esecuzione: conservare
la chiave farebbe s\`i che la riprova si riconosca come duplicato di se stessa e
riproduca l'esito fallito. (Per questo il record ha catturato la chiave alla
costruzione.)

Secondo, il rilascio dello slot di concorrenza \`e idempotente per tentativo. Nella
prima versione era un insieme di tentativi rilasciati nel record; oggi \`e
\code{DispatchLease.release()}, che rilascia una volta sola. Senza questa
protezione, un esito che arriva due volte --- risposta HTTP \emph{e} callback ---
rilascerebbe due slot avendone occupato uno, e il contatore delle richieste in
volo scenderebbe sotto il valore reale, con effetti diretti sull'autoscaler e sul
controllore di concorrenza che leggono quel contatore.

Terzo, la riprova pu\`o fallire perch\'e non c'\`e modo di ripubblicarla (coda piena,
lavori pendenti al massimo, timer rifiutato), e in tal caso l'esecuzione viene
completata con l'errore \emph{originale} anzich\'e con un errore di coda: il
chiamante riceve la ragione per cui la funzione \`e fallita, non la ragione per cui
la piattaforma ha rinunciato a riprovare. Il codice riverifica che l'esecuzione sia
ancora parcheggiata su quel tentativo prima di concluderla, perch\'e una scadenza
amministrativa pu\`o averla conclusa nel frattempo.

\section{Limitazione di frequenza}

Il limitatore \`e ancora una finestra fissa di un secondo, ma ha cambiato sia
posizione sia realizzazione. La posizione: non \`e pi\`u chiamato da
\code{InvocationService}, dopo la decodifica HTTP e JSON, ma da un \code{WebFilter}
(\code{RateLimitWebFilter}) ristretto ai due suffissi \code{:invoke} e \code{:enqueue}, che
rifiuta \emph{prima} di leggere il corpo: nel commento, ogni rifiuto pagava prima
tutto il costo di decodifica, proprio quando la piattaforma ha meno da spendere.
Il rifiuto svuota comunque il corpo, altrimenti Reactor Netty non riceve il segnale
di richiesta consumata e pu\`o chiudere la connessione invece di riusarla in
keep-alive, un costo maggiore di ci\`o che il rifiuto risparmia; risponde \code{429}
con \code{Retry-After: 1}. La realizzazione: finestra e conteggio stanno in un \emph{unico} \code{AtomicLong}
(33 bit per il secondo epoch, 31 per il conteggio), cos\`i che il passaggio di
finestra e il reset del contatore siano una sola transizione CAS. La versione
precedente, con due contatori atomici separati, aveva una corsa che il commento
descrive: rifiutava a torto richieste della nuova finestra e permetteva a un reset
di cancellare incrementi fatti dopo il passaggio.

\begin{lstlisting}[language=Java]
public boolean allow() {
    long max = maxPerSecond;
    long now = epochSecondSource.getAsLong();     // letto una volta per decisione
    while (true) {
        long state = windowAndCount.get();
        long window = state >>> COUNT_BITS, count = state & COUNT_MASK;
        if (now > window) {                       // nuova finestra: la reclama a 1
            if (windowAndCount.compareAndSet(state, encode(now, 1L))) return 1L <= max;
            continue;
        }
        if (count >= max) return false;
        if (windowAndCount.compareAndSet(state, state + 1L)) return true;
    }
}
\end{lstlisting}

Il difetto noto della finestra fissa non \`e cambiato, e il commento della classe
lo dichiara: a cavallo del secondo si pu\`o ammettere fino al doppio del limite, e
non \`e un token bucket. La scelta \`e coerente con il ruolo del limitatore in
questo sistema --- una protezione presente sul percorso ma neutra per default
(\cref{cap:03-requisiti-principi}) --- e la sua imprecisione \`e irrilevante a un
milione di richieste al secondo. Il limite \`e \code{volatile} e riletto a ogni
decisione: l'endpoint di configurazione a caldo lo modifica con effetto alla
decisione successiva, senza sincronizzazione aggiuntiva.

\section{Metriche}

\code{Metrics} incapsula la registrazione Micrometer. La struttura ha una
propriet\`a che sorprende alla prima lettura: i contatori per funzione sono
creati alla registrazione della funzione e \emph{rimossi} alla sua
cancellazione, e ogni metodo di incremento verifica prima che i contatori
esistano ancora.

\begin{lstlisting}[language=Java]
public void success(String function) {
    FunctionMeters metersFor = metersOrNull(function);
    if (metersFor != null) { metersFor.success().increment(); }
}
\end{lstlisting}

Il percorso comune --- una funzione registrata e viva --- non prende lock. Lo
prendeva sempre, ed era un monitor \emph{globale} condiviso da tutte le
funzioni, su un percorso che ogni invocazione attraversa sei volte: il
commento riporta 223~ns per operazione con un thread e 2\,638~ns con otto (il
costo per operazione cresceva con i thread invece di restare piatto, che \`e la
firma di una serializzazione). Il lock resta sul percorso lento, dove serve:
prima registrazione e corsa con la rimozione. Un incremento che afferra i
contatori un istante prima della rimozione va perso, ed \`e un prezzo accettato.

La rimozione dei contatori alla cancellazione evita che il registro Prometheus
cresca indefinitamente in un sistema in cui le funzioni sono create e distrutte da
script di test. Ha per\`o una conseguenza sperimentale che il progetto ha
imparato per via empirica: una verifica che interroghi le metriche \emph{dopo}
aver cancellato la funzione non trova nulla. Il \cref{cap:22-metodo-sperimentale}
riporta come gli scenari di validazione sono stati corretti di conseguenza.
Le metriche di carico per funzione dei moduli (\code{function\_queue\_depth} e
compagne) seguono la stessa regola, ma con \emph{drenaggio}: i contatori vengono
rimossi solo dopo che il lavoro in volo della generazione \`e tornato
(\cref{cap:09-async-queue}).

Le esecuzioni relative a funzioni gi\`a rimosse vengono contabilizzate su un
registro separato e non esportato (\code{removed\_function\_*}), cos\`i che la
loro esistenza non produca n\'e serie orfane n\'e eccezioni. Il conteggio dei
campioni \`e legato alla generazione: un completamento la cui funzione \`e stata
rimossa e ri-registrata dopo l'ammissione non attribuisce il proprio campione ai
contatori della nuova registrazione (\code{isCurrentGeneration}, I7).
Tre contatori per percorso (\code{admitted}, \code{refused}, \code{replayed})
distinguono la porta da cui l'invocazione \`e entrata, e sono nomi \emph{nuovi} e non
un'etichetta sui contatori esistenti: cinque di quelli alimentano anelli di
controllo (la latenza il governatore di concorrenza, i dispatch l'autoscaler, in
volo e profondit\`a l'HPA di Kubernetes), e dividere una serie che un anello legge
cambia ci\`o che l'anello vede.

\chapter{Modalit\`a di esecuzione}
\label{cap:07-modalita-esecuzione}

Una funzione registrata su \nf{} dichiara \emph{dove} vive attraverso il campo
\code{executionMode}, che ammette tre valori. La distinzione non riguarda
il protocollo di invocazione --- che \`e HTTP in tutti i casi in cui esiste un
processo separato --- ma la responsabilit\`a sul ciclo di vita del processo che
esegue il codice.

\begin{figure}[htbp]
  \centering
  \resizebox{0.95\textwidth}{!}{% nanofaas execution modes (LOCAL / POOL / DEPLOYMENT) and the cold-start vs
% warm-start latency contrast that motivates warm pools.
\documentclass[border=8pt]{standalone}
\usepackage{nanofaas-figs}
\begin{document}
\begin{tikzpicture}[nf]

  % ============ (a) dispatch targets per mode ============
  \node[core, minimum width=30mm] (disp) {PoolDispatcher};

  \node[neutral, below=13mm of disp, xshift=-52mm, minimum width=30mm] (local)
    {\textbf{LOCAL}\\{\scriptsize in-process}};
  \node[neutral, below=13mm of disp, minimum width=30mm] (pool)
    {\textbf{POOL}\\{\scriptsize warm pod pool}};
  \node[neutral, below=13mm of disp, xshift=52mm, minimum width=34mm] (dep)
    {\textbf{DEPLOYMENT}\\{\scriptsize backend provider}};

  \draw[flowc] (disp) -- (local) node[tag,midway,left=0.5mm]{test};
  \draw[flowc] (disp) -- (pool);
  \draw[flowc] (disp) -- (dep);

  % targets
  \node[backend, below=9mm of local, minimum width=30mm] (localt)
    {handler in JVM\\{\scriptsize no container}};
  \node[backend, below=9mm of pool, minimum width=34mm] (poolt)
    {pre-warmed pods\\{\scriptsize OpenWhisk-style}};
  \node[backend, below=9mm of dep, xshift=-11mm, minimum width=26mm] (k8s)
    {Kubernetes\\{\scriptsize Deploy+Svc}};
  \node[backend, below=9mm of dep, xshift=19mm, minimum width=26mm] (clocal)
    {container-local\\{\scriptsize Docker CLI}};

  \draw[flowc] (local) -- (localt);
  \draw[flowc] (pool) -- (poolt);
  \draw[flowc] (dep.south) -- (k8s.north);
  \draw[flowc] (dep.south) -- (clocal.north);

  % ============ (b) cold-start vs warm-start timeline ============
  \newcommand{\phase}[6]{% x1,x2,y,fill,draw,label
    \draw[fill=#4, draw=#5, line width=0.5pt, rounded corners=1pt]
      (#1,#3-0.3) rectangle (#2,#3+0.3);
    \node[font=\sffamily\scriptsize, text=nfink] at ({(#1+#2)/2},#3) {#6};}

  \coordinate (t0) at (-5.2,-7.9);   % origin of the timeline block
  \begin{scope}[shift={(t0)}]
    \node[clabel, anchor=west] at (-0.2,1.95)
      {Cold start vs.\ warm start \textnormal{\footnotesize— first request latency}};

    % axis
    \draw[flowc] (0,-1.5) -- (11.2,-1.5) node[tag, right]{time};

    % cold track (DEPLOYMENT, first hit)
    \node[tag, anchor=east, align=right] at (-0.15,0.35)
      {DEPLOYMENT\\(cold)};
    \phase{0}{1.4}{0.35}{nfslatebg}{nfslate}{schedule}
    \phase{1.4}{5.0}{0.35}{nfamberbg}{nfamber}{pull image}
    \phase{5.0}{6.8}{0.35}{nfgreenbg}{nfgreen}{start ctr}
    \phase{6.8}{9.0}{0.35}{nfpurplebg}{nfpurple}{app ready}
    \phase{9.0}{10.6}{0.35}{nfbluebg}{nfblue}{serve}

    % warm track (POOL / warm pool)
    \node[tag, anchor=east, align=right] at (-0.15,-0.65)
      {POOL /\\warm pool};
    \phase{0}{0.7}{-0.65}{nfgreenbg}{nfgreen}{ready}
    \phase{0.7}{2.3}{-0.65}{nfbluebg}{nfblue}{serve}

    % cold-start penalty bracket
    \draw[decorate, decoration={brace, amplitude=4pt, raise=1pt}, draw=nfred]
      (9.0,0.72) -- (0.0,0.72);
    \node[tag, text=nfred] at (4.5,1.15) {cold-start penalty};
    \draw[dashed, draw=nfslate] (0.7,-1.0) -- (0.7,-1.5);
    \draw[dashed, draw=nfslate] (9.0,0.0) -- (9.0,-1.5);
  \end{scope}

\end{tikzpicture}
\end{document}
}
  \caption{Le tre modalit\`a di esecuzione e il contrasto fra avvio a freddo e a
  caldo.}
  \label{fig:modalita}
\end{figure}

\section{\code{LOCAL}: esecuzione in-process}

\code{LocalDispatcher} \`e sei righe di codice:

\begin{lstlisting}[language=Java]
@Component
public class LocalDispatcher implements Dispatcher {
    @Override
    public CompletableFuture<DispatchResult> dispatch(InvocationTask task) {
        return CompletableFuture.completedFuture(
                DispatchResult.warm(InvocationResult.success(task.request().input())));
    }
}
\end{lstlisting}

Restituisce l'input come output, immediatamente, dichiarandosi caldo. \`E un eco.

Il suo scopo \`e sperimentale, non applicativo. Fornisce un percorso di
invocazione la cui componente di esecuzione ha costo nullo e varianza nulla, e
quindi rende misurabile il costo del control plane in isolamento: qualunque
latenza osservata su una funzione \code{LOCAL} \`e attribuibile interamente alla
piattaforma. \`E anche il modo in cui la maggior parte dei test di integrazione
esercita il percorso completo senza avviare container.

\section{\code{EXTERNAL}: inoltro a un endpoint esistente}

In modalit\`a \code{EXTERNAL} la funzione \`e ospitata fuori dal control plane,
che ne conosce soltanto l'URL (\code{endpointUrl}) e vi inoltra le invocazioni.
Non ne gestisce n\'e la creazione, n\'e le repliche, n\'e la distruzione. \`E,
letteralmente, un passthrough.

Il codice ne trae due conseguenze. La prima \`e che l'assenza di
\code{endpointUrl} \`e un errore rilevato al momento del dispatch, non della
registrazione:

\begin{lstlisting}[language=Java]
String endpoint = task.functionSpec().endpointUrl();
if (endpoint == null || endpoint.isBlank()) {
    return CompletableFuture.completedFuture(DispatchResult.warm(
            InvocationResult.error("EXTERNAL_ENDPOINT_MISSING",
                                   "endpointUrl is required for EXTERNAL mode")));
}
\end{lstlisting}

La seconda \`e che il blocco di scaling non viene risolto per le funzioni
\code{EXTERNAL} (\cref{cap:06-nucleo}): il control plane non pu\`o
ragionevolmente decidere il numero di repliche di un deployment che non
possiede.

La modalit\`a \code{EXTERNAL} \`e quella che il servizio di riferimento
\code{warm-echo} usa, ed \`e il veicolo principale degli esperimenti in cui il
runtime della funzione dev'essere tenuto costante mentre si variano le politiche
del control plane.

\section{\code{DEPLOYMENT}: intento gestito}

\code{DEPLOYMENT} \`e la modalit\`a completa: il control plane si assume la
responsabilit\`a di far esistere le repliche. \`E anche il default in assenza di
dichiarazione esplicita.

La sua realizzazione segue un principio: \code{DEPLOYMENT} \`e un
\emph{intento}, non un meccanismo. Il control plane non sa creare deployment;
sa dichiarare che ne serve uno e delegarne la realizzazione a un provider.

\subsection{L'interfaccia del provider}

\begin{lstlisting}[language=Java]
public interface ManagedDeploymentProvider {
    String backendId();
    boolean isAvailable();
    boolean supports(FunctionSpec spec);
    ProvisionResult provision(FunctionSpec spec);
    void deprovision(String functionName);
    void setReplicas(String functionName, int replicas);
    int getReadyReplicas(String functionName);
    default ReplicaStatus getReplicaStatus(String functionName) { ... }
}
\end{lstlisting}

L'interfaccia \`e piccola e i suoi primi tre metodi sono i pi\`u interessanti,
perch\'e non riguardano il provisioning ma la \emph{selezione}. Un provider deve
poter dichiarare di esistere (\code{backendId}), di essere utilizzabile in questo
momento (\code{isAvailable} --- il provider Kubernetes non lo \`e senza un
kubeconfig valido) e di poter servire una specifica particolare
(\code{supports}).

Il metodo \code{provision} restituisce un \code{ProvisionResult} che contiene
l'endpoint risultante e i metadati degli oggetti creati. Il ritorno di un
endpoint anzich\'e di un handle opaco \`e ci\`o che consente al dispatcher di
trattare uniformemente \code{DEPLOYMENT} ed \code{EXTERNAL} dopo il
provisioning: una volta che l'endpoint esiste, invocarlo \`e la stessa
operazione.

\subsection{La risoluzione del backend}

\code{DeploymentProviderResolver} sceglie il provider con una precedenza a tre
livelli: suggerimento esplicito per registrazione, default di piattaforma
(\code{nanofaas.deployment.default-backend}), e infine risoluzione implicita.

La risoluzione implicita \`e la parte in cui il progetto compie una scelta
difendibile ma non ovvia:

\begin{lstlisting}[language=Java]
if (candidates.isEmpty()) {
    throw new IllegalStateException("No managed deployment provider available for function ...");
}
if (candidates.size() > 1) {
    throw new IllegalStateException("Ambiguous managed deployment provider selection for function '"
            + spec.name() + "': " + ids);
}
return candidates.getFirst();
\end{lstlisting}

Con due provider entrambi disponibili e capaci, il resolver \textbf{fallisce}
anzich\'e scegliere. Una piattaforma con un ordine di preferenza implicito ---
``prova Kubernetes, altrimenti Docker'' --- funzionerebbe, e produrrebbe
deployment su un backend diverso da quello atteso ogni volta che l'ambiente
cambia sotto i piedi dell'operatore. Un'ambiguit\`a dichiarata \`e preferibile a
una preferenza nascosta. L'operatore risolve fissando
\code{default-backend}, che \`e un atto esplicito.

\subsection{Degradazione da \code{DEPLOYMENT} a \code{EXTERNAL}}

Esiste un caso in cui il provisioning fallisce ma la registrazione riesce:

\begin{lstlisting}[language=Java]
public ProvisionResult resolveAndProvision(FunctionSpec spec, String backendHint) {
    try {
        return resolve(spec, backendHint).provision(spec);
    } catch (IllegalStateException ex) {
        if (spec.endpointUrl() != null && !spec.endpointUrl().isBlank()) {
            return new ProvisionResult(spec.endpointUrl(), null,
                    ExecutionMode.EXTERNAL, ex.getMessage(), Map.of());
        }
        throw ex;
    }
}
\end{lstlisting}

Se la funzione dichiara \emph{anche} un \code{endpointUrl}, un fallimento del
provisioning la degrada a \code{EXTERNAL} anzich\'e respingerla. \`E la
degradazione che consente allo stesso manifesto di funzione di essere usato su
un cluster Kubernetes e su una macchina di sviluppo priva di backend gestito.

La degradazione \`e \emph{osservabile}, ed \`e questo a renderla accettabile
secondo il criterio del \cref{cap:03-requisiti-principi}. Il tipo
\code{DeploymentMetadata} conserva contemporaneamente la modalit\`a richiesta,
quella effettiva e la ragione della differenza:

\begin{lstlisting}[language=Java]
public record DeploymentMetadata(
        ExecutionMode requestedExecutionMode,
        ExecutionMode effectiveExecutionMode,
        String deploymentBackend,
        String degradationReason,
        String effectiveEndpointUrl,
        Map<String, String> deploymentObjects
) { ... }
\end{lstlisting}

L'API di ispezione della funzione restituisce entrambe le modalit\`a e la
ragione, cos\`i che un operatore che si aspettava un deployment Kubernetes possa
scoprire di averne ottenuto un passthrough --- e perch\'e --- senza leggere i
log.

\section{Risveglio da zero repliche}

Con un deployment scalato a zero, un'invocazione arriva prima che esista una
replica capace di servirla. Il nucleo affronta il caso con
\code{DeploymentWakeUpGate}: prima del dispatch verso l'endpoint, se la funzione
\`e gestita e non ha repliche pronte, la richiesta attende che ne compaia una.

Il meccanismo \`e coordinato: pi\`u invocazioni concorrenti sulla stessa funzione
condividono la stessa attesa anzich\'e emettere ciascuna la propria richiesta di
scale-up. \`E il minimo indispensabile per evitare che una raffica su una
funzione fredda si traduca in una raffica di comandi al backend.

\begin{damisurare}
La distribuzione del tempo di risveglio da zero repliche \`e strumentata
(\code{function\_cold\_start\_ms}) ma non \`e stata caratterizzata
sistematicamente: le campagne di misura del \cref{cap:23-prestazioni-base}
partono da repliche gi\`a calde. \todomis{distribuzione del tempo di risveglio
per backend Kubernetes e container locale}
\end{damisurare}

\section{Confronto}

\begin{table}[htbp]
\centering
\caption{Le tre modalit\`a a confronto.}
\label{tab:modalita}
\small
\begin{tabularx}{\textwidth}{@{}lXXX@{}}
\toprule
 & \textbf{\code{LOCAL}} & \textbf{\code{EXTERNAL}} & \textbf{\code{DEPLOYMENT}} \\
\midrule
Processo funzione & nessuno & preesistente & creato dal control plane \\
Ciclo di vita & --- & del proprietario & del control plane \\
Scaling repliche & --- & non applicabile & \code{autoscaler} o backend \\
Cold start & assente & fuori dal controllo & osservato e misurato \\
Validazione immagine & no & no & s\`i, dal validatore del backend \\
Uso previsto & misura del control plane, test & runtime fissato negli esperimenti & carichi realistici \\
\bottomrule
\end{tabularx}
\end{table}

La riga ``uso previsto'' riassume la logica del progetto: le tre modalit\`a non
sono tre livelli di completezza ma tre strumenti sperimentali. \code{LOCAL}
isola il control plane, \code{EXTERNAL} lo mette a confronto con un runtime che
non cambia fra un esperimento e l'altro, \code{DEPLOYMENT} chiude il ciclo e
include l'infrastruttura nella misura.

\chapter{Il sistema a moduli}
\label{cap:08-sistema-moduli}

Il \cref{cap:03-requisiti-principi} ha enunciato la regola che governa la
divisione fra nucleo e moduli: il binario prodotto senza moduli non deve
contenerli. Questo capitolo descrive il meccanismo che la realizza, ne discute i
limiti, e presenta il criterio con cui il progetto tratta le dipendenze fra
moduli.

\section{Perch\'e a tempo di build e non a runtime}

L'approccio convenzionale a un sistema estensibile \`e l'attivazione
condizionale: tutte le funzionalit\`a sono compilate nel binario e un flag di
configurazione decide quali istanziare. Spring Boot lo supporta nativamente con
\code{@ConditionalOnProperty}, ed \`e la scelta che quasi ogni applicazione fa.

\nf{} non la fa, per tre ragioni.

\paragraph{La superficie del binario.} Un modulo disattivato \`e comunque
presente: le sue classi sono nel classpath, le sue dipendenze transitive sono
risolte, il suo codice \`e analizzato dall'ottimizzatore. Per la compilazione
in immagine nativa (\cref{cap:20-build-packaging}) questo \`e sostanziale: il
provider Kubernetes trascina il client Fabric8 e Vert.x, il provider container
trascina il client Docker Java. Un'immagine nativa che contenga entrambi \`e
significativamente pi\`u grande e pi\`u lenta da compilare di una che contenga
solo ci\`o che verr\`a usato.

\paragraph{L'onest\`a della misura.} Un esperimento che confronta il costo del
control plane con e senza controllo di concorrenza deve confrontare due binari
diversi, non lo stesso binario con un flag diverso. Nel secondo caso le classi
inattive contribuiscono comunque all'impronta di memoria, al tempo di
inizializzazione del contesto Spring e alla pressione sulla cache di istruzioni.
Il \cref{cap:24-footprint} misura configurazioni, e questo richiede che le
configurazioni siano artefatti distinti.

\paragraph{L'assenza di rami condizionali.} \`E l'argomento pi\`u forte. Se i
moduli fossero attivabili a runtime, il codice del nucleo dovrebbe contenere test
sulla loro presenza. Con la selezione a build time, il nucleo dipende solo da
interfacce di cui esiste sempre un'implementazione, e il percorso di invocazione
non contiene un solo \code{if} che riguardi la configurazione modulare.

\section{Il meccanismo}

\subsection{Selezione}

La selezione avviene attraverso una propriet\`a Gradle o una variabile d'ambiente:

\begin{lstlisting}[language=bash]
./gradlew :control-plane:bootJar -PcontrolPlaneModules=none
./gradlew :control-plane:bootJar -PcontrolPlaneModules=async-queue,sync-queue
./gradlew :control-plane:bootJar -PcontrolPlaneModules=all
NANOFAAS_CONTROL_PLANE_MODULES=all ./gradlew :control-plane:bootJar
\end{lstlisting}

I moduli disponibili sono scoperti scandendo le sottodirectory di
\code{platform/modules/} che contengono uno script di build --- non esiste un
elenco da mantenere allineato. La selezione \`e poi validata con tre regole:

\begin{lstlisting}[language=Java, caption={Validazione della selezione in
\code{settings.gradle}.}]
if (requestedModules.contains('none')) {
    if (requestedModules.size() > 1) {
        throw new GradleException("Module selector 'none' cannot be combined with other values: ...")
    }
    requestedModules.clear()
}
if (requestedModules.contains('all')) {
    requestedModules = new LinkedHashSet<String>(availableModules)
}
def unknownModules = requestedModules.findAll { !availableModules.contains(it) }
if (!unknownModules.isEmpty()) {
    throw new GradleException("Unknown control-plane modules: ${unknownModules}. Available: ${availableModules}")
}
\end{lstlisting}

Un nome sconosciuto fa fallire la build. \`E la scelta corretta per un selettore
di questo tipo: ignorare silenziosamente \code{-PcontrolPlaneModules=async-quee}
produrrebbe un binario privo della coda asincrona e un operatore convinto del
contrario.

I moduli selezionati entrano nel classpath come dipendenze \code{runtimeOnly}.
La scelta \`e significativa: il codice del control plane non pu\`o compilare
contro una classe di modulo, perch\'e a tempo di compilazione i moduli non sono
sul suo compile classpath. Il compilatore \`e cos\`i l'esecutore della regola di
isolamento, non un revisore.

\subsection{Caricamento}

A runtime, i moduli si presentano attraverso il \code{ServiceLoader} di Java.
Ciascuno pubblica un file
\code{META-INF/services/it.unimib.datai.nanofaas.common.controlplane.ControlPlaneModule}
contenente il nome della propria classe di implementazione, che \`e tipicamente
di dieci righe:

\begin{lstlisting}[language=Java]
public final class ConcurrencyControlModule implements ControlPlaneModule {
    @Override
    public Set<Class<?>> configurationClasses() {
        return Set.of(ConcurrencyControlConfiguration.class);
    }
}
\end{lstlisting}

Il ponte fra \code{ServiceLoader} e Spring \`e un \code{DeferredImportSelector}:

\begin{lstlisting}[language=Java, caption={\code{ControlPlaneModuleImportSelector}.}]
public final class ControlPlaneModuleImportSelector implements DeferredImportSelector {
    @Override
    public String[] selectImports(AnnotationMetadata importingClassMetadata) {
        ClassLoader classLoader = Thread.currentThread().getContextClassLoader();
        Set<String> moduleConfigurations = new LinkedHashSet<>();
        ServiceLoader.load(ControlPlaneModule.class, classLoader)
                .forEach(module -> module.configurationClasses().stream()
                        .filter(Objects::nonNull)
                        .map(Class::getName)
                        .forEach(moduleConfigurations::add));
        return moduleConfigurations.toArray(String[]::new);
    }
}
\end{lstlisting}

Il commento nel sorgente ne spiega la ragione d'essere: \emph{``Ensures optional
module \code{@Configuration} classes are imported for every bootstrap path,
including tests that do not execute \code{ControlPlaneApplication.main()}.''} Un
caricamento eseguito nel \code{main} funzionerebbe in produzione e non nei test
di integrazione, che avviano il contesto per altra via. Legandolo a
un'annotazione \code{@Import} sulla classe applicativa, ogni percorso di
avvio che usi quella classe carica i moduli.

Va notato che i package dei moduli \textbf{non} sono soggetti a component
scanning: ogni modulo registra i propri bean con \code{@Bean} espliciti nella
propria classe di configurazione. La scansione automatica avrebbe reso il
contenuto del contesto dipendente dai nomi dei package, che \`e esattamente il
tipo di accoppiamento implicito che il sistema modulare vuole evitare.

\section{Le dipendenze fra moduli}

L'interfaccia \code{ControlPlaneModule} non consente a un modulo di dichiarare
di averne bisogno di un altro. \`E una lacuna apparente, e il progetto l'ha
mantenuta di proposito: l'unica dipendenza legittima fra moduli \`e attraverso
un'interfaccia del nucleo, e in tal caso il modulo dipende dall'interfaccia,
non dal modulo che la implementa.

Rimane per\`o il caso in cui l'implementazione neutra dell'interfaccia rende il
modulo dipendente \emph{inerte}, e questo caso ha bisogno di una risposta.

\subsection{Il caso fatale: \code{concurrency-control}}

Il modulo \code{concurrency-control} legge lo stato delle code da
\code{ScalingMetricsSource} e vi \emph{riscrive} i limiti che calcola. Il nucleo
registra sempre un'implementazione no-op di quell'interfaccia, quindi la
condizione \code{@ConditionalOnBean} \`e soddisfatta anche quando nessun modulo
ne fornisce una reale. Il costruttore della configurazione codifica quindi il
controllo che l'annotazione non pu\`o fare:

\begin{lstlisting}[language=Java]
/**
 * Refuses to start when no module supplies real queue state.
 *
 * <p>{@code @ConditionalOnBean} above cannot catch this: the core always registers a no-op
 * source, so the condition holds even when nothing produces one. The governor would then run
 * against a source reporting depth 0 and in-flight 0 forever, and -- worse -- every limit it
 * computed would be written into that same no-op and enforced by nobody. The module would be
 * entirely inert while its metrics claimed otherwise.
 */
public ConcurrencyControlConfiguration(ScalingMetricsSource metricsSource) {
    if (!metricsSource.enabled()) {
        throw new IllegalStateException(
                "concurrency-control needs a module that supplies queue state (async-queue); "
                        + "without one the governor reads zeroes and the limits it computes "
                        + "are enforced by nothing. Add async-queue to the module selection.");
    }
}
\end{lstlisting}

Il metodo \code{enabled()} sull'interfaccia \`e ci\`o che rende possibile la
distinzione. Un'implementazione no-op \`e indistinguibile da una reale che
osserva un sistema scarico --- entrambe riportano zero --- a meno che non
dichiari di essere no-op. \`E una tecnica generale: quando un default neutro
pu\`o essere confuso con un'osservazione legittima, deve poter essere
interrogato sulla propria natura.

Il messaggio di errore nomina il modulo mancante e l'azione correttiva. In un
sistema in cui la configurazione avviene a tempo di build, un errore di
configurazione si manifesta all'avvio ed \`e l'unica occasione per dire
all'operatore che cosa fare.

\subsection{Il caso non fatale: \code{autoscaler}}

Lo stesso modulo \code{autoscaler} legge la stessa sorgente, ma non fallisce.
La differenza \`e che la sua metrica principale --- \code{rps} --- proviene da
un contatore Micrometer indipendente dalle code. Senza \code{async-queue}
l'autoscaler resta pienamente funzionante sulla metrica di default, e sono le
due metriche derivate dalle code (\code{queue\_depth} e \code{in\_flight}) a
diventare cieche. Il modulo emette un avviso al primo uso di una di esse.

Il criterio che distingue i due casi \`e quello enunciato nel
\cref{cap:03-requisiti-principi}: degradazione osservabile, s\`i; degradazione
silenziosa, no. \code{autoscaler} degrada parzialmente e lo dichiara;
\code{concurrency-control} degraderebbe totalmente e in silenzio, quindi non
degrada affatto.

\section{Isolamento verificato}

La regola secondo cui un modulo non deve dipendere da un altro modulo \`e
verificata automaticamente con ArchUnit~\cite{archunit}, e ciascuno dei nove
moduli porta la propria classe di verifica. La regola \`e in vigore senza
eccezioni: al momento della stesura nessun modulo dichiara una deroga.

La formulazione della regola merita una nota, perch\'e definisce con precisione
che cosa sia una dipendenza illegittima. Un modulo non pu\`o dipendere da
\emph{altri moduli}; pu\`o invece dipendere liberamente dalle interfacce del
nucleo, che \`e il meccanismo previsto. Cos\`i i provider di deployment
implementano \code{ImageValidator}, dichiarata nel package \code{registry} del
control plane (\cref{cap:15-deployment-provider}), senza violare nulla: il
riferimento va verso il nucleo, non lateralmente.

Il costo di questa formulazione \`e che una dipendenza fra moduli non pu\`o
essere introdotta per errore ma pu\`o esserlo per decisione: il commento nel
codice prescrive che ogni coppia sanzionata sia aggiunta come riga esplicita con
la propria giustificazione. Nessuna \`e stata aggiunta finora.

\section{Il catalogo dei moduli}

\begin{table}[htbp]
\centering
\caption{I nove moduli opzionali, con dimensione del codice di produzione e
capitolo di riferimento.}
\label{tab:moduli}
\small
\begin{tabularx}{\textwidth}{@{}lrXl@{}}
\toprule
\textbf{Modulo} & \textbf{LOC} & \textbf{Ruolo} & \textbf{Cap.} \\
\midrule
\code{async-queue} & 682 & Code per funzione e scheduler del percorso asincrono
& \ref{cap:09-async-queue} \\
\code{sync-queue} & 905 & Ammissione e backpressure sul percorso sincrono
& \ref{cap:10-sync-queue} \\
\code{autoscaler} & 754 & Scaling interno del numero di repliche
& \ref{cap:11-autoscaler} \\
\code{concurrency-control} & 1\,314 & Governo del limite di concorrenza per
funzione & \ref{cap:12-concurrency-control} \\
\code{offload} & 230 & Proxy condizionale verso un'istanza remota
& \ref{cap:13-offload} \\
\code{runtime-config} & 540 & Riconfigurazione a caldo e API di
amministrazione & \ref{cap:14-runtime-config} \\
\code{k8s-deployment-provider} & 1\,202 & Backend di deployment Kubernetes
& \ref{cap:15-deployment-provider} \\
\code{container-deployment-provider} & 1\,166 & Backend di deployment su runtime
container locale & \ref{cap:15-deployment-provider} \\
\code{build-metadata} & 43 & Endpoint diagnostico sulla build
& \ref{cap:16-moduli-diagnostici} \\
\midrule
\textbf{Totale moduli} & \textbf{6\,836} & & \\
\textit{Nucleo, per confronto} & \textit{5\,065} & & \\
\bottomrule
\end{tabularx}
\end{table}

Il rapporto fra le due ultime righe \`e l'osservazione con cui vale la pena
chiudere il capitolo. La somma dei moduli supera il nucleo che estendono: la
maggior parte del codice di \nf{} \`e codice che si pu\`o non compilare. Se il
meccanismo di selezione modulare fosse mal progettato, questo sarebbe un
problema di manutenzione; essendo verificato dal compilatore e dalle regole
architetturali, \`e invece la misura di quanto la piattaforma sia riducibile.

\section{Il default \`e \code{all}}

Un'ultima nota, di metodo pi\`u che di architettura. In assenza di selezione
esplicita, ogni invocazione di Gradle --- compresa quella dei test --- usa
\code{all}. La ragione \`e che la suite di test deve esercitare la
configurazione che viene distribuita, non una configurazione parziale che
nessuno esegue.

La configurazione priva di moduli non resta per\`o senza copertura: la classe
\code{CoreOnlyApiTest} verifica le propriet\`a che valgono \emph{solo} in
assenza di moduli --- accodatore no-op, \code{501} su \code{:enqueue} --- e si
auto-esclude quando i moduli sono presenti. Poich\'e la selezione modulare \`e
risolta a tempo di configurazione di Gradle e non pu\`o variare all'interno di
una singola build, quel test richiede una propria invocazione con
\code{-PcontrolPlaneModules=none}. \`E un costo reale sulla catena di
integrazione continua, ed \`e il prezzo diretto della scelta di selezionare a
build time anzich\'e a runtime.


\part{I moduli}
\chapter{Il modulo \code{async-queue}}
\label{cap:09-async-queue}

\code{async-queue} \`e il modulo che rende possibile l'invocazione asincrona, e
--- meno ovviamente --- \`e anche l'unico che fornisce alla piattaforma
un'osservazione reale dello stato delle code. Il primo ruolo \`e quello che il
suo nome dichiara; il secondo \`e ci\`o che lo rende una dipendenza per
\code{concurrency-control} e una dipendenza parziale per \code{autoscaler}.

Sono 682 righe di codice in nove classi.

\section{Struttura}

\begin{center}
\small
\begin{adjustbox}{max width=\textwidth}
\begin{tabular}{@{}lp{9cm}@{}}
\toprule
\textbf{Classe} & \textbf{Ruolo} \\
\midrule
\code{FunctionQueueState} & Stato di una singola funzione: coda limitata,
contatore delle richieste in volo, limiti di concorrenza. \\
\code{QueueManager} & Mappa da nome di funzione a stato, registrazione delle
metriche, notifica del lavoro disponibile. \\
\code{Scheduler} & Ciclo a thread singolo che consuma dalle code e dispaccia. \\
\code{QueueBackedEnqueuer} & Implementazione di \code{InvocationEnqueuer}. \\
\code{QueueBackedMetricsSource} & Implementazione di \code{ScalingMetricsSource}. \\
\code{QueueConcurrencyControlMetrics} & Metriche descrittive del controllore
attivo per funzione. \\
\code{WorkSignaler} & Interfaccia a un metodo per la notifica. \\
\bottomrule
\end{tabular}
\end{adjustbox}
\end{center}

\section{Lo stato di una funzione}

\code{FunctionQueueState} incapsula quattro grandezze: la coda dei compiti
(\code{ArrayBlockingQueue} di dimensione fissa), il contatore delle richieste in
volo, il limite di concorrenza \emph{configurato} e il limite di concorrenza
\emph{effettivo}.

La distinzione fra i due limiti \`e la chiave dell'intero modulo. Il limite
configurato \`e ci\`o che l'utente ha dichiarato nella specifica della funzione.
Il limite effettivo \`e ci\`o che il modulo \code{concurrency-control}, se
presente, ha deciso di applicare in questo momento. In assenza di quel modulo i
due coincidono sempre.

\subsection{Acquisizione dello slot}

L'ammissione al dispatch \`e un'operazione atomica di test-e-incremento:

\begin{lstlisting}[language=Java, caption={\code{FunctionQueueState.tryAcquireSlot()}.}]
/**
 * Atomically checks if dispatch is allowed and increments inFlight if so.
 * This prevents race conditions where multiple threads could both pass
 * canDispatch() check and then both increment, exceeding the concurrency limit.
 */
public boolean tryAcquireSlot() {
    while (true) {
        int current = inFlight.get();
        if (current >= effectiveConcurrency) {
            return false;
        }
        if (inFlight.compareAndSet(current, current + 1)) {
            return true;
        }
        // CAS failed, another thread modified - retry
    }
}
\end{lstlisting}

L'alternativa ingenua --- verificare \code{canDispatch()} e poi incrementare ---
\`e corretta solo con un unico consumatore. Il modulo ne ha uno solo oggi (lo
scheduler \`e a thread singolo), ma il metodo \`e esposto anche al percorso
sincrono attraverso \code{InvocationEnqueuer.tryAcquireSlot}, che gira su thread
arbitrari. La primitiva atomica \`e quindi necessaria, non difensiva.

\subsection{Il rilascio con pavimento a zero}

Il rilascio \`e simmetrico ma con una protezione aggiuntiva:

\begin{lstlisting}[language=Java]
private void decrementInFlightNonNegative() {
    while (true) {
        int current = inFlight.get();
        if (current == 0) { return; }
        if (inFlight.compareAndSet(current, current - 1)) { return; }
    }
}
\end{lstlisting}

Un decremento su un contatore gi\`a a zero viene ignorato anzich\'e produrre un
valore negativo. La condizione non dovrebbe verificarsi --- il
\cref{cap:06-nucleo} descrive il meccanismo che rende idempotente il rilascio
per tentativo --- ma se si verificasse, un contatore negativo si propagherebbe
direttamente nella metrica \code{function\_inFlight}, che \`e la sorgente su cui
decidono sia l'autoscaler sia il controllore di concorrenza. Un errore di
contabilit\`a in questo punto non produce un malfunzionamento locale: produce
decisioni di politica sbagliate a valle, per un tempo indeterminato. Il pavimento
a zero \`e la contenzione minima di un errore che si propaga.

\subsection{Interazione fra limite configurato ed effettivo}

I due metodi che scrivono i limiti codificano una semantica che vale la pena
esplicitare:

\begin{lstlisting}[language=Java]
public void concurrency(int concurrency) {
    int previousConfigured = this.configuredConcurrency;
    int normalized = Math.max(1, concurrency);
    this.configuredConcurrency = normalized;
    if (effectiveConcurrency == previousConfigured || effectiveConcurrency > normalized) {
        // Preserve fixed-mode semantics: effective limit tracks configured limit
        // (and clamps down on shrink).
        effectiveConcurrency = normalized;
    }
}

public void setEffectiveConcurrency(int effectiveConcurrency) {
    int clamped = Math.max(1, effectiveConcurrency);
    if (clamped > configuredConcurrency) { clamped = configuredConcurrency; }
    this.effectiveConcurrency = clamped;
}
\end{lstlisting}

Due invarianti emergono. Il primo: \textbf{il limite effettivo non pu\`o mai
superare quello configurato}. Un controllore adattivo pu\`o restringere ma non
allargare oltre ci\`o che l'utente ha autorizzato; la specifica della funzione
resta un tetto. Il secondo: un aggiornamento del limite configurato trascina
quello effettivo \emph{solo} se questo non era stato scostato da un controllore,
oppure se il nuovo tetto \`e pi\`u basso. In pratica: cambiare la configurazione
non cancella una decisione adattiva in corso, ma un abbassamento del tetto la
vincola immediatamente.

\section{Lo scheduler}

Lo scheduler \`e un ciclo su thread singolo. La sua struttura \`e determinata da
un requisito: non deve sondare le code inattive.

\begin{lstlisting}[language=Java]
private void loop() {
    while (running.get()) {
        String functionName = activeFunctions.poll(500, TimeUnit.MILLISECONDS);
        if (functionName != null) {
            enqueuedFunctions.remove(functionName);
            processFunction(functionName);
        }
        ...
    }
}
\end{lstlisting}

La coda \code{activeFunctions} contiene i nomi delle funzioni che hanno lavoro
pendente; l'insieme \code{enqueuedFunctions} evita che lo stesso nome vi compaia
pi\`u volte. Un accodamento notifica lo scheduler solo se il nome non era gi\`a
segnalato. Il risultato \`e che il costo del ciclo \`e proporzionale al numero di
funzioni \emph{attive}, non al numero di funzioni registrate: registrare mille
funzioni inattive non costa nulla.

Il timeout di 500\,ms sulla \code{poll} \`e ci\`o che rende il ciclo terminabile:
senza di esso l'arresto richiederebbe un'interruzione del thread bloccato.

\subsection{Il batch limitato e l'equit\`a}

\begin{lstlisting}[language=Java]
private static final int MAX_BATCH_PER_FUNCTION = 2;

private void processFunction(String functionName) {
    ...
    int dispatched = 0;
    while (running.get() && dispatched < MAX_BATCH_PER_FUNCTION && state.tryAcquireSlot()) {
        InvocationTask task = state.poll();
        if (task == null) { state.releaseSlot(); break; }
        dispatched++;
        SchedulerDispatchSupport.dispatchWithFailureCleanup(
                task, () -> invocationService.dispatch(task), state::releaseSlot, log);
    }
    if (state.queued() > 0 && dispatched > 0) {
        signalWork(functionName);
    }
}
\end{lstlisting}

Ogni visita a una funzione dispaccia al massimo due compiti, poi la funzione
viene rimessa in fondo alla coda dei nomi attivi se ha ancora lavoro. \`E un
round-robin con quanto pari a due.

Il quanto limitato \`e ci\`o che garantisce l'equit\`a fra funzioni. Senza di
esso, una funzione con mille compiti accodati e slot disponibili terrebbe il
thread dello scheduler finch\'e non li avesse dispacciati tutti, e una funzione
con un solo compito accodata subito dopo attenderebbe l'intero drenaggio. Con il
quanto, la seconda attende al pi\`u due dispatch.

Il valore due, e non uno, riduce il costo di rotazione --- ogni ciclo comporta
un'operazione su due strutture concorrenti --- senza degradare
apprezzabilmente l'equit\`a. \`E un compromesso, e il progetto non ne documenta
una taratura sperimentale.

\begin{damisurare}
L'effetto di \code{MAX\_BATCH\_PER\_FUNCTION} sull'equit\`a fra funzioni
co-residenti e sul throughput aggregato non \`e stato misurato. Il
\cref{cap:25-valutazione-concorrenza} riporta misure di co-tenancy che
\emph{includono} l'effetto di questo parametro senza isolarlo.
\todomis{sensibilit\`a del throughput e della latenza al quanto di scheduling}
\end{damisurare}

\subsection{Pulizia in caso di fallimento del dispatch}

Il dispatch \`e affidato a un helper del nucleo, \code{dispatchWithFailureCleanup},
a cui viene passato il rilascio dello slot come azione di compensazione. Il
motivo \`e che lo slot \`e stato acquisito \emph{prima} del dispatch: se il
dispatch solleva prima di aver stabilito una catena di completamento, nessuno
rilascer\`a lo slot, e la funzione perderebbe permanentemente una unit\`a di
concorrenza. Ripetuto, il fenomeno porterebbe il limite effettivo a zero senza
che alcuna metrica lo segnali come anomalia --- il contatore delle richieste in
volo apparirebbe semplicemente sempre saturo.

\section{Le tre metriche esportate}

\code{QueueManager} registra tre gauge per ciascuna funzione al momento della
sua creazione:

\begin{center}
\small
\begin{adjustbox}{max width=\textwidth}
\begin{tabular}{@{}lll@{}}
\toprule
\textbf{Metrica} & \textbf{Sorgente} & \textbf{Consumatore} \\
\midrule
\code{function\_queue\_depth} & \code{state::queued} & autoscaler,
concurrency-control, operatore \\
\code{function\_inFlight} & \code{state::inFlight} & autoscaler,
concurrency-control \\
\code{function\_effective\_concurrency} & \code{state::effectiveConcurrency} &
osservazione della politica \\
\bottomrule
\end{tabular}
\end{adjustbox}
\end{center}

La terza \`e diagnostica e non alimenta alcuna decisione: rende osservabile
\emph{che cosa il controllore ha deciso}, distinguendolo da ci\`o che l'utente
aveva configurato. Senza di essa, osservare l'effetto di una politica adattiva
richiederebbe di dedurlo dal comportamento.

I meter sono rimossi quando la funzione viene rimossa, con la conseguenza gi\`a
discussa nel \cref{cap:06-nucleo} sugli scenari di verifica che interrogano le
metriche dopo la cancellazione.

\section{I due punti di estensione implementati}

\subsection{\code{QueueBackedEnqueuer}}

Implementa \code{InvocationEnqueuer} delegando a \code{QueueManager}. La sua
esistenza \`e ci\`o che fa passare \code{enabled()} da \code{false} a
\code{true}, e quindi ci\`o che fa rispondere \code{202} anzich\'e \code{501} a
\code{POST :enqueue}.

\subsection{\code{QueueBackedMetricsSource}}

Implementa \code{ScalingMetricsSource}, ed \`e la classe da cui dipende l'intera
Parte III di questo documento. Fornisce quattro operazioni: leggere la
profondit\`a della coda, leggere le richieste in volo, \emph{scrivere} il limite
di concorrenza effettivo, e registrare il modo del controllore attivo.

La terza \`e quella che conta. Un modulo di controllo della concorrenza calcola
un limite; \`e la scrittura attraverso questa interfaccia a rendere quel limite
effettivo, perch\'e da quel momento \code{tryAcquireSlot} lo confronta con il
contatore delle richieste in volo. \`E letteralmente l'unico punto in cui una
politica di concorrenza tocca il comportamento del sistema, e questo spiega
perch\'e \code{concurrency-control} rifiuti di avviarsi senza
\code{async-queue}: la scrittura verso un'implementazione no-op non fallisce, non
segnala nulla, e non ha effetto.

Il metodo \code{enabled()} distingue questa implementazione dal no-op del nucleo,
e nel modulo restituisce costantemente \code{true}.

\section{Il ruolo del modulo nel percorso sincrono}

Un'ultima nota, perch\'e il nome del modulo pu\`o trarre in inganno. La catena di
ammissione descritta nel \cref{cap:04-architettura-control-plane} usa la coda
asincrona anche per le invocazioni \emph{sincrone}, quando la coda sincrona non
\`e presente:

\begin{lstlisting}[language=Java]
private void admitLocally(ExecutionRecord executionRecord) {
    if (syncQueueGateway.enabled()) {
        syncQueueGateway.enqueueOrThrow(executionRecord.task());
    } else if (enqueuer.enabled()) {
        InvocationEnqueueSupport.enqueueOrThrow(enqueuer, metrics, executionRecord);
    } else {
        completionHandler.dispatch(executionRecord.task());
    }
}
\end{lstlisting}

In quella configurazione una richiesta sincrona viene accodata sulla coda della
funzione, il chiamante attende il futuro di completamento, e lo scheduler la
dispaccia quando uno slot si libera. Il risultato \`e che il limite di
concorrenza per funzione \`e applicato anche al traffico sincrono, il che \`e
esattamente ci\`o che si desidera. La differenza rispetto alla coda sincrona
(\cref{cap:10-sync-queue}) \`e che qui non esiste ammissione: una richiesta
sincrona che trova la coda piena riceve \code{429} da \code{QueueFullException}
senza che nessuno abbia stimato quanto avrebbe dovuto attendere.

\chapter{Il modulo \code{sync-queue}}
\label{cap:10-sync-queue}

Il modulo \code{sync-queue} inserisce sul percorso sincrono una coda esplicita
con controllo di ammissione. \`E l'applicazione diretta del principio di
SEDA~\cite{welsh2001seda}: il punto in cui le richieste si accumulano dev'essere
noto e strumentato, anzich\'e disperso nei buffer del sistema operativo o nella
lunghezza delle connessioni pendenti.

Il modulo risponde a una domanda che il \code{async-queue} non pone: \emph{questa
richiesta ha senso accettarla?} Una coda senza ammissione accetta finch\'e ha
posto e rifiuta quando \`e piena, il che significa che l'ultima richiesta
accettata prima del riempimento attende quanto l'intera coda impiega a
drenare. Con l'ammissione, il rifiuto avviene prima, sulla base di una stima.

Sono 905 righe in dieci classi.

\section{Le due domande dell'ammissione}

\code{SyncQueueAdmissionController} \`e la classe pi\`u corta e pi\`u
significativa del modulo:

\begin{lstlisting}[language=Java]
public SyncQueueAdmissionResult evaluate(String functionName, int depth, Instant now) {
    if (depth >= maxDepth) {
        return SyncQueueAdmissionResult.rejected(SyncQueueRejectReason.DEPTH,
                                                 Double.POSITIVE_INFINITY);
    }
    double estWaitSeconds = estimator.estimateWaitSeconds(functionName, depth, now);
    long maxWaitSeconds = configSource.syncQueueMaxEstimatedWait().toSeconds();
    if (configSource.syncQueueAdmissionEnabled()
            && (maxWaitSeconds == 0 || estWaitSeconds > maxWaitSeconds)) {
        return SyncQueueAdmissionResult.rejected(SyncQueueRejectReason.EST_WAIT, estWaitSeconds);
    }
    return SyncQueueAdmissionResult.accepted(estWaitSeconds);
}
\end{lstlisting}

Le due condizioni sono qualitativamente diverse.

\paragraph{\code{DEPTH}} \`e un vincolo di risorsa. La coda ha una capacit\`a
finita e non pu\`o accettare oltre. Il rifiuto \`e incondizionato e non dipende
dalla configurazione dell'ammissione: anche disattivando il controllo di
ammissione la coda continua a rifiutare quando \`e piena.

\paragraph{\code{EST\_WAIT}} \`e un vincolo di qualit\`a del servizio. La coda
avrebbe posto, ma la richiesta attenderebbe troppo a lungo. Il rifiuto \`e una
scelta di politica, disattivabile.

La distinzione \`e resa visibile al chiamante attraverso l'header
\code{X-Queue-Reject-Reason}, ed \`e la stessa distinzione su cui il modulo
\code{offload} basa la propria attivazione reattiva
(\cref{cap:13-offload}): entrambe le condizioni sono \emph{pressione}, e
giustificano una deviazione verso un'istanza remota, mentre un
\code{TIMEOUT} in coda non lo \`e.

\section{La stima dell'attesa}

\subsection{Il modello}

\code{WaitEstimator} applica la forma pi\`u diretta della legge di
Little~\cite{little1961proof}. Se il sistema smaltisce $\mu$ richieste al secondo
e ne ha $Q$ in coda, una richiesta che si accodasse adesso attenderebbe
approssimativamente

\[
  W_{\text{stim}} = \frac{Q}{\mu}.
\]

Il throughput $\mu$ \`e stimato contando i dispatch effettivi in una finestra
scorrevole:

\begin{lstlisting}[language=Java]
public double estimateWaitSeconds(String functionName, int queueDepth, Instant now) {
    if (queueDepth <= 0) { return 0.0; }
    ThroughputSnapshot perFunction = snapshot(perFunctionEvents.get(functionName), now);
    if (perFunction.samples() >= perFunctionMinSamples && perFunction.throughput() > 0) {
        return queueDepth / perFunction.throughput();
    }
    ThroughputSnapshot global = snapshot(globalEvents, now);
    if (global.throughput() <= 0) { return Double.POSITIVE_INFINITY; }
    return queueDepth / global.throughput();
}
\end{lstlisting}

\subsection{Le tre decisioni della stima}

\paragraph{Prima: per funzione, se ci sono abbastanza campioni.} Funzioni diverse
hanno tempi di servizio diversi, quindi la stima per funzione \`e la sola
significativa. Ma una stima costruita su pochi campioni \`e rumorosa, e un
rumore su questa grandezza si traduce direttamente in rifiuti ingiustificati. Il
parametro \code{perFunctionMinSamples} (default 50) fissa la soglia sotto la
quale la stima per funzione non \`e considerata affidabile.

\paragraph{Seconda: altrimenti, globale.} Sotto la soglia si ripiega sul
throughput aggregato della piattaforma. \`E una stima peggiore --- mescola
funzioni con costi diversi --- ma \`e disponibile prima. Il ripiego \`e una
scelta pragmatica: preferire una stima grossolana a nessuna stima.

\paragraph{Terza: senza campioni, attesa infinita.} \`E il caso pi\`u
interessante e ha una conseguenza operativa che il progetto ha appreso per via
empirica. Con \code{globalThroughput = 0} la stima \`e
\code{Double.POSITIVE\_INFINITY}, che eccede qualunque soglia, e quindi
\emph{ogni} richiesta viene rifiutata con \code{429}.

Questo \`e corretto in produzione --- un sistema che non ha ancora dispacciato
nulla non pu\`o promettere alcun tempo di attesa --- e distruttivo nei test: una
suite che avvia il control plane e invia la prima richiesta la vede respinta con
\code{429} senza alcuna congestione reale. La soluzione adottata \`e che i test
di API disattivano esplicitamente la coda sincrona
(\code{sync-queue.enabled=false}), il che \`e registrato come vincolo noto
dell'impianto di test (\cref{cap:26-qualita-verifica}).

\`E anche un esempio del genere di comportamento che una piattaforma di ricerca
deve documentare: la condizione ``nessuna storia'' non \`e un caso limite raro,
\`e lo stato in cui il sistema si trova ogni volta che viene avviato.

\subsection{La finestra}

Il default \`e una finestra di 30 secondi. Il calcolo del throughput divide il
numero di campioni per la durata della finestra, non per l'intervallo effettivo
fra il primo e l'ultimo campione:

\begin{lstlisting}[language=Java]
double seconds = Math.max(1.0, window.toSeconds());
return new ThroughputSnapshot(samples, samples / seconds);
\end{lstlisting}

Questo \emph{sottostima} il throughput durante i primi 30 secondi di attivit\`a,
perch\'e divide per una finestra che il sistema non ha ancora riempito. Una
sottostima del throughput produce una sovrastima dell'attesa, e quindi un
comportamento pi\`u conservativo --- pi\`u rifiuti --- nella fase iniziale. La
direzione dell'errore \`e quindi sicura, ed \`e presumibilmente il motivo per cui
la formulazione \`e stata mantenuta; il progetto non lo documenta esplicitamente.

\section{Lo scheduler sincrono e il forward progress}

Il consumatore della coda \`e \code{SyncScheduler}, e la sua funzione
\code{tickOnce()} risolve un problema che una coda FIFO ingenua non risolve.

\subsection{Il problema della testa bloccata}

Con una singola coda condivisa da tutte le funzioni, la richiesta in testa
appartiene a una funzione che potrebbe non avere slot di concorrenza liberi. Un
consumatore FIFO stretto si bloccherebbe su di essa, e nessuna richiesta di
nessun'altra funzione verrebbe dispacciata --- \emph{head-of-line blocking}, con
la variante particolarmente sgradevole che il blocco dura quanto il tempo di
servizio della funzione bloccante.

\subsection{La soluzione: scansione selettiva}

\begin{lstlisting}[language=Java, caption={Il cuore di \code{tickOnce()}.}]
SyncQueueItem item = queue.findReadyMatching(now,
        task -> enqueuer.hasAvailableSlot(task.functionName()));
if (item == null) {
    if (queue.peekReady(now) == null) {
        blockedBackoffMs = tickMs;
        queue.awaitWork(tickMs);          // coda vuota: attendi lavoro
    } else {
        queue.rotateReadyScanWindow(now); // tutti bloccati: ruota e ritenta
        pause.accept(currentBlockedBackoff());
    }
    return;
}
String functionName = item.task().functionName();
if (!enqueuer.tryAcquireSlot(functionName)) {
    queue.rotateReadyItem(item, now);     // slot perso in corsa: rimetti in fondo
    pause.accept(currentBlockedBackoff());
    return;
}
...
\end{lstlisting}

Lo scheduler non prende la testa: \emph{cerca} il primo elemento la cui funzione
ha uno slot disponibile. La ricerca \`e limitata a una finestra di 64 elementi
(\code{POLL\_READY\_MATCHING\_SCAN\_LIMIT}), il che mantiene il costo per tick
limitato indipendentemente dalla profondit\`a della coda.

Se nessun elemento nella finestra \`e servibile, la finestra viene \emph{ruotata}
--- i suoi elementi vanno in fondo --- cos\`i che il tick successivo esamini
elementi diversi. La rotazione \`e ci\`o che rende la scansione limitata
comunque equa: nessun elemento resta indefinitamente fuori dalla finestra
esaminata.

Il codice tratta anche il caso di corsa: fra il momento in cui un elemento \`e
selezionato perch\'e la sua funzione aveva uno slot e il momento in cui lo slot
viene acquisito, un altro consumatore pu\`o averlo preso. In tal caso l'elemento
viene ruotato anzich\'e scartato.

\subsection{Il backoff}

\begin{lstlisting}[language=Java]
private volatile long tickMs = 2;
private volatile long blockedBackoffMs = tickMs;

private long currentBlockedBackoff() {
    long current = blockedBackoffMs;
    blockedBackoffMs = Math.min(current * 2, 50);
    return current;
}
\end{lstlisting}

Il tick base \`e di 2\,ms. Quando lo scheduler trova lavoro ma non riesce a
dispacciarlo --- ossia quando ogni funzione in coda \`e satura --- attende con
backoff esponenziale fino a un tetto di 50\,ms. Il backoff viene azzerato appena
un dispatch riesce.

I due valori sono un compromesso esplicito fra reattivit\`a e consumo di CPU. A
2\,ms il ciclo aggiunge al pi\`u due millisecondi alla latenza di una richiesta
che attende uno slot appena liberato, il che \`e dello stesso ordine del tempo di
servizio delle funzioni di riferimento del progetto. Senza backoff, un sistema
saturo farebbe girare il thread dello scheduler a vuoto 500 volte al secondo.

\begin{damisurare}
Il contributo del tick di 2\,ms alla latenza end-to-end non \`e stato isolato
sperimentalmente. Il \cref{cap:25-valutazione-concorrenza} riporta tempi di
attesa dell'ordine di 40\,ms contro tempi di servizio di 5\,ms, in cui il
contributo del tick sarebbe minoritario ma non trascurabile.
\todomis{contributo del periodo di tick alla latenza a bassa concorrenza}
\end{damisurare}

\section{Timeout in coda}

Un elemento pu\`o scadere mentre attende. La scadenza \`e governata da
\code{sync-queue.max-queue-wait} (default 2\,s) e viene verificata a ogni
ispezione della coda: \code{peekReady}, \code{findReadyMatching} e le rotazioni
rimuovono gli elementi scaduti e li completano con esito \code{QUEUE\_TIMEOUT}.

Il coordinatore reattivo traduce quell'esito in un rifiuto con ragione
\code{TIMEOUT}:

\begin{lstlisting}[language=Java]
if (result.error() != null && "QUEUE_TIMEOUT".equals(result.error().code())) {
    throw new SyncQueueRejectedException(SyncQueueRejectReason.TIMEOUT,
                                         syncQueueGateway.retryAfterSeconds());
}
\end{lstlisting}

Il chiamante riceve quindi \code{429} con \code{X-Queue-Reject-Reason: timeout},
distinguibile dai due rifiuti di ammissione. La distinzione conta: un rifiuto per
\code{DEPTH} o \code{EST\_WAIT} avviene \emph{prima} che la richiesta occupi la
coda, un \code{TIMEOUT} avviene dopo che l'ha occupata per l'intera durata
consentita. Solo i primi due sono trattati come pressione dal modulo di offload.

\section{Rimozione di una funzione}

Quando una funzione viene cancellata mentre ha richieste in coda, il modulo
esegue tre azioni: registra il nome in un insieme di funzioni rimosse, drena
dalla coda tutti gli elementi che le appartengono completandoli con errore
\code{FUNCTION\_REMOVED}, e cancella lo stato dell'estimatore e delle metriche.

L'insieme delle funzioni rimosse \`e consultato anche in ingresso, due volte: una
prima del lock sulla coda e una dentro. La doppia verifica gestisce la corsa fra
una registrazione in corso e una cancellazione, e ha il costo di una lettura su
un insieme concorrente, trascurabile rispetto alla sezione critica che protegge.

\section{Metriche esportate}

\begin{center}
\small
\begin{adjustbox}{max width=\textwidth}
\begin{tabular}{@{}lll@{}}
\toprule
\textbf{Metrica} & \textbf{Tipo} & \textbf{Significato} \\
\midrule
\code{sync\_queue\_depth} & gauge & Profondit\`a, globale (\code{function=""}) e per funzione. \\
\code{sync\_queue\_wait\_seconds} & gauge & Attesa stimata, globale e per funzione. \\
\code{sync\_queue\_admitted\_total} & counter & Ammissioni. \\
\code{sync\_queue\_rejected\_total} & counter & Rifiuti di ammissione. \\
\code{sync\_queue\_timedout\_total} & counter & Scadenze in coda. \\
\bottomrule
\end{tabular}
\end{adjustbox}
\end{center}

La documentazione operativa del progetto fornisce una regola di lettura che vale
la pena riportare perch\'e sintetizza il rapporto fra le due grandezze:
\emph{``un valore alto di \code{sync\_queue\_rejected\_total\{function\}} con
profondit\`a bassa indica un rifiuto per attesa stimata, non un esaurimento della
capacit\`a''}. \`E il modo in cui un operatore distingue una piattaforma che ha
finito la coda da una che sta proteggendo il proprio SLO.

\section{Configurazione}

\begin{center}
\small
\begin{adjustbox}{max width=\textwidth}
\begin{tabular}{@{}llp{6.4cm}@{}}
\toprule
\textbf{Chiave} & \textbf{Default} & \textbf{Effetto} \\
\midrule
\code{enabled} & \code{true} & Attiva il modulo. \\
\code{admission-enabled} & \code{true} & Attiva il rifiuto per \code{EST\_WAIT}
(quello per \code{DEPTH} resta). \\
\code{max-depth} & 200 & Capacit\`a della coda. \\
\code{max-estimated-wait} & 2\,s & Soglia di rifiuto per attesa stimata. \\
\code{max-queue-wait} & 2\,s & Scadenza in coda. \\
\code{retry-after-seconds} & 2 & Valore dell'header \code{Retry-After}. \\
\code{throughput-window} & 30\,s & Finestra della stima di throughput. \\
\code{per-function-min-samples} & 50 & Campioni per fidarsi della stima per
funzione. \\
\bottomrule
\end{tabular}
\end{adjustbox}
\end{center}

Tutti questi parametri, tranne \code{max-depth} e i due che governano la stima,
sono modificabili a caldo attraverso il modulo \code{runtime-config}
(\cref{cap:14-runtime-config}): la lettura avviene infatti attraverso
\code{SyncQueueConfigSource} anzich\'e direttamente dalle propriet\`a. \`E il
motivo per cui \code{evaluate()} interroga il config source a ogni valutazione
anzich\'e leggere un campo finale.

\chapter{Il modulo \code{autoscaler}}
\label{cap:11-autoscaler}

Il modulo \code{autoscaler} decide \emph{quante repliche} deve avere una
funzione gestita. \`E la controparte, e la controfigura, del modulo trattato nel
capitolo successivo: entrambi regolano la capacit\`a offerta a una funzione, ma
agiscono su leve diverse e --- questo \`e il punto architetturale di \nf{} ---
sono sostituibili indipendentemente.

Sono circa 850 righe in undici classi.

\section{Il posto dell'autoscaler nel sistema}

Il campo \code{ScalingStrategy} nella specifica di una funzione ammette tre
valori, e definisce chi possiede la decisione:

\begin{center}
\small
\begin{adjustbox}{max width=\textwidth}
\begin{tabular}{@{}lp{9.4cm}@{}}
\toprule
\textbf{Strategia} & \textbf{Chi decide} \\
\midrule
\code{INTERNAL} & Questo modulo. Il control plane campiona le proprie metriche e
comanda il numero di repliche al provider di deployment. \\
\code{HPA} & Kubernetes. Il control plane crea le risorse e si astiene; la
decisione \`e dell'Horizontal Pod Autoscaler~\cite{k8shpa}. \\
\code{NONE} & Nessuno. Il numero di repliche resta quello impostato
manualmente. \\
\bottomrule
\end{tabular}
\end{adjustbox}
\end{center}

Il default per una funzione in modalit\`a \code{DEPLOYMENT} \`e \code{INTERNAL}
(\cref{cap:06-nucleo}). La strategia \code{NONE} non \`e una lacuna ma uno
strumento sperimentale: consente di fissare il numero di repliche mentre si
studia il comportamento del controllore di concorrenza, isolando le due
dinamiche che altrimenti interagirebbero.

\section{La decisione}

\subsection{La formula}

\code{ScalingDecisionCalculator} implementa la stessa formula dell'HPA di
Kubernetes: il rapporto fra valore osservato e valore obiettivo, moltiplicato
per il numero corrente di repliche.

\[
  R_{\text{racc}} = \left\lceil \max_{m \in M}
    \frac{v_m}{t_m} \cdot R_{\text{corr}} \right\rceil,
  \qquad
  R_{\text{des}} = \operatorname{clamp}(R_{\text{racc}}, R_{\min}, R_{\max})
\]

dove $M$ \`e l'insieme delle metriche configurate, $v_m$ il valore osservato e
$t_m$ l'obiettivo. Con pi\`u metriche vince la pi\`u esigente: il massimo dei
rapporti, non la media. \`E la scelta corretta per un insieme di vincoli che
devono valere tutti. Un obiettivo assente o non numerico ricade su 50.

Un dettaglio del codice attuale cambia il significato di $R_{\text{corr}}$: il
moltiplicatore \`e il numero di repliche \emph{pronte} (\emph{serving}), non quello
richiesto al deployment. Il rapporto misura quanto il carico eccede la capacit\`a
che sta effettivamente servendo; moltiplicarlo per repliche non ancora pronte
sovrastimerebbe la capacit\`a e frenerebbe la risposta a un picco.

\subsection{Una separazione che sembra ridondante}

\begin{lstlisting}[language=Java]
// Kept apart on purpose: the clamp used to overwrite the recommendation in
// place, so a decision pinned at maxReplicas was indistinguishable from a
// metric that had stopped rising.
int recommendedReplicas = (int) Math.ceil(maxRatio * normalizedCurrentReplicas);
int desiredReplicas =
        Math.clamp(recommendedReplicas, scaling.minReplicas(), scaling.maxReplicas());
\end{lstlisting}

Il valore raccomandato e quello desiderato sono conservati entrambi nel record
di decisione, anche se solo il secondo viene applicato. Il commento del record
lo dice esplicitamente: \`e una distinzione che l'HPA di Kubernetes non espone,
perch\'e pubblica soltanto il valore limitato. La ragione \`e
diagnostica: una funzione ferma a dieci repliche pu\`o esserlo perch\'e la
domanda si \`e stabilizzata oppure perch\'e la domanda continua a crescere ma il
tetto \`e stato raggiunto. Sono due situazioni che richiedono azioni opposte
dall'operatore, e sovrascrivendo la raccomandazione con il valore limitato
diventano indistinguibili.

\`E un esempio del principio ricorrente in questo progetto: \emph{conservare la
distinzione fra ci\`o che il sistema voleva e ci\`o che ha potuto}. Lo stesso
principio governa i due limiti di concorrenza del \cref{cap:09-async-queue} e i
due campi \code{requested}/\code{effective} della modalit\`a di esecuzione nel
\cref{cap:07-modalita-esecuzione}.

\section{Le tre metriche}

\subsection{\code{queue\_depth} e \code{in\_flight}}

Entrambe sono lette da una \code{WorkloadMetricsSource}, interfaccia del progetto
\code{workload-metrics} con quattro letture (profondit\`a, in volo, concorrenza
effettiva, backlog smaltibile). L'implementazione non appartiene a un modulo di
coda: \`e \code{EngineWorkloadMetricsSource}, nel control plane, sopra il motore di
scheduling composto (\cref{cap:09-async-queue}). La profondit\`a \`e il numero di
prenotazioni che la funzione occupa nel motore, mantenuto incrementalmente e
letto in tempo costante, in modo che uno scrape non scandisca mai il backlog; il
commento nel sorgente dichiara anche che differisce di al pi\`u $+1$, globalmente,
dalla definizione ``accodato e non ancora reclamato'' dei vecchi scheduler, e che
n\'e il rapporto dell'autoscaler n\'e la soglia del governatore possono essere spostati
da una lettura di $\pm 1$. Il valore in volo \`e quello dei lease di capacit\`a
della funzione. Sono metriche di \emph{stato}: descrivono la situazione
istantanea del sistema.

\subsection{\code{rps}: un caso a parte}

La terza metrica non viene dalla sorgente ma dalle osservazioni sulle
invocazioni (\code{InvocationObservations}), un'interfaccia dell'SPI implementata
da \code{Metrics}, che restituisce uno snapshot con il totale cumulativo dei
dispatch \emph{e la generazione della funzione} a cui appartiene. La lettura la
trasforma in derivata:

\begin{lstlisting}[language=Java]
var sample = observations.snapshot(functionName);
if (sample.generation() == null) { lastDispatchSamples.remove(functionName); return 0; }
CounterSample current = new CounterSample(sample.generation(), sample.dispatched(), now);
CounterSample previous = lastDispatchSamples.put(functionName, current);
if (previous == null || !previous.generation().equals(current.generation())) { return 0.0; }
double deltaCount = current.count() - previous.count();
long deltaMs = current.epochMs() - previous.epochMs();
if (deltaCount <= 0.0 || deltaMs <= 0L) { return 0.0; }
return deltaCount / (deltaMs / 1000.0);
\end{lstlisting}

Il totale \`e cumulativo, e usarne il valore grezzo come carico sarebbe un errore
--- un sistema attivo da un'ora avrebbe un ``carico'' enorme e crescente. La
lettura ne calcola quindi la differenza fra campioni successivi divisa per il
tempo trascorso. Restituisce zero al primo campione, perch\'e non esiste un
precedente, e \emph{al cambio di generazione}: una funzione cancellata e
ri-registrata riparte con contatori nuovi, e sottrarre il totale della vecchia
generazione da quello della nuova produrrebbe una derivata priva di senso
(negativa, quindi azzerata, o spuria). La versione precedente leggeva un
\code{Counter} Micrometer per nome e non aveva questo confine.

La documentazione operativa del progetto lo enuncia esplicitamente: \emph{``The
raw cumulative counter value is not used directly as load.''} \`E un'avvertenza
che vale la pena avere scritta, perch\'e l'errore \`e comune e le sue
conseguenze --- scale-up monotono e irreversibile --- sono spettacolari.

\subsection{Che cosa succede senza un modulo di coda}

\begin{lstlisting}[language=Java]
@ConditionalOnBean({WorkloadMetricsSource.class, FunctionCatalogView.class})
public class AutoscalerConfiguration { ... }
\end{lstlisting}

L'autoscaler \`e condizionato all'esistenza di una \code{WorkloadMetricsSource}, e
questa esiste solo se il motore di scheduling esiste, cio\`e se almeno un modulo di
coda fornisce una strategia. Senza, il modulo \emph{non viene creato}, e non c'\`e pi\`u
la degradazione parziale (``\code{rps} funziona, il resto no'') che le versioni
precedenti segnalavano con un avviso emesso una volta. Il commento nel sorgente
dichiara la lezione di un incidente sullo stesso meccanismo: la condizione
elencava anche il \code{MeterRegistry}, valutata prima che la sua definizione
esistesse, falliva e portava con s\'e l'intero modulo \emph{in silenzio} --- nessun
bean, nessun \code{InternalScaler}, nessuna riga di log --- e una campagna ha
proseguito con 340 richieste al secondo contro una soglia di 100 restando su una
replica. Il registro \`e oggi iniettato nei bean e non figura nella condizione,
cos\`i che una sua assenza sarebbe rumorosa. Resta per\`o un caso silenzioso, che
il commento di \code{EngineWorkloadMetricsSource} nomina: senza la sorgente,
autoscaler e governatore di concorrenza sono entrambi disattivati senza errore.
Chi vuole l'autoscaler deve quindi selezionare almeno un modulo di coda.

\section{Isteresi: i cooldown}

Un controllore reattivo che agisce a ogni campione oscilla. \`E il fenomeno che
in teoria del controllo si chiama \emph{hunting}: la correzione modifica la
grandezza osservata, che genera la correzione opposta.

Il modulo lo previene con due cooldown asimmetrici:

\begin{lstlisting}[language=Java]
private static final long SCALE_UP_COOLDOWN_MS = 30_000;
private static final long SCALE_DOWN_COOLDOWN_MS = 60_000;
\end{lstlisting}

L'asimmetria \`e il punto. Scalare in su \`e un'operazione a basso rischio: si
paga capacit\`a inutilizzata per qualche minuto. Scalare in gi\`u \`e ad alto
rischio: se la domanda ritorna, le richieste pagano un cold start. Quindi il
sistema sale rapidamente e scende con riluttanza --- lo stesso principio che
governa le finestre di stabilizzazione dell'HPA di Kubernetes, dove il default di
scale-down \`e di cinque minuti.

I valori 30\,s e 60\,s sono costanti nel codice, non configurabili
(\code{ScalingCooldownTracker}); sono configurabili solo l'intervallo di
polling e i limiti di default di repliche (\code{nanofaas.scaling.*}). \`E una
scelta di semplicit\`a che il progetto potrebbe voler rivedere: sono i due
parametri che pi\`u probabilmente un esperimento vorrebbe variare.

\begin{damisurare}
L'effetto dei due cooldown sulla stabilit\`a del numero di repliche e sul costo
in capacit\`a inutilizzata non \`e stato caratterizzato. Le campagne del
\cref{cap:23-prestazioni-base} riportano il numero di repliche di picco per
release ma non la traiettoria temporale.
\todomis{traiettoria del numero di repliche sotto carico a gradino, al variare
dei cooldown}
\end{damisurare}

\section{Rollout in stallo e protezione dal risveglio}

Il ciclo confronta la raccomandazione non con le repliche pronte ma con quelle
\emph{gi\`a richieste} al deployment, e lo fa leggendo richieste e pronte dallo
\emph{stesso} snapshot, perch\'e una decisione non sia presa su un target che nel
frattempo si \`e mosso. Ne segue un problema che il codice risolve con un
componente proprio. Se la raccomandazione scende sotto il richiesto mentre il
rollout sta ancora raggiungendo il richiesto, tornare indietro subito
annullerebbe uno scale-up che sta funzionando (il commento rimanda a un test di
regressione, \code{InternalScalerDesiredVsReadyRegressionTest}); ma non pu\`o
essere una protezione senza fine, perch\'e se alcune repliche non diventano mai
pronte un vero scale-down resterebbe bloccato per sempre su un target fantasma.
\code{ScalingProgressTracker} decide quando ``ancora in salita'' diventa
``bloccato'': il rollout \`e in stallo se il richiesto non \`e raggiunto e il numero di
pronte non \`e avanzato per una finestra intera di 60\,s.

Lo scale-down segue quindi due strade: il \emph{segnale di scale-down} esplicito
(le repliche pronte superano gi\`a ci\`o che il carico richiede, e non attende il
completamento del rollout), oppure lo stallo. In entrambi i casi passa per
\code{scaleDownIfUnprotected} del coordinatore di risveglio, che rifiuta se la
generazione ha un lease di risveglio attivo (\cref{sec:06-riprova}): non si toglie
capacit\`a a una funzione che una richiesta sta aspettando. La lettura delle
repliche \`e un'osservazione non bloccante con tre esiti (fresca, vecchia,
assente); senza alcuna lettura il ciclo salta la funzione, perch\'e trattarla come
zero repliche scalerebbe un deployment sano sulla base di una GET fallita.

\section{Inferenza del cold start}

\code{ColdStartTracker} risolve un problema di osservabilit\`a: un runtime
vecchio, o scritto in un linguaggio il cui SDK non implementa la convenzione,
non emette l'header \code{X-Cold-Start}, e il control plane non saprebbe che
un'invocazione ha pagato un'inizializzazione.

Il tracker ricorda gli eventi di scale-up recenti e ne inferisce il cold start:

\begin{lstlisting}[language=Java]
/**
 * Returns true if the function recently scaled up from 0 replicas
 * (within the expiry window), indicating a likely cold start.
 */
public boolean isPotentialColdStart(String functionName) {
    ScaleUpEvent event = recentScaleUps.get(functionName);
    if (event == null) { return false; }
    long ageMs = Instant.now().toEpochMilli() - event.timestamp().toEpochMilli();
    if (ageMs > EXPIRY_MS) { recentScaleUps.remove(functionName); return false; }
    return event.fromReplicas() == 0;
}
\end{lstlisting}

\`E un'euristica, e il codice non lo nasconde --- il metodo si chiama
\code{isPotentialColdStart}, non \code{isColdStart}. La finestra di 30 secondi \`e
generosa, e il criterio \code{fromReplicas() == 0} \`e conservativo: uno scale-up
da una a due repliche non \`e considerato, anche se la seconda replica sar\`a
fredda.

L'euristica pu\`o sbagliare in entrambe le direzioni: attribuisce cold start a
richieste servite da una replica gi\`a calda entro trenta secondi da un
risveglio, e non lo attribuisce a richieste servite da una replica nuova
aggiunta a un insieme non vuoto. \`E accettabile per il ruolo che ha --- un
ripiego per runtime che non collaborano --- e diventa problematica se qualcuno
ne usasse i risultati come dato sperimentale.

\begin{damisurare}
Il tasso di errore dell'inferenza rispetto all'header \code{X-Cold-Start}
riportato dai runtime che lo implementano \`e misurabile confrontando le due
sorgenti sulla stessa esecuzione, e non \`e stato misurato.
\todomis{accordo fra cold start inferito e cold start dichiarato}
\end{damisurare}

\section{Il ciclo di scaling}

\code{InternalScaler} \`e un \code{SmartLifecycle} con un esecutore schedulato a
periodo fisso (default 5\,s, \code{nanofaas.scaling.poll-interval-ms}). A ogni tick
scandisce le funzioni registrate, e per ciascuna con strategia \code{INTERNAL},
deployment gestito e generazione corrente calcola la decisione, verifica il
cooldown pertinente, e comanda il nuovo numero di repliche al coordinatore di
deployment. Un'eccezione su una funzione non ferma le altre.

Il modulo si disattiva completamente in assenza di un coordinatore di deployment:

\begin{lstlisting}[language=Java]
if (deploymentCoordinator == null) {
    log.info("InternalScaler disabled: no ManagedReplicaControl available");
    return;
}
\end{lstlisting}

\`E la degradazione corretta: senza un backend di deployment non esistono
repliche da scalare, e un ciclo che gira senza poter agire sarebbe puro consumo.
Anche qui la condizione \`e dichiarata nel log all'avvio anzich\'e lasciata
dedurre.

Il periodo di 5\,s \`e sostanzialmente pi\`u lungo dei tempi dell'ammissione, e la
differenza \`e istruttiva: l'ammissione deve reagire in millisecondi perch\'e
governa singole richieste, lo scaling reagisce in secondi perch\'e la sua azione
--- creare un pod --- richiede secondi. Un autoscaler che campionasse a pochi
millisecondi produrrebbe decisioni su uno stato che non ha ancora recepito
l'azione precedente.

\section{Metriche di decisione}

\code{ScalingDecisionMetrics} esporta la decisione stessa, in quattro gauge per
funzione: repliche raccomandate, desiderate, un indicatore di limitazione
(\code{function\_scaling\_limited}) e il rapporto massimo moltiplicato per 1\,000
(\code{function\_scaling\_ratio\_milli}, intero perch\'e un rapporto di 3,65 non
sia riportato come 3 o 4). \code{TargetLoadMetrics} esporta invece gli obiettivi
configurati (\code{gateway\_service\_target\_load}) in una forma compatibile con
OpenFaaS. Rende
possibile ricostruire \emph{a posteriori} perch\'e il sistema abbia scelto un
certo numero di repliche, senza dover correlare i log con le metriche di stato.

La documentazione operativa fornisce la lettura combinata: \emph{``Internal
autoscaling for \code{rps} uses the delta between successive
\code{function\_dispatch\_total} samples divided by elapsed sample time''}, ed \`e
la stessa quantit\`a che l'operatore ricalcolerebbe da PromQL. Esportare la
grandezza su cui il sistema ha effettivamente deciso, anzich\'e lasciarla
ricostruire, evita la classe di errori in cui l'analisi post-mortem calcola
qualcosa di leggermente diverso da ci\`o che il controllore aveva visto.

\section{Rapporto con il controllo della concorrenza}

I due moduli agiscono su leve diverse:

\begin{center}
\small
\begin{adjustbox}{max width=\textwidth}
\begin{tabular}{@{}lll@{}}
\toprule
 & \textbf{\code{autoscaler}} & \textbf{\code{concurrency-control}} \\
\midrule
Leva & numero di repliche & limite di concorrenza per funzione \\
Latenza dell'azione & secondi (creazione pod) & immediata (variabile in memoria) \\
Periodo di decisione & 5\,s & \multicolumn{1}{l}{\`e nel \cref{cap:12-concurrency-control}} \\
Costo dell'errore & capacit\`a sprecata o cold start & attesa o sotto-utilizzazione \\
\bottomrule
\end{tabular}
\end{adjustbox}
\end{center}

Il modo \code{STATIC\_PER\_POD} del controllore di concorrenza \`e l'unico punto
in cui i due si toccano: calcola il limite come \emph{repliche pronte $\times$
obiettivo per replica}, e quindi il suo risultato si muove quando l'autoscaler
agisce. \`E un accoppiamento voluto e unidirezionale --- il controllore di
concorrenza legge lo stato delle repliche, l'autoscaler non legge i limiti di
concorrenza.

Che i due possano essere selezionati indipendentemente \`e il fatto
architettonico che rende possibili gli esperimenti del
\cref{cap:25-valutazione-concorrenza}: fissando la strategia di scaling a
\code{NONE} si ottiene un sistema in cui l'unica variabile \`e il limite di
concorrenza, e le due dinamiche non si confondono.

\chapter{Il modulo \code{concurrency-control}}
\label{cap:12-concurrency-control}

Questo \`e il modulo su cui il progetto ha investito di pi\`u, ed \`e il capitolo
con il maggior contenuto originale del documento. Il modulo decide, per ciascuna
funzione e a intervalli regolari, quante invocazioni possono essere
contemporaneamente in volo. \`E 1\,314 righe in quindici classi --- pi\`u di un
quarto del codice di tutti i moduli messi insieme --- e implementa cinque
politiche, tre delle quali adattive.

Il capitolo procede cos\`i: prima il problema e la sua struttura analitica, poi
le cinque politiche in ordine di sofisticazione crescente, infine il ciclo che
le esegue. I risultati sperimentali che hanno motivato le scelte sono riassunti
qui dove servono a giustificare una decisione di progetto, e riportati per esteso
nel \cref{cap:25-valutazione-concorrenza}.

\section{Il problema}

\subsection{Che cosa fa un limite di concorrenza}

In \nf{}, il limite di concorrenza di una funzione \`e il valore contro cui
\code{FunctionQueueState.tryAcquireSlot()} confronta il contatore delle richieste
in volo (\cref{cap:09-async-queue}). Una richiesta che non ottiene lo slot resta
in coda.

Ne segue che il limite governa due grandezze contemporaneamente, e questa
duplicit\`a \`e la fonte di tutte le difficolt\`a che seguono.

\paragraph{Primo: l'utilizzazione.} Un limite troppo basso lascia le repliche
inattive mentre le richieste attendono. Un limite troppo alto le sovraccarica,
allungando il tempo di servizio di ciascuna richiesta.

\paragraph{Secondo: l'ammissione.} Una funzione pu\`o trattenere al pi\`u
$C + Q$ richieste, dove $C$ \`e il limite e $Q$ la dimensione della coda; oltre
quella soglia rifiuta. Il limite \`e quindi anche, insieme alla coda, la
finestra di ammissione.

La seconda propriet\`a \`e stata scoperta sperimentalmente, non progettata, ed \`e
costata un fallimento --- descritto nella \cref{sec:sojourn} --- in cui un
controllore ottimizzava correttamente la prima dimenticando la seconda, e
rifiutava il 4\% del traffico offerto.

\subsection{La struttura analitica}

Sia $\lambda$ il tasso di arrivo, $S(C)$ il tempo di servizio medio con limite
$C$, $W(C)$ il tempo di attesa in coda e $T(C) = W(C) + S(C)$ il tempo di
sojourn --- ci\`o che il chiamante osserva.

Dalla legge di Little~\cite{little1961proof} applicata al solo servizio, con il
sistema saturo, il throughput \`e $\mu(C) = C / S(C)$. Poich\'e le repliche
condividono CPU e memoria, $S$ cresce con $C$: pi\`u richieste contemporanee
significa pi\`u contesa, quindi

\[
  \frac{\mathrm{d}S}{\mathrm{d}C} > 0 \qquad \text{per ogni } C > 0 .
\]

Questa disuguaglianza ha una conseguenza immediata e distruttiva per una classe
intera di controllori:

\begin{quote}
\textbf{Un controllore che minimizza $S(C)$ non ha ottimo interno.} Poich\'e $S$
\`e monotona crescente, il gradiente punta in discesa ovunque, e una regola che
riduce il limite quando la latenza cresce cammina fino al proprio limite
inferiore.
\end{quote}

Il tempo di attesa ha il comportamento opposto: la coda drena pi\`u in fretta
quando pi\`u richieste sono servite in parallelo, quindi $W$ decresce con $C$. La
somma ha un minimo interno. La \cref{fig:sojourn} illustra la relazione.

\begin{figure}[htbp]
  \centering
  \resizebox{0.85\textwidth}{!}{% Schematic: why service time has no interior optimum and sojourn time does.
% Illustrative shapes, NOT measured data — the caption must say so.
\documentclass[border=8pt]{standalone}
\usepackage{nanofaas-figs}
\usepackage{pgfplots}
\pgfplotsset{compat=1.18}
\begin{document}
\begin{tikzpicture}
\begin{axis}[
    width=11cm, height=7cm,
    xlabel={limite di concorrenza $C$},
    ylabel={tempo (unit\`a arbitrarie)},
    xmin=1, xmax=10, ymin=0, ymax=13,
    domain=1:10, samples=200,
    axis lines=left, font=\sffamily\small,
    legend style={draw=nfslate!50, font=\sffamily\footnotesize, at={(0.98,0.98)}, anchor=north east},
    tick label style={font=\sffamily\footnotesize},
    grid=major, grid style={nfslate!20},
]
  % service time: monotonically increasing with concurrency
  \addplot[nfblue, line width=1.1pt] {2 + 0.55*x};
  \addlegendentry{$S(C)$ --- tempo di servizio}
  % wait: decreasing, ~ backlog / drain rate
  \addplot[nfamber, line width=1.1pt] {18/x};
  \addlegendentry{$W(C)$ --- tempo di attesa}
  % sojourn = sum, interior minimum
  \addplot[nfgreen, line width=1.5pt] {2 + 0.55*x + 18/x};
  \addlegendentry{$T(C)=W+S$ --- sojourn}
  % marker at the minimum: d/dC (0.55C + 18/C) = 0  ->  C = sqrt(18/0.55) ~ 5.72
  \addplot[only marks, mark=*, mark size=2.4pt, nfgreen]
      coordinates {(5.72,11.29)};
  \node[font=\sffamily\footnotesize, text=nfgreen, anchor=south]
      at (axis cs:5.72,11.5) {ginocchio};
  \draw[dashed, nfslate] (axis cs:5.72,0) -- (axis cs:5.72,11.29);
\end{axis}
\end{tikzpicture}
\end{document}
}
  \caption{Schema qualitativo (\emph{non} dati misurati). Il tempo di servizio
  cresce monotonamente con il limite di concorrenza e non ammette minimo interno;
  il tempo di attesa decresce; la loro somma --- il tempo di sojourn --- ha un
  ginocchio. Un controllore che decide da $S$ non ha nulla da trovare; uno che
  decide da $T$ s\`i.}
  \label{fig:sojourn}
\end{figure}

\subsection{Che cosa il control plane misura davvero}

La metrica \code{function\_latency\_ms} \`e misurata dal control plane fra
l'istante di dispatch e quello di completamento. \`E quindi il \emph{tempo di
servizio}, comprensivo del round trip HTTP, e non include l'attesa in coda.

La misura riportata nel \cref{cap:25-valutazione-concorrenza} quantifica lo
scarto: sotto il profilo a raffica il tempo di attesa medio \`e stato di
37--43\,ms contro un tempo di servizio di circa 5\,ms. La grandezza da cui i
primi controllori decidevano \`e circa \textbf{un settimo} di ci\`o che il
chiamante sperimenta.

La conseguenza \`e stata misurata e non solo temuta: in un confronto iniziale il
modo \code{BUDGETED} ha ridotto il tempo di servizio del 31\% su una funzione e
del 44\% sull'altra --- entrambe ampiamente dentro un SLO di 10\,ms --- e ha
\emph{peggiorato} il p95 end-to-end, perch\'e l'attesa creata dal suo stesso
limite ha pi\`u che assorbito il guadagno. Un limite di concorrenza non elimina
lavoro: lo sposta nel buffer.

\section{Le cinque politiche}

\begin{table}[htbp]
\centering
\caption{Le cinque politiche di controllo della concorrenza.}
\label{tab:modi}
\small
\begin{tabularx}{\textwidth}{@{}lXll@{}}
\toprule
\textbf{Modo} & \textbf{Regola} & \textbf{Segnale} & \textbf{Ambito} \\
\midrule
\code{FIXED} & Il limite configurato, invariato. & nessuno & per funzione \\
\code{STATIC\_PER\_POD} & repliche pronte $\times$ obiettivo per replica. &
repliche & per funzione \\
\code{ADAPTIVE\_PER\_POD} & Gradiente di latenza contro il minimo storico della
funzione. & servizio & per funzione \\
\code{BUDGETED} & Fabbisogno per l'SLO, servito da un budget di piattaforma con
equit\`a max-min pesata. & servizio, throughput & \textbf{globale} \\
\code{SOJOURN} & Modello: guardia di ammissione, legge di Little su tasso
predetto, orizzonte di drenaggio. & sojourn, servizio, coda & per funzione \\
\bottomrule
\end{tabularx}
\end{table}

Tutte le politiche sono soggette a due vincoli invalicabili, applicati a valle:
il limite effettivo non pu\`o superare quello configurato nella specifica
(\cref{cap:09-async-queue}), e non pu\`o scendere sotto il pavimento
configurato. La specifica della funzione resta il tetto e il pavimento; le
politiche si muovono all'interno.

\section{\code{FIXED} e \code{STATIC\_PER\_POD}}

\code{FIXED} \`e il default e non fa nulla: il limite \`e quello configurato. \`E
la configurazione di riferimento contro cui ogni altra viene confrontata.

\code{STATIC\_PER\_POD} \`e ventotto righe:

\begin{lstlisting}[language=Java]
int replicas = Math.max(1, readyReplicas);
int targetPerPod = control.targetInFlightPerPod() == null
        ? 1 : Math.max(1, control.targetInFlightPerPod());
long desired = (long) replicas * targetPerPod;
if (desired > configured) { return configured; }
return Math.max(1, (int) desired);
\end{lstlisting}

L'idea \`e che il limite sensato non \`e una propriet\`a della funzione ma della
capacit\`a disponibile: se ogni replica pu\`o servire bene due richieste
contemporanee e le repliche pronte sono cinque, il limite \`e dieci. \`E la
stessa nozione di \code{containerConcurrency} di Knative, con la differenza che
qui \`e il control plane a moltiplicare, avendo visibilit\`a sul numero di
repliche.

Il modo \`e utile e ha un limite noto: assume che l'obiettivo per replica sia
noto all'operatore. Determinarlo richiede un esperimento --- che \`e esattamente
ci\`o che le politiche adattive cercano di evitare.

\section{\code{ADAPTIVE\_PER\_POD}: il gradiente di latenza}

\subsection{Il segnale}

Il commento di classe enuncia la scelta del segnale con precisione:

\begin{quote}\itshape
Il segnale non \`e deliberatamente n\'e la profondit\`a di coda n\'e
l'utilizzazione. Una funzione che gira alla massima utilizzazione con una coda
vuota \`e sana, non satura, quindi nessuna delle due pu\`o localizzare il punto in
cui la concorrenza aggiuntiva inizia a costare throughput. Un tempo di servizio
crescente a lavoro costante \`e esattamente quel punto.
\end{quote}

L'osservazione \`e corretta e vale la pena isolarla, perch\'e distingue questa
famiglia di controllori da quelli basati su soglie di utilizzazione. Il segnale
\`e la \emph{degradazione} rispetto al miglior tempo di servizio osservato:

\[
  d = \frac{\bar{S}_{\text{int}} - S_{\text{base}}}{\bar{S}_{\text{int}}}
  \;\in\; [0, 1),
\]

dove $\bar{S}_{\text{int}}$ \`e la media dell'intervallo appena concluso e
$S_{\text{base}}$ il minimo storico. \`E la costruzione di TCP
Vegas~\cite{brakmo1995vegas} trasposta alle chiamate RPC, come nella libreria di
Netflix~\cite{netflixlimits}.

\subsection{La regola}

\begin{lstlisting}[language=Java]
if (degradation >= 0) {
    if (degradation >= highThreshold) {
        if (nowEpochMs - state.lastDecreaseEpochMs() >= downCooldown) {
            target = Math.max(minTarget, target - 1);
            state.lastDecreaseEpochMs(nowEpochMs);
        }
    } else if (degradation <= lowThreshold
            && nowEpochMs - state.lastIncreaseEpochMs() >= upCooldown) {
        target = Math.min(maxTarget, target + 1);
        state.lastIncreaseEpochMs(nowEpochMs);
    }
}
\end{lstlisting}

Passi unitari, con banda morta fra le due soglie (default 0{,}15 e 0{,}5) e
cooldown asimmetrici (30\,s in su, 60\,s in gi\`u). Il valore $-1$ della
degradazione codifica ``nessuna invocazione completata in questo intervallo'', e
in quel caso il limite non si muove: \emph{un intervallo vuoto non \`e evidenza di
niente}.

\subsection{La deriva della baseline}

\begin{lstlisting}[language=Java]
/**
 * ponytail: the baseline creeps up 0.2%/tick so a single lucky-fast interval cannot pin it
 * forever, while a real 2x degradation still stands out for minutes. Replace with a windowed
 * minimum if functions turn out to have long-period service-time seasonality.
 */
private static final double BASELINE_DRIFT_PER_TICK = 1.002;
\end{lstlisting}

Il minimo storico puro ha un difetto noto: un singolo intervallo fortunato ---
una finestra in cui il sistema era scarico --- fissa un riferimento che il
sistema non raggiunger\`a pi\`u, e da quel momento la degradazione appare
permanentemente alta. La deriva dello 0{,}2\% per tick fa dimenticare lentamente
i valori estremi. \`E, come il commento riconosce, la versione economica di un
minimo su finestra scorrevole.

\subsection{Perch\'e non basta}

Due difetti, entrambi misurati.

\paragraph{Non ha ottimo interno.} Il riferimento \`e il minimo osservato della
funzione stessa, e per la monotonia di $S(C)$ il gradiente punta in discesa
ovunque. Il commento di \code{SloDemandEstimator} lo registra: \emph{``a
controller that chases its own minimum has no interior optimum to find [\ldots]
the limit walks to its floor, which is what two runtimes were measured doing''}.

\paragraph{Decide su una risorsa condivisa senza vedere i vicini.} Con pi\`u
funzioni sullo stesso control plane, il limite di una si muove perch\'e un'altra
\`e arrivata, mentre il proprio carico non \`e cambiato. \`E il fenomeno misurato
nella campagna di co-tenancy (\cref{cap:25-valutazione-concorrenza}): la funzione
stabile perde il 28\% di throughput e guadagna il 36\% di latenza quando la
vicina compare. Il controllore interpreta il degrado come proprio e reagisce, ma
la causa \`e altrove.

I due difetti motivano rispettivamente i due modi successivi.

\section{\code{BUDGETED}: la concorrenza come risorsa di piattaforma}

\subsection{L'idea}

Il commento di classe la enuncia:

\begin{quote}\itshape
Ogni funzione dichiara ci\`o di cui ha bisogno, e una sola allocazione risponde a
tutte. [\ldots] Una decisione per funzione non pu\`o rispettare un budget
condiviso, perch\'e nessuna funzione sa che cosa stanno chiedendo le altre.
\end{quote}

La domanda viene quindi spezzata in due, e le due parti sono risolte da due
componenti distinti:

\begin{enumerate}
  \item \emph{Quanto serve a questa funzione} per rispettare il proprio SLO al
  carico che sta vedendo? Risponde \code{SloDemandEstimator}, per funzione.
  \item \emph{Quanto pu\`o dare la piattaforma}, dato che il budget \`e finito e
  altre funzioni chiedono? Risponde \code{ConcurrencyBudgetAllocator}, una volta
  per tick su tutte le richieste.
\end{enumerate}

Per costruzione la somma delle concessioni non pu\`o eccedere il budget, e le
funzioni non devono scoprirsi a vicenda attraverso l'interferenza.

\subsection{La stima del fabbisogno}

La regola base \`e il gradiente della libreria Netflix, con l'SLO come obiettivo
anzich\'e il minimo auto-misurato:

\[
  g = \operatorname{clamp}\!\left(\frac{L_{\text{target}}}{L_{\text{oss}}},\,
      0{,}5,\, 1\right),
  \qquad
  C_{\text{des}} = C \cdot g + \sqrt{C}.
\]

Tre elementi meritano commento.

\paragraph{L'obiettivo \`e un SLO, non un minimo osservato.} \`E la correzione
diretta del primo difetto di \code{ADAPTIVE\_PER\_POD}. Contro un obiettivo
fisso il gradiente \`e maggiore di 1 finch\'e la funzione \`e dentro il proprio
SLO, e minore di 1 solo quando la promessa \`e effettivamente violata. Esiste
quindi un punto di equilibrio interno.

\paragraph{Il pavimento a 0{,}5 sul gradiente.}
\begin{lstlisting}[language=Java]
/**
 * Floor on the gradient. Without it a single slow interval — a cold start, a GC pause — would
 * halve the limit repeatedly; with it the worst case is halving per tick, which recovers.
 */
private static final double MIN_GRADIENT = 0.5;
\end{lstlisting}
Limita il danno di un intervallo anomalo. Il caso peggiore diventa un
dimezzamento per tick, da cui il sistema si riprende, anzich\'e un crollo.

\paragraph{Il termine $\sqrt{C}$.} \`E margine deliberato. Un limite
dimensionato esattamente sul tasso di arrivo corrente non lascia nulla per una
raffica: il primo arrivo sopra la media si accoda e, se la coda \`e piena, viene
rifiutato. Per la legge di Little la concorrenza necessaria ad assorbire
$\lambda$ a tempo di servizio $S$ \`e $\lambda S$; il margine \`e ci\`o che sta
sopra, e cresce in modo sublineare cos\`i che un limite grande non porti uno
spreco proporzionalmente grande.

\subsection{Il tetto al ginocchio}

Il gradiente da solo non basta, e la ragione \`e misurata:

\begin{quote}\itshape
Un controllore guidato dall'SLO spende ogni millisecondo di latenza che gli viene
concesso: detto che pu\`o impiegare 10\,ms, \`e cresciuto fino a quattro volte la
concorrenza del modo a gradiente sullo stesso carico e ha consegnato il 6\% di
richieste in meno, perch\'e nulla nella regola chiedeva se la concorrenza
aggiuntiva comprasse dei completamenti.
\end{quote}

La correzione \`e un tetto alla crescita calcolato anzich\'e cercato:

\[
  C_{\text{ginocchio}} = \lambda_{\max} \cdot S_{\min}
  \;+\; \sqrt{\lambda_{\max} \cdot S_{\min}},
\]

ossia la legge di Little applicata ai due estremi che la funzione ha dimostrato
di saper raggiungere. Una volta che il throughput satura, $\lambda_{\max}$ smette
di crescere e il tetto smette di crescere con esso, qualunque cosa l'SLO
consentirebbe ancora.

Il commento identifica correttamente la parentela: \emph{``This is BBR's shape
rather than Vegas': track the two extremes separately and target their product,
instead of reading one degraded signal.''} Vegas legge un segnale degradato --- il
ritardo --- e ne deduce lo stato; qui si tracciano separatamente il miglior tempo
di servizio e il miglior throughput e se ne prende il prodotto.

Gli estremi decadono verso il presente all'1\% per tick, per la stessa ragione
per cui la baseline del modo precedente deriva: un minimo assoluto non si riprende
mai da un intervallo fortunato, e un massimo assoluto non ammette mai che la
macchina sia diventata pi\`u lenta.

Un'asimmetria importante \`e codificata esplicitamente:

\begin{lstlisting}[language=Java]
// Growth stops at the knee even when the SLO would allow more. Shrinking is never
// capped: an SLO breach is a broken promise, and no throughput argument outranks it.
if (desired > current) {
    desired = Math.min(desired, knee(spec.name(), observedLatencyMs, throughputRps));
}
\end{lstlisting}

Il tetto vincola solo la crescita. La contrazione non \`e mai limitata, perch\'e
una violazione dell'SLO \`e una promessa infranta e nessun argomento di
throughput la sopravanza. \`E una gerarchia di obiettivi dichiarata nel codice,
non implicita.

\subsection{L'allocazione: equit\`a max-min pesata}

L'allocatore implementa l'equit\`a max-min pesata~\cite{bertsekas1992datanetworks}:

\begin{enumerate}[leftmargin=*]
  \item Assegna a ciascuna funzione il proprio pavimento.
  \item Distribuisce il residuo proporzionalmente ai pesi, ma nessuna funzione
  riceve pi\`u di quanto ha chiesto.
  \item Ripete finch\'e resta budget e qualcuno \`e ancora insoddisfatto.
  \item Se le quote arrotondano a zero, distribuisce il residuo a unit\`a singole
  anzich\'e ciclare a vuoto.
\end{enumerate}

Il commento giustifica la scelta rispetto all'alternativa ovvia --- scalare ogni
richiesta dello stesso fattore --- con due propriet\`a:

\begin{quote}\itshape
Una funzione che chiede meno della propria quota non viene mai tagliata, quindi
una funzione piccola non \`e punita per aver condiviso la piattaforma con una
grande. E nulla resta sul tavolo mentre qualcuno \`e ancora a corto: lo
scalamento proporzionale lascerebbe inattiva capacit\`a che una funzione affamata
potrebbe usare.
\end{quote}

\subsection{I pavimenti prima di tutto}

Il trattamento del caso in cui i soli pavimenti eccedono il budget \`e la
decisione pi\`u interessante dell'allocatore:

\begin{quote}\itshape
Una funzione a cui viene concesso zero non pu\`o servire, e una piattaforma che
affama una funzione fino al silenzio totale per soddisfarne una vicina ha preso
una decisione di ammissione fingendo di prenderne una di scheduling. Se i soli
pavimenti eccedono il budget, il budget \`e per definizione troppo piccolo per le
funzioni registrate, e ogni pavimento viene concesso: l'alternativa \`e rompere
una promessa che la piattaforma fa a ogni tenant per tenere in ordine
l'aritmetica.
\end{quote}

Il budget viene quindi sforato. \`E la scelta corretta: un budget \`e uno
strumento di ripartizione, non un vincolo fisico, e quando diventa
insoddisfacibile il fatto da segnalare \`e che \`e mal dimensionato, non che una
funzione debba smettere di esistere.

\subsection{Il budget di piattaforma}

\begin{lstlisting}[language=Java]
/**
 * A platform capacity statement, not something to infer from load: derived from observed
 * throughput it would grow under load and shrink when idle, which is the opposite of what a
 * budget is for. Unset, it scales with the cores the control plane can see, which is a guess —
 * but a guess that at least tracks the machine rather than a constant.
 */
public int totalBudgetOrDefault() {
    if (totalBudget != null && totalBudget > 0) { return totalBudget; }
    return Math.max(8, Runtime.getRuntime().availableProcessors() * 4);
}
\end{lstlisting}

\`E lo stesso principio dell'SLO di default (\cref{cap:06-nucleo}): una grandezza
che esprime un vincolo non deve essere derivata dall'osservazione, o cesser\`a di
vincolare esattamente quando serve.

\section{\code{SOJOURN}: decidere dalla latenza del chiamante}
\label{sec:sojourn}

Il modo \code{SOJOURN} \`e il pi\`u recente e il pi\`u istruttivo, perch\'e la
sua documentazione conserva due progetti falliti insieme a quello adottato.

\subsection{Primo fallimento: il gradiente sulla latenza end-to-end}

L'idea immediata --- alimentare la regola a gradiente esistente con la latenza
end-to-end anzich\'e con il tempo di servizio --- \`e instabile, e \`e stata
respinta prima dell'implementazione. La regola riduce il limite quando la latenza
eccede l'obiettivo; un limite pi\`u piccolo drena la coda pi\`u lentamente, quindi
l'attesa cresce, quindi il limite si riduce ancora, fino al pavimento. \`E un
anello di retroazione positiva.

\subsection{Secondo fallimento: la ricerca del minimo}

Poich\'e $T(C)$ ha un minimo interno (\cref{fig:sojourn}), cercarlo per passi
successivi --- muovi, confronta con l'ultima volta, prosegui se \`e migliorato,
inverti se \`e peggiorato --- \`e ben posto in linea di principio. \`E stato
implementato e falsificato dalla misura:

\begin{center}
\small
\begin{adjustbox}{max width=\textwidth}
\begin{tabular}{@{}lrr@{}}
\toprule
 & \textbf{\code{ADAPTIVE\_PER\_POD}} & \textbf{\code{SOJOURN}, a ricerca} \\
\midrule
richieste rifiutate & 32 (0{,}004\%) & \textbf{36\,181 (4{,}275\%)} \\
limite medio [min--max] & 7{,}8 [4--8] & 4{,}0 [1--8] \\
\bottomrule
\end{tabular}
\end{adjustbox}
\end{center}

La causa \`e un problema di attribuzione:

\begin{quote}\itshape
Non pu\`o distinguere ``la mia mossa ha aiutato'' da ``il carico \`e calato''. In
un avvallamento ha attribuito il sollievo naturale alla propria decisione di
restringere, ed era ancora piccolo quando il carico \`e tornato. Confrontare due
istanti presuppone che il mondo sia rimasto fermo fra di essi.
\end{quote}

Il dettaglio pi\`u interessante \`e che il punto di lavoro scelto non era
sbagliato: a limite 4 la ricerca consegnava il 96{,}5\% del throughput con met\`a
del tempo di servizio. Ci\`o che ha mancato \`e la seconda funzione del limite ---
essere la finestra di ammissione --- e dimensionarlo dalla sola latenza ha fatto
perdere il 4\% del traffico.

\subsection{Il progetto adottato: calcolare anzich\'e confrontare}

Il controllore adottato non confronta due intervalli. Calcola, in tre parti con
un ordine di autorit\`a stretto.

\paragraph{1. La guardia di ammissione.} Ha la precedenza assoluta:

\begin{lstlisting}[language=Java]
private static final double HIGH_WATER = 0.7;

if (underQueuePressure(queueDepth, bounds)) {
    // Never below what is already in flight while the buffer is filling: lowering the limit
    // here shrinks the admission window and refuses the very traffic that is queueing.
    needed = Math.max(needed, Math.max(inFlight + 1, bounds.floor()));
}
\end{lstlisting}

Quando la coda supera il 70\% della propria capacit\`a, il limite non pu\`o
scendere, qualunque cosa l'aritmetica della latenza preferisca. Il commento \`e
esplicito sul fatto che questa non \`e una preferenza tarabile: \emph{``refusing
traffic is worse than being slow, and that ordering is not a tuning
parameter''}.

\paragraph{2. La legge di Little su un tasso predetto.}

\[
  C_{\text{servizio}} = \hat{\lambda} \cdot S,
  \qquad
  C_{\text{drenaggio}} = \frac{Q \cdot S}{H},
  \qquad
  C = \left\lceil C_{\text{servizio}} + C_{\text{drenaggio}} \right\rceil
\]

dove $\hat{\lambda}$ \`e il tasso di arrivo \emph{predetto}, $Q$ la profondit\`a
di coda corrente e $H$ l'orizzonte di drenaggio concesso.

\paragraph{3. La promessa come orizzonte.} Quando il sojourn osservato eccede
l'obiettivo, l'orizzonte di drenaggio si accorcia in proporzione:

\[
  H = \max\left(1\,\mathrm{s},\; 5\,\mathrm{s} \cdot
      \frac{T_{\text{target}}}{T_{\text{oss}}}\right).
\]

Un orizzonte pi\`u corto richiede pi\`u concorrenza per smaltire lo stesso
arretrato. La logica \`e che l'attesa \`e l'unica parte della latenza del
chiamante che un limite di concorrenza possa effettivamente accorciare.

\subsection{La predizione del carico}

La predizione $\hat{\lambda}$ \`e ci\`o che elimina il problema di attribuzione:
il controllore smette di confrontare istanti e inizia a calcolare da un modello.
\code{LoadForecast} implementa il metodo lineare di Holt con trend smorzato:

\[
  \ell_t = \alpha \lambda_t + (1-\alpha)(\ell_{t-1} + \phi b_{t-1}),
  \quad
  b_t = \beta(\ell_t - \ell_{t-1}) + (1-\beta)\phi b_{t-1},
  \quad
  \hat{\lambda}_{t+1} = \ell_t + \phi b_t,
\]

con $\alpha = 0{,}5$, $\beta = 0{,}3$, $\phi = 0{,}8$.

La scelta \`e giustificata cos\`i:

\begin{quote}\itshape
Una media mobile \`e in ritardo su una rampa, e una rampa \`e dove la coda si
costruisce: quando una stima in ritardo se ne accorge, il buffer \`e gi\`a pieno.
Il termine di trend vede la pendenza; lo smorzamento impedisce di portare quella
pendenza oltre la cima della rampa, dove un'estrapolazione non smorzata
sovrastima di pi\`u.
\end{quote}

Il primo campione produce livello senza trend: \emph{``una lettura \`e un livello
e nessun trend; una pendenza richiede due punti, e inventarne uno farebbe
predire una rampa al primo intervallo di ogni funzione''}.

\subsection{Il risultato}

Sotto arrivi aperti --- richieste programmate dall'orologio anzich\'e da un
generatore a ciclo chiuso --- il confronto con \code{ADAPTIVE\_PER\_POD} sullo
stesso carico:

\begin{center}
\small
\begin{adjustbox}{max width=\textwidth}
\begin{tabular}{@{}lrrr@{}}
\toprule
 & \textbf{\code{ADAPTIVE\_PER\_POD}} & \textbf{\code{SOJOURN}} & \\
\midrule
servite & 549\,856 & \textbf{567\,210} & $+3{,}2\%$ \\
rifiutate & 110\,117 (16{,}7\%) & \textbf{92\,785 (14{,}1\%)} & $-15{,}7\%$ \\
p95 end-to-end & 120{,}48\,ms & \textbf{107{,}68\,ms} & $-10{,}6\%$ \\
p99 end-to-end & 157{,}67\,ms & \textbf{133{,}82\,ms} & $-15{,}1\%$ \\
attesa in coda & 35{,}61\,ms & 31{,}03\,ms & $-12{,}9\%$ \\
tempo di servizio & 1{,}69\,ms & 2{,}00\,ms & $+18\%$ \\
limite medio & 3{,}9 & 2{,}8 & \\
\bottomrule
\end{tabular}
\end{adjustbox}
\end{center}

Lo scambio \`e andato nella direzione prevista: il tempo di servizio peggiora del
18\%, l'attesa migliora del 13\%, e poich\'e l'attesa \`e venti volte il tempo di
servizio il chiamante ci guadagna. Il dettaglio istruttivo \`e che \code{SOJOURN}
ottiene questo con \emph{meno} concorrenza media: non vince aggiungendo
capacit\`a ma collocandola nel tempo, che \`e ci\`o per cui la predizione era
stata introdotta.

Il \cref{cap:25-valutazione-concorrenza} riporta i caveat --- una sola esecuzione
per modo, entrambi i sistemi in sovraccarico --- e il
\cref{cap:27-discussione} li tratta come limiti alla validit\`a.

\section{Il governor: il ciclo che esegue le politiche}

\code{ConcurrencyGovernor} \`e il componente che a ogni tick (default 5\,s)
raccoglie le osservazioni, invoca la politica pertinente per ciascuna funzione, e
pubblica i limiti attraverso \code{ScalingMetricsSource}.

Ha una struttura obbligata dalla differenza di ambito fra le politiche. I modi
\code{STATIC\_PER\_POD} e \code{ADAPTIVE\_PER\_POD} decidono per funzione e sono
gestiti da \code{ConcurrencyControlCoordinator} in un ciclo. Il modo
\code{BUDGETED} decide su tutte le funzioni insieme e riceve quindi la lista
completa delle osservazioni in una sola chiamata. Il modo \code{SOJOURN} torna a
decidere per funzione, ma legge tre segnali anzich\'e uno.

\`E il motivo per cui l'interfaccia dei controllori non \`e uniforme: un modo che
alloca una risorsa condivisa non pu\`o essere ricondotto alla stessa firma di uno
che decide in isolamento senza inventare uno stato globale nascosto.

\subsection{L'intervallo come unit\`a di osservazione}

Tutti i controllori adattivi leggono i timer cumulativi di Micrometer e li
convertono in medie dell'intervallo appena concluso, attraverso
\code{LatencyIntervals}:

\begin{lstlisting}[language=Java]
/**
 * Turns a cumulative service-time timer into the mean of the interval just ended.
 *
 * <p>The cumulative value answers "since startup", which flattens exactly the change a controller
 * reacts to: a function that has served a million fast requests and is now slow still reads fast.
 * Only the difference between two readings describes the present.</p>
 */
\end{lstlisting}

\`E lo stesso errore, e la stessa correzione, gi\`a incontrati per la metrica
\code{rps} dell'autoscaler (\cref{cap:11-autoscaler}). Che il progetto abbia
dovuto risolverlo due volte in due moduli \`e un piccolo indizio a favore di una
astrazione condivisa che non esiste.

\section{Configurazione}

\begin{center}
\small
\begin{adjustbox}{max width=\textwidth}
\begin{tabular}{@{}llp{6cm}@{}}
\toprule
\textbf{Chiave} & \textbf{Default} & \textbf{Significato} \\
\midrule
\code{nanofaas.concurrency-control.poll-interval-ms} & 5\,000 & Periodo del tick. \\
\code{...default-target-in-flight-per-pod} & 2 & Obiettivo iniziale. \\
\code{...total-budget} & $\max(8, 4 \times \text{core})$ & Budget globale per
\code{BUDGETED}. \\
\midrule
\multicolumn{3}{@{}l@{}}{\emph{Per funzione, in \code{scalingConfig.concurrencyControl}:}} \\
\code{mode} & \code{FIXED} & \\
\code{targetInFlightPerPod} & 2 & \code{STATIC}/\code{ADAPTIVE}. \\
\code{minTargetInFlightPerPod} & 1 & Pavimento. \\
\code{maxTargetInFlightPerPod} & 8 & Tetto. \\
\code{upscaleCooldownMs} & 30\,000 & \code{ADAPTIVE}. \\
\code{downscaleCooldownMs} & 60\,000 & \code{ADAPTIVE}. \\
\code{highLoadThreshold} & 0{,}5 & \code{ADAPTIVE}. \\
\code{lowLoadThreshold} & 0{,}15 & \code{ADAPTIVE}. \\
\code{targetLatencyMs} & 250 & SLO: di servizio per \code{BUDGETED},
end-to-end per \code{SOJOURN}. \\
\code{weight} & 1{,}0 & Peso nell'allocazione \code{BUDGETED}. \\
\bottomrule
\end{tabular}
\end{adjustbox}
\end{center}

Va notato che \code{targetLatencyMs} cambia significato fra i due modi che lo
usano: in \code{BUDGETED} \`e un SLO sul tempo di servizio, in \code{SOJOURN}
sull'end-to-end. \`E un'ambiguit\`a reale nel modello di configurazione, mitigata
dal fatto che ciascun modo interpreta il campo coerentemente con la propria
semantica dichiarata, e non risolta.

\section{Sintesi}

\begin{table}[htbp]
\centering
\caption{Le cinque politiche: che problema risolve ciascuna, e quale lascia
aperto.}
\label{tab:modi-sintesi}
\small
\begin{tabularx}{\textwidth}{@{}lXX@{}}
\toprule
\textbf{Modo} & \textbf{Risolve} & \textbf{Lascia aperto} \\
\midrule
\code{FIXED} & Nulla; \`e il riferimento. & Tutto. \\
\code{STATIC\_PER\_POD} & Il limite segue la capacit\`a disponibile. &
L'obiettivo per replica va determinato sperimentalmente. \\
\code{ADAPTIVE\_PER\_POD} & Trova il proprio punto senza taratura. & Non ha
ottimo interno; ignora i vicini. \\
\code{BUDGETED} & La risorsa condivisa \`e allocata una volta sola, con
equit\`a. & Decide ancora dal tempo di servizio. \\
\code{SOJOURN} & Decide da ci\`o che il chiamante osserva; non confronta
istanti. & Mai osservato fuori dal sovraccarico; guardia al 70\% non tarata. \\
\bottomrule
\end{tabularx}
\end{table}

Le due colonne di destra sono la struttura del capitolo: ogni modo esiste perch\'e
il precedente ha lasciato aperto qualcosa, e ciascuno lascia aperto qualcosa a
sua volta. La progressione non \`e terminata --- l'ultima riga della tabella
elenca due questioni che il \cref{cap:28-lavori-futuri} riprende.

\chapter{Il modulo \code{offload}}
\label{cap:13-offload}

Il modulo \code{offload} consente a un'istanza di \nf{} di inoltrare
un'invocazione sincrona a un'altra istanza remota anzich\'e eseguirla localmente.
Lo scenario di riferimento \`e edge$\to$cloud: un nodo periferico con capacit\`a
limitata che, quando non riesce ad ammettere una richiesta, la delega a
un'istanza pi\`u capiente~\cite{satya2017edge, shi2016edge}.

\`E fra i moduli pi\`u piccoli che facciano qualcosa di sostanziale: 470 righe in
tre classi (configurazione, propriet\`a, gateway), di cui 411 nel gateway; una
parte consistente di queste \`e un percorso di compatibilit\`a per i consumatori
che forniscono un registro di metriche anzich\'e il contratto del nucleo. La
brevit\`a \`e conseguenza diretta di due rinunce, entrambe deliberate, che il
capitolo argomenta, e del fatto che la decisione di scaricare non vive nel
modulo: il modulo implementa \code{OffloadGateway} (un'interfaccia del contratto
del nucleo, con un'implementazione nulla di default) e fa la chiamata remota; chi
decide \`e il nucleo.

\begin{figure}[htbp]
  \centering
  \resizebox{0.95\textwidth}{!}{% nanofaas offload: an edge control plane transparently proxies sync invocations
% to a cloud control plane under pressure (or by eager per-function policy).
% Background-layer containers are declared AFTER all blur-shadow nodes (see the
% ordering note in fig-architecture.tex).
\documentclass[border=8pt]{standalone}
\usepackage{nanofaas-figs}
\begin{document}
\begin{tikzpicture}[nf, node distance=6mm]

  % ---------- client ----------
  \node[client, minimum width=20mm, minimum height=16mm] (client) {Client};

  % ---------- edge control plane (inner nodes) ----------
  \node[edgebox, right=24mm of client, minimum width=34mm] (sq)
    {sync-queue admission\\{\scriptsize DEPTH $\cdot$ EST\_WAIT}};
  \node[edgebox, below=13mm of sq, minimum width=34mm] (epods)
    {warm pods\\{\scriptsize local dispatch}};
  \node[edgebox, right=13mm of sq, minimum width=24mm] (proxy)
    {offload\\proxy};

  % ---------- cloud control plane (inner nodes) ----------
  \node[cloudbox, right=54mm of proxy, minimum width=30mm] (cpods)
    {warm pods\\{\scriptsize cloud dispatch}};

  % ---------- failure semantics (shadowed; before the containers) ----------
  \node[danger, above=8mm of proxy, minimum width=40mm, font=\sffamily\scriptsize] (fail)
    {no local fallback\\remote failure $\Rightarrow$ \texttt{502 / 504}};

  % ---------- trigger legend (shadowed; before the containers) ----------
  \node[neutral, anchor=north west, align=left, font=\sffamily\scriptsize,
        minimum height=0mm, text width=88mm] (leg)
    at ($(epods.south west)+(-2mm,-15mm)$)
    {\textbf{Offload trigger} (single hop, no fan-out):\\
     \;$\bullet$\; \texttt{EAGER} — per-function policy (\texttt{offload.enabled=true})\\
     \;$\bullet$\; \texttt{DEPTH} / \texttt{EST\_WAIT} — sync-queue backpressure rejection\\
     Counters: \texttt{nanofaas\_offload\_total\{function,trigger\}},
     \texttt{nanofaas\_offload\_failure\_total}};

  % ---------- grouping containers (background layer, AFTER all shadowed nodes) ----------
  \begin{scope}[on background layer]
    \node[container, draw=nfblue, line width=1.5pt, fill=nfblue!9, fit=(sq)(epods)(proxy)] (edge) {};
    \node[container, draw=nfgreen, line width=1.5pt, fill=nfgreen!12, fit=(cpods)] (cloud) {};
  \end{scope}
  \node[clabel, text=nfblue, anchor=south west] at ($(edge.north west)+(1mm,0.8mm)$)
    {Edge control plane};
  \node[clabel, text=nfgreen, anchor=south west] at ($(cloud.north west)+(1mm,0.8mm)$)
    {Cloud control plane};

  % ---------- request flow ----------
  \draw[flow] (client.east) -- (client.east-|edge.west)
    node[tag, midway, above] {\texttt{:invoke}};
  \draw[flow] (sq.south) -- (epods.north)
    node[tag, midway, right=0.5mm, fill=white, inner sep=1pt] {admit};
  \draw[flow] (sq.east) -- (proxy.west)
    node[tag, midway, above, align=center] {reject /\\eager};

  % the offload hop, highlighted
  \draw[offloadarr] (proxy.east) -- (proxy.east-|cloud.west)
    node[text=nfamber, font=\sffamily\scriptsize, midway, above=0.6mm, align=center]
      {transparent proxy $\cdot$ single hop\\\texttt{X-NanoFaaS-Offloaded}};

  % ---------- responses ----------
  \draw[ret] (cpods.south) to[out=-90,in=-90,looseness=0.55]
    node[tagc, pos=0.5, below] {response — transparent to the client} (client.south);

  % ---------- failure link ----------
  \draw[flowc, draw=nfred, dashed] (fail.south) -- (proxy.north);

\end{tikzpicture}
\end{document}
}
  \caption{Offload trasparente edge$\to$cloud: attivazione eager per politica di
  funzione oppure reattiva su rifiuto della coda sincrona; salto singolo, nessun
  fallback locale.}
  \label{fig:offload}
\end{figure}

\section{Quando si scarica}

L'attivazione \`e decisa in \code{ReactiveInvocationCoordinator}, nel nucleo, sul
percorso di ammissione sincrona (\cref{cap:04-architettura-control-plane}), e ha
due innesti. In entrambi i casi vale una precondizione, valutata per prima: la
richiesta non deve provenire a sua volta da un salto di offload, e il gateway deve
essere attivo.

\subsection{Eager: la politica della funzione}

\begin{lstlisting}[language=Java]
@Override
public boolean shouldOffloadEagerly(FunctionSpec spec) {
    OffloadPolicy policy = spec.offload();
    return policy != null
            && !Boolean.FALSE.equals(policy.enabled())
            && OffloadPolicy.MODE_ALWAYS.equalsIgnoreCase(policy.mode())
            && targetUrl(spec) != null;
}
\end{lstlisting}

Una funzione con \code{offload.mode = "always"} viene sempre inoltrata, senza
tentare l'esecuzione locale. \`E la dichiarazione, da parte dell'operatore, che
quella funzione non ha senso al bordo --- per esempio perch\'e richiede risorse
che il nodo periferico non ha.

\subsection{Pressure: il rifiuto come segnale}

\begin{lstlisting}[language=Java]
try {
    admitLocally(executionRecord);
} catch (SyncQueueRejectedException ex) {
    OffloadTrigger trigger = pressureTrigger(ex.reason());
    if (offloadable && trigger != null && offloadGateway.shouldOffloadOnPressure(spec)) {
        startOffload(executionRecord, spec, trigger, context, offloadedTarget);
        return;
    }
    throw ex;
}

private static OffloadTrigger pressureTrigger(SyncQueueRejectReason reason) {
    return switch (reason) {
        case DEPTH -> OffloadTrigger.DEPTH;
        case EST_WAIT -> OffloadTrigger.EST_WAIT;
        case TIMEOUT -> null;
    };
}
\end{lstlisting}

Questa \`e la parte concettualmente interessante. La decisione di scaricare
\emph{non} \`e presa da un modello di costo che confronti la latenza locale
attesa con quella remota, come nella formulazione classica del problema di
offloading. \`E presa dal \emph{rifiuto} della coda di ammissione locale.

L'eccezione \code{SyncQueueRejectedException} \`e lanciata \emph{soltanto} dal
gateway della coda sincrona. Ne segue una conseguenza pratica che il nome del
modulo non lascia intuire: l'offload reattivo funziona solo quando l'ammissione
sincrona passa per la coda sincrona, cio\`e nel profilo \code{SYNC\_QUEUE}
(\cref{sec:09-profili}). Nei profili \code{FUNCTION\_QUEUE} e \code{DIRECT} una
funzione che trova coda o slot esauriti riceve il \code{429} del nucleo e nulla
viene scaricato: resta operativo il solo offload eager. Con entrambi i moduli di
coda nell'artefatto il profilo derivato \`e \code{FUNCTION\_QUEUE}, quindi
l'offload reattivo richiede in quel caso di dichiarare esplicitamente il profilo.

Il vantaggio \`e che non richiede alcuna calibrazione preventiva: la coda
sincrona sa gi\`a stimare l'attesa e sa gi\`a decidere che \`e eccessiva
(\cref{cap:10-sync-queue}); l'offload riusa quella decisione anzich\'e
ripeterla con parametri propri. Il costo \`e che la decisione \`e binaria e
tardiva: non c'\`e modo di scaricare \emph{preventivamente} una frazione del
traffico per evitare la congestione, e non c'\`e alcuna stima della latenza
remota nella decisione.

\subsection{Perch\'e \code{TIMEOUT} non \`e pressione}

Il \code{switch} mappa esplicitamente \code{TIMEOUT} a \code{null}. Le prime due
ragioni di rifiuto avvengono \emph{prima} che la richiesta occupi la coda: la
piattaforma dichiara di non poterla accettare, e inoltrarla altrove \`e una
risposta sensata. \code{TIMEOUT} avviene \emph{dopo} che la richiesta ha
occupato la coda per l'intera durata concessa: il budget temporale \`e gi\`a
esaurito, e iniziare adesso una chiamata remota consumerebbe altro tempo per
consegnare comunque una risposta tardiva.

\section{Le due rinunce}

\subsection{Salto singolo}

Ogni richiesta inoltrata porta l'header \code{X-NanoFaaS-Offload-Hop: 1}, e il
control plane ricevente lo interpreta come divieto di ulteriore inoltro:

\begin{lstlisting}[language=Java]
boolean offloadable = !context.offloadedHop() && offloadGateway.enabled();
\end{lstlisting}

La ragione \`e la terminazione. Con $n$ istanze configurate in catena o in ciclo,
una richiesta potrebbe rimbalzare indefinitamente, e ogni rimbalzo consuma parte
del budget temporale del chiamante originale senza avvicinarsi a un'esecuzione.
Un contatore di salti con limite $k > 1$ sarebbe pi\`u generale e richiederebbe
di decidere $k$; il valore 1 \`e l'unica scelta che non richiede una decisione.

\subsection{Nessun fallback locale}

Se la chiamata remota fallisce, l'esecuzione fallisce. Il chiamante riceve
\code{502 Bad Gateway} o \code{504 Gateway Timeout}. Non viene ritentata
localmente.

Il commento nel codice del nucleo motiva la scelta sul piano meccanico:

\begin{lstlisting}[language=Java]
// Bypasses the local queue entirely: no local concurrency slots are consumed,
// so completion goes through the offload-specific path (no slot release, no retry).
\end{lstlisting}

Una richiesta inoltrata non ha mai acquisito uno slot locale. Un fallback la
farebbe rientrare nella coda locale --- la stessa che l'aveva rifiutata, e che
nel frattempo non si \`e svuotata --- riproducendo il rifiuto con il budget
temporale ormai consumato. Il fallback trasformerebbe un \code{502} rapido in un
\code{429} tardivo, che \`e strettamente peggiore.

C'\`e anche una ragione di onest\`a diagnostica: con il fallback, un'istanza
remota irraggiungibile si manifesterebbe come un aumento di latenza locale
anzich\'e come un errore, e la configurazione errata resterebbe invisibile.

\section{La semantica degli errori}

La distinzione fra ``la funzione remota ha fallito'' e ``l'inoltro ha fallito''
\`e centrale, ed \`e implementata con cura:

\begin{lstlisting}[language=Java]
// A marked response is the function's own answer, whatever its status —
// read it down the same body-parsing path as a plain 2xx, before any
// status-based branching (a marker-bearing 404 is a function decision,
// not "unregistered function").
if (functionDecided || response.statusCode().is2xxSuccessful()) {
    return response.bodyToMono(InvocationResponse.class)
            .map(this::toResult)
            .switchIfEmpty(Mono.error(new OffloadFailedException(target, false,
                    "empty response body from remote " + target)));
}
if (response.statusCode().value() == 404) {
    return response.releaseBody().then(Mono.error(new OffloadFailedException(target, false,
            "function '" + task.functionName() + "' not registered on remote " + target)));
}
\end{lstlisting}

Il caso critico \`e il \code{404}. Un \code{404} \emph{senza} marcatore significa
che la funzione non \`e registrata sull'istanza remota: \`e un errore di
configurazione della piattaforma, e va segnalato come fallimento di offload
(\code{502}), non propagato al chiamante come se la sua richiesta fosse
malformata. Un \code{404} \emph{con} marcatore \code{X-NanoFaaS-Function-Status}
\`e invece la risposta deliberata della funzione remota, e va propagato tale e
quale.

L'ordine dei rami \`e quindi essenziale: il marcatore viene consultato prima di
qualunque diramazione sul codice di stato.

\subsection{Che cosa attraversa il salto}

Il gateway inoltra la richiesta originale come envelope JSON verso
\code{/v1/functions/\{nome\}:invoke} dell'istanza remota, e vi aggiunge il
marcatore di salto e, se presenti, gli header di tracciamento
(\code{X-Trace-Id}, \code{traceparent}, \code{tracestate}). Il control plane
ricevente ricostruisce gli header applicativi della richiesta dal trasporto HTTP
del salto, quindi il gateway copia gli header applicativi del chiamante come veri
header HTTP: senza questa copia la funzione remota li perderebbe.

La copia \`e difensiva. Sono esclusi: l'inquadramento del trasporto che il client
HTTP possiede (\code{host}, \code{content-length}, \code{transfer-encoding} e
soprattutto \code{content-type}, che deve continuare a descrivere l'envelope JSON
e non il corpo originale del chiamante, e con esso la negoziazione e la codifica
del contenuto), gli header che il gateway imposta da s\'e o che il ricevente
lega a parametri dedicati (identificativi di esecuzione e di tentativo, timeout,
chiave di idempotenza, marcatore di salto, tracciamento), gli header
\emph{hop-by-hop} definiti da RFC~9110 e ogni header \code{proxy-*}. L'elenco
rispecchia di proposito quello del controller ricevente: inoltrare un nome che il
ricevente scarterebbe non serve, e un valore applicativo non deve poter fingersi
un header di controllo.

\subsection{Il margine di timeout}

\begin{lstlisting}[language=Java]
// ponytail: fixed margin so the gateway's remote budget expires before the
// coordinator's local wait, making 504 (not a local "timeout") deterministic
private static final long TIMEOUT_MARGIN_MS = 50;
...
long timeoutMs = Math.max(1, timeoutBudgetMs - TIMEOUT_MARGIN_MS);
\end{lstlisting}

Due timeout corrono in parallelo: quello locale del coordinatore, che attende il
futuro di completamento, e quello della chiamata remota. Se scadessero
contemporaneamente, l'esito dipenderebbe da una corsa e il chiamante riceverebbe
talvolta \code{504} e talvolta una risposta di timeout locale. Il margine fisso
di 50\,ms garantisce che il timeout remoto scatti per primo, e quindi che la
diagnosi sia deterministica.

Il margine ha una condizione che il codice non enuncia: al gateway viene passato
come budget il timeout della \emph{specifica} della funzione, mentre l'attesa
locale usa l'eventuale timeout indicato dal chiamante per quella richiesta. Se il
chiamante chiede un timeout inferiore a quello della specifica, l'attesa locale
scade prima e il margine non ordina pi\`u i due eventi; la determinazione del
\code{504} vale, in quel caso, solo per l'attesa di specifica.
\begin{damisurare}
\todomis{esito osservato per una richiesta offload con timeout per-richiesta inferiore a quello della specifica}
\end{damisurare}

\`E un dettaglio piccolo con una propriet\`a importante: rende \emph{riproducibile}
un comportamento in condizioni di errore, che \`e la classe di comportamenti pi\`u
difficile da testare e la pi\`u facile da rendere non deterministica per
distrazione.

\subsection{Chi possiede l'input}

Un'esecuzione scaricata non acquisisce alcun \code{DispatchLease}, ma tiene comunque
una risorsa: l'input fisico della richiesta. La chiamata remota ne \`e proprietaria
fino al proprio segnale terminale, non fino a quando la vista locale \`e ancora
annullabile. Una cancellazione amministrativa conclude soltanto la vista locale
dell'esecuzione; il publisher della chiamata remota pu\`o ignorare la cancellazione
e conservare la richiesta, quindi l'input viene chiuso solo quando la chiamata
termina, in qualunque modo. \`E la stessa regola di propriet\`a che il nucleo applica
agli altri percorsi (\cref{cap:04-architettura-control-plane}).

L'esito di un'esecuzione scaricata \`e definitivo: il completamento passa per un
percorso dedicato del nucleo che non ritenta e non rilascia slot, e un fallimento
dell'inoltro conclude l'esecuzione con il codice \code{OFFLOAD\_FAILED} o, se
scaduto il budget, \code{OFFLOAD\_TIMEOUT}. Il chiamante li vede come \code{502} e
\code{504}.

\section{Osservabilit\`a}

\begin{center}
\small
\begin{adjustbox}{max width=\textwidth}
\begin{tabular}{@{}lll@{}}
\toprule
\textbf{Segnale} & \textbf{Tipo} & \textbf{Etichette} \\
\midrule
\code{nanofaas.offload} & counter & \code{function}, \code{trigger}
(\code{eager} / \code{depth} / \code{est\_wait}) \\
\code{nanofaas.offload.failure} & counter & \code{function} \\
\code{X-NanoFaaS-Offloaded} & header & URL dell'istanza remota, anche sulle
risposte \code{502} e \code{504} \\
\bottomrule
\end{tabular}
\end{adjustbox}
\end{center}

L'header di risposta \`e ci\`o che rende l'offload \emph{trasparente ma non
invisibile}: il chiamante non deve fare nulla di diverso, ma se lo desidera pu\`o
sapere che la sua richiesta \`e stata servita altrove e da chi. \`E la stessa
distinzione fra trasparenza e opacit\`a che il \cref{cap:07-modalita-esecuzione}
applica alla degradazione da \code{DEPLOYMENT} a \code{EXTERNAL}.

I due contatori sono legati alla \emph{generazione} della funzione attiva
all'inizio dell'offload. Il tentativo \`e contato alla sottoscrizione della
chiamata e non quando la pipeline viene costruita, e un fallimento \`e contato solo
per un tentativo gi\`a contato; un tentativo appartenente a una generazione ritirata
scarta silenziosamente i propri aggiornamenti, cos\`i che una funzione cancellata
non veda risorgere la serie (invariante I7 del ciclo di vita).

L'etichetta \code{trigger} sul contatore consente di distinguere le tre cause di
attivazione, il che \`e necessario per interpretare il dato: un'alta frequenza di
\code{eager} \`e una politica che funziona, un'alta frequenza di \code{depth} \`e
un nodo periferico sotto-dimensionato.

\section{Configurazione}

\begin{center}
\small
\begin{adjustbox}{max width=\textwidth}
\begin{tabular}{@{}llp{6.6cm}@{}}
\toprule
\textbf{Chiave} & \textbf{Default} & \textbf{Effetto} \\
\midrule
\code{nanofaas.offload.enabled} & \code{true} & Interruttore generale; il gateway
resta comunque inerte senza un target. \\
\code{nanofaas.offload.target-url} & --- & URL base dell'istanza remota. \\
\code{nanofaas.offload.pressure-enabled} & \code{true} & Attiva l'offload
reattivo; quello eager per funzione resta comunque. \\
\midrule
\multicolumn{3}{@{}l@{}}{\emph{Per funzione, nel blocco \code{offload}:}} \\
\code{enabled} & --- & \code{false} esclude la funzione. \\
\code{targetUrl} & --- & Sovrascrive il target globale; la sua sola presenza
attiva l'offload per quella funzione. \\
\code{mode} & \code{pressure} & \code{always} per l'offload eager. \\
\bottomrule
\end{tabular}
\end{adjustbox}
\end{center}

Il default \code{enabled = true} in assenza di un target \`e apparentemente
incoerente ed \`e invece corretto: il modulo \`e attivo ma non ha dove inoltrare,
quindi non fa nulla. Configurare un target \`e l'unico atto necessario, e non
occorre ricordarsi anche di attivare l'interruttore. Il commento nel codice lo
enuncia: \emph{``even when true the gateway stays inert until a target URL is
configured''}.

\section{Ci\`o che il modulo non fa}

\begin{itemize}[leftmargin=*]
  \item \textbf{Non modella il costo.} La formulazione classica del problema di
  offloading valuta se il trasferimento convenga, dati il costo di rete e il
  guadagno di esecuzione. Qui la decisione \`e binaria e derivata dal rifiuto
  locale.

  \item \textbf{Non sceglie fra pi\`u destinazioni.} Il target \`e uno, globale
  o per funzione. Non esiste bilanciamento fra istanze remote.

  \item \textbf{Non scarica il percorso asincrono.} Solo \code{:invoke} \`e
  soggetto a offload. L'invocazione asincrona ha gi\`a una coda che assorbe la
  pressione, e il chiamante non sta attendendo.
\end{itemize}

\begin{damisurare}
Il beneficio dell'offload sotto pressione non \`e stato caratterizzato
sperimentalmente: non esiste una campagna che confronti latenza end-to-end e
tasso di rifiuto di un nodo periferico con e senza offload attivo, al variare
della latenza di rete verso il target. \`E l'esperimento naturale per questo
modulo e non \`e stato eseguito.
\todomis{p95 end-to-end e tasso di rifiuto con e senza offload, al variare del
RTT verso l'istanza remota}
\end{damisurare}

\chapter{Il modulo \code{runtime-config}}
\label{cap:14-runtime-config}

Il modulo \code{runtime-config} consente di modificare a caldo un sottoinsieme
dei parametri della piattaforma, senza riavviare il processo. Sono 540 righe in
dieci classi.

Il modulo ha una motivazione strettamente sperimentale. Un esperimento che
confronta due configurazioni deve, per essere valido, mantenere costante tutto il
resto: stesso processo, stessa JVM riscaldata, stesse code popolate, stesso stato
del garbage collector. Riavviare il control plane fra una configurazione e
l'altra reintroduce esattamente le variabili che l'esperimento intende
escludere.

\section{Che cosa \`e modificabile a caldo}

Sei parametri, esattamente:

\begin{lstlisting}[language=Java]
public record RuntimeConfigSnapshot(
        long revision,
        int rateMaxPerSecond,
        boolean syncQueueEnabled,
        boolean syncQueueAdmissionEnabled,
        Duration syncQueueMaxEstimatedWait,
        Duration syncQueueMaxQueueWait,
        int syncQueueRetryAfterSeconds
) { ... }
\end{lstlisting}

Il criterio di inclusione \`e la reversibilit\`a senza effetti sullo stato
esistente. Il limite di frequenza \`e un intero letto a ogni richiesta:
cambiarlo ha effetto dalla richiesta successiva e non invalida nulla. Le soglie
della coda sincrona sono lette a ogni valutazione di ammissione (\cref{cap:10-sync-queue}):
lo stesso.

Sono invece esclusi \code{max-depth} della coda sincrona --- ridurla
richiederebbe di decidere che cosa fare degli elementi in eccesso --- e la
selezione dei moduli, che per costruzione avviene a tempo di build
(\cref{cap:08-sistema-moduli}).

\section{Il meccanismo: snapshot immutabile e revisione}

Lo stato \`e un unico record immutabile in un \code{AtomicReference}. Ogni
aggiornamento ne produce uno nuovo con revisione incrementata. I consumatori
leggono sempre l'ultimo:

\begin{lstlisting}[language=Java]
@Override
public Duration syncQueueMaxEstimatedWait() {
    return current.get().syncQueueMaxEstimatedWait();
}
\end{lstlisting}

La propriet\`a che questa costruzione garantisce \`e che una singola valutazione
di ammissione veda una configurazione \emph{coerente}: legge uno snapshot, e i
campi di quello snapshot appartengono tutti alla stessa revisione. Con sei
variabili volatili indipendenti, un aggiornamento concorrente potrebbe far s\`i
che una valutazione legga la nuova soglia di attesa e il vecchio flag di
ammissione --- una combinazione che nessuno ha mai richiesto.

\subsection{Il lock ottimistico}

L'aggiornamento richiede al chiamante di dichiarare la revisione che si aspetta:

\begin{lstlisting}[language=Java]
public RuntimeConfigSnapshot update(long expectedRevision, RuntimeConfigPatch patch) {
    while (true) {
        RuntimeConfigSnapshot snapshot = current.get();
        if (snapshot.revision() != expectedRevision) {
            throw new RevisionMismatchException(expectedRevision, snapshot.revision());
        }
        RuntimeConfigSnapshot candidate = snapshot.applyPatch(patch);
        if (current.compareAndSet(snapshot, candidate)) {
            return candidate;
        }
        // CAS failed because another thread updated; retry loop will re-check revision
    }
}
\end{lstlisting}

\`E il pattern \emph{compare-and-swap} classico, con il controllo di revisione
che lo trasforma in un lock ottimistico visibile al chiamante. Un aggiornamento
basato su uno stato obsoleto riceve \code{409 Conflict} con la revisione
corrente, anzich\'e sovrascrivere in silenzio una modifica concorrente.

L'aggiornamento \`e un \code{PATCH} con semantica ``i campi nulli restano
invariati''. La combinazione di patch parziale e revisione attesa \`e ci\`o che
rende sicure due modifiche concorrenti su campi diversi: la seconda vede il
conflitto, rilegge, e riapplica.

\section{L'API di amministrazione}

\begin{center}
\small
\begin{adjustbox}{max width=\textwidth}
\begin{tabular}{@{}llp{6.6cm}@{}}
\toprule
\textbf{Metodo} & \textbf{Percorso} & \textbf{Semantica} \\
\midrule
\code{GET} & \code{/v1/admin/runtime-config} & Snapshot corrente con revisione. \\
\code{POST} & \code{/v1/admin/runtime-config/validate} & Prova la patch senza
applicarla. \\
\code{PATCH} & \code{/v1/admin/runtime-config} & Applica, con
\code{expectedRevision} obbligatoria. \\
\bottomrule
\end{tabular}
\end{adjustbox}
\end{center}

\subsection{L'endpoint di validazione}

L'esistenza di \code{/validate} merita una nota. Applica la patch a una copia
dello snapshot corrente, ne verifica la validit\`a e riporta gli errori, senza
modificare nulla. \`E lo stesso codice che il \code{PATCH} esegue prima di
applicare.

Il suo valore \`e per gli strumenti automatici: uno script di esperimento pu\`o
verificare che la propria matrice di configurazioni sia interamente valida prima
di iniziare, anzich\'e scoprire al terzo punto della matrice che una combinazione
verr\`a rifiutata --- avendo gi\`a consumato il tempo dei primi due.

\subsection{La sequenza del \code{PATCH}}

L'ordine delle operazioni codifica una scelta di sicurezza:

\begin{enumerate}[leftmargin=*]
  \item \code{expectedRevision} presente, altrimenti \code{400};
  \item la patch \`e sintatticamente valida (durate parsabili, ecc.), altrimenti
  \code{400} con nome del campo e valore offensivo;
  \item lo snapshot risultante \`e semanticamente valido, altrimenti \code{422};
  \item la revisione corrisponde, altrimenti \code{409} con la revisione reale;
  \item lo snapshot viene sostituito atomicamente;
  \item l'applicatore propaga ai componenti che hanno bisogno di essere
  informati, con rollback in caso di fallimento (\code{503}).
\end{enumerate}

La validazione precede la sostituzione: uno snapshot invalido non entra mai nello
stato. Il rollback al punto 6 copre il caso residuo in cui la propagazione
fallisca dopo che lo stato \`e cambiato.

\subsection{Il metodo sincronizzato}

\begin{lstlisting}[language=Java]
@PatchMapping
public synchronized ResponseEntity<Object> patch(@RequestBody PatchRequest request) { ... }
\end{lstlisting}

Il metodo del controller \`e \code{synchronized}, il che serializza tutte le
modifiche di configurazione. \`E ridondante rispetto al CAS del servizio per la
correttezza dello stato, ma serializza anche la fase di applicazione, che il CAS
non copre: due patch concorrenti potrebbero altrimenti propagare i propri effetti
in ordine invertito rispetto a quello in cui hanno vinto il CAS.

\`E anche una violazione del principio non bloccante del
\cref{cap:03-requisiti-principi}, ed \`e accettabile perch\'e riguarda un
endpoint amministrativo a bassissima frequenza. Vale la pena notare che il
codice non lo giustifica, e che il lettore deve dedurre da s\'e che il costo \`e
confinato.

\section{L'applicatore}

\code{RuntimeConfigApplier} propaga lo snapshot ai componenti mutabili in ordine
deterministico:

\begin{lstlisting}[language=Java]
/**
 * Applies the new snapshot to all mutable runtime components.
 * Sync queue fields are read directly from {@link RuntimeConfigService} by consumers,
 * so only components with dedicated setters need explicit application here.
 */
public void apply(RuntimeConfigSnapshot snapshot, RuntimeConfigSnapshot previous,
                  RuntimeConfigService configService) { ... }
\end{lstlisting}

Il commento rivela che esistono due modalit\`a di propagazione:

\begin{description}[style=nextline, leftmargin=0pt]
  \item[\textbf{Pull}] --- i consumatori della coda sincrona leggono dallo
  snapshot a ogni uso. Non c'\`e nulla da propagare; la modifica ha effetto
  immediatamente e senza coordinamento.
  \item[\textbf{Push}] --- il limitatore di frequenza possiede una propria
  variabile interna, e l'applicatore deve chiamarne il setter.
\end{description}

La modalit\`a pull \`e evidentemente preferibile e il modulo la usa dove
possibile; la push esiste per un componente del nucleo che \`e antecedente al
modulo e che il modulo non deve costringere a cambiare struttura. \`E una
concessione ragionevole al principio di non far dipendere il nucleo dai moduli.

\section{Osservabilit\`a}

\begin{center}
\small
\begin{adjustbox}{max width=\textwidth}
\begin{tabular}{@{}lll@{}}
\toprule
\textbf{Metrica} & \textbf{Tipo} & \textbf{Significato} \\
\midrule
\code{controlplane\_runtime\_config\_revision} & gauge & Revisione corrente. \\
\code{controlplane\_runtime\_config\_updates\_total} & counter &
Aggiornamenti, con etichetta \code{status}. \\
\code{controlplane\_runtime\_config\_apply\_duration\_seconds} & timer &
Durata dell'applicazione. \\
\bottomrule
\end{tabular}
\end{adjustbox}
\end{center}

Il gauge sulla revisione \`e lo strumento pi\`u utile del gruppo in contesto
sperimentale: esportando la revisione nelle stesse serie temporali delle metriche
di prestazione, un cambio di configurazione diventa un gradino osservabile sul
grafico, e l'attribuzione di una variazione di latenza al cambio che l'ha causata
non richiede di correlare a mano i log con le misure.

\section{Disattivato per default}

\begin{lstlisting}[language=yamlnf]
nanofaas:
  admin:
    runtime-config:
      enabled: false
\end{lstlisting}

L'API espone operazioni di scrittura su una piattaforma priva di autenticazione
(\cref{cap:03-requisiti-principi}). Il default disattivato fa s\`i che
l'esposizione sia una decisione consapevole dell'operatore anzich\'e una
propriet\`a ereditata dalla scelta dei moduli. \`E l'unico modulo che, pur essendo
compilato nel binario, resta inerte finch\'e non viene attivato esplicitamente.

\chapter{I provider di deployment}
\label{cap:15-deployment-provider}

Due moduli implementano l'interfaccia \code{ManagedDeploymentProvider}
(\cref{cap:07-modalita-esecuzione}) e forniscono alla modalit\`a
\code{DEPLOYMENT} un meccanismo concreto: \code{k8s-deployment-provider}
(1\,202 righe) e \code{container-deployment-provider} (1\,166 righe). Insieme
sono un terzo del codice di tutti i moduli.

Che siano due, e che siano intercambiabili, \`e la prova che l'astrazione
introdotta nel \cref{cap:07-modalita-esecuzione} non \`e speculativa: il progetto
la esercita con due realizzazioni sostanzialmente diverse.

\section{\code{k8s-deployment-provider}}

Il provider Kubernetes traduce una \code{FunctionSpec} in tre oggetti: un
\code{Deployment}, un \code{Service} e --- se la strategia di scaling \`e
\code{HPA} --- un \code{HorizontalPodAutoscaler}. Usa il client
Fabric8.

\subsection{Convenzioni di denominazione ed etichette}

I nomi degli oggetti sono derivati dal nome della funzione con prefisso
\code{fn-}, e ogni oggetto porta due etichette: \code{app=nanofaas} e
\code{function=<nome>}. La seconda \`e anche il selettore del Deployment.

Il template dei pod porta tre annotazioni per lo scraping Prometheus:

\begin{lstlisting}[language=Java]
.addToAnnotations(NanofaasDeploymentConstants.ANNOTATION_SCRAPE, "true")
.addToAnnotations(NanofaasDeploymentConstants.ANNOTATION_PATH, "/metrics")
.addToAnnotations(NanofaasDeploymentConstants.ANNOTATION_PORT, "8080")
\end{lstlisting}

Le metriche lato funzione sono quindi raccolte senza che l'operatore configuri
nulla. \`E una piccola decisione con una conseguenza sperimentale grande: le
misure del \cref{cap:25-valutazione-concorrenza} che confrontano il tempo di
servizio misurato dal control plane con quello misurato dentro la funzione sono
possibili solo perch\'e la seconda sorgente esiste per default.

\subsection{Le variabili d'ambiente riservate}

\begin{lstlisting}[language=Java]
private static final Set<String> RESERVED_ENV = Set.of(
        "FUNCTION_NAME", "WARM", "TIMEOUT_MS", "EXECUTION_MODE", "WATCHDOG_CMD", "CALLBACK_URL"
);
\end{lstlisting}

Sono il contratto fra control plane e runtime della funzione
(\cref{cap:17-runtime-funzione}). Il builder le inietta e impedisce alla
specifica dell'utente di sovrascriverle: \code{CALLBACK\_URL} in particolare \`e
l'indirizzo a cui il runtime invia il completamento asincrono, e un valore
sbagliato produrrebbe esecuzioni che non terminano mai --- il caso che il terzo
orizzonte di eviction dello \code{ExecutionStore} deve poi ripulire
(\cref{cap:06-nucleo}).

\subsection{Il deadlock dello scale-to-zero con HPA}

Il commento pi\`u istruttivo del modulo:

\begin{lstlisting}[language=Java]
if (replicas == 0 && hpaOwnsScaling(spec.scalingConfig())) {
    // Kubernetes 1.36 only lets an HPA scale a target up from zero once it
    // has parked that target at zero itself (status condition ScaledToZero).
    // A Deployment created at zero therefore deadlocks: "scaling is disabled
    // since the replica count of the target is zero", forever. Start at one
    // and let the HPA scale it down, which is what stamps that condition.
    replicas = 1;
}
\end{lstlisting}

Una funzione con \code{minReplicas = 0} e strategia \code{HPA} verrebbe creata a
zero repliche, e l'HPA rifiuterebbe di risvegliarla perch\'e non \`e stato lui a
portarla a zero. Il rimedio \`e creare il Deployment a una replica e lasciare che
sia l'HPA a scalarla a zero, il che imprime la condizione di stato che poi gli
consente di risalire.

\`E il tipo di dettaglio che vale la pena riportare in un documento accademico
per una ragione metodologica: una piattaforma FaaS costruita sopra Kubernetes
eredita i vincoli di Kubernetes, e alcuni di questi non sono documentati come
vincoli ma si manifestano come comportamenti. Il costo dell'astrazione
``deployment gestito'' include l'obbligo di conoscerli.

\subsection{Validazione delle immagini}

\code{KubernetesImageValidator} verifica che un'immagine esista e sia scaricabile
\emph{prima} che la registrazione della funzione venga accettata. Il metodo \`e
diretto: crea un pod di prova, ne osserva lo stato del container, e cerca le
condizioni di attesa che segnalano un fallimento di pull
(\code{ErrImagePull}, \code{ImagePullBackOff}), con un timeout di 20 secondi e
polling a 500\,ms.

\`E un controllo costoso --- crea e distrugge un pod --- eseguito una sola volta
per registrazione. L'alternativa sarebbe scoprire l'immagine inesistente alla
prima invocazione, quando il chiamante sta gi\`a attendendo e la diagnosi \`e
molto meno diretta. La registrazione fallisce con un errore che nomina il
problema.

\section{\code{container-deployment-provider}}

Il provider locale crea container su un runtime compatibile con Docker. \`E
pensato per lo sviluppo e --- soprattutto --- per gli esperimenti su una singola
macchina, dove interporre Kubernetes fra il control plane e le funzioni
aggiungerebbe una variabile senza aggiungere realismo.

\subsection{Struttura}

\begin{center}
\small
\begin{adjustbox}{max width=\textwidth}
\begin{tabular}{@{}lp{8.6cm}@{}}
\toprule
\textbf{Componente} & \textbf{Ruolo} \\
\midrule
\code{ContainerRuntimeAdapter} & Astrazione sul runtime; due implementazioni,
via CLI e via client Docker Java. \\
\code{EphemeralPortAllocator} & Assegna una porta libera dell'host a ciascun
container. \\
\code{HttpEndpointProbe} & Attende che il container risponda sull'endpoint di
salute. \\
\code{RoundRobinFunctionProxy} & Espone un endpoint stabile che bilancia sulle
repliche. \\
\code{DockerImageValidator} & Verifica l'esistenza dell'immagine. \\
\bottomrule
\end{tabular}
\end{adjustbox}
\end{center}

\subsection{Il proxy: perch\'e serve}

Kubernetes fornisce un \code{Service}, ossia un nome stabile che bilancia sui
pod. Un runtime container locale non ha nulla di equivalente: ogni container ha
la propria porta effimera sull'host. Se il control plane conoscesse quelle porte
direttamente, ogni scale-up e ogni scale-down cambierebbero l'endpoint della
funzione.

\code{RoundRobinFunctionProxy} colma la differenza con un piccolo server HTTP ---
costruito sul \code{HttpServer} del JDK, senza dipendenze aggiuntive --- che
espone \code{/invoke} e \code{/health} su una porta stabile e distribuisce le
richieste sui backend correnti con round robin. L'insieme dei backend \`e un
riferimento atomico che il provider aggiorna a ogni cambio di scala.

\`E, in sostanza, una reimplementazione minimale del \code{Service} di
Kubernetes, ed \`e la componente che rende i due provider \emph{equivalenti dal
punto di vista del control plane}: entrambi restituiscono un URL stabile.

\subsection{Il \code{cpuset}: una decisione derivata da una misura}

Il commento su \code{ContainerLocalProperties} \`e uno dei pi\`u densi del
progetto:

\begin{quote}\itshape
\code{cpuset}: le CPU dell'host a cui ogni container di funzione gestita \`e
vincolato, per esempio \code{"0-3"}. Un limite di CPU per funzione limita ciascun
container separatamente, il che su una macchina con core liberi significa che le
funzioni non competono mai davvero --- misurato: due funzioni con quattro CPU
ciascuna su un host a undici core hanno condiviso un control plane e si sono a
malapena influenzate a vicenda. Vincolarle agli stessi core \`e ci\`o che rende la
capacit\`a della piattaforma una quantit\`a che devono dividersi, che \`e la
premessa su cui \`e costruito il budget di concorrenza.
\end{quote}

Questa \`e una funzionalit\`a nata da un requisito \emph{sperimentale}, non
operativo. Il modo \code{BUDGETED} del \cref{cap:12-concurrency-control} presume
che la concorrenza sia una risorsa scarsa condivisa; su un host con core in
eccesso quella premessa \`e falsa e il modo non pu\`o essere valutato. Il
\code{cpuset} fabbrica la scarsit\`a necessaria a rendere l'esperimento
significativo.

\`E anche il parametro che ha consentito di \emph{escludere} la contesa di CPU
come spiegazione del cross-talk fra funzioni co-residenti: vincolando entrambe
agli stessi quattro core, i risultati sono rimasti identici a tre cifre
significative (\cref{cap:25-valutazione-concorrenza}).

\subsection{Readiness}

\code{HttpEndpointProbe} interroga l'endpoint di salute del container con un
timeout di 20 secondi e un intervallo di 250\,ms. Un container che non diventa
pronto entro il timeout fa fallire il provisioning.

I valori sono deliberatamente pi\`u aggressivi delle sonde di Kubernetes:
l'obiettivo qui non \`e la robustezza in produzione ma il tempo di ciclo di un
esperimento, in cui il provisioning avviene decine di volte.

\section{Chi possiede la validazione delle immagini}

Un punto architetturale gi\`a anticipato nel \cref{cap:03-requisiti-principi} e
che qui trova la sua giustificazione completa: la validazione delle immagini
\emph{non} \`e un modulo autonomo. Ciascun provider possiede il proprio
validatore e lo attiva quando \`e selezionato come backend.

Il motivo \`e che l'operazione non \`e la stessa. Verificare un'immagine tramite
Kubernetes significa creare un pod e osservare le condizioni di attesa del
container; verificarla tramite un demone Docker locale significa interrogarne
l'indice delle immagini o tentare un pull. Un modulo generico avrebbe dovuto
reintrodurre al proprio interno la distinzione fra backend che l'astrazione dei
provider esiste per eliminare, e avrebbe dovuto essere selezionato
coerentemente con essi --- un vincolo di coerenza in pi\`u nella selezione
modulare, per nessun guadagno.

Resta il fatto che l'interfaccia \code{ImageValidator} vive nel package
\code{registry} del control plane anzich\'e in \code{platform/common}, dove
stanno gli altri contratti condivisi (\cref{cap:05-contratti-dati}). Non \`e una
violazione architetturale --- i moduli possono dipendere dal nucleo, e la regola
ArchUnit vieta soltanto le dipendenze \emph{fra} moduli
(\cref{cap:26-qualita-verifica}) --- ma \`e una collocazione incoerente con il
criterio dichiarato altrove, e il posto naturale dell'interfaccia \`e il modulo
dei contratti.

\section{Confronto}

\begin{table}[htbp]
\centering
\caption{I due provider di deployment a confronto.}
\label{tab:provider}
\small
\begin{tabularx}{\textwidth}{@{}lXX@{}}
\toprule
 & \textbf{\code{k8s}} & \textbf{\code{container-local}} \\
\midrule
Backend & API Kubernetes (Fabric8) & runtime Docker-compatibile \\
Endpoint stabile & \code{Service} & proxy round robin interno \\
Scaling & \code{replicas} del Deployment o HPA & creazione/distruzione
container \\
Validazione immagine & pod di prova, condizioni di attesa & indice del demone \\
Isolamento CPU & \code{resources.limits} & \code{resources.limits} +
\code{cpuset} \\
Dipendenze pesanti & Fabric8, Vert.x & client Docker Java \\
Uso previsto & carichi realistici su cluster & esperimenti su singola macchina \\
\bottomrule
\end{tabularx}
\end{table}

L'ultima riga \`e la ragione per cui entrambi esistono. Il provider Kubernetes
\`e quello che rende \nf{} una piattaforma FaaS reale; quello locale \`e quello
che rende gli esperimenti eseguibili in un ciclo abbastanza breve da poterli
ripetere.

\chapter{Moduli diagnostici e superficie minima}
\label{cap:16-moduli-diagnostici}

Questo capitolo chiude la Parte III con due argomenti che il resto non ha
collocato: il modulo pi\`u piccolo del progetto, e la configurazione che non ha
moduli affatto.

\section{\code{build-metadata}}

Quarantatr\'e righe in tre classi. Espone un endpoint,
\code{/modules/build-metadata}, che riporta l'identit\`a della build in
esecuzione.

La sua utilit\`a \`e proporzionale alla combinatoria che il sistema modulare
introduce. Un binario di \nf{} \`e definito da: la versione del codice, la
selezione dei moduli, il tipo di build (JVM o immagine nativa), il livello di
ottimizzazione, il garbage collector. Il \cref{cap:24-footprint} confronta
configurazioni che differiscono \emph{solo} per l'ultimo di questi assi. In un
tale contesto, l'assenza di un modo diretto per interrogare un'istanza in
esecuzione sulla propria identit\`a \`e una fonte reale di errore sperimentale:
il rischio non \`e teorico, ed \`e quello di attribuire a una politica una
differenza dovuta al fatto che una delle due istanze non era quella che si
credeva.

Che sia un modulo, e non un endpoint del nucleo, \`e coerente: la
diagnostica sulla build serve a chi conduce esperimenti, non alla piattaforma.
Un deployment di produzione pu\`o non compilarla.

\section{La configurazione senza moduli}

\subsection{Che cosa resta}

Compilando con \code{-PcontrolPlaneModules=none} si ottiene una piattaforma con
le seguenti propriet\`a:

\begin{center}
\small
\begin{adjustbox}{max width=\textwidth}
\begin{tabular}{@{}lp{8.4cm}@{}}
\toprule
\textbf{Capacit\`a} & \textbf{Stato} \\
\midrule
Registrazione e ciclo di vita & Completa. \\
Invocazione sincrona & Completa; dispatch diretto, senza coda n\'e ammissione. \\
Invocazione asincrona & \code{501 Not Implemented}. \\
Idempotenza, retry, timeout & Complete. \\
Modalit\`a \code{LOCAL} ed \code{EXTERNAL} & Complete. \\
Modalit\`a \code{DEPLOYMENT} & Non disponibile: nessun provider. \\
Limite di concorrenza & Nessuna applicazione: senza code non c'\`e chi conti gli
slot. \\
Metriche & Complete per le grandezze del nucleo. \\
\bottomrule
\end{tabular}
\end{adjustbox}
\end{center}

\`E ancora una piattaforma FaaS secondo la definizione operativa del
\cref{cap:02-stato-dell-arte}? Soddisfa l'obbligo di registrazione, quello di
instradamento, e --- in modalit\`a \code{EXTERNAL} --- opera su unit\`a di calcolo
che qualcun altro ha attivato. Non soddisfa attivazione e ritiro. \`E quindi una
piattaforma FaaS \emph{ridotta}, come il \cref{cap:03-requisiti-principi}
prometteva, non una piattaforma rotta: ogni operazione che accetta la porta a
termine correttamente, e le operazioni che non pu\`o svolgere le rifiuta
esplicitamente.

\subsection{\code{501}, non \code{503}}

La scelta del codice di stato per \code{:enqueue} \`e deliberata. \code{503
Service Unavailable} significa ``riprova pi\`u tardi''; \code{501 Not
Implemented} significa ``questa capacit\`a non esiste in questo binario''. Un
client che riceve \code{503} riprova e continuer\`a a fallire; uno che riceve
\code{501} sa che deve cambiare qualcosa.

\subsection{La copertura di test}

La configurazione senza moduli ha una classe di test dedicata,
\code{CoreOnlyApiTest}, che verifica ci\`o che vale \emph{solo} in sua assenza:
che l'accodatore sia il no-op e che l'API asincrona risponda \code{501}. La
classe si auto-esclude quando i moduli sono presenti.

Il costo \`e reale. Poich\'e la selezione modulare \`e risolta a tempo di
configurazione di Gradle e non pu\`o variare all'interno di una singola build,
quella classe richiede una propria invocazione:

\begin{lstlisting}[language=bash]
./gradlew test                                   # tutti i moduli (default: all)
./gradlew test -PcontrolPlaneModules=none        # la configurazione minima
\end{lstlisting}

Due invocazioni della suite per ogni verifica completa. \`E il prezzo diretto
della scelta di selezionare a build time anzich\'e a runtime
(\cref{cap:08-sistema-moduli}), ed \`e utile averlo quantificato: la
modularit\`a a costo zero sul percorso critico non \`e a costo zero sulla catena
di integrazione continua.

\subsection{Perch\'e il default \`e \code{all}}

Ogni invocazione di Gradle priva di selezione esplicita usa \code{all}, test
compresi. La ragione \`e che la suite deve esercitare la configurazione che viene
distribuita. Una suite che girasse per default su una configurazione parziale
verificherebbe un artefatto che nessuno esegue, e le interazioni fra moduli ---
che sono la parte pi\`u fragile di un sistema componibile --- resterebbero
scoperte.

La combinazione delle due scelte, default \code{all} e test dedicato per
\code{none}, copre i due estremi dello spazio delle configurazioni. Le
configurazioni intermedie --- per esempio \code{async-queue} senza
\code{sync-queue} --- non hanno copertura dedicata.

\begin{damisurare}
Lo spazio delle configurazioni modulari ha $2^9 = 512$ punti, di cui alcuni non
validi per la dipendenza fra \code{concurrency-control} e \code{async-queue}. La
suite ne verifica due. Una verifica di costruibilit\`a su tutte le combinazioni
valide \`e realizzabile e non \`e stata eseguita.
\todomis{costruibilit\`a e avviabilit\`a di tutte le combinazioni modulari valide}
\end{damisurare}


\part{Runtime, SDK, strumenti}
\chapter{Il runtime di funzione}
\label{cap:17-runtime-funzione}

I capitoli precedenti hanno descritto il control plane. Questo attraversa il
confine e descrive che cosa gira dentro il container di una funzione, quale
contratto lo lega al control plane, e perch\'e quel contratto ha la forma che ha.

\section{Anatomia del container}

Un container gestito da \nf{} contiene due strati.

\begin{center}
\small
\begin{adjustbox}{max width=\textwidth}
\begin{tabular}{@{}lp{9cm}@{}}
\toprule
\textbf{Strato} & \textbf{Ruolo} \\
\midrule
\code{watchdog} & Entrypoint del container. Binario Rust statico, circa 2\,MB,
distribuito su un'immagine \code{FROM scratch}. Possiede il ciclo di vita del
processo utente. \\
\emph{runtime} & Il processo dell'utente: un'applicazione Spring Boot, un
servizio FastAPI, un binario Go, uno script, qualunque cosa risponda a uno dei
tre protocolli. \\
\bottomrule
\end{tabular}
\end{adjustbox}
\end{center}

La separazione risolve un problema pratico. Il control plane vuole parlare un
solo protocollo HTTP con ogni funzione; gli autori di funzioni vogliono scrivere
in qualunque linguaggio, incluso uno che non ha un server HTTP. Il watchdog
adatta il secondo al primo.

\section{I tre protocolli di runtime}

L'enumerazione \code{RuntimeMode} (\cref{cap:05-contratti-dati}) dichiara come il
watchdog deve parlare con il processo figlio.

\begin{center}
\small
\begin{adjustbox}{max width=\textwidth}
\begin{tabular}{@{}lp{8.8cm}@{}}
\toprule
\textbf{Modo} & \textbf{Protocollo} \\
\midrule
\code{HTTP} & Il watchdog sonda \code{GET /health} finch\'e il runtime \`e
pronto, poi fa da reverse proxy su \code{POST /invoke} verso
\code{127.0.0.1:8081}. \\
\code{STDIO} & Il watchdog scrive il payload JSON su stdin del processo, legge
la risposta JSON da stdout, e ne raccoglie l'uscita. \\
\code{FILE} & Il watchdog scrive \code{/tmp/input.json}, avvia il processo con
le variabili \code{INPUT\_FILE} e \code{OUTPUT\_FILE}, e legge
\code{/tmp/output.json}. \\
\bottomrule
\end{tabular}
\end{adjustbox}
\end{center}

I tre modi coprono uno spettro di invasivit\`a decrescente. \code{HTTP} richiede
un server nel processo utente ed \`e il pi\`u efficiente, perch\'e il processo
resta vivo fra le invocazioni. \code{STDIO} non richiede nulla oltre alla
capacit\`a di leggere e scrivere JSON, ma paga la creazione di un processo per
invocazione. \code{FILE} \`e il ripiego per binari che non hanno un'interfaccia
migliore.

Il punto importante \`e che i tre modi sono \emph{indistinguibili dal control
plane}: qualunque sia il modo, il control plane invia una \code{POST} al
watchdog e riceve una risposta. La differenza \`e interamente confinata dietro
il watchdog.

\section{Warm e one-shot}

Il watchdog opera in due regimi.

In modalit\`a \code{DEPLOYMENT} resta vivo fra le invocazioni (container caldo):
espone il proprio server HTTP su \code{:8080}, esegue \code{WATCHDOG\_CMD} per
invocazione nei modi \code{STDIO} e \code{FILE}, e fa da reverse proxy verso il
runtime interno nel modo \code{HTTP}. Il control plane invia un
\code{InvocationRequest} come corpo e legge la risposta direttamente --- nessuna
callback per invocazione.

Il regime one-shot, precedente, avvia il processo una volta per invocazione e ne
riporta l'esito con una callback verso \code{CALLBACK\_URL}. \`E il regime che
paga un'inizializzazione per ogni richiesta, e la sua esistenza \`e ci\`o che
rende \code{coldStart} e \code{initDurationMs} campi significativi in
\code{ExecutionStatus}.

\subsection{Le variabili d'ambiente del contratto}

\begin{center}
\small
\begin{adjustbox}{max width=\textwidth}
\begin{tabular}{@{}llp{5.8cm}@{}}
\toprule
\textbf{Variabile} & \textbf{Regime} & \textbf{Significato} \\
\midrule
\code{FUNCTION\_NAME} & warm & Nome logico. \\
\code{WARM} & warm & \code{true}. \\
\code{TIMEOUT\_MS} & entrambi & Budget per invocazione. \\
\code{EXECUTION\_MODE} & entrambi & \\
\code{WATCHDOG\_CMD} & entrambi & Comando del processo utente. \\
\code{RUNTIME\_URL} & warm & \code{http://127.0.0.1:8081/invoke}. \\
\code{EXECUTION\_ID} & one-shot & Identificatore; nel warm arriva come header
\code{X-Execution-Id}. \\
\code{CALLBACK\_URL} & one-shot & Endpoint di completamento. \\
\code{INVOCATION\_PAYLOAD} & one-shot & Il payload, nell'ambiente. \\
\code{TRACE\_ID} & one-shot & \\
\bottomrule
\end{tabular}
\end{adjustbox}
\end{center}

Sono le stesse variabili che il builder Kubernetes protegge dalla sovrascrittura
da parte dell'utente (\cref{cap:15-deployment-provider}).

Il passaggio dal regime one-shot al warm si legge nella tabella: nel warm
l'identificatore di esecuzione si sposta dall'ambiente agli header, il che \`e
obbligato --- un processo che serve molte invocazioni non pu\`o riceverne
l'identit\`a in una variabile d'ambiente fissata all'avvio.

\section{Il contratto della risposta}

\subsection{\code{HandlerResponse}}

Un handler pu\`o restituire un valore semplice, che diventa l'output, oppure una
busta:

\begin{lstlisting}[language=Java]
/**
 * Optional response envelope a handler may return instead of a plain value to control
 * the HTTP status code, a safe set of response headers, and whether {@code output} is
 * base64-encoded binary. Its mere presence as the handler's return value is itself the
 * signal that the response was function-decided, not a platform error.
 */
public record HandlerResponse(Object output, int statusCode, Map<String, String> headers, String encoding) { ... }
\end{lstlisting}

La frase decisiva \`e l'ultima: \emph{la sola presenza della busta \`e il segnale
che la risposta \`e stata decisa dalla funzione}. \`E ci\`o che consente al
control plane di distinguere un \code{404} applicativo --- ``la risorsa che hai
chiesto alla mia funzione non esiste'' --- da un \code{404} di piattaforma ---
``questa funzione non \`e registrata''. La distinzione viaggia sul filo come
header \code{X-NanoFaaS-Function-Status} ed \`e interpretata sia dal dispatcher
esterno (\cref{cap:06-nucleo}) sia dal gateway di offload
(\cref{cap:13-offload}).

Senza questo meccanismo, il tasso di errore misurato da una piattaforma FaaS \`e
ambiguo: mescola i fallimenti della piattaforma con le risposte non-2xx
deliberate delle funzioni. Il \cref{cap:21-osservabilita} riprende il punto.

\subsection{La lista bianca degli header}

Una funzione non pu\`o impostare header arbitrari sulla risposta:

\begin{lstlisting}[language=Java]
public static final Set<String> ALLOWED_RESPONSE_HEADERS = Set.of(
        "content-type", "location", "cache-control", "etag",
        "content-disposition", "content-language", "retry-after", "vary");
\end{lstlisting}

Otto header, tutti descrittivi del contenuto o della sua cacheabilit\`a. Sono
esclusi gli header hop-by-hop, quelli che governano la connessione, e quelli che
il control plane usa per il proprio protocollo. Un handler che potesse impostare
\code{X-NanoFaaS-Function-Status} potrebbe mentire sulla natura della propria
risposta; uno che potesse impostare \code{Transfer-Encoding} potrebbe corrompere
la connessione.

Il filtro contiene due dettagli che vale la pena isolare, perch\'e riguardano
classi di bug reali.

\paragraph{Collisioni per maiuscole.}
\begin{lstlisting}[language=Java]
/**
 * At most one entry survives per header name, compared case-insensitively: HTTP header names
 * are case-insensitive, so emitting both {@code Content-Type} and {@code content-type} would put
 * two colliding entries on the response. The first occurrence in iteration order wins and keeps
 * its original casing — that casing is part of the public {@code InvocationResponse.headers}
 * contract and callers read it.
 */
\end{lstlisting}

\paragraph{Il locale turco.}
\begin{lstlisting}[language=Java]
// Locale.ROOT, not the default locale: under a Turkish/Azerbaijani default,
// 'I' folds to dotless 'i' (U+0131) and the allow-list match would silently fail.
String lowerKey = entry.getKey().toLowerCase(Locale.ROOT);
\end{lstlisting}

In turco la minuscola di \code{I} \`e \code{\i} senza punto, e
\code{"Location".toLowerCase()} su una JVM con locale turco non produce
\code{"location"}. Il confronto con la lista bianca fallirebbe, e l'header
sarebbe silenziosamente scartato --- su alcune macchine e non su altre. Lo stesso
accorgimento compare nel controller per il valore dell'header di ragione del
rifiuto (\cref{cap:04-architettura-control-plane}).

Sono dettagli, ma sono i dettagli che distinguono un contratto che funziona da
uno che funziona sulla macchina dello sviluppatore.

\section{Il watchdog}

Il watchdog \`e scritto in Rust, compilato staticamente con musl, e distribuito
su un'immagine \code{FROM scratch} di circa 2\,MB. Il sorgente principale \`e di
circa 1\,500 righe.

La scelta di Rust e della compilazione statica \`e coerente con il ruolo. Il
watchdog \`e presente in \emph{ogni} container di funzione, quindi la sua
impronta si moltiplica per il numero di repliche di ogni funzione; ed \`e
l'entrypoint, quindi il suo tempo di avvio si somma a ogni cold start. Un
binario statico da 2\,MB su un'immagine vuota \`e ci\`o che rende quei due costi
trascurabili.

Le sue responsabilit\`a: leggere la configurazione dall'ambiente, avviare il
processo figlio, attendere che diventi pronto, applicare il timeout per
invocazione, inoltrare la richiesta secondo il modo configurato, inviare la
callback nel regime one-shot, e gestire i segnali di terminazione.

\begin{damisurare}
Il costo del watchdog sul percorso di invocazione --- l'ulteriore salto di
reverse proxy nel modo \code{HTTP} --- non \`e stato isolato. \`E incluso in
tutte le misure di \code{function\_latency\_ms} riportate nel
\cref{cap:23-prestazioni-base}, ma non \`e stato confrontato con un'invocazione
diretta al runtime.
\todomis{latenza aggiunta dal reverse proxy del watchdog}
\end{damisurare}

\section{Il servizio di riferimento \code{warm-echo}}

\code{services/java/warm-echo} \`e un servizio che incorpora l'SDK Java e si
limita a restituire l'input. Sono ventiquattro righe.

La sua funzione nel progetto \`e sperimentale: fornisce un runtime \emph{reale},
con un vero salto HTTP e un vero processo JVM, ma con un tempo di esecuzione
dell'handler prossimo a zero. Il confronto con la modalit\`a \code{LOCAL}
(\cref{cap:07-modalita-esecuzione}) isola cos\`i il costo del trasporto dal costo
del control plane.

Il progetto insiste su un punto: \code{warm-echo} \`e distribuito attraverso la
normale API di registrazione, non installato come infrastruttura condivisa. La
distinzione conta perch\'e mantiene il servizio di riferimento soggetto alle
stesse politiche --- code, limiti di concorrenza, retry --- di qualunque altra
funzione. Un servizio di riferimento privilegiato misurerebbe un percorso che
nessun'altra funzione percorre.

\chapter{Gli SDK multi-linguaggio}
\label{cap:18-sdk}

\nf{} fornisce SDK per sei ambienti: Java, Java-lite, Python, Go, JavaScript e
un adattatore per eseguibili generici (\code{exec}/bash). Questo capitolo li
attraversa e --- pi\`u interessante --- descrive il meccanismo con cui il
progetto verifica che parlino tutti lo stesso contratto.

\section{Il contratto minimo}

Un SDK di funzione ha tre obblighi: consentire all'autore di dichiarare un
handler, esporlo secondo uno dei tre protocolli di runtime
(\cref{cap:17-runtime-funzione}), e tradurre il valore restituito in una
risposta conforme.

L'interfaccia lato Java \`e un solo metodo:

\begin{lstlisting}[language=Java]
public interface FunctionHandler {
    Object handle(InvocationRequest request);
}
\end{lstlisting}

Ogni SDK realizza questo contratto con l'idioma del proprio linguaggio.

\section{Java: due SDK, due filosofie}

\subsection{\code{sdks/java}}

L'SDK completo (1\,607 righe) \`e costruito su Spring Boot. La dichiarazione di
un handler \`e un'annotazione:

\begin{lstlisting}[language=Java]
/**
 * Marks a Spring bean as a Nanofaas function handler.
 *
 * <p>Spring component scanning, not user code, discovers this annotation and registers the
 * handler in the application context. The runtime then resolves the single active handler bean
 * from that context when the control plane calls {@code /invoke}.</p>
 */
@Target(ElementType.TYPE) @Retention(RetentionPolicy.RUNTIME) @Documented @Component
public @interface NanofaasFunction {
    String value() default "";
}
\end{lstlisting}

L'autore della funzione ottiene iniezione delle dipendenze, configurazione, e
l'intero ecosistema Spring. Paga in impronta di memoria e tempo di avvio.

\subsection{\code{sdks/java-lite}}

L'SDK leggero (1\,454 righe) rinuncia a Spring. Costruisce direttamente sul
\code{HttpServer} del JDK, con tre gestori --- \code{InvokeHandler},
\code{HealthHandler}, \code{MetricsHandler} --- e nessun contenitore di
inversione del controllo.

Il confronto fra i due, misurato sulla stessa funzione
(\cref{cap:24-footprint}), \`e la giustificazione della loro coesistenza:

\begin{center}
\small
\begin{adjustbox}{max width=\textwidth}
\begin{tabular}{@{}lrr@{}}
\toprule
 & \textbf{\code{word-stats-java}} & \textbf{\code{word-stats-java-lite}} \\
\midrule
build & Spring Boot su JVM & binario nativo GraalVM \\
SDK & \code{sdks/java} & \code{sdks/java-lite} \\
residente, a riposo & 133\,MiB & 2{,}7\,MiB \\
residente, a caldo & 160\,MiB & 51{,}6\,MiB \\
tempo di servizio & 5{,}57\,ms & 2{,}15\,ms \\
\bottomrule
\end{tabular}
\end{adjustbox}
\end{center}

Il fattore cinquanta a riposo \`e la differenza fra un framework e la sua
assenza; il fattore tre a caldo mostra che il carico di lavoro effettivo
riavvicina i due (\cref{cap:24-footprint}).

\begin{quote}
\textbf{Avvertenza sperimentale.} La documentazione di progetto registra che i
commenti degli scenari hanno per un certo tempo trattato le due funzioni
\code{word-stats} come intercambiabili. Non lo sono: differiscono per SDK, per
tipo di build e per tempo di servizio. \emph{Qualunque esperimento che le tratti
come equivalenti \`e confuso.} \`E il tipo di errore che una piattaforma con
molte varianti dello stesso artefatto rende facile commettere, ed \`e riportato
qui perch\'e la sua scoperta ha invalidato analisi precedenti.
\end{quote}

\section{Python}

L'SDK Python \`e costruito su FastAPI (1\,407 righe di sorgente). La
dichiarazione dell'handler \`e un decoratore:

\begin{lstlisting}[language=Python]
@nanofaas_function
def handle(input):
    return {"echo": input}
\end{lstlisting}

Il decoratore registra la funzione in una variabile a livello di modulo, con la
regola che un solo handler pu\`o essere attivo e che una seconda applicazione
sovrascrive la prima. Sono supportati sia handler sincroni sia \code{async}.

La scelta di una registrazione globale anzich\'e di un registro esplicito \`e la
forma pi\`u idiomatica in Python e la meno cerimoniosa; ha il difetto di rendere
l'ordine degli import significativo, che \`e il prezzo abituale di questa
tecnica.

\section{Go, JavaScript, exec}

L'SDK Go (1\,914 righe) incorpora un runtime HTTP nel binario della
funzione ed espone un proprio endpoint \code{/metrics} Prometheus, con contatori
lato runtime per invocazioni, durata dell'handler, cold start e callback
asincrone scartate. \`E l'unico SDK che strumenta autonomamente il proprio lato
del confine, il che lo rende lo strumento pi\`u diretto per separare il tempo
speso dentro l'handler da quello speso nel trasporto.

L'SDK JavaScript (1\,531 righe di TypeScript) espone un runtime HTTP analogo per Node.

L'adattatore \code{exec} non \`e un SDK ma il modo \code{STDIO} del watchdog: un
eseguibile che legge JSON da stdin e scrive JSON su stdout \`e una funzione
valida senza libreria alcuna. \`E la via d'accesso per gli script bash e per i
binari preesistenti.

\section{Il contratto di saturazione}
\label{sec:18-saturazione}

Un runtime che accetta lavoro senza un limite \`e un runtime che, sotto carico,
consuma la memoria del contenitore finch\'e il kernel non lo termina. Il progetto
ha quindi fissato, in un unico documento (\code{sdks/runtime-contract/README.md}),
la politica con cui \emph{ogni} runtime di funzione deve rifiutare, e ne ha
fatto un contratto verificabile allo stesso modo della parit\`a funzionale.

Il contratto vale sia per le chiamate dirette a \code{POST /invoke} sia per i
dispacci del piano di controllo. L'ammissione \`e \emph{fail-fast}: il runtime
riserva capacit\`a di conteggio e di byte per ingresso, handler e callback
\emph{prima} di avviare l'handler o di produrre una copia serializzata di grandi
dimensioni. Una callback richiesta dall'invocazione fa parte dell'ammissione: un
runtime non pu\`o accettare il lavoro e scartare in silenzio la sua callback
terminale.

\begin{center}
\small
\begin{adjustbox}{max width=\textwidth}
\begin{tabular}{@{}llp{6.6cm}@{}}
\toprule
\textbf{Esito} & \textbf{Codice} & \textbf{Quando} \\
\midrule
429 & \code{RUNTIME\_HANDLER\_SATURATED} & ammissione dell'handler satura \\
429 & \code{RUNTIME\_CALLBACK\_SATURATED} & capacit\`a di callback satura \\
413 & \code{RUNTIME\_INPUT\_TOO\_LARGE} & ingresso oltre il limite di byte \\
500 & \code{RUNTIME\_OUTPUT\_TOO\_LARGE} & uscita scoperta \emph{dopo} l'esecuzione;
l'handler \`e partito, l'uscita non viene trattenuta e una callback d'errore
limitata \`e comunque dovuta \\
503 & \code{RUNTIME\_STOPPING} & richiesta accettata dopo l'inizio dell'arresto \\
504 & \code{HANDLER\_TIMEOUT} & attesa dell'handler scaduta, a livello di
tentativo (non il \code{408} per singolo attendente) \\
\bottomrule
\end{tabular}
\end{adjustbox}
\end{center}

Gli errori hanno forma \code{\{"error":\{"code":...,"message":...\}\}} e quelli
ritentabili portano \code{Retry-After: 1}. Tre regole completano la politica.
Ogni attesa \`e finita e configurabile. Ogni riserva di conteggio e di byte resta
posseduta finch\'e il lavoro rappresentato e il carico trattenuto non sono
\emph{fisicamente} spariti: scadere l'attesa HTTP non prova che un handler non
cooperativo si sia fermato, e per questo il runtime Go, per esempio, documenta
che un handler scaduto continua a occupare il proprio posto. Infine i runtime non
ridispacciano mai: la decisione di ritentare appartiene al piano di controllo
(\cref{cap:06-nucleo}), che la esercita mantenendo stabile
\code{X-Execution-Id} e incrementando \code{X-Dispatch-Attempt}.

Il contratto \`e eseguibile. Le casistiche sono in un solo file,
\code{saturation-wire-corpus.json}, e ciascuna implementazione ha un adattatore
di test sottile che lo legge; un validatore Python
(\code{validate\_saturation\_wire\_corpus.py}) verifica la coerenza interna del
corpus --- riferimenti, azioni e barriere, esiti esatti, contatori finali a
regime --- e ha una propria suite di mutazioni, cos\`i che il corpus non possa
diventare un oracolo indipendente e incoerente. L'ordinamento fra azioni usa
barriere nominate, non attese temporizzate: il tempo trascorso non \`e evidenza
di ordine. Gli adattatori sono nella suite nativa di ciascun linguaggio; Java e
Java-lite deserializzano il modello tipizzato completo, Go rifiuta i campi
sconosciuti.

I limiti hanno valori predefiniti finiti e sono configurabili per variabile
d'ambiente. Il runtime Go e quello Python condividono lo stesso vocabolario
(\code{NANOFAAS\_MAX\_CONCURRENT\_HANDLERS}, 32 predefiniti;
\code{NANOFAAS\_MAX\_INPUT\_BYTES} e \code{NANOFAAS\_MAX\_OUTPUT\_BYTES}, 1\,MiB;
\code{NANOFAAS\_MAX\_PENDING\_CALLBACKS}, 128;
\code{NANOFAAS\_MAX\_CALLBACK\_PAYLOAD\_BYTES}, 2\,MiB;
\code{NANOFAAS\_MAX\_PENDING\_CALLBACK\_BYTES}, 16\,MiB; timeout di lettura del
corpo, di tentativo di callback e di arresto di 5\,s), e valori oltre i massimi
di produzione ricadono sui predefiniti prima che venga allocato alcunch\'e. Un
dettaglio che mostra il criterio: la capacit\`a di callback effettiva \`e il
minimo fra il limite di conteggio e il rapporto fra il tetto aggregato di byte e
il tetto per singola callback, perch\'e il runtime riserva un'intera quota di
callback prima di avviare l'handler.

\`E la stessa filosofia dell'ArchUnit del piano di controllo
(\cref{sec:26-archunit}) applicata ai runtime: un limite che non \`e imposto e
verificato \`e un'intenzione. Il divario iniziale era reale ---
l'inventario del documento registra, per ciascun SDK, code e attivit\`a
\emph{senza} un tetto di ammissione (Java e Java-lite con un thread virtuale per
richiesta, Python con la coda illimitata di \code{to\_thread}, Go con una
goroutine per richiesta) --- e la sua chiusura \`e avvenuta per fasi che il
documento numera (P16a fissa politica e inventario, P16b implementa i limiti su
tutti i runtime, P17 il tracciamento del lavoro fisico in Python, P18 le risorse
di arresto di Java-lite). Non \`e stata misurata qui l'impronta di questi limiti
sul percorso caldo.

\section{La parit\`a fra SDK, e come viene verificata}

\subsection{Il problema}

Sei implementazioni indipendenti dello stesso contratto divergono. Non
divergono in modo spettacolare --- ciascuna \`e testata --- ma nei dettagli:
come viene serializzato un valore nullo, che cosa accade con un input vuoto, come
viene rappresentato un errore. Divergenze di questo tipo sono invisibili finch\'e
qualcuno non confronta due SDK sulla stessa funzione, che \`e esattamente ci\`o
che un esperimento comparativo fa.

\subsection{La soluzione: fixture condivise}

Il progetto definisce tre \emph{famiglie} di funzioni di riferimento ---
\code{word-stats}, \code{json-transform}, \code{roman-numeral} --- e per ciascuna
un file di fixture JSON con input e output attesi. Ogni implementazione consuma
lo stesso file dalla propria suite di test nativa.

\begin{center}
\small
\begin{adjustbox}{max width=\textwidth}
\begin{tabular}{@{}lccc@{}}
\toprule
\textbf{SDK} & \textbf{\code{word-stats}} & \textbf{\code{json-transform}} &
\textbf{\code{roman-numeral}} \\
\midrule
Java & s\`i & s\`i & s\`i \\
Java Lite & s\`i & s\`i & s\`i \\
Go & s\`i & s\`i & s\`i \\
Python & s\`i & s\`i & s\`i \\
JavaScript & s\`i & s\`i & s\`i \\
exec/bash & s\`i & s\`i & s\`i \\
\bottomrule
\end{tabular}
\end{adjustbox}
\end{center}

Il cancello completo si esegue con un comando:

\begin{lstlisting}[language=bash]
./functions/contract-tests/run.sh
\end{lstlisting}

Diciotto funzioni, sei SDK, tre famiglie. \`E il meccanismo che converte
un'affermazione di parit\`a in una verifica.

\subsection{Fixture di correttezza e fixture di carico}

Il progetto distingue due tipi di fixture, e la distinzione \`e metodologicamente
significativa:

\begin{itemize}[leftmargin=*]
  \item \code{correctness.json} contiene input \emph{e} output attesi. Verifica
  la semantica.
  \item \code{performance-*.json} contiene input soltanto, e \emph{omette
  deliberatamente} gli output attesi.
\end{itemize}

L'omissione \`e voluta. Una fixture di carico serve a fissare l'input dei
benchmark cos\`i che due esecuzioni siano confrontabili; verificare l'output
durante un benchmark aggiungerebbe alla misura il costo del confronto e
introdurrebbe una variabile nel percorso caldo. Separare i due ruoli evita che
uno strumento serva a due scopi facendone bene nessuno.

\section{Il catalogo delle funzioni di esempio}

Il repository contiene funzioni di esempio in cinque linguaggi:

\begin{center}
\small
\begin{adjustbox}{max width=\textwidth}
\begin{tabular}{@{}lp{8.4cm}@{}}
\toprule
\textbf{Linguaggio} & \textbf{Funzioni} \\
\midrule
Java & \code{word-stats}, \code{json-transform}, \code{roman-numeral},
\code{figlet}, \code{qr-code}, \code{handler-envelope},
\code{binary-envelope}, pi\`u tre varianti \code{-lite} \\
Python & le stesse tre famiglie, pi\`u \code{qr-code},
\code{handler-envelope}, \code{mlimage} \\
Go, JavaScript & le tre famiglie, pi\`u \code{qr-code} e
\code{handler-envelope} \\
bash & le stesse cinque, pi\`u \code{figlet} \\
\bottomrule
\end{tabular}
\end{adjustbox}
\end{center}

Le funzioni non sono soltanto dimostrative. Ciascuna copre un aspetto del
contratto: \code{word-stats} \`e il carico di riferimento delle campagne di
misura; \code{handler-envelope} esercita la busta di risposta
(\cref{cap:17-runtime-funzione}); \code{binary-envelope} esercita la codifica
base64; \code{qr-code} e \code{figlet} producono output non testuali;
\code{mlimage} \`e l'unica con un tempo di esecuzione sostanziale e un'immagine
pesante, e serve a esercitare il cold start reale.

\chapter{CLI e scaffolding}
\label{cap:19-cli-scaffolding}

Due strumenti stanno fra l'utente e la piattaforma: la CLI \code{nanofaas}, che
parla con un control plane, e \code{fn-init}, che genera lo scheletro di una
nuova funzione. Un terzo strumento --- quello che provvisiona l'infrastruttura
--- vive in un repository separato ed \`e trattato nel
\cref{cap:22-metodo-sperimentale}.

\section{Il confine fra gli strumenti}

La documentazione del progetto enuncia la divisione senza ambiguit\`a:

\begin{quote}\itshape
\code{nanofaas} \`e un client HTTP neutro rispetto al backend per un control
plane NanoFaaS. Non provvisiona infrastruttura, non chiama Helm o \code{kubectl},
e non ispeziona risorse specifiche del provider.
\end{quote}

La separazione ha una conseguenza di progetto: la CLI non ha alcuna dipendenza
da Kubernetes, da Docker o da qualunque backend, e la sua unica dipendenza
esterna \`e HTTP. \`E per questo che pu\`o essere compilata in un binario nativo
di dimensioni contenute.

\begin{center}
\small
\begin{adjustbox}{max width=\textwidth}
\begin{tabular}{@{}lp{8.2cm}@{}}
\toprule
\textbf{Obiettivo} & \textbf{Strumento} \\
\midrule
Costruire, registrare, invocare o ispezionare una funzione & CLI
\code{nanofaas} \\
Provvisionare una VM, installare k3s/Helm, distribuire immagini, eseguire E2E &
strumento di provisioning esterno \\
Eseguire il control plane in sviluppo & Gradle / Spring Boot \\
\bottomrule
\end{tabular}
\end{adjustbox}
\end{center}

\section{La CLI}

2\,426 righe di sorgente Java (\code{wc -l} su \code{clients/cli/src/main}, 2026-09-28), costruita su Picocli. La struttura dei comandi:

\begin{center}
\small
\begin{adjustbox}{max width=\textwidth}
\begin{tabular}{@{}ll@{}}
\toprule
\textbf{Comando} & \textbf{Ruolo} \\
\midrule
\code{fn apply} & Registra o aggiorna una funzione da un manifesto YAML. \\
\code{fn list} / \code{fn get} / \code{fn delete} & Ciclo di vita. \\
\code{fn test} & Invoca con un input di prova e confronta output e stato HTTP atteso. \\
\code{invoke} & Invocazione sincrona. \\
\code{enqueue} & Invocazione asincrona. \\
\code{exec get} & Stato di un'esecuzione. \\
\code{deploy} & Costruisce l'immagine e registra in un passo. \\
\bottomrule
\end{tabular}
\end{adjustbox}
\end{center}

La configurazione risolve in ordine: opzione da riga di comando, variabile
d'ambiente, file di configurazione (\code{\textasciitilde/.config/nanofaas/config.yaml}).
\`E la precedenza convenzionale, e l'assenza di sorprese \`e essa stessa una
propriet\`a desiderabile in uno strumento che compare dentro script di
esperimento.

\subsection{Due build, due scopi}

\begin{lstlisting}[language=bash]
./gradlew :nanofaas-cli:installDist      # distribuzione JVM
./gradlew :nanofaas-cli:nativeCompile    # eseguibile nativo GraalVM
\end{lstlisting}

La build JVM \`e quella di sviluppo. Quella nativa esiste perch\'e la CLI \`e
invocata dentro cicli: uno scenario di validazione la chiama decine di volte, e
il tempo di avvio della JVM --- dell'ordine del secondo --- si somma a ogni
chiamata. Un eseguibile nativo parte in decine di millisecondi.

\`E lo stesso argomento che il \cref{cap:20-build-packaging} applica al control
plane, ma con una differenza: per un processo di lunga durata il tempo di avvio
\`e ammortizzato e ci\`o che conta \`e il throughput a regime; per uno strumento a
riga di comando il tempo di avvio \emph{\`e} il tempo di esecuzione.

\subsection{Lo smoke test nativo}

\begin{lstlisting}[language=bash]
./gradlew :nanofaas-cli:nativeSmoke
\end{lstlisting}

Verifica help, versione e un comando HTTP contro uno stub locale. \`E una
verifica breve con uno scopo preciso: la compilazione nativa pu\`o riuscire e
produrre un binario che fallisce a runtime per una configurazione di riflessione
mancante --- l'errore caratteristico di GraalVM
(\cref{cap:20-build-packaging}). Un test che esercita almeno un percorso di
serializzazione reale \`e ci\`o che distingue ``ha compilato'' da ``funziona''.

Il rischio non \`e teorico. Il file di configurazione della riflessione della
CLI nominava undici tipi, e quattro comandi si rompevano nel binario nativo
mentre i loro test JVM passavano (\code{fn replicas get}, \code{control-plane
info}, \code{fn update -f} e \code{--{}-config <file>}: errori di lettura o di
serializzazione JSON e YAML). La correzione registra ogni \code{record} della CLI
pi\`u la classe di configurazione YAML (29 voci) e, soprattutto, aggiunge un test
(\code{NativeReflectConfigCoverageTest}) che \emph{scansiona} il codice
compilato alla ricerca dei \code{record} invece di elencarli: un tipo nuovo non
pu\`o pi\`u essere dimenticato. Con la vecchia configurazione il test nomina
esattamente i tipi mancanti. Dopo la correzione, secondo il messaggio del
commit, la CLI nativa e quella JVM producono lo stesso output per ogni comando
esercitato.

\subsection{Che cosa verifica \code{fn test}}

\code{fn test} passava un caso soltanto su uno stato 2xx, quindi un input che
descrive un percorso d'errore non poteva mai passare: un 422 deciso dalla
funzione arriva come HTTP 422 con stato di esecuzione \code{success} e il corpo
d'errore come output. Ogni scaffold di \code{fn-init} risponde 422 a un input
mancante, e dunque la CLI bocciava il \code{missing-input.json} dello scaffold
stesso appena generato. Ora un input pu\`o dichiarare \code{expectedStatusCode}
(chiave e predefinito 200 gi\`a usati dal corpus \code{correctness.json}): il caso
passa se lo stato HTTP coincide e l'output coincide con l'atteso, mentre un
errore di piattaforma continua a fallire. \`E la stessa distinzione fra errore
della funzione ed errore della piattaforma che il contratto della risposta
(\cref{cap:17-runtime-funzione}) fissa lato runtime.

\section{\code{fn-init}}

Genera lo scheletro di una nuova funzione: struttura di directory, file di
build, handler di esempio, Dockerfile, manifesto. \`E scritto in Python e vive
sotto \code{tools/}.

Il progetto ha registrato su questo strumento una regressione istruttiva: la
generazione dello scaffold JavaScript aveva perso la capacit\`a di incorporare
l'SDK, e la perdita non era stata rilevata perch\'e il test corrispondente era
rimasto orfano --- verificava una funzionalit\`a che non veniva pi\`u esercitata.
\`E il modo caratteristico in cui uno strumento di scaffolding degrada: il codice
generato continua a esistere e a sembrare corretto, e solo chi lo usa scopre che
manca qualcosa.

\begin{damisurare}
Non esiste una verifica che il codice generato da \code{fn-init} sia
effettivamente costruibile e registrabile per ogni linguaggio supportato. Un
cancello di questo tipo --- generare, costruire, registrare, invocare, per
ciascun template --- \`e realizzabile con gli strumenti gi\`a presenti.
\todomis{costruibilit\`a e invocabilit\`a degli scaffold generati, per linguaggio}
\end{damisurare}

\section{Il manifesto di funzione}

Il formato che \code{fn apply} consuma \`e la serializzazione YAML di
\code{FunctionSpec} (\cref{cap:05-contratti-dati}). Un esempio con le politiche
discusse nella Parte III:

\begin{lstlisting}[language=yamlnf, caption={Manifesto con controllo di
concorrenza a budget e offload reattivo.}]
name: word-stats
image: nanofaas/word-stats:1.0
executionMode: DEPLOYMENT
runtimeMode: HTTP
timeoutMs: 5000
concurrency: 8
queueSize: 100
resources:
  requests: { cpu: 0.5, memoryMiB: 256 }
  limits:   { cpu: 2.0, memoryMiB: 512 }
scalingConfig:
  strategy: INTERNAL
  minReplicas: 1
  maxReplicas: 6
  metrics:
    - { type: queue_depth, target: "5" }
  concurrencyControl:
    mode: BUDGETED
    targetLatencyMs: 10
    weight: 2.0
    minTargetInFlightPerPod: 1
    maxTargetInFlightPerPod: 8
offload:
  mode: pressure
  targetUrl: http://cloud-host:8080
\end{lstlisting}

Il manifesto \`e il punto in cui tutte le politiche descritte nella Parte III si
incontrano, ed \`e utile leggerlo come indice: \code{scalingConfig.strategy}
seleziona l'autoscaler (\cref{cap:11-autoscaler}),
\code{concurrencyControl.mode} il controllore di concorrenza
(\cref{cap:12-concurrency-control}), \code{offload} la politica di deviazione
(\cref{cap:13-offload}), e \code{queueSize} dimensiona la coda che l'ammissione
protegge (\cref{cap:10-sync-queue}).

\chapter{Build, packaging, immagini}
\label{cap:20-build-packaging}

Un control plane che si vuole misurabile deve essere anche \emph{producibile in
varianti confrontabili}. Questo capitolo descrive la catena di build di \nf{},
gli assi lungo cui un binario pu\`o variare, e le decisioni di packaging che il
progetto motiva con misure.

\section{Gli assi di variazione}

Un artefatto di \nf{} \`e determinato da cinque scelte indipendenti:

\begin{center}
\small
\begin{adjustbox}{max width=\textwidth}
\begin{tabular}{@{}lp{8.2cm}@{}}
\toprule
\textbf{Asse} & \textbf{Valori} \\
\midrule
Selezione dei moduli & \code{none}, \code{all}, o una lista
(\cref{cap:08-sistema-moduli}) \\
Tipo di build & JVM (Spring Boot fat jar) oppure immagine nativa GraalVM \\
Livello di ottimizzazione nativa & \code{-Os} (dimensione) o \code{-O3}
(velocit\`a, default) \\
Garbage collector nativo & \code{serial} (default Community) o \code{G1}
(Oracle GraalVM) \\
Limite di memoria del container & assente o fissato \\
\bottomrule
\end{tabular}
\end{adjustbox}
\end{center}

Il \cref{cap:24-footprint} misura l'effetto degli ultimi quattro. Che il
progetto li abbia resi tutti selezionabili da riga di comando \`e la condizione
che rende quella misura possibile.

\section{La build JVM}

Java 25, Spring Boot 4.1, Gradle con toolchain risolta automaticamente.
L'immagine OCI \`e costruita su Distroless con una tecnica che vale la pena
descrivere: il runtime Java non \`e un JRE standard ma uno costruito su misura
con \code{jlink}.

\begin{lstlisting}[language=bash, caption={Estratto dal \code{Dockerfile} del
control plane.}]
RUN jar xf app.jar BOOT-INF/lib BOOT-INF/classes && \
    jdeps --ignore-missing-deps --multi-release 25 --recursive \
      --print-module-deps --class-path 'BOOT-INF/lib/*' \
      BOOT-INF/classes app.jar > modules.txt && \
    jlink --add-modules "$(cat modules.txt),jdk.unsupported,jdk.crypto.ec,jdk.httpserver" \
      --strip-debug --no-header-files --no-man-pages --compress=zip-9 \
      --output /jre
\end{lstlisting}

\code{jdeps} determina i moduli JDK effettivamente raggiunti dall'applicazione,
e \code{jlink} ne assembla un runtime minimale. Tre moduli sono aggiunti a mano,
e il commento spiega perch\'e: \code{jdk.unsupported} (Netty usa
\code{sun.misc.Unsafe}), \code{jdk.crypto.ec} (le suite TLS a curve ellittiche
sono caricate via \code{ServiceLoader}) e \code{jdk.httpserver} (il proxy del
provider container). Sono tre casi in cui l'analisi statica di \code{jdeps} non
vede una dipendenza reale, perch\'e la dipendenza \`e risolta a runtime.

\`E lo stesso problema che la compilazione nativa affronta su scala molto
maggiore, e la sua comparsa gi\`a qui, in forma ridotta, \`e istruttiva.

\subsection{I flag della JVM}

\begin{lstlisting}[language=bash]
"-XX:+UseSerialGC", "-XX:MaxRAMPercentage=70", "-XX:TieredStopAtLevel=1",
"-Xss256k", "-Dspring.jmx.enabled=false", ...
\end{lstlisting}

Ciascuno rinuncia a qualcosa in cambio di impronta. \code{TieredStopAtLevel=1}
ferma la compilazione JIT al primo livello: il codice risultante \`e pi\`u lento
ma il compilatore lavora meno e la CodeHeap resta piccola. \code{-Xss256k}
riduce lo stack per thread da 1\,MB al quarto --- con ventitr\'e thread, come
misurato nel \cref{cap:24-footprint}, sono 17\,MB risparmiati.

Va detto che questa configurazione \`e ottimizzata per l'impronta e non per il
throughput di picco, e che le misure di throughput riportate nel
\cref{cap:24-footprint} sono state raccolte su una configurazione diversa. \`E
un'incoerenza da tenere presente nel leggerle.

\begin{damisurare}
Il costo in throughput di \code{TieredStopAtLevel=1} sul control plane non \`e
stato misurato. \`E l'unico flag della lista il cui effetto potrebbe essere
sostanziale sul percorso caldo.
\todomis{throughput con compilazione JIT completa contro C1 soltanto}
\end{damisurare}

\section{La build nativa}

\subsection{Perch\'e}

Tre motivi, in ordine di importanza per questo progetto.

\emph{Il tempo di avvio.} Da 1{,}585\,s a 0{,}072\,s misurati sul control plane
(\cref{cap:24-footprint}). Per un processo di lunga durata \`e ininfluente; per
la CLI \`e la differenza fra usabile e no; per un esperimento che avvia e ferma
il control plane decine di volte \`e tempo di ciclo.

\emph{L'impronta a riposo.} Da 133\,MiB a 25{,}6\,MiB su una funzione Spring.

\emph{L'eliminazione della fase di riscaldamento.} Un'immagine nativa \`e
compilata in anticipo e non ha una curva di riscaldamento JIT. Per un esperimento
questo elimina una variabile: non c'\`e da decidere quanto carico scartare prima
di iniziare a misurare.

\subsection{Il prezzo: la riflessione}

La compilazione anticipata richiede di conoscere staticamente ci\`o che verr\`a
caricato. Ogni uso della riflessione, ogni proxy dinamico, ogni risorsa caricata
per nome deve essere dichiarato. Il progetto lo fa con classi di \emph{runtime
hints} dedicate:

\begin{center}
\small
\begin{adjustbox}{max width=\textwidth}
\begin{tabular}{@{}ll@{}}
\toprule
\textbf{Classe} & \textbf{Copre} \\
\midrule
\code{CaffeineRuntimeHints} & La cache di idempotenza. \\
\code{VertxRuntimeHints} & Vert.x, trascinato dal client Fabric8. \\
\code{DockerJavaRuntimeHints} & Il client Docker Java. \\
\code{GcMetricsConfiguration} & I binder di metriche del collector. \\
\bottomrule
\end{tabular}
\end{adjustbox}
\end{center}

\`E significativo che tre delle quattro riguardino dipendenze dei \emph{moduli}
e non del nucleo. La compilabilit\`a nativa \`e un vincolo che si propaga alle
dipendenze transitive, ed \`e uno degli argomenti a favore della selezione
modulare a tempo di build (\cref{cap:08-sistema-moduli}): un binario che non
include il provider Kubernetes non ha bisogno degli hint di Vert.x.

Il rischio caratteristico \`e che un hint mancante non fa fallire la
compilazione: produce un binario che fallisce a runtime, sul primo percorso che
usa quella riflessione. \`E la ragione per cui esiste lo smoke test nativo della
CLI (\cref{cap:19-cli-scaffolding}).

\subsection{Il livello di ottimizzazione}

Il progetto usa \code{-O3} per default, e il commento nel file di build riporta
la misura che ha motivato il cambio da \code{-Os}:

\begin{center}
\small
\begin{adjustbox}{max width=\textwidth}
\begin{tabular}{@{}lrr@{}}
\toprule
 & \textbf{\code{-Os}} & \textbf{\code{-O3}} \\
\midrule
throughput & 8\,398 rps & \textbf{9\,251 rps} $(+10{,}2\%)$ \\
p95 & 15{,}64\,ms & 14{,}79\,ms \\
massimo & 148{,}7\,ms & 127{,}2\,ms \\
avvio & 0{,}173\,s & \textbf{0{,}068\,s} \\
residente a riposo & 120{,}6\,MiB & \textbf{39{,}3\,MiB} \\
dimensione immagine & 195\,MB & 397\,MB \\
\bottomrule
\end{tabular}
\end{adjustbox}
\end{center}

Il risultato \`e controintuitivo: il livello che ottimizza per la \emph{velocit\`a}
produce anche un'impronta a riposo tre volte pi\`u piccola, a fronte di
un'immagine doppia. Il \cref{cap:24-footprint} discute il meccanismo ---
l'inlining --- e la condizione sotto cui il guadagno svanisce.

Il commento nel build file chiude con l'avvertenza corretta: \emph{``The gain is
conditional [\ldots] under a 1GB limit the two were identical at $\sim$3\,490
requests per second, because there the bottleneck is collection rather than code
quality.''}

\subsection{Il garbage collector}

\begin{lstlisting}[language=Java]
// -PnativeGc=G1 needs Oracle GraalVM on Linux; Community offers only
// 'serial' and 'epsilon', and rejects anything else at build time.
// Measured motivation: the serial collector spent 32% of wall-clock
// time collecting under sustained load, with pauses reaching 1.9s.
if (project.hasProperty('nativeGc')) {
    binary.buildArgs.add("--gc=${project.property('nativeGc')}")
}
\end{lstlisting}

\`E la decisione di packaging pi\`u consequenziale del progetto, ed \`e trattata
in dettaglio nel \cref{cap:24-footprint}. In sintesi: GraalVM Community offre
solo un collector seriale, che su questo carico spende il 32\% del tempo di
parete a raccogliere; G1, disponibile in Oracle GraalVM, elimina il problema. La
selezione resta opt-in (\code{NATIVE\_GC=G1}) perch\'e Oracle GraalVM ha una
licenza diversa.

\subsection{Il budget di memoria del compilatore}

\begin{lstlisting}[language=Java]
// -PnativeBuildMemory=6g bounds the BUILDER's heap, not the built
// image's. native-image sizes its own heap from the memory it can
// see, which is the whole machine even when most of it is already
// spoken for: on a 12GB VM also running k3s, a control plane and
// Prometheus, it was SIGKILLed by the OOM killer after 9m50s of
// work. A build that dies that way wastes the ten minutes and says
// "exit 137"; one that is told its budget either fits or fails fast.
\end{lstlisting}

L'osservazione riguarda una classe generale di problemi: uno strumento che
dimensiona le proprie risorse dalla memoria \emph{visibile} anzich\'e da quella
\emph{disponibile} si comporta male su una macchina condivisa. Il parametro di
parallelismo esiste per lo stesso motivo, con il compromesso opposto: meno
lavoratori costano tempo di parete e risparmiano memoria di picco.

Entrambi sono decisioni che riguardano una macchina su cui la build compete con
il carico contro cui verr\`a poi misurata --- una situazione specifica di un
progetto sperimentale, e non affrontata dalla documentazione degli strumenti.

\section{Una dipendenza esclusa}

\begin{lstlisting}[language=Java]
configurations.all {
    // HTTP/3/QUIC support that reactor-netty pulls in automatically; nothing in this
    // codebase enables HTTP/3. Drops netty-codec-http3 plus its 5 platform-native quic jars.
    exclude group: 'io.netty', module: 'netty-codec-http3'
}
\end{lstlisting}

Sei archivi rimossi, di cui cinque contenenti codice nativo per piattaforma
specifica. Nel contesto della compilazione nativa e delle immagini Distroless,
una dipendenza inutilizzata che porta librerie native non \`e soltanto peso: \`e
superficie su cui gli hint di runtime potrebbero servire e non ci sono.

Il progetto utilizza il plugin \code{dependency-analysis} per individuare
sistematicamente dipendenze non usate o non dichiarate
(\cref{cap:26-qualita-verifica}).

\section{La matrice degli artefatti}

\begin{center}
\small
\begin{adjustbox}{max width=\textwidth}
\begin{tabular}{@{}lll@{}}
\toprule
\textbf{Artefatto} & \textbf{Build JVM} & \textbf{Build nativa} \\
\midrule
control plane & \code{Dockerfile} + Distroless & \code{native-java-image.sh control-plane} \\
\code{warm-echo} & \code{Dockerfile} + Distroless & \code{native-java-image.sh warm-echo} \\
funzioni Java & per funzione & \code{native-java-image.sh <nome>} \\
CLI & \code{installDist} & \code{nativeCompile} \\
watchdog & --- & Rust statico musl, \code{FROM scratch} \\
\bottomrule
\end{tabular}
\end{adjustbox}
\end{center}

Un unico comando costruisce e verifica l'intera matrice nativa:

\begin{lstlisting}[language=bash]
scripts/native-build.sh
\end{lstlisting}

Le immagini native non usano i buildpack di Spring Boot ma un
\code{Dockerfile} che copia il binario su Distroless. La scelta \`e coerente con
il resto: un buildpack decide da s\'e i flag di compilazione, e questo progetto
ha bisogno che quei flag siano una variabile sperimentale esplicita.


\part{Osservabilit\`a e valutazione sperimentale}
\chapter{Osservabilit\`a}
\label{cap:21-osservabilita}

In una piattaforma pensata come strumento di misura, l'osservabilit\`a non \`e un
corredo operativo: \`e l'apparato sperimentale. Ogni risultato riportato nella
Parte V \`e stato ottenuto leggendo le metriche descritte in questo capitolo, e
diverse decisioni di progetto della Parte III sono state prese sulla base di
grandezze che qui sono definite.

\section{Il modello}

Le metriche sono Micrometer esportate in formato Prometheus~\cite{prometheus} su
\code{/actuator/prometheus}, porta \code{8081}. La separazione dalla porta
dell'API (\cref{cap:04-architettura-control-plane}) fa s\`i che lo scraping non
condivida code n\'e limiti di ammissione con il traffico misurato.

La quasi totalit\`a delle metriche \`e etichettata per funzione. Il ciclo di vita
dei meter segue quello della funzione: creati alla registrazione, rimossi alla
cancellazione (\cref{cap:06-nucleo}).

\section{Il catalogo}

\subsection{Nucleo}

\begin{center}
\small
\begin{adjustbox}{max width=\textwidth}
\begin{tabular}{@{}llp{6.2cm}@{}}
\toprule
\textbf{Metrica} & \textbf{Tipo} & \textbf{Significato} \\
\midrule
\code{function\_enqueue\_total} & counter & Accodamenti. \\
\code{function\_dispatch\_total} & counter & Dispatch verso il runtime.
Sorgente della metrica \code{rps} dell'autoscaler. \\
\code{function\_success\_total} & counter & Esiti positivi. \\
\code{function\_error\_total} & counter & Esiti negativi. \\
\code{function\_retry\_total} & counter & Riprove. \\
\code{function\_latency\_ms} & timer & \textbf{Tempo di servizio}: da dispatch a
completamento. \\
\code{function\_e2e\_latency\_ms} & timer & Tempo di sojourn: da ammissione a
completamento. \\
\code{function\_queue\_wait\_ms} & timer & Attesa in coda. \\
\code{function\_cold\_start\_ms} & timer & Durata dell'inizializzazione. \\
\code{function\_init\_duration\_ms} & timer & Come sopra, dal runtime. \\
\bottomrule
\end{tabular}
\end{adjustbox}
\end{center}

La distinzione fra le prime tre metriche temporali \`e la pi\`u importante
dell'intero capitolo, ed \`e il fondamento del \cref{cap:12-concurrency-control}:

\[
  \underbrace{\code{function\_e2e\_latency\_ms}}_{\text{sojourn}}
  \;=\;
  \underbrace{\code{function\_queue\_wait\_ms}}_{\text{attesa}}
  \;+\;
  \underbrace{\code{function\_latency\_ms}}_{\text{servizio}} .
\]

Che le tre esistano separatamente \`e ci\`o che ha reso possibile stabilire che
la seconda \`e sei volte la terza sotto carico, e quindi che un controllore che
decide dalla terza sta ottimizzando la grandezza sbagliata.

\subsection{Code}

\begin{center}
\small
\begin{adjustbox}{max width=\textwidth}
\begin{tabular}{@{}llp{6.2cm}@{}}
\toprule
\textbf{Metrica} & \textbf{Tipo} & \textbf{Fornita da} \\
\midrule
\code{function\_queue\_depth} & gauge & \code{async-queue} \\
\code{function\_inFlight} & gauge & \code{async-queue} \\
\code{function\_effective\_concurrency} & gauge & \code{async-queue} \\
\code{sync\_queue\_depth} & gauge & \code{sync-queue}, globale e per funzione \\
\code{sync\_queue\_wait\_seconds} & gauge & \code{sync-queue}, attesa stimata \\
\code{sync\_queue\_admitted\_total} & counter & \code{sync-queue} \\
\code{sync\_queue\_rejected\_total} & counter & \code{sync-queue} \\
\code{sync\_queue\_timedout\_total} & counter & \code{sync-queue} \\
\bottomrule
\end{tabular}
\end{adjustbox}
\end{center}

\subsection{Politiche}

\begin{center}
\small
\begin{adjustbox}{max width=\textwidth}
\begin{tabular}{@{}llp{6cm}@{}}
\toprule
\textbf{Metrica} & \textbf{Tipo} & \textbf{Fornita da} \\
\midrule
\code{function\_target\_inflight\_per\_pod} & gauge & \code{concurrency-control} \\
\code{function\_concurrency\_controller\_mode} & gauge &
\code{concurrency-control} \\
\code{nanofaas\_offload\_total} & counter & \code{offload}, con etichetta
\code{trigger} \\
\code{nanofaas\_offload\_failure\_total} & counter & \code{offload} \\
\code{controlplane\_runtime\_config\_revision} & gauge & \code{runtime-config} \\
\bottomrule
\end{tabular}
\end{adjustbox}
\end{center}

Nessuna di queste alimenta una decisione: sono tutte \emph{descrittive della
decisione}. \`E una categoria che vale la pena isolare concettualmente, perch\'e
risponde a una domanda che le metriche di stato non possono risolvere: dato un
comportamento osservato, quale politica era attiva e che cosa aveva deciso.

\section{I due profili}

\begin{lstlisting}[language=Java]
private static final Set<String> ADVANCED_METRICS = Set.of(
        "function_cold_start_total", "function_warm_start_total",
        "function_effective_concurrency", "function_target_inflight_per_pod",
        "function_concurrency_controller_mode", "gateway_service_target_load",
        "controlplane_runtime_config_revision", ...);
\end{lstlisting}

Il profilo \code{basic} (default) esporta il nucleo. Il profilo
\code{advanced} aggiunge le metriche descrittive delle politiche e --- pi\`u
significativamente --- attiva gli istogrammi per percentili sui quattro timer
per funzione.

Il costo che i due profili arbitrano \`e quello degli istogrammi. Un timer
Micrometer senza istogramma mantiene conteggio, somma e massimo: tre valori. Con
gli istogrammi mantiene una serie di secchi, e ogni secchio \`e una serie
temporale in Prometheus. Per una piattaforma con molte funzioni la differenza
nella cardinalit\`a esportata \`e sostanziale.

La conseguenza pratica \`e che i percentili --- p95, p99 --- sono disponibili
soltanto sotto \code{advanced}. Poich\'e i risultati della Parte V sono espressi
prevalentemente in percentili, tutte le campagne sperimentali girano su quel
profilo.

Il profilo \`e validato all'avvio: un valore diverso da \code{basic} o
\code{advanced} fa fallire il contesto con un messaggio esplicito, anzich\'e
ripiegare in silenzio su un default.

\section{Le letture canoniche}

La documentazione operativa del progetto raccoglie le interrogazioni PromQL di
riferimento. Le pi\`u istruttive sono quelle che \emph{combinano} due metriche,
perch\'e ciascuna combinazione codifica una diagnosi.

\begin{lstlisting}[language=bash, caption={Tasso di errore per funzione.}]
rate(function_error_total{function="word-stats"}[5m])
  / (rate(function_success_total{function="word-stats"}[5m])
     + rate(function_error_total{function="word-stats"}[5m])) * 100
\end{lstlisting}

\begin{lstlisting}[language=bash, caption={Pressione di ammissione sincrona.}]
rate(sync_queue_rejected_total{function="word-stats"}[5m])
  / rate(sync_queue_admitted_total{function="word-stats"}[5m])
\end{lstlisting}

\begin{lstlisting}[language=bash, caption={Carico effettivo per replica: il
segnale su cui l'autoscaler decide.}]
rate(function_dispatch_total{function="word-stats"}[1m])
  / max by (function) (kube_deployment_status_replicas{deployment="fn-word-stats"})
\end{lstlisting}

\subsection{La guida alla lettura della contesa}

Il progetto documenta quattro regole di interpretazione che meritano di essere
riportate, perch\'e sono il tipo di conoscenza che di norma resta implicita:

\begin{itemize}[leftmargin=*]
  \item \code{sync\_queue\_depth} crescente con \code{function\_dispatch\_total}
  piatto significa che l'ammissione riesce pi\`u in fretta di quanto gli slot si
  riaprano.
  \item \code{sync\_queue\_rejected\_total} alto con profondit\`a bassa indica un
  rifiuto per attesa stimata, non un esaurimento di capacit\`a.
  \item \code{function\_dispatch\_total} che cresce mentre successi ed errori
  restano indietro indica latenza di completamento, non iniquit\`a dello
  scheduler.
  \item Profondit\`a persistente con poche richieste in volo indica una funzione
  limitata dagli slot; profondit\`a persistente con molte richieste in volo
  indica un runtime lento.
\end{itemize}

L'ultima coppia \`e la pi\`u utile e non \`e ovvia: la stessa profondit\`a di coda
ha due cause opposte, distinguibili solo confrontandola con il contatore delle
richieste in volo.

\section{Il tasso di errore \`e ambiguo, e il progetto lo disambigua}

Un problema ricorrente nel misurare piattaforme FaaS: il tasso di errore
mescola i fallimenti della piattaforma con le risposte non-2xx deliberate delle
funzioni. Una funzione che restituisce correttamente \code{404} per una risorsa
inesistente non \`e un errore della piattaforma, ma da fuori \`e indistinguibile.

Il meccanismo della busta di risposta (\cref{cap:17-runtime-funzione}) risolve
l'ambiguit\`a alla radice: una risposta decisa dall'handler porta
\code{X-NanoFaaS-Function-Status: true} ed \`e contabilizzata come successo dalla
piattaforma, qualunque sia il suo codice di stato. \code{function\_error\_total}
conta quindi soltanto i fallimenti che la piattaforma si attribuisce.

\`E una precondizione per il gate sperimentale del
\cref{cap:23-prestazioni-base}, che richiede tasso di errore esattamente zero:
un tale gate sarebbe inapplicabile se il contatore includesse le risposte
applicative negative.

\section{Log e tracing}

I log sono strutturati con cinque campi: \code{executionId},
\code{functionName}, \code{traceId}, \code{attempt}, \code{status}. Il tracing
\`e la propagazione di \code{X-Trace-Id} lungo l'intera catena, con supporto per
i contesti W3C (\code{traceparent}, \code{tracestate}) sul percorso di offload.

Una regola \`e enunciata esplicitamente e vale la pena riportarla:

\begin{quote}\itshape
Il runtime non sintetizza identificatori di esecuzione segnaposto per il
logging. Se n\'e l'header della richiesta n\'e \code{EXECUTION\_ID} sono
disponibili, il contesto diagnostico resta vuoto per \code{executionId}.
\end{quote}

Un identificatore inventato riempirebbe il campo e renderebbe la correlazione
\emph{apparentemente} possibile ma sbagliata. Un campo vuoto \`e onesto.

\subsection{La consegna delle callback \`e limitata e pu\`o essere scartata}

Un dettaglio con conseguenze sperimentali. La consegna della callback di
completamento dal runtime Java \`e asincrona e limitata: \code{/invoke}
restituisce senza attenderla, e quando il dispatcher \`e saturo la callback viene
\emph{scartata} con un avviso. Il runtime Go si comporta allo stesso modo e
incrementa una metrica di callback scartate.

Una callback scartata produce un'esecuzione che non raggiunge mai uno stato
terminale --- esattamente il caso che il terzo orizzonte di eviction dello
\code{ExecutionStore} deve ripulire (\cref{cap:06-nucleo}). Che il runtime Go
esponga un contatore dedicato e quello Java soltanto un avviso nel log \`e
un'asimmetria che rende il fenomeno misurabile su un SDK e non sull'altro.

Le riprove della callback sono limitate ai fallimenti ritentabili: errori di
trasporto, \code{408}, \code{429} e \code{5xx}. Gli altri \code{4xx} sono
considerati permanenti. \`E la distinzione corretta --- un \code{400} sulla
callback significa che il control plane ha rifiutato il corpo, e ripeterlo non
lo far\`a accettare.

\section{Salute}

Due sonde: \code{/actuator/health/liveness} e
\code{/actuator/health/readiness}, entrambe sulla porta di management. Sono la
base delle \code{livenessProbe} e \code{readinessProbe} del Deployment del
control plane.

\section{Copertura contro le regressioni}

Il progetto adotta una posizione metodologica che vale la pena riportare per
esteso:

\begin{quote}\itshape
Il repository include test strutturali di regressione sul percorso caldo anzich\'e
microbenchmark assoluti. Asseriscono progresso e comportamenti sensibili
all'allocazione, come il riuso del replay, il forward progress della coda
sincrona dietro una testa bloccata, e l'equit\`a asincrona fra funzioni calde e
fredde. [\ldots] Nel tarare il comportamento delle code, si preferisca
preservare quelle garanzie strutturali anzich\'e inseguire un numero fisso di
temporizzazione locale. Le temporizzazioni assolute sono sensibili all'ambiente,
le garanzie di equit\`a e riuso non lo sono.
\end{quote}

\`E la risposta corretta a un problema noto dei test di prestazione: un test che
asserisce ``questa operazione impiega meno di $X$ millisecondi'' fallisce sulla
macchina di un collega, e viene disattivato. Un test che asserisce ``dietro una
testa bloccata, un'altra funzione riesce comunque a progredire'' cattura la
propriet\`a che interessava, ed \`e vero su qualunque macchina.

\chapter{Metodo sperimentale}
\label{cap:22-metodo-sperimentale}

I capitoli seguenti riportano misure. Questo descrive come sono state ottenute,
e --- pi\`u importante --- quali propriet\`a dell'impianto rendono confrontabili
due numeri raccolti in momenti diversi.

\section{Il principio: comparabilit\`a nel tempo, non fra piattaforme}

Il \cref{cap:02-stato-dell-arte} ha osservato che confrontare piattaforme FaaS
misurate da persone diverse su hardware diverso \`e notoriamente inaffidabile.
Il metodo di \nf{} non tenta quel confronto. Persegue invece la comparabilit\`a
\emph{intra}-piattaforma: fissati un profilo hardware e un carico, il risultato
di una release \`e confrontato con quello della release precedente sullo stesso
profilo, e la deviazione \`e trattata come segnale di regressione.

\`E un obiettivo pi\`u ristretto e pi\`u difendibile. Non consente di affermare
che \nf{} sia pi\`u veloce di Knative; consente di affermare che una certa
modifica ha cambiato il throughput del $x\%$ su un profilo dichiarato, che \`e
l'affermazione di cui il progetto ha effettivamente bisogno.

\section{L'impianto: uno strumento separato}

L'orchestrazione degli esperimenti vive in un repository separato. La
separazione \`e deliberata e ha una giustificazione di metodo: uno strumento che
provvisiona macchine, installa Kubernetes, distribuisce immagini ed esegue
carichi \`e un sistema a s\'e, con dipendenze proprie, e includerlo nel
repository della piattaforma renderebbe la piattaforma pi\`u difficile da
costruire per chi vuole solo usarla.

Il confine \`e quello enunciato nel \cref{cap:19-cli-scaffolding}: la CLI \`e un
client HTTP e non provvisiona nulla; lo strumento di provisioning non conosce il
protocollo delle funzioni.

\subsection{Scenari dichiarativi}

Un esperimento \`e descritto da uno scenario YAML e da un file di ambiente. Lo
scenario dichiara \emph{che cosa} accade --- provvisiona, installa, registra,
carica, verifica --- e l'ambiente dichiara \emph{dove}.

\begin{lstlisting}[language=bash]
# esecuzione
./nanolab.sh run packages/nanolab/scenarios-v2/deployment-lifecycle-k8s.yaml \
    --environment packages/nanolab/environments/multipass.yaml
# anteprima senza eseguire
./nanolab.sh plan packages/nanolab/scenarios-v2/deployment-lifecycle-k8s.yaml \
    --environment packages/nanolab/environments/local.yaml
\end{lstlisting}

La validazione end-to-end della piattaforma appartiene a questo strumento e non
al repository di \nf{}: la suite di quest'ultimo si ferma ai test unitari e di
integrazione, e il ciclo di vita completo --- registrazione, invocazione,
verifiche HTTP e Kubernetes --- \`e uno scenario
(\cref{cap:26-qualita-verifica}). Le famiglie di scenari sono declinate per
backend di esecuzione: lo stesso ciclo di vita esiste per il backend container,
per containerd e per Kubernetes (\code{deployment-lifecycle-container},
\code{-containerd}, \code{-k8s}), cos\`i che un cambiamento di backend non
richieda di riscrivere l'esperimento; al 2026-09-28 la cartella
\code{scenarios-v2} contiene 70 file.

L'esistenza di una modalit\`a \code{plan} \`e significativa: un esperimento che
provvisiona infrastruttura a pagamento non dovrebbe scoprire un errore di
configurazione dopo dieci minuti di provisioning.

\subsection{Ambienti disponibili}

\begin{center}
\small
\begin{adjustbox}{max width=\textwidth}
\begin{tabular}{@{}lp{7.6cm}@{}}
\toprule
\textbf{Ambiente} & \textbf{Uso} \\
\midrule
locale (Docker) & Validazione del container, cicli rapidi. \\
Multipass & VM locale con k3s: percorso Kubernetes completo senza costo. \\
Azure & Profilo fissato per le misure pubblicate; VM dedicate per stack e
generatore di carico. \\
Proxmox & Ambiente su hardware proprio. \\
\bottomrule
\end{tabular}
\end{adjustbox}
\end{center}

\section{Il generatore di carico}

Il carico \`e generato con k6~\cite{k6}. Due propriet\`a dell'impianto meritano
attenzione perch\'e hanno influenzato i risultati.

\subsection{Corpora di payload versionati}

L'input dei benchmark non \`e generato casualmente ma letto da corpora
versionati (\code{functions/test-data/<famiglia>/performance-<profilo>.json}),
con tre profili di dimensione. Due esecuzioni della stessa configurazione
ricevono quindi esattamente lo stesso input.

\`E la stessa distinzione discussa nel \cref{cap:18-sdk}: i corpora di
prestazione omettono deliberatamente gli output attesi, perch\'e verificarli
durante un benchmark aggiungerebbe alla misura il costo del confronto.

\subsection{Ciclo chiuso e ciclo aperto}

Questa \`e la distinzione metodologicamente pi\`u importante del capitolo, e la
sua rilevanza \`e stata scoperta a caro prezzo.

In un generatore a \textbf{ciclo chiuso}, ogni utente virtuale mantiene una sola
richiesta in volo e ne invia un'altra solo dopo aver ricevuto la risposta. Il
numero di richieste nel sistema \`e quindi fissato dal generatore. In un
generatore a \textbf{ciclo aperto}, gli arrivi sono programmati dall'orologio
indipendentemente da quanto il sistema stia impiegando a rispondere.

La differenza \`e cruciale per valutare un controllore di concorrenza:

\begin{quote}\itshape
Sotto il profilo a raffica ogni utente virtuale tiene una richiesta, quindi il
numero nel sistema \`e fissato dal generatore e la profondit\`a di coda \`e
\code{VU} $-$ \code{limite}. Due conseguenze, entrambe misurate. Portare il
limite da 4 a 8 ha accorciato di quattro una coda profonda 99 --- nessun
controllore pu\`o muovere l'attesa di molto. E la coda si assesta a 95 su 100 per
costruzione, permanentemente sopra qualunque soglia di guardia, quindi la
guardia di ammissione scatta a ogni tick e il modello non parla mai.
\end{quote}

La conseguenza \`e enunciata senza mezzi termini: \emph{``\`E per questo che due
controllori che leggevano segnali completamente diversi sono finiti entro il 2\%
l'uno dall'altro su questo profilo. La somiglianza era una propriet\`a
dell'impianto.''}

\`E un risultato metodologico che va oltre questo progetto. Un carico a ciclo
chiuso non pu\`o porre al sistema la domanda a cui un controllore di ammissione
dovrebbe rispondere, perch\'e retroaziona sulla stessa grandezza che il
controllore manipola. Il \cref{cap:25-valutazione-concorrenza} riporta i due
insiemi di risultati, con e senza questa correzione.

\section{Profili hardware fissati}

Ogni record di prestazione porta l'identificatore del profilo su cui \`e stato
ottenuto:

\begin{lstlisting}[language=bash]
azure-d8s-v5+d2s-v5-amd64-native-loadtest-v1
\end{lstlisting}

L'identificatore codifica: provider, tipo di VM dello stack, tipo di VM del
generatore, architettura, tipo di build del control plane, e una versione del
profilo stesso. Un confronto fra record di profili diversi \`e vietato per
costruzione.

Il generatore di carico gira su una VM \emph{separata} da quella dello stack. \`E
una precauzione elementare e spesso omessa: un generatore co-residente compete
per CPU con il sistema misurato, e il suo contributo alla latenza osservata
cresce proprio quando il sistema \`e sotto pressione.

\section{Il record di release}

Ogni release produce un record JSON versionato in
\code{docs/performance/releases/}. Contiene:

\begin{center}
\small
\begin{adjustbox}{max width=\textwidth}
\begin{tabular}{@{}lp{7.6cm}@{}}
\toprule
\textbf{Sezione} & \textbf{Contenuto} \\
\midrule
\code{aggregates} & Throughput, tasso di errore, p50/p95/p99, cold start,
repliche di picco, CPU e heap di picco del control plane, attesa in coda. \\
\code{profile} & L'identificatore del profilo e le sue componenti. \\
\code{runCount} & Numero di esecuzioni la cui mediana \`e riportata (3). \\
\code{sourceCommit} & Il commit esatto da cui l'artefatto \`e stato costruito. \\
\code{imageDigests} & Digest SHA-256 di \emph{ogni} immagine coinvolta. \\
\code{thresholds} & Le soglie di regressione applicate. \\
\bottomrule
\end{tabular}
\end{adjustbox}
\end{center}

I digest sono la parte che rende il record realmente riproducibile: non
identificano un'etichetta mobile ma il contenuto esatto delle immagini. Un
record di \nf{} v0.19.0 elenca oltre centoventi digest --- il control plane in
quattro varianti, il watchdog, e le funzioni di riferimento in cinque linguaggi
per due architetture.

\section{Il gate di regressione}

Tre esecuzioni per release, mediane confrontate con il record pi\`u recente dello
stesso profilo. Le soglie:

\begin{center}
\small
\begin{adjustbox}{max width=\textwidth}
\begin{tabular}{@{}lrl@{}}
\toprule
\textbf{Soglia} & \textbf{Valore} & \textbf{Significato} \\
\midrule
\code{errorRateMax} & 0{,}3\,\% & Tasso di errore massimo assoluto. \\
\code{p95MaxIncreasePercent} & 15\,\% & Peggioramento massimo del p95 rispetto
al record precedente. \\
\code{throughputMaxLossPercent} & 10\,\% & Perdita massima di throughput. \\
\bottomrule
\end{tabular}
\end{adjustbox}
\end{center}

Le due soglie relative sono generose --- 15\% e 10\% --- e la ragione \`e onesta:
con tre esecuzioni la varianza campionaria non consente di distinguere
differenze pi\`u piccole. Una soglia stretta produrrebbe falsi allarmi a ogni
release, e il gate verrebbe disattivato. Il \cref{cap:27-discussione} riprende
il punto come limite alla sensibilit\`a del metodo.

Accanto alle soglie numeriche esiste un gate strutturale, gi\`a citato nel
\cref{cap:21-osservabilita}: una modifica che regredisca le garanzie
strutturali del percorso caldo --- riuso del replay, forward progress della coda
sincrona, equit\`a asincrona --- viene respinta anche se i numeri assoluti
sembrano accettabili.

\section{Un errore di calibrazione, riportato}

La documentazione del progetto conserva un errore che vale la pena riportare
perch\'e \`e generale:

\begin{quote}\itshape
La prima esecuzione di questo confronto ha fissato il budget a 8 mentre il carico
raggiungeva un picco di 105 contro una coda di 100. Sotto contesa ciascuna
funzione teneva circa 4 posti di 104 contro 105 arrivi, quindi la coda traboccava
per aritmetica in un modo e non nell'altro, producendo 86\,302 rifiuti che non
dicevano nulla sul controllore.
\end{quote}

E la regola che ne \`e stata derivata:

\begin{quote}\itshape
Il budget deve stare dentro una finestra: a o sopra \emph{funzioni $\times$
tetto} nulla \`e mai scarso e il modo degenera in controllo per funzione; sotto
\emph{funzioni $\times$ (picco $-$ coda)} i rifiuti sono fabbricati. Il picco
offerto deve essere derivato dal limite pi\`u piccolo che qualunque modo sotto
confronto conceder\`a, o le due esecuzioni non stanno misurando la stessa cosa.
\end{quote}

\`E un vincolo che non compare in nessuna documentazione di benchmarking
serverless, e che discende dalla struttura del sistema misurato: quando la
grandezza controllata determina anche la finestra di ammissione, il carico
offerto non pu\`o essere scelto indipendentemente dal controllore.

\section{Riproducibilit\`a: che cosa \`e garantito e che cosa no}

\begin{center}
\small
\begin{adjustbox}{max width=\textwidth}
\begin{tabular}{@{}lp{7.8cm}@{}}
\toprule
\textbf{Elemento} & \textbf{Stato} \\
\midrule
Codice sorgente & Fissato dal commit nel record. \\
Immagini & Fissate dai digest nel record. \\
Input del carico & Fissato dai corpora versionati. \\
Profilo hardware & Fissato dall'identificatore; la VM \`e ricreata a ogni
esecuzione. \\
Prestazioni della VM & \textbf{Non garantite}: due VM Azure dello stesso tipo
possono avere vicini diversi. \\
Numero di ripetizioni & 3, che \`e poco. \\
\bottomrule
\end{tabular}
\end{adjustbox}
\end{center}

Le ultime due righe sono i limiti reali. Il \cref{cap:27-discussione} li tratta
per esteso, ma \`e utile enunciarli qui, accanto al metodo che pretendono di
qualificare: le misure riportate nei prossimi tre capitoli sono \emph{mediane di
tre esecuzioni su infrastruttura condivisa}, non stime con intervallo di
confidenza.

\chapter{Prestazioni di base}
\label{cap:23-prestazioni-base}

Questo capitolo riporta le prestazioni della piattaforma sul carico di
riferimento, attraverso sei release consecutive, sul profilo hardware fissato
descritto nel \cref{cap:22-metodo-sperimentale}.

\section{La serie}

Tutte le misure che seguono sono mediane di tre esecuzioni sul profilo
\code{azure-d8s-v5+d2s-v5-amd64-native-loadtest-v1}: control plane compilato
nativamente per AMD64 su VM \code{Standard\_D8s\_v5}, generatore di carico k6 su
VM separata \code{Standard\_D2s\_v5}, scenario \code{autoscaling-cycle-k8s}.

\begin{table}[htbp]
\centering
\caption{Serie di prestazioni per release. Mediane di tre esecuzioni sullo
stesso profilo.}
\label{tab:serie}
\small
\begin{adjustbox}{max width=\textwidth}
\begin{tabular}{@{}lrrrr@{}}
\toprule
\textbf{Versione} & \textbf{Throughput (req/s)} & \textbf{Tasso errore} &
\textbf{p95 (ms)} & \textbf{Repliche di picco} \\
\midrule
0.17.0  & 241{,}90 & 0{,}000 & 6{,}50  & 3 \\
0.17.1  & 269{,}07 & 0{,}000 & 5{,}77  & 3 \\
v0.18.1 & 239{,}13 & 0{,}000 & 6{,}61  & 4 \\
v0.18.2 & 249{,}34 & 0{,}000 & 5{,}13  & 3 \\
v0.18.3 & 232{,}41 & 0{,}000 & 5{,}13  & 4 \\
v0.19.0 & \textbf{300{,}48} & 0{,}000 & \textbf{11{,}26} & 5 \\
\bottomrule
\end{tabular}
\end{adjustbox}
\end{table}

\begin{figure}[htbp]
  \centering
  \resizebox{0.95\textwidth}{!}{% Release performance series, from docs/performance/history.md (medians of 3 runs
% on profile azure-d8s-v5+d2s-v5-amd64-native-loadtest-v1).
\documentclass[border=8pt]{standalone}
\usepackage{nanofaas-figs}
\usepackage{pgfplots}
\pgfplotsset{compat=1.18}
\begin{document}
\begin{tikzpicture}
\begin{axis}[
    width=13cm, height=7cm,
    axis y line*=left, axis x line*=bottom,
    xlabel={release}, ylabel={throughput (req/s)},
    ylabel style={nfblue}, y tick label style={nfblue},
    symbolic x coords={0.17.0,0.17.1,v0.18.1,v0.18.2,v0.18.3,v0.19.0},
    xtick=data, ymin=0, ymax=340,
    font=\sffamily\small, tick label style={font=\sffamily\footnotesize},
    legend style={draw=none, font=\sffamily\footnotesize, at={(0.02,0.98)}, anchor=north west},
    ybar, bar width=16pt,
]
  \addplot[fill=nfbluebg, draw=nfblue, line width=0.8pt] coordinates {
    (0.17.0,241.904) (0.17.1,269.069) (v0.18.1,239.130)
    (v0.18.2,249.341) (v0.18.3,232.407) (v0.19.0,300.483)};
  \addlegendentry{throughput}
\end{axis}
\begin{axis}[
    width=13cm, height=7cm,
    axis y line*=right, axis x line=none,
    ylabel={p95 (ms)}, ylabel style={nfamber}, y tick label style={nfamber},
    symbolic x coords={0.17.0,0.17.1,v0.18.1,v0.18.2,v0.18.3,v0.19.0},
    ymin=0, ymax=13, font=\sffamily\small,
    tick label style={font=\sffamily\footnotesize},
    legend style={draw=none, font=\sffamily\footnotesize, at={(0.02,0.86)}, anchor=north west},
]
  \addplot[nfamber, line width=1.3pt, mark=*, mark size=2.2pt] coordinates {
    (0.17.0,6.496) (0.17.1,5.772) (v0.18.1,6.607)
    (v0.18.2,5.125) (v0.18.3,5.131) (v0.19.0,11.256)};
  \addlegendentry{p95}
\end{axis}
\end{tikzpicture}
\end{document}
}
  \caption{Throughput (barre, asse sinistro) e p95 (linea, asse destro) per
  release. Dati da \code{docs/performance/history.md}.}
  \label{fig:serie}
\end{figure}

\section{Che cosa la serie mostra}

\subsection{Il tasso di errore \`e esattamente zero}

Su tutte e sei le release, su tre esecuzioni ciascuna. Non ``prossimo a zero'':
zero. \`E il gate rigido del flusso di rilascio, e la sua tenuta \`e il
risultato pi\`u solido del capitolo --- l'unico che non dipende dalla
sensibilit\`a del metodo, perch\'e distingue fra zero e non-zero anzich\'e fra
due valori vicini.

Va ricordato (\cref{cap:21-osservabilita}) che questo contatore esclude le
risposte non-2xx deliberate delle funzioni, grazie al meccanismo della busta.
Un gate a zero sarebbe inapplicabile senza quella disambiguazione.

\subsection{Il throughput oscilla in una banda, poi sale}

Le prime cinque release stanno fra 232 e 269 richieste al secondo: una banda del
$\pm 7\%$ attorno a 250. Con tre esecuzioni per punto, \emph{questa banda non
\`e distinguibile dal rumore}. Nessuna delle differenze fra 0.17.0, v0.18.1,
v0.18.2 e v0.18.3 pu\`o essere attribuita a una modifica del codice sulla base
di questi dati.

La v0.19.0 esce dalla banda: 300{,}5 req/s, ossia $+29\%$ sulla release
precedente e $+12\%$ sul massimo precedente. Questa differenza \`e maggiore di
qualunque oscillazione osservata prima, e supera la soglia del gate.

\subsection{Il p95 della v0.19.0 \`e il dato da spiegare}

Nella stessa release il p95 passa da 5{,}13\,ms a 11{,}26\,ms: pi\`u che
raddoppiato, e ben oltre la soglia di regressione del 15\%. Il numero di
repliche di picco sale da 4 a 5.

L'interpretazione che i dati sostengono \`e che la v0.19.0 abbia spostato il
sistema in un punto di lavoro diverso, non che sia peggiorata. Un throughput di
300 req/s su cinque repliche significa un carico per replica confrontabile con
quello precedente; il p95 pi\`u alto \`e coerente con un sistema che opera pi\`u
vicino alla propria capacit\`a. Il record di dettaglio della v0.19.0 lo conferma:

\begin{center}
\small
\begin{adjustbox}{max width=\textwidth}
\begin{tabular}{@{}lr@{}}
\toprule
\textbf{Grandezza} & \textbf{v0.19.0} \\
\midrule
p50 & 4{,}43\,ms \\
p95 & 11{,}26\,ms \\
p99 & 21{,}05\,ms \\
attesa in coda (media) & 0{,}82\,ms \\
cold start osservati & 5 \\
CPU di picco del control plane & 0{,}664 core \\
heap di picco del control plane & 199{,}5\,MB \\
\bottomrule
\end{tabular}
\end{adjustbox}
\end{center}

Il p50 di 4{,}43\,ms \`e in linea con le release precedenti: il tipico non \`e
peggiorato, si \`e allungata la coda della distribuzione. L'attesa in coda media
\`e sotto il millisecondo, quindi la coda non \`e la causa.

\begin{damisurare}
L'attribuzione del salto di p95 della v0.19.0 non \`e stata condotta. I dati
disponibili escludono l'attesa in coda come causa dominante e sono coerenti con
un punto di lavoro pi\`u vicino alla saturazione, ma non lo stabiliscono. Un
confronto a parit\`a di throughput offerto fra v0.18.3 e v0.19.0 lo
deciderebbe.
\todomis{p95 di v0.18.3 e v0.19.0 a parit\`a di carico offerto}
\end{damisurare}

\subsection{L'osservabilit\`a del control plane}

Due grandezze del record riguardano il control plane e non le funzioni: CPU di
picco 0{,}664 core e heap di picco 199{,}5\,MB. Sono la risposta quantitativa
alla domanda del \cref{cap:01-introduzione} su quanto costi un control plane
minimale: meno di un core e circa 200\,MB di heap per sostenere 300 richieste al
secondo.

Va sottolineato che si tratta di \emph{heap}, non di memoria residente. Il
\cref{cap:24-footprint} mostra che per un processo Spring su JVM la memoria
residente \`e diverse volte l'heap, perch\'e la maggior parte non \`e dati.

\section{Gli SLO indicativi}

Il progetto dichiara obiettivi di lavoro, esplicitamente non contrattuali:

\begin{center}
\small
\begin{adjustbox}{max width=\textwidth}
\begin{tabular}{@{}lp{7.6cm}@{}}
\toprule
\textbf{Obiettivo} & \textbf{Stato} \\
\midrule
Tasso di errore 0\% & Rispettato in tutte e sei le release; \`e un gate rigido. \\
p95 $\leq$ 10\,ms & Rispettato fino a v0.18.3; \textbf{violato} in v0.19.0
(11{,}26\,ms). \\
Throughput $\geq$ 200 req/s & Rispettato in tutte le release. \\
Cold start & Nessun budget fisso; si traccia la distribuzione. \\
\bottomrule
\end{tabular}
\end{adjustbox}
\end{center}

La scelta di non fissare un budget per il cold start ma di tracciarne la
distribuzione \`e coerente con l'analisi del \cref{cap:02-stato-dell-arte}: il
cold start su \nf{} \`e dominato dal tempo di creazione del pod da parte di
Kubernetes e dall'avvio del runtime, entrambi fuori dal controllo del control
plane. Un budget su una grandezza che non si governa \`e un obiettivo che si
manca senza poter reagire.

\begin{quote}
\textbf{Nota sulla documentazione.} Il documento \code{docs/slo.md} riporta la
serie fino a v0.18.3 e non include la v0.19.0, che \`e presente in
\code{docs/performance/history.md} e nel record di release. La discrepanza \`e
segnalata qui perch\'e la violazione dell'obiettivo di p95 \`e visibile solo nella
seconda fonte. \`E un esempio del genere di deriva che colpisce la
documentazione mantenuta a mano, gi\`a osservato per \code{openapi.yaml} nel
\cref{cap:05-contratti-dati}.
\end{quote}

\section{Limiti di questa serie}

\begin{enumerate}[leftmargin=*]
  \item \textbf{Tre esecuzioni per punto.} Non consentono di stimare una
  varianza. Le differenze inferiori al 10\% non sono interpretabili.

  \item \textbf{Infrastruttura condivisa.} Due VM Azure dello stesso tipo non
  hanno necessariamente le stesse prestazioni. Parte dell'oscillazione osservata
  fra le prime cinque release potrebbe essere questa.

  \item \textbf{Un solo scenario.} Tutta la serie usa
  \code{autoscaling-cycle-k8s}. Non esistono serie comparabili per altri profili
  di carico o altre famiglie di funzioni.

  \item \textbf{Il confronto \`e a coppie consecutive.} Il gate confronta con il
  record precedente, quindi una deriva lenta che resti sotto soglia a ogni passo
  non viene rilevata. La serie completa nella \cref{tab:serie} \`e il rimedio
  manuale, ma non \`e automatizzato.
\end{enumerate}

Il quarto punto merita enfasi perch\'e \`e un difetto strutturale del metodo, non
dei dati: un gate che confronta solo con l'immediato precedente ammette una
regressione monotona del 9\% per release indefinitamente.

\begin{damisurare}
Non esiste una verifica automatica contro un riferimento assoluto o contro la
migliore release storica. La \cref{tab:serie} rende la deriva visibile a chi la
legge, e nient'altro.
\todomis{gate contro un riferimento storico oltre al precedente immediato}
\end{damisurare}

\chapter{Footprint: JVM, nativo, memoria, GC}
\label{cap:24-footprint}

Questo capitolo risponde alla quarta domanda di ricerca del
\cref{cap:01-introduzione}: qual \`e il costo reale di un control plane, e come
cambia al variare della catena di build. \`E il capitolo con la maggiore densit\`a
di risultati negativi e di ipotesi rigettate, ed \`e riportato con quelle
incluse.

\section{Dove va la memoria di una JVM Spring}

Il punto di partenza \`e una funzione, non il control plane: \code{word-stats}
Java su Spring Boot, a 190\,MiB di memoria residente.

\begin{center}
\small
\begin{adjustbox}{max width=\textwidth}
\begin{tabular}{@{}lrr@{}}
\toprule
\textbf{Area} & \textbf{Dimensione} & \textbf{Quota} \\
\midrule
\textbf{heap in uso} & \textbf{22{,}4\,MB} & \textbf{12\,\%} \\
Metaspace (9\,375 classi caricate) & 53{,}0\,MB & 28\,\% \\
Compressed class space & 6{,}6\,MB & 3\,\% \\
CodeHeap (codice compilato JIT) & 13{,}5\,MB & 7\,\% \\
23 stack di thread & $\sim$23\,MB & 12\,\% \\
\emph{resto} (buffer diretti, GC, nativo) & $\sim$71\,MB & 37\,\% \\
\bottomrule
\end{tabular}
\end{adjustbox}
\end{center}

\begin{quote}
Lo heap --- i \emph{dati} --- \`e il 12\% del totale. Il resto \`e il runtime che
descrive se stesso: metadati delle classi, codice compilato, e i profili
conservati per ricompilarlo. Spring \`e costoso in memoria perch\'e carica
migliaia di classi e perch\'e una JVM porta con s\'e un compilatore, non perch\'e
allochi molto.
\end{quote}

Questa osservazione ha una conseguenza pratica immediata: ridurre l'uso di
memoria di un'applicazione Spring ottimizzando le allocazioni agisce sul 12\% del
problema. Le leve efficaci sono altre --- caricare meno classi, compilare in
anticipo, ridurre gli stack.

\`E anche la giustificazione quantitativa dei flag JVM del
\cref{cap:20-build-packaging}: \code{-Xss256k} agisce sui 23\,MB degli stack. Sui 13{,}5\,MB della CodeHeap
agiva \code{TieredStopAtLevel=1}, che il default di oggi non contiene pi\`u:
la latenza che C1 faceva perdere sul percorso caldo superava il risparmio
(\cref{cap:20-build-packaging}).

\section{Compilare nativamente}

\subsection{Sulla funzione}

\begin{center}
\small
\begin{adjustbox}{max width=\textwidth}
\begin{tabular}{@{}lrr@{}}
\toprule
 & \textbf{JVM} & \textbf{nativo} \\
\midrule
residente a riposo & 133\,MiB & \textbf{25{,}6\,MiB} \\
avvio & 1{,}046\,s & \textbf{0{,}066\,s} \\
residente sotto carico & 160\,MiB & 56\,MiB \\
\bottomrule
\end{tabular}
\end{adjustbox}
\end{center}

Il fattore cinque a riposo e il fattore sedici sull'avvio sono quanto ci si
aspetta dalla compilazione anticipata: sparisce il compilatore, spariscono i
metadati delle classi caricate dinamicamente.

Il dato interessante \`e l'ultimo. Sotto carico, la funzione Spring compilata
nativamente (56\,MiB) e la funzione \emph{lite} nativa (51{,}6\,MiB) quasi
convergono, pur differendo di un fattore dieci a riposo (25{,}6 contro
2{,}7\,MiB). L'interpretazione: \emph{il working set a caldo \`e dominato
dall'allocazione per richiesta, non dal framework}.

\`E un risultato che ridimensiona un'intuizione comune. La scelta di un framework
leggero fa una differenza enorme sull'impronta a riposo --- che conta per una
piattaforma che tiene molte funzioni idle --- e una differenza modesta
sull'impronta sotto carico.

\subsection{Sul control plane: il primo risultato \`e sbagliato}

La prima misura sembr\`o mostrare che il control plane nativo dimezzava sia
memoria sia CPU. Aveva dimezzato il \emph{lavoro}: 415\,184 richieste contro
865\,692, con p99 a 1\,696\,ms.

\`E l'errore di misura pi\`u istruttivo del progetto: il consumo di risorse per
unit\`a di tempo non \`e una metrica di efficienza se il sistema non sta
consegnando lo stesso lavoro.

\section{Il garbage collector dell'immagine nativa}

\subsection{La misura in isolamento}

Stesso backend, stesso carico, unica differenza la build del control plane:

\begin{center}
\small
\begin{adjustbox}{max width=\textwidth}
\begin{tabular}{@{}lrrr@{}}
\toprule
 & \textbf{JVM} & \textbf{nativo (serial)} & \textbf{nativo (G1)} \\
\midrule
throughput & 8\,205 rps & 4\,983 rps & \textbf{8\,085 rps} \\
mediana & 5{,}98\,ms & \textbf{4{,}71\,ms} & 6{,}08\,ms \\
p95 & 16{,}06\,ms & 27{,}95\,ms & 16{,}51\,ms \\
\textbf{massimo} & 92\,ms & \textbf{1{,}77\,s} & 123\,ms \\
avvio & 1{,}585\,s & 0{,}192\,s & \textbf{0{,}072\,s} \\
\bottomrule
\end{tabular}
\end{adjustbox}
\end{center}

Due osservazioni.

La \textbf{mediana della build nativa \`e migliore di quella della JVM}. Fra una
pausa e l'altra il codice compilato in anticipo \`e veloce; la perdita \`e
interamente sulla coda.

Il \textbf{massimo della build seriale \`e 1{,}77 secondi}, venti volte il
massimo della JVM. \`E la firma di un sistema che si ferma, non di uno che \`e
lento.

Il costo di CPU per richiesta era 1{,}62\,ms contro 1{,}32 --- la penalizzazione
ordinaria della compilazione anticipata, il 23\%, ben lontana dallo spiegare un
dimezzamento. E il sistema usava il 265\% di un 1\,100\% disponibile mentre
consegnava met\`a del lavoro: \emph{nulla era saturo, qualcosa si fermava}.

\subsection{La causa}

Corretto un difetto nel binder delle metriche del collector --- \code{JvmGcMetrics}
richiede notifiche che SubstrateVM non emette, e riportava uno zero costante ---
la causa \`e risultata diretta:

\begin{center}
\small
\begin{adjustbox}{max width=\textwidth}
\begin{tabular}{@{}lrr@{}}
\toprule
 & \textbf{JVM} & \textbf{nativo (serial)} \\
\midrule
raccolte complete & 2, $\sim$50\,ms ciascuna & \textbf{25, 883\,ms ciascuna} \\
tempo totale di raccolta & 6{,}6\,s & \textbf{22{,}1\,s su 60} \\
quota del tempo di parete & \textbf{0{,}9\,\%} & \textbf{31{,}8\,\%} \\
\bottomrule
\end{tabular}
\end{adjustbox}
\end{center}

Quasi un terzo del tempo speso a raccogliere.

\subsection{Tre ipotesi rigettate}

Il progetto ha testato e scartato tre spiegazioni alternative, e le riporta:

\begin{itemize}[leftmargin=*]
  \item \textbf{Limitare lo heap.} \code{-Xmx512m} ha peggiorato molto la
  situazione: 1\,543 rps e timeout a 30\,s.
  \item \textbf{Promozione prematura.} Ingrandire la generazione giovane non ha
  cambiato nulla. \emph{Il primo tentativo di questo test non era valido}:
  \code{-XX:MaxNewSize} \`e limitato da
  \code{-XX:MaximumYoungGenerationSizePercent}, il cui default \`e 10, quindi il
  flag non faceva nulla.
  \item \textbf{Buffer heap di Netty.} Entrambe le build preferiscono buffer
  diretti, entrambe usano l'allocatore adattivo, a nessuna delle due \`e
  disponibile \code{Unsafe}. Identiche.
\end{itemize}

Il secondo caso \`e riportato con la propria invalidazione. \`E buona pratica
sperimentale e vale la pena notarla: un test che non fa ci\`o che crede di fare
produce un risultato ``nessun effetto'' indistinguibile da quello vero.

\subsection{Il meccanismo}

La pausa scala con la \emph{dimensione} dello heap, non con i dati vivi:
563\,ms per uno heap di 512\,MB contro un insieme vivo di decine di megabyte. \`E
il comportamento di un collector che percorre l'intero spazio --- che \`e ci\`o
che fa l'unico collector utilizzabile di GraalVM Community: \code{--gc} accetta
\code{serial} ed \code{epsilon}, e nient'altro.

G1, disponibile in Oracle GraalVM, elimina il problema: 98{,}5\% del throughput
della JVM, coda a 123\,ms, avvio in 0{,}07\,s. Poich\'e Oracle GraalVM ha licenza
GFTC anzich\'e GPL, la build resta opt-in tramite \code{NATIVE\_GC=G1}.

\begin{damisurare}
\emph{Perch\'e} una raccolta completa costi 883\,ms con un insieme vivo piccolo
non \`e stato confermato. La pausa traccia la dimensione dello heap anzich\'e i
dati vivi, ma il meccanismo non \`e stato stabilito.
\todomis{meccanismo della pausa del collector seriale di SubstrateVM}
\end{damisurare}

\section{Osservare il collector: un problema di strumento}

Sotto G1 l'immagine nativa non registra alcun \code{GarbageCollectorMXBean},
quindi n\'e Micrometer n\'e il binder a polling riportano nulla --- sulla build
che si comporta meglio. \code{--enable-monitoring=jfr,jvmstat,jmxserver} non li
riporta in vita; \`e stato provato e non funziona.

JFR funziona, per una via indiretta. I tipi di evento specifici del collector
(\code{jdk.G1GarbageCollection}, \code{jdk.GCHeapSummary}, la famiglia
\code{GCPhase*}) sono dichiarati nella registrazione ed emettono zero eventi. Le
pause sono registrate come \textbf{operazioni della VM}, ciascuna con la propria
durata. Da una registrazione di 90\,s sotto carico:

\begin{center}
\small
\begin{adjustbox}{max width=\textwidth}
\begin{tabular}{@{}lrrrr@{}}
\toprule
\textbf{Operazione} & \textbf{Conteggio} & \textbf{Totale} & \textbf{Media} &
\textbf{Massimo} \\
\midrule
\code{G1 wrapper} & 43 & 835{,}1\,ms & 19{,}4\,ms & 97{,}6\,ms \\
\code{Collect for allocation} & 31 & 651{,}1\,ms & 21{,}0\,ms & 94{,}5\,ms \\
\code{Try init concurrent mark} & 2 & 37{,}3\,ms & 18{,}6\,ms & 22{,}5\,ms \\
\midrule
\textbf{totale} & \textbf{76} & \textbf{1\,524\,ms su 90\,s} & &
\textbf{1{,}69\,\%} \\
\bottomrule
\end{tabular}
\end{adjustbox}
\end{center}

Il risultato conferma la sezione precedente per una via indipendente ---
1{,}69\% del tempo di parete contro il 31{,}8\% del collector seriale --- e
fornisce pi\`u di quanto gli MXBean potrebbero: una durata per singola pausa, e
quindi il massimo, che i contatori interrogati periodicamente non danno neppure
dove funzionano.

\begin{center}
\small
\begin{adjustbox}{max width=\textwidth}
\begin{tabular}{@{}lll@{}}
\toprule
\textbf{Build} & \textbf{Conteggi e tempo totale} & \textbf{Pausa massima} \\
\midrule
JVM & Micrometer, nativamente & s\`i, dalle notifiche \\
nativo, serial & il binder a polling & no \\
nativo, G1 & \textbf{solo JFR} & \textbf{s\`i, per pausa} \\
\bottomrule
\end{tabular}
\end{adjustbox}
\end{center}

La conclusione \`e precisa: \emph{G1 non \`e inosservabile, \`e inosservabile
attraverso Prometheus}. Il dato richiede una registrazione estratta e analizzata
fuori banda anzich\'e un gauge da raccogliere. Adeguato per un'indagine, non per
un allarme.

\`E anche un limite reale dell'apparato sperimentale di questo progetto, che
altrove poggia interamente su Prometheus.

\section{Il livello di ottimizzazione}

\begin{center}
\small
\begin{adjustbox}{max width=\textwidth}
\begin{tabular}{@{}lrr@{}}
\toprule
 & \textbf{\code{-Os}} & \textbf{\code{-O3}} \\
\midrule
throughput & 8\,398 rps & \textbf{9\,251 rps} $(+10{,}2\%)$ \\
p95 & 15{,}64\,ms & 14{,}79\,ms \\
massimo & 148{,}7\,ms & 127{,}2\,ms \\
avvio & 0{,}173\,s & \textbf{0{,}068\,s} \\
residente a riposo & 120{,}6\,MiB & \textbf{39{,}3\,MiB} \\
dimensione immagine & 195\,MB & 397\,MB \\
\bottomrule
\end{tabular}
\end{adjustbox}
\end{center}

A \code{-O3} la build nativa non si limita a eguagliare la JVM: la supera ---
9\,251 contro 8\,205 richieste al secondo, con un avvio ventitr\'e volte pi\`u
rapido.

\subsection{Perch\'e l'immagine cresce}

Dal rapporto di build: l'area di codice triplica, da 47{,}78\,MB a 166{,}59\,MB,
a partire da \emph{meno} unit\`a di compilazione --- 84\,694 contro 105\,256.

Quell'accoppiamento \`e la firma dell'inlining. I metodi vengono copiati nei
chiamanti anzich\'e chiamati, quindi le unit\`a spariscono mentre ogni
sopravvissuta porta copie di ci\`o che ha assorbito; lo srotolamento dei cicli
aggiunge il resto. Stesso programma, pi\`u codice macchina --- che \`e anche
perch\'e \`e pi\`u veloce: nessun costo di chiamata, e ogni copia specializzata
per il proprio sito.

\subsection{Dove cade il costo, e dove non cade il guadagno}

Due qualificazioni, entrambe misurate.

Con un tetto di heap fissato, entrambi i livelli hanno raggiunto lo stesso picco
di 744\,MiB residenti, perch\'e il codice \`e mappato dal file anzich\'e copiato
nello heap: i 200\,MB in pi\`u sono un costo per il registry e per il disco, non
per il nodo.

Il guadagno di throughput \`e \emph{condizionale}: sotto un limite di 1\,GB le
due build sono risultate identiche a circa 3\,490 richieste al secondo, perch\'e
l\`i il collo di bottiglia \`e la raccolta e non la qualit\`a del codice.
\code{-O3} compra velocit\`a solo dove la memoria non \`e gi\`a il vincolo.

\subsection{Un avvertimento su PGO}

Il log di build riporta \code{PGO: off} a \code{-Os} e
\code{PGO: ML-inferred} a \code{-O3}: Oracle GraalVM applica profili inferiti da
modello quando ottimizza per velocit\`a. Parte del guadagno \`e quindi gi\`a una
forma di guida per profilo. Un eventuale PGO esplicito va misurato contro
\code{-O3}, non contro \code{-Os}, o gli verr\`a accreditato ci\`o che
\code{-O3} ha gi\`a raccolto.

\section{I limiti di memoria cambiano la risposta}

Tutte le cifre precedenti sono senza limite di memoria del container, dove ogni
build cresce finch\'e il suo collector non vede ragione di fermarsi.

\begin{table}[htbp]
\centering
\caption{Throughput e latenza sotto limite di memoria del container.}
\label{tab:memlimit}
\small
\begin{adjustbox}{max width=\textwidth}
\begin{tabular}{@{}llrrrr@{}}
\toprule
\textbf{Limite} & \textbf{Build} & \textbf{Throughput} & \textbf{p95} &
\textbf{Max} & \textbf{Memoria usata} \\
\midrule
500\,MB & JVM & 1\,515 rps & 324\,ms & 1{,}99\,s & 455 / 500 \\
        & nativo serial & 1\,370 rps & 32\,ms & 30{,}0\,s & 489 / 500 \\
        & \textbf{nativo G1, default} & \textbf{520 rps} & 61\,ms & 30{,}0\,s &
        \textbf{145 / 500} \\
        & nativo G1, \code{-Xmx350m} & 1\,445 rps & 96\,ms & 30{,}0\,s & 378 / 500 \\
\midrule
1\,GB & JVM & 2\,450 rps & 46\,ms & 1{,}34\,s & 805 / 1024 \\
      & \textbf{nativo G1, \code{-Xmx700m}} & \textbf{3\,089 rps} & 44\,ms &
      \textbf{135\,ms} & 736 / 1024 \\
\midrule
2\,GB & JVM & 2\,712 rps & 47{,}5\,ms & 451\,ms & 727 / 2048 \\
      & nativo G1, \code{-Xmx1400m} & 3\,400 rps & 39{,}5\,ms & 151\,ms & 1\,201 / 2048 \\
\midrule
nessuno & JVM & 8\,205 rps & 16{,}1\,ms & 92\,ms & 1{,}8\,GiB \\
        & nativo G1 & 8\,085 rps & 16{,}5\,ms & 123\,ms & 2{,}7\,GiB \\
\bottomrule
\end{tabular}
\end{adjustbox}
\end{table}

\subsection{G1 sotto un limite piccolo \`e una trappola}

Il tetto di default di G1 \`e il 25\% della RAM visibile. Rispetta il cgroup ---
il che \`e corretto --- ma poi si confina a 125\,MB di un container da 500\,MB,
lascia inutilizzato il 71\% della memoria, e crolla a 520 rps. Fissare
esplicitamente \code{-Xmx} quasi lo triplica.

\begin{quote}
\textbf{Con un limite di memoria, una build nativa G1 deve ricevere uno heap
esplicito.}
\end{quote}

\subsection{A 500\,MB non ci sta nulla}

Tutte le build perdono da cinque a sei volte il throughput senza vincolo, e
quelle native iniziano ad andare in timeout. Il rimedio non \`e un collector
migliore ma meno concorrenza offerta.

\subsection{Il ginocchio}

Il ginocchio della curva \`e fra 1 e 2\,GB. A 1\,GB con G1 e
\code{-Xmx700m} si ottengono 3\,089 rps con una coda di 135\,ms e nessun
fallimento, in un terzo della memoria che l'esecuzione senza vincoli aveva
occupato.

\subsection{Una configurazione da correggere}

Il control plane distribuito via Docker Compose fissa
\code{-XX:MaxRAMPercentage=70} senza \code{mem\_limit}, il che su un host da
12{,}5\,GB autorizza uno heap da 8{,}77\,GB. Fino a che non venga fissato un
limite, ogni cifra di memoria proveniente da quel percorso descrive
\emph{appetito}, non fabbisogno.

\`E la conclusione generale del capitolo, ed \`e metodologica prima che tecnica:
\textbf{misure di memoria senza un limite descrivono quanto un runtime \`e
disposto a prendere, non quanto gli serve.}

\section{Sintesi}

\begin{table}[htbp]
\centering
\caption{Quando conviene ciascuna build.}
\label{tab:build-scelta}
\small
\begin{tabularx}{\textwidth}{@{}lXX@{}}
\toprule
\textbf{Build} & \textbf{Conviene quando} & \textbf{Da evitare quando} \\
\midrule
JVM & memoria abbondante; si vuole osservabilit\`a del collector via Prometheus &
il tempo di avvio conta \\
nativo, serial & si tollerano pause al secondo & sotto carico sostenuto: 32\% del
tempo speso a raccogliere \\
nativo, G1, \code{-O3} & throughput e avvio contano; si accetta la licenza
Oracle & sotto limiti di memoria stretti senza \code{-Xmx} esplicito \\
\bottomrule
\end{tabularx}
\end{table}

\chapter{Valutazione dei controllori di concorrenza}
\label{cap:25-valutazione-concorrenza}

Questo capitolo riporta per esteso le campagne che hanno prodotto i risultati
citati nel \cref{cap:12-concurrency-control}. \`E il capitolo pi\`u ricco di
risultati negativi del documento, ed \`e organizzato in modo che ciascuno preceda
la correzione che ha motivato.

Tutte le misure sono di agosto 2026 e provengono dall'impianto descritto nel
\cref{cap:22-metodo-sperimentale}. Dove un'affermazione \`e stata fatta e poi
confutata dalla misura, il documento conserva la confutazione anzich\'e
l'affermazione.

\section{Il confronto fra i due controllori per funzione}

\subsection{Impianto}

Due funzioni su un solo control plane, vincolate agli stessi quattro core, coda
di 100 ciascuna, tetto per funzione 8, budget totale \code{BUDGETED} pari a 12.
Carico a ciclo chiuso con picco di 103 richieste offerte per funzione su tre
finestre: una funzione da sola, poi le due in antifase, poi le due in fase.
Entrambi i modi hanno ricevuto esattamente lo stesso carico.

\subsection{Risultati aggregati}

\begin{center}
\small
\begin{adjustbox}{max width=\textwidth}
\begin{tabular}{@{}lrr@{}}
\toprule
 & \textbf{\code{ADAPTIVE\_PER\_POD}} & \textbf{\code{BUDGETED}} \\
\midrule
richieste & 899\,684 & 893\,841 \\
rifiutate & 7 (0{,}001\%) & 33 (0{,}004\%) \\
p95 end-to-end & 80{,}2\,ms & 79{,}7\,ms \\
p99 end-to-end & 90{,}2\,ms & 87{,}2\,ms \\
\bottomrule
\end{tabular}
\end{adjustbox}
\end{center}

I due modi sono entro l'1\% l'uno dall'altro end-to-end. Preso da solo, questo
risultato direbbe che il modo a budget non porta alcun beneficio. \`E la
conclusione sbagliata, per due ragioni distinte che le sezioni seguenti
sviluppano.

\subsection{La concorrenza concessa, per finestra}

\begin{center}
\small
\begin{adjustbox}{max width=\textwidth}
\begin{tabular}{@{}lrrrrrr@{}}
\toprule
\textbf{Finestra} & \textbf{AD. java} & \textbf{AD. lite} & \textbf{somma} &
\textbf{BU. java} & \textbf{BU. lite} & \textbf{somma} \\
\midrule
una sola & 7{,}8 & 4{,}1 (idle) & --- & 8{,}0 & 1{,}2 (idle) & --- \\
antifase & 8{,}0 & 5{,}0 & 13{,}0 & 8{,}0 & 2{,}9 & \textbf{10{,}9} \\
in fase & 7{,}9 & 5{,}5 & 13{,}4 & 8{,}0 & 3{,}7 & \textbf{11{,}7} \\
\bottomrule
\end{tabular}
\end{adjustbox}
\end{center}

\code{BUDGETED} mantiene il totale sotto il proprio budget di 12 e sposta
capacit\`a fra le funzioni al muoversi del carico. \code{ADAPTIVE\_PER\_POD} non
ha alcun meccanismo per fare n\'e l'una n\'e l'altra cosa, perch\'e i suoi due
controllori indipendenti non sanno l'uno dell'esistenza dell'altro. \`E la
propriet\`a per cui il modo esiste, ed \`e visibile qui e non negli aggregati.

\subsection{L'asimmetria non \`e iniquit\`a}

Le due funzioni ricevono concorrenza molto diversa, e questo \emph{non} \`e un
difetto dell'allocatore. Sono runtime diversi: \code{lite} serve in 2{,}15\,ms
contro 5{,}57\,ms, quindi la legge di Little le assegna un limite pi\`u piccolo.
A 1\,493 richieste al secondo le servono circa 3{,}2 in volo, e oltre non compra
completamenti.

\code{BUDGETED} le ha concesso 2{,}6 in media contro i 4{,}9 di
\code{ADAPTIVE} \emph{a parit\`a di throughput}, e il suo tempo di servizio \`e
sceso del 37\%. Met\`a della concorrenza, lo stesso lavoro svolto.

\subsection{Un errore di calibrazione}

Il \cref{cap:22-metodo-sperimentale} riporta l'errore che ha invalidato la prima
esecuzione di questo confronto e la regola che ne \`e stata derivata. Va
richiamato qui perch\'e riguarda direttamente l'interpretabilit\`a della tabella
precedente: con budget 8 anzich\'e 12, la coda traboccava per aritmetica in un
modo e non nell'altro, producendo 86\,302 rifiuti che non dicevano nulla sui
controllori.

\section{La coda \`e dove sta la latenza}

\code{function\_latency\_ms} \`e misurata dal control plane fra dispatch e
completamento: \`e il round trip, non il tempo di esecuzione interno alla
funzione. Il governor decide da essa.

Sotto il profilo a raffica l'attesa media in coda \`e stata di \textbf{37--43\,ms}
contro un tempo di servizio di circa \textbf{5\,ms}. La grandezza da cui il
controllore si orienta \`e circa \textbf{un settimo} di ci\`o che il chiamante
sperimenta.

La conseguenza \`e misurabile, non teorica. In un confronto iniziale
\code{BUDGETED} ha ridotto il tempo di servizio del 31\% su una funzione e del
44\% sull'altra --- entrambe molto dentro un SLO di 10\,ms --- e ha
\textbf{peggiorato} il p95 end-to-end, perch\'e l'attesa creata dal suo stesso
limite ha pi\`u che assorbito il guadagno.

\begin{quote}
Un limite di concorrenza non rimuove lavoro: lo sposta nel buffer.
\end{quote}

\section{Il ciclo chiuso non poteva porre la domanda}

Prima di poter valutare un controllore che decide dal sojourn, \`e stato
necessario correggere l'impianto.

Sotto il profilo \code{burst} ogni utente virtuale tiene una richiesta, quindi
il numero nel sistema \`e fissato dal generatore e la profondit\`a di coda \`e
\code{VU} $-$ \code{limite}. Due conseguenze, entrambe misurate:

\begin{itemize}[leftmargin=*]
  \item Portare il limite da 4 a 8 ha accorciato di quattro una coda profonda 99.
  Nessun controllore pu\`o muovere l'attesa in modo apprezzabile.
  \item La coda si assesta a 95 su 100 per costruzione, permanentemente sopra
  qualunque soglia di guardia. La guardia di ammissione del modo \code{SOJOURN}
  scatta a ogni tick e il modello non parla mai: la riesecuzione legge
  \code{8 [8--8]} per tutta la sua durata.
\end{itemize}

\begin{quote}
\`E per questo che due controllori che leggevano segnali completamente diversi
sono finiti entro il 2\% l'uno dall'altro su questo profilo. La somiglianza era
una propriet\`a dell'impianto.
\end{quote}

\`E il risultato metodologico centrale del capitolo, e la ragione per cui il
confronto aggregato della prima sezione va letto con cautela: quel confronto
girava su un impianto incapace di distinguere.

\section{Decidere dalla latenza del chiamante: due fallimenti e un risultato}

\subsection{Fallimento 1: il gradiente sul sojourn}

Alimentare la regola a gradiente esistente con la latenza end-to-end \`e
instabile, e \`e stato scartato prima dell'implementazione. La regola riduce il
limite quando la latenza eccede l'obiettivo; un limite pi\`u piccolo drena la
coda pi\`u lentamente, quindi l'attesa cresce, quindi si riduce ancora, fino al
pavimento.

\subsection{Fallimento 2: la ricerca del minimo}

Il sojourn ha un minimo interno dove il tempo di servizio non ce l'ha, quindi la
ricerca \`e ben posta in linea di principio. \`E stata implementata e falsificata:

\begin{center}
\small
\begin{adjustbox}{max width=\textwidth}
\begin{tabular}{@{}lrr@{}}
\toprule
 & \textbf{\code{ADAPTIVE\_PER\_POD}} & \textbf{\code{SOJOURN}, a ricerca} \\
\midrule
rifiutate & 32 (0{,}004\%) & \textbf{36\,181 (4{,}275\%)} \\
limite [min--max] & 7{,}8 [4--8] & 4{,}0 [1--8] \\
\bottomrule
\end{tabular}
\end{adjustbox}
\end{center}

Il punto di lavoro scelto non era il problema: a limite 4 consegnava il 96{,}5\%
del throughput con met\`a del tempo di servizio. Il problema \`e che \emph{il
limite fa due lavori}: insieme alla coda \`e anche la finestra di ammissione,
perch\'e una funzione trattiene \code{limite + queueSize} e rifiuta il resto.
Dimensionarlo dalla sola latenza ha fatto perdere il 4\% del traffico.

La causa profonda \`e un problema di attribuzione: la ricerca non pu\`o
distinguere ``la mia mossa ha aiutato'' da ``il carico \`e calato''. In un
avvallamento ha attribuito il sollievo naturale alla propria decisione di
restringere, ed era ancora piccola quando il carico \`e tornato.

\subsection{Il risultato: calcolare da un modello}

Sotto arrivi aperti, con picco sopra la capacit\`a misurata, stesso carico,
stessi core, stesse code, entrambi i modi per funzione:

\begin{table}[htbp]
\centering
\caption{\code{ADAPTIVE\_PER\_POD} contro \code{SOJOURN} sotto arrivi aperti.}
\label{tab:sojourn}
\small
\begin{adjustbox}{max width=\textwidth}
\begin{tabular}{@{}lrrl@{}}
\toprule
 & \textbf{\code{ADAPTIVE\_PER\_POD}} & \textbf{\code{SOJOURN}} & \\
\midrule
offerte & 659\,973 & 659\,995 & identiche per costruzione \\
\textbf{servite} & 549\,856 & \textbf{567\,210} & \textbf{$+3{,}2\%$} \\
rifiutate & 110\,117 (16{,}7\%) & \textbf{92\,785 (14{,}1\%)} &
\textbf{$-15{,}7\%$} \\
\textbf{p95 end-to-end} & 120{,}48\,ms & \textbf{107{,}68\,ms} &
\textbf{$-10{,}6\%$} \\
\textbf{p99 end-to-end} & 157{,}67\,ms & \textbf{133{,}82\,ms} &
\textbf{$-15{,}1\%$} \\
attesa in coda & 35{,}61\,ms & 31{,}03\,ms & $-12{,}9\%$ \\
tempo di servizio & 1{,}69\,ms & 2{,}00\,ms & $+18\%$ \\
limite medio & 3{,}9 & 2{,}8 & \\
\bottomrule
\end{tabular}
\end{adjustbox}
\end{table}

Lo scambio \`e andato nella direzione prevista: il tempo di servizio peggiora del
18\%, l'attesa migliora del 13\%, e poich\'e l'attesa \`e venti volte il tempo di
servizio il chiamante ci guadagna. L'obiettivo --- stesso throughput, stessi
rifiuti, p95 end-to-end inferiore --- era stato scritto nello scenario
\emph{prima} dell'esecuzione, e tutti e tre sono stati superati.

Il dettaglio istruttivo \`e che \code{SOJOURN} ha ottenuto questo con
\emph{meno} concorrenza media. Non vince aggiungendo capacit\`a ma collocandola
nel tempo, che \`e ci\`o per cui la predizione era stata introdotta.

\subsection{Caveat}

Riportati dal progetto e da prendere sul serio:

\begin{itemize}[leftmargin=*]
  \item Una sola esecuzione per modo, non ripetuta. Differenze del 10--15\% sono
  grandi ma non stabilite.
  \item Entrambi i sistemi stanno scartando il 14--16\% del carico, perch\'e il
  picco a ciclo aperto eccede genuinamente la capacit\`a. Il confronto \`e fra
  due sistemi sovraccarichi, non fra due sistemi comodi.
  \item \code{dropped\_iterations} \`e stato 21 e 0, quindi il generatore non
  stava rinunciando e i rifiuti sono della piattaforma.
\end{itemize}

Il secondo caveat \`e il pi\`u limitante: \code{SOJOURN} non \`e mai stato
osservato scegliere un limite in condizioni comode.

\section{Co-tenancy: cross-talk senza contesa}

\subsection{Il fenomeno}

Due funzioni su un control plane, una a carico costante mentre l'altra va e
viene. Quando la vicina arriva, la funzione stabile perde il 28\% del proprio
throughput e guadagna il 36\% di latenza.

\subsection{L'ipotesi ovvia \`e sbagliata}

La spiegazione immediata era la contesa di CPU. \`E stato aggiunto
\code{nanofaas.container-local.cpuset} per vincolare entrambe le funzioni agli
stessi quattro core, e l'esecuzione \`e stata ripetuta:

\begin{center}
\small
\begin{adjustbox}{max width=\textwidth}
\begin{tabular}{@{}lrr@{}}
\toprule
 & \textbf{4 core propri ciascuna} & \textbf{core 0--3 condivisi} \\
\midrule
java da sola & 1\,762 rps / 4{,}57\,ms & 1\,762 rps / 4{,}57\,ms \\
java con la vicina & 1\,274 rps / 6{,}23\,ms & 1\,274 rps / 6{,}23\,ms \\
combinato durante la sovrapposizione & 2\,769 rps & 2\,751 rps \\
\bottomrule
\end{tabular}
\end{adjustbox}
\end{center}

Identiche a tre cifre significative. Vincolare entrambe le funzioni a quattro
core condivisi costa lo 0{,}6\% del throughput combinato, che \`e rumore.

I quattro core non erano mai il vincolo attivo: a circa 2\,760 richieste al
secondo ciascuna richiesta pu\`o costare al massimo 1{,}45\,ms di CPU, e
l'obiettivo \`e raggiunto. Il cross-talk viene da ci\`o che le due funzioni
condividono in \emph{entrambe} le configurazioni --- l'unico control plane ---
non dai core.

\subsection{Perch\'e il \code{cpuset} resta}

Pur non essendo la spiegazione cercata, la funzionalit\`a \`e stata mantenuta:
rende il budget di concorrenza una divisione di una capacit\`a che
\emph{esiste}. Senza di essa, su una macchina con core in eccesso, la premessa
del modo \code{BUDGETED} \`e falsa e il modo non pu\`o essere valutato
(\cref{cap:15-deployment-provider}).

\begin{damisurare}
Il meccanismo esatto del cross-talk attraverso il control plane condiviso non \`e
stato isolato. I candidati --- contesa sul thread dello scheduler, sull'event
loop, sulle strutture concorrenti delle code --- sono distinguibili con la
profilazione, e non sono stati distinti.
\todomis{origine del cross-talk fra funzioni co-residenti nel control plane}
\end{damisurare}

\section{Un confondente scoperto: i due \code{word-stats}}

Registrato dal progetto perch\'e i commenti degli scenari hanno affermato il
contrario per un certo periodo:

\begin{center}
\small
\begin{adjustbox}{max width=\textwidth}
\begin{tabular}{@{}lrr@{}}
\toprule
 & \textbf{\code{word-stats-java}} & \textbf{\code{word-stats-java-lite}} \\
\midrule
build & Spring Boot su JVM & binario nativo GraalVM \\
SDK & \code{sdks/java} & \code{sdks/java-lite} \\
residente, a riposo & 133\,MiB & 2{,}7\,MiB \\
residente, a caldo & 160\,MiB & 51{,}6\,MiB \\
tempo di servizio & 5{,}57\,ms & 2{,}15\,ms \\
\bottomrule
\end{tabular}
\end{adjustbox}
\end{center}

\begin{quote}
Qualunque esperimento che le tratti come intercambiabili \`e confuso.
\end{quote}

\`E riportato in questo capitolo, e non solo nel \cref{cap:18-sdk}, perch\'e
riguarda direttamente l'interpretazione della prima sezione: l'asimmetria fra le
due funzioni nella tabella delle concessioni \`e in larga parte questa
differenza di runtime, non una scelta dell'allocatore.

\section{Che cosa non \`e stabilito}

Il progetto enumera esplicitamente i limiti dei propri risultati. Sono riportati
qui integralmente perch\'e sono parte del contributo:

\begin{enumerate}[leftmargin=*]
  \item Perch\'e una raccolta completa costi 883\,ms quando l'insieme vivo \`e
  piccolo. La pausa traccia la dimensione dello heap anzich\'e i dati vivi, ma il
  meccanismo non \`e stato confermato.
  \item Se i pesi di \code{BUDGETED} facciano qualcosa di utile. Entrambe le
  funzioni hanno ricevuto SLO e pesi identici di proposito, quindi i livelli di
  servizio differenziati restano non verificati.
  \item Se il risultato di \code{SOJOURN} sotto arrivi aperti regga. \`E una sola
  esecuzione non ripetuta per modo, su un banco in cui entrambi i sistemi
  scartavano il 14--16\% del carico offerto.
  \item Che cosa faccia \code{SOJOURN} quando il sistema \emph{non} \`e
  sovraccarico. Ogni sua esecuzione finora \`e stata a capacit\`a o oltre, quindi
  il modo non \`e mai stato osservato scegliere un limite in condizioni comode.
  \item Se la soglia di guardia al 70\% sia corretta. \`E stata scelta, non
  misurata, e sul profilo a ciclo chiuso ha scavalcato completamente il modello.
\end{enumerate}

Il punto 2 merita enfasi: l'equit\`a max-min \emph{pesata} \`e presentata nel
\cref{cap:12-concurrency-control} come uno dei contributi del lavoro, e i pesi
non sono mai stati esercitati con valori diversi. Ci\`o che \`e stato validato \`e
l'allocazione a pesi uguali.

\section{Sintesi}

\begin{table}[htbp]
\centering
\caption{Le tre campagne e i loro esiti.}
\label{tab:campagne}
\small
\begin{tabularx}{\textwidth}{@{}lXX@{}}
\toprule
\textbf{Campagna} & \textbf{Risultato} & \textbf{Forza dell'evidenza} \\
\midrule
\code{BUDGETED} contro \code{ADAPTIVE\_PER\_POD} & Pari end-to-end; il primo con
met\`a della concorrenza sulla funzione che non ne ha bisogno, e totale sotto
budget. & Buona sugli aggregati; l'impianto a ciclo chiuso non poteva
distinguere i segnali. \\
\code{SOJOURN} contro \code{ADAPTIVE\_PER\_POD} & $+3{,}2\%$ servite,
$-15{,}7\%$ rifiuti, $-10{,}6\%$ p95, con \emph{meno} concorrenza. & Una sola
esecuzione per modo, entrambi in sovraccarico. \\
Co-tenancy & Cross-talk reale ($-28\%$ throughput), \emph{non} contesa di CPU
(differenza 0{,}6\%). & Solida: la manipolazione del \code{cpuset} \`e un
controllo diretto. \\
\bottomrule
\end{tabularx}
\end{table}

Il terzo risultato \`e il pi\`u solido dei tre, ed \`e negativo: esclude
un'ipotesi anzich\'e stabilirne una. \`E anche il pi\`u utile per il seguito del
lavoro, perch\'e restringe lo spazio di ricerca sul meccanismo del cross-talk al
control plane condiviso.


\part{Qualit\`a, discussione, chiusura}
\chapter{Qualit\`a del software e verifica}
\label{cap:26-qualita-verifica}

In un progetto che si propone come strumento di misura, la verifica non \`e
igiene ma parte del contributo: un risultato ottenuto su una piattaforma di cui
non si conosce lo stato non \`e un risultato. Questo capitolo descrive
l'apparato di verifica di \nf{}, dalle unit\`a alle regole architetturali.

\section{Le proporzioni}

\begin{center}
\small
\begin{adjustbox}{max width=\textwidth}
\begin{tabular}{@{}lr@{}}
\toprule
\textbf{Grandezza} & \textbf{Valore} \\
\midrule
Righe di codice di produzione (Java) & 16\,487 \\
Righe di codice di test (Java) & 31\,664 \\
Rapporto test / produzione & $\approx 1{,}9$ \\
Classi di test & 235 \\
Soglia minima di copertura (JaCoCo) & 85\,\% \\
\bottomrule
\end{tabular}
\end{adjustbox}
\end{center}

Quasi due righe di test per ogni riga di produzione. Il rapporto \`e alto ma non
sorprendente per un progetto la cui superficie \`e piccola e le cui interazioni
--- fra moduli, fra code, fra politiche --- sono numerose: il codice combinatorio
\`e nei test, non nella produzione.

\section{I livelli}

\begin{center}
\small
\begin{adjustbox}{max width=\textwidth}
\begin{tabular}{@{}lp{8.4cm}@{}}
\toprule
\textbf{Livello} & \textbf{Caratteristiche} \\
\midrule
Unit\`a & JUnit 5, classi \code{*Test.java}, nessuna infrastruttura. \\
Integrazione & Contesto Spring completo con \code{WebTestClient}; server mock
Fabric8 per Kubernetes. \\
Container & Testcontainers~\cite{testcontainers}: richiede Docker. \\
Kubernetes & Richiede \code{KUBECONFIG}. \\
Contratto multi-SDK & \code{functions/contract-tests/run.sh}
(\cref{cap:18-sdk}). \\
Watchdog & Test unitari Rust pi\`u integrazione locale per HTTP, STDIO, FILE,
callback, metriche e ciclo di vita warm. \\
E2E di scenario & Di propriet\`a dello strumento esterno
(\cref{cap:22-metodo-sperimentale}). \\
\bottomrule
\end{tabular}
\end{adjustbox}
\end{center}

Va notato che i test del watchdog non richiedono n\'e Docker n\'e una VM: bastano
Rust, Python 3 e alcune utilit\`a a riga di comando. \`E una scelta che mantiene
verificabile in locale il componente presente in ogni container di funzione.

\section{Le due invocazioni della suite}

Come discusso nel \cref{cap:16-moduli-diagnostici}, la selezione modulare \`e
risolta a tempo di configurazione di Gradle e non pu\`o variare all'interno di
una build. La verifica completa richiede quindi due invocazioni:

\begin{lstlisting}[language=bash]
./gradlew test                              # default: tutti i moduli
./gradlew test -PcontrolPlaneModules=none   # la configurazione minima
\end{lstlisting}

\section{Le regole architetturali}

Il progetto codifica sei regole con ArchUnit~\cite{archunit}. Cinque riguardano
il nucleo e una i moduli.

\subsection{Le regole del nucleo}

\begin{description}[style=nextline, leftmargin=0pt]
  \item[\textbf{R1 --- Assenza di cicli.}] I package di primo livello del nucleo
  non devono formare cicli fra loro.
\begin{lstlisting}[language=Java]
slices().matching("..controlplane.(*)..").should().beFreeOfCycles();
\end{lstlisting}

  \item[\textbf{R2 --- \code{api} \`e l'unico punto di ingresso.}] Nessuna classe
  fuori dal package \code{api} deve dipendere da esso. La direzione \`e
  unidirezionale: l'API conosce il servizio, il servizio non conosce l'API.

  \item[\textbf{R3 --- Il nucleo non dipende dai moduli.}]
\begin{lstlisting}[language=Java]
noClasses().that().resideInAPackage("..controlplane..")
        .should().dependOnClassesThat().resideInAPackage("..modules..")
        .as("the core must work without optional modules (SPI)");
\end{lstlisting}
  \`E la regola che rende \emph{verificata} la promessa del
  \cref{cap:03-requisiti-principi}. Senza di essa, l'affermazione ``il nucleo
  funziona senza moduli'' sarebbe un'intenzione; con essa \`e una propriet\`a che
  la build accerta.

  \item[\textbf{R4 --- La direzione dei livelli.}] I package \code{dispatch},
  \code{execution} e \code{deployment} non devono dipendere da \code{service}. La
  gerarchia \`e \code{api} $\to$ \code{service} $\to$ \code{dispatch}, e i livelli
  bassi non risalgono.

  \item[\textbf{R6 --- Il namespace appartiene al nucleo.}] Ogni classe nel
  package \code{it.unimib.datai.nanofaas.controlplane} deve avere il proprio
  file di classe sotto \code{platform/control-plane/}.
\end{description}

L'ultima \`e la pi\`u insolita e vale la pena soffermarcisi. Verifica non una
dipendenza ma una \emph{corrispondenza fra nome logico e collocazione fisica}:
impedisce che un modulo dichiari una classe dentro il namespace del nucleo per
aggirare le altre regole. Senza di essa, R3 sarebbe eludibile spostando il
package anzich\'e la dipendenza.

L'implementazione ha dovuto adattarsi a un cambiamento della libreria:

\begin{lstlisting}[language=Java]
// ArchUnit 1.4.1 dropped SourceCodeLocation's source-path accessor, so read the
// source URI instead; its path points under platform/control-plane/.
String path = item.getSource().map(Source::getUri).map(URI::getPath).orElse("");
\end{lstlisting}

\subsection{La regola dei moduli}

Ciascuno dei nove moduli porta la propria classe di verifica, con la stessa
regola:

\begin{lstlisting}[language=Java]
// R5': isolamento SPI. Un'estensione tra moduli è una decisione esplicita:
// aggiungere una riga not(resideInAPackage(...)) con commento per ogni coppia sanzionata
// (ArchUnit 1.4 ha rimosso ignoreDependency dal fluent API di base).
@ArchTest
static final ArchRule does_not_depend_on_other_modules =
        noClasses()
                .that().resideInAPackage("..modules.offload..")
                .should().dependOnClassesThat(
                        resideInAPackage("..modules..")
                                .and(DescribedPredicate.not(resideInAPackage("..modules.offload.."))))
                .as("module must not depend on other modules (SPI isolation)");
\end{lstlisting}

Il divieto \`e sulle dipendenze \emph{laterali}. Dipendere dalle interfacce del
nucleo \`e il meccanismo previsto e resta libero: \`e cos\`i che i provider di
deployment implementano \code{ImageValidator} e che
\code{concurrency-control} scrive attraverso \code{ScalingMetricsSource}.

Il commento prescrive anche la procedura in caso di deroga: aggiungere una riga
esplicita con la propria giustificazione, cos\`i che ogni eccezione sia visibile
nel codice della regola. Al momento della stesura nessuna deroga \`e stata
aggiunta.

\subsection{Che cosa questo apparato compra}

Le sei regole insieme rendono \emph{eseguibile} l'architettura descritta nella
Parte II. Il \cref{cap:08-sistema-moduli} ha osservato che il compilatore \`e gi\`a
un esecutore parziale --- i moduli entrano come dipendenze
\code{runtimeOnly}, quindi il nucleo non pu\`o compilare contro di essi. Le
regole ArchUnit coprono ci\`o che il compilatore non vede: i cicli, la direzione
dei livelli, e la corrispondenza fra namespace e modulo fisico.

\section{Analisi statica}

\subsection{SonarQube}

L'analisi \`e su richiesta, non in integrazione continua:
\code{scripts/sonar.sh} avvia un container SonarQube effimero, analizza Java,
Python e Rust, e stampa il conteggio dei problemi aperti per linguaggio.

La rinuncia al \emph{quality gate} \`e argomentata esplicitamente:

\begin{quote}\itshape
Il server \`e effimero (nessun volume, nessuna storia), quindi SonarQube tratta
l'intera base di codice come ``nuovo codice'' a ogni esecuzione. Il gate di
default fallisce ogni volta che esiste un problema preesistente, il che
significa che un verdetto del gate sarebbe un FAIL costante e privo di
informazione.
\end{quote}

\`E la stessa logica che governa le soglie del gate di regressione
(\cref{cap:22-metodo-sperimentale}): un controllo che fallisce sempre viene
disattivato, e un controllo disattivato non controlla nulla. Meglio riportare un
numero che un verdetto inutile.

La conseguenza operativa \`e che la linea di base \`e zero problemi aperti, e
quindi \emph{qualunque} problema rilevato \`e nuovo. \`E un criterio pi\`u
severo di un gate sul nuovo codice, e sostenibile solo perch\'e la base di
codice \`e piccola.

\subsection{Analisi delle dipendenze}

Il progetto applica il plugin \code{dependency-analysis} a livello di settings,
che segnala dipendenze dichiarate e non usate, usate e non dichiarate, e
dipendenze che potrebbero cambiare configurazione.

Due osservazioni dall'esperienza registrata dal progetto: cinque dipendenze di
test sono state effettivamente rimosse a seguito dell'analisi, e gli starter di
Spring producono sistematicamente falsi positivi --- uno starter \`e per
costruzione una dipendenza che non si usa direttamente.

L'esclusione di \code{netty-codec-http3} descritta nel
\cref{cap:20-build-packaging} \`e il risultato pi\`u tangibile di questa linea di
lavoro: sei archivi rimossi, di cui cinque con codice nativo per piattaforma.

\subsection{Copertura}

JaCoCo su ogni sottoprogetto, con soglia minima all'85\%. I report XML
alimentano anche l'analisi SonarQube.

\section{Test di regressione strutturali}

La posizione del progetto sui test di prestazione \`e stata riportata nel
\cref{cap:21-osservabilita} e vale la pena richiamarla qui perch\'e appartiene a
questo capitolo: il repository contiene test \emph{strutturali} sul percorso
caldo anzich\'e microbenchmark assoluti. Asseriscono il riuso del replay, il
forward progress della coda sincrona dietro una testa bloccata, e l'equit\`a
asincrona fra funzioni calde e fredde.

\begin{quote}\itshape
Le temporizzazioni assolute sono sensibili all'ambiente, le garanzie di equit\`a
e riuso non lo sono.
\end{quote}

Ciascuno dei tre corrisponde a una propriet\`a discussa nella Parte III: il
replay al meccanismo di idempotenza (\cref{cap:06-nucleo}), il forward progress
alla scansione selettiva dello scheduler sincrono (\cref{cap:10-sync-queue}),
l'equit\`a al quanto di scheduling del modulo asincrono
(\cref{cap:09-async-queue}). Sono, in sostanza, i tre invarianti che i capitoli
precedenti hanno argomentato, resi eseguibili.

\section{Trappole note dell'impianto di test}

Il progetto ha accumulato alcune conoscenze operative che vale la pena
registrare, perch\'e ciascuna corrisponde a una propriet\`a reale del sistema.

\begin{description}[style=nextline, leftmargin=0pt]
  \item[\textbf{La coda sincrona respinge la prima richiesta.}] Senza storia di
  throughput la stima dell'attesa \`e infinita, e ogni richiesta viene respinta
  con \code{429} (\cref{cap:10-sync-queue}). I test di API devono quindi
  disattivare esplicitamente la coda sincrona. Non \`e un difetto: \`e il
  comportamento corretto in produzione che si manifesta in modo scomodo in
  laboratorio.

  \item[\textbf{Le metriche spariscono con la funzione.}] I meter per funzione
  sono rimossi alla cancellazione (\cref{cap:06-nucleo}), quindi una verifica
  che interroghi le metriche dopo aver cancellato la funzione non trova nulla.
  Gli scenari di validazione sono stati corretti mantenendo viva l'ultima
  funzione fino a dopo il controllo.

  \item[\textbf{Aggiungere un componente a un \code{record} rompe i test.}] Il
  costruttore canonico cambia firma e ogni istanziazione va aggiornata. \`E la
  ragione dei costruttori di compatibilit\`a descritti nel
  \cref{cap:05-contratti-dati}.

  \item[\textbf{Modificare i contratti condivisi richiede la suite intera.}] Una
  modifica a \code{platform/common} o a un SDK si propaga oltre il control plane;
  eseguire soltanto \code{:control-plane:test} non basta.
\end{description}

\section{Limiti dell'apparato}

\begin{enumerate}[leftmargin=*]
  \item \textbf{Lo spazio delle configurazioni modulari \`e coperto ai soli
  estremi} (\cref{cap:16-moduli-diagnostici}): \code{all} e \code{none}. Le
  configurazioni intermedie non hanno copertura dedicata.

  \item \textbf{L'analisi statica non \`e in integrazione continua.} \`E su
  richiesta, quindi la sua efficacia dipende dalla disciplina di chi sviluppa.

  \item \textbf{Gli scaffold generati non sono verificati end-to-end}
  (\cref{cap:19-cli-scaffolding}).
\end{enumerate}

I tre limiti hanno una struttura comune: riguardano tutti \emph{artefatti
generati o dichiarati} --- una configurazione di build, un file di
specifica, uno scaffold --- anzich\'e codice eseguito. \`E la categoria che
l'apparato di test copre peggio, ed \`e una lacuna riconoscibile e colmabile con
gli strumenti gi\`a presenti.

\chapter{Discussione e minacce alla validit\`a}
\label{cap:27-discussione}

Questo capitolo separa ci\`o che il lavoro stabilisce da ci\`o che suggerisce e
da ci\`o che non tocca affatto. \`E scritto per essere letto da chi volesse
contestare i risultati: elenca in anticipo le obiezioni che riterrei fondate.

\section{Che cosa i dati stabiliscono}

\begin{description}[style=nextline, leftmargin=0pt]
  \item[\textbf{S1. Il cross-talk fra funzioni co-residenti non \`e contesa di
  CPU.}] \`E il risultato pi\`u solido, perch\'e ottenuto per manipolazione diretta:
  vincolando entrambe le funzioni agli stessi quattro core, il throughput
  combinato \`e cambiato dello 0{,}6\%, mentre la funzione stabile continuava a
  perdere il 28\% del proprio throughput all'arrivo della vicina. Il controllo \`e
  pulito e l'ipotesi \`e falsificata.

  \item[\textbf{S2. Un controllore che decide dal tempo di servizio non ha
  ottimo interno.}] \`E stabilito analiticamente --- $S(C)$ \`e monotona crescente
  su una risorsa condivisa, quindi il gradiente punta in discesa ovunque --- e
  confermato per osservazione su due runtime, entrambi visti camminare fino al
  proprio pavimento.

  \item[\textbf{S3. Il tempo di servizio \`e una frazione minoritaria della
  latenza del chiamante sotto carico.}] Misurato: attesa 37--43\,ms contro
  servizio $\approx 5$\,ms. \`E un rapporto di sette a uno, molto oltre qualunque
  incertezza di misura.

  \item[\textbf{S4. Un carico a ciclo chiuso non pu\`o distinguere politiche di
  ammissione.}] Segue dalla struttura del generatore --- il numero in sistema \`e
  fissato dal generatore, la profondit\`a di coda \`e \code{VU} $-$ \code{limite}
  --- ed \`e stato osservato: la guardia di ammissione scattava a ogni tick e il
  modello non parlava mai.

  \item[\textbf{S5. Il collector seriale di GraalVM Community spende circa un
  terzo del tempo di parete a raccogliere, sotto questo carico.}] 22{,}1\,s su
  60, con 25 raccolte complete da 883\,ms. Confermato per due vie indipendenti
  (contatori corretti e registrazione JFR).

  \item[\textbf{S6. Misure di memoria senza limite del container descrivono
  appetito.}] Dimostrato dalla differenza fra le esecuzioni vincolate e non
  vincolate: 1{,}8--2{,}7\,GiB senza limite contro 736\,MiB con limite a 1\,GB e
  throughput superiore.

  \item[\textbf{S7. Il tasso di errore \`e zero su sei release.}] \`E un gate
  rigido e la sua tenuta \`e verificata su diciotto esecuzioni.
\end{description}

\section{Che cosa i dati suggeriscono senza stabilire}

\begin{description}[style=nextline, leftmargin=0pt]
  \item[\textbf{D1. \code{SOJOURN} \`e superiore a \code{ADAPTIVE\_PER\_POD} sotto
  arrivi aperti.}] I numeri sono netti --- $+3{,}2\%$ servite, $-15{,}7\%$
  rifiuti, $-10{,}6\%$ p95 --- ma provengono da \emph{una sola esecuzione per
  modo}. Differenze del 10--15\% sono grandi, e con $n = 1$ non sono
  distinguibili da una fluttuazione dell'ambiente.

  \item[\textbf{D2. \code{BUDGETED} concede meno concorrenza a parit\`a di
  lavoro.}] Il dato \`e coerente e spiegato dalla legge di Little, ma il
  confronto \`e stato eseguito su un impianto a ciclo chiuso che, per S4, non
  poteva distinguere i segnali dei due controllori. La propriet\`a osservata
  --- il totale sotto budget, la capacit\`a spostata fra funzioni --- \`e reale;
  la sua conseguenza sulla latenza del chiamante non \`e stata misurata su un
  impianto capace di rilevarla.

  \item[\textbf{D3. Il salto di p95 della v0.19.0 \`e uno spostamento del punto
  di lavoro, non una regressione.}] L'attesa in coda sotto il millisecondo e il
  p50 invariato lo rendono plausibile. Non \`e stabilito.
\end{description}

\section{Che cosa non \`e stabilito}

Il progetto enumera esplicitamente cinque questioni aperte, riportate nel
\cref{cap:25-valutazione-concorrenza}. Vale la pena isolarne due, perch\'e
toccano contributi dichiarati nel \cref{cap:01-introduzione}.

\paragraph{I pesi di \code{BUDGETED} non sono mai stati esercitati.} L'equit\`a
max-min \emph{pesata} \`e presentata come contributo, e ogni esecuzione ha usato
pesi identici. Ci\`o che \`e stato validato \`e l'allocazione a pesi uguali, che
\`e il caso degenere. La capacit\`a di servire livelli di servizio differenziati
--- l'argomento principale per cui il peso esiste --- \`e non verificata.

\paragraph{\code{SOJOURN} non \`e mai stato osservato fuori dal sovraccarico.}
Ogni sua esecuzione \`e stata a capacit\`a o oltre. Il modo \`e progettato per
trovare un ginocchio, e il ginocchio esiste solo dove il sistema ha margine.
Il suo comportamento nel regime per cui \`e stato progettato non \`e stato
osservato.

\section{Minacce alla validit\`a}

\subsection{Validit\`a interna}

\begin{description}[style=nextline, leftmargin=0pt]
  \item[\textbf{Numerosit\`a campionaria.}] Tre esecuzioni per punto nelle serie
  di release, una nelle campagne di concorrenza. Non consentono di stimare una
  varianza. Le soglie del gate (10\% e 15\%) sono state scelte per essere
  robuste a questa limitazione, il che significa che il gate \`e cieco alle
  regressioni pi\`u piccole.

  \item[\textbf{Il confondente dei due \code{word-stats}.}] Scoperto in corso
  d'opera. Le due funzioni differiscono per SDK, tipo di build e tempo di
  servizio (5{,}57 contro 2{,}15\,ms). Analisi precedenti alla scoperta sono
  invalidate; quelle riportate qui ne tengono conto, ma la scoperta stessa
  suggerisce che altri confondenti dello stesso tipo possano esistere.

  \item[\textbf{L'errore di calibrazione del budget.}] La prima esecuzione del
  confronto fra controllori ha prodotto 86\,302 rifiuti fabbricati
  dall'aritmetica del budget. La regola derivata --- il picco offerto deve
  essere derivato dal limite pi\`u piccolo che qualunque modo conceder\`a --- \`e
  ora nota, ma \`e stata scoperta dopo un'esecuzione, non prima.

  \item[\textbf{Un test invalido riportato come tale.}] L'ipotesi della
  promozione prematura \`e stata testata con un flag che, a causa di un
  parametro percentuale con default 10, non faceva nulla. L'esito ``nessun
  effetto'' era indistinguibile da quello vero. \`E stato rilevato e corretto;
  la classe di errore resta possibile altrove.
\end{description}

\subsection{Validit\`a esterna}

\begin{description}[style=nextline, leftmargin=0pt]
  \item[\textbf{Un solo profilo hardware.}] Tutta la serie di release \`e su
  \code{azure-d8s-v5+d2s-v5-amd64-native-loadtest-v1}. Non esistono serie
  comparabili su altri profili, e nulla nei dati consente di estrapolare.

  \item[\textbf{Infrastruttura condivisa.}] Due VM Azure dello stesso tipo
  possono avere vicini diversi. Parte dell'oscillazione fra le prime cinque
  release potrebbe essere questa e non il codice.

  \item[\textbf{Un solo scenario di carico.}] \code{autoscaling-cycle-k8s}, con
  corpora di payload fissati. Profili di arrivo diversi --- code pesanti,
  raffiche correlate, payload molto grandi --- non sono stati esplorati.

  \item[\textbf{Poche funzioni contemporanee.}] Le campagne di co-tenancy usano
  due funzioni. Il modo \code{BUDGETED} \`e progettato per una piattaforma
  multi-tenant, e con due tenant l'allocatore max-min opera nel proprio caso pi\`u
  semplice.
\end{description}

\subsection{Validit\`a di costrutto}

\begin{description}[style=nextline, leftmargin=0pt]
  \item[\textbf{La latenza end-to-end \`e misurata dal control plane, non dal
  chiamante.}] \code{function\_e2e\_latency\_ms} parte dall'ammissione. Il tempo
  che una richiesta passa nella rete e nei buffer prima di essere ammessa non \`e
  incluso. Per un chiamante remoto la grandezza rilevante \`e maggiore.

  \item[\textbf{Il tasso di errore esclude le risposte applicative negative.}]
  Si veda il \cref{cap:21-osservabilita}. \`E la definizione corretta per un gate di
  piattaforma, e va tenuta presente confrontando con piattaforme che contano
  diversamente.

  \item[\textbf{Il cold start \`e in parte inferito.}] Il tracker
  dell'autoscaler lo deduce da un evento di scale-up da zero entro trenta
  secondi (\cref{cap:11-autoscaler}), quando il runtime non lo dichiara.
  L'accordo fra inferenza e dichiarazione non \`e stato misurato.
\end{description}

\section{Quando \nf{} \`e la scelta sbagliata}

\begin{itemize}[leftmargin=*]
  \item \textbf{Se l'indisponibilit\`a ha un costo.} Un solo processo di
  controllo, nessuna replicazione, stato in memoria: un riavvio perde il
  registro e interrompe le invocazioni in volo.
  \item \textbf{Se la rete non \`e fidata.} Nessuna autenticazione, nessuna
  autorizzazione.
  \item \textbf{Se servono garanzie di durabilit\`a sull'asincrono.} Le code
  sono in memoria; un riavvio le perde.
  \item \textbf{Se il carico \`e dominato dal cold start.} \nf{} lo misura e non
  contribuisce a ridurlo; le piattaforme citate nel
  \cref{cap:02-stato-dell-arte} hanno meccanismi che \nf{} non ha.
  \item \textbf{Se servono composizione di funzioni, trigger da sorgenti di
  eventi, versioning o traffic splitting.} Non esistono.
\end{itemize}

\section{Riflessione sui vincoli di progetto}

Il \cref{cap:03-requisiti-principi} ha sostenuto che le tre rinunce fondamentali
non fossero debito ma metodo. A valle dei risultati, il bilancio \`e
disomogeneo.

Il vincolo del \textbf{processo singolo} ha pagato. La brevit\`a del percorso
critico \`e ci\`o che ha reso attribuibile una differenza di alcuni millisecondi
di attesa a una politica di ammissione; e --- ironicamente --- \`e anche il
sospettato principale del cross-talk che S1 non ha saputo spiegare altrimenti.
Il vincolo ha prodotto tanto lo strumento di misura quanto il fenomeno da
misurare.

Il vincolo dello \textbf{stato in memoria} ha pagato in parte. Ha eliminato l'I/O
dal percorso critico come promesso, ma ha imposto tre orizzonti di ritenzione
nello \code{ExecutionStore} e un intero apparato di gestione della dimenticanza.
Il costo \`e stato pi\`u alto di quanto il capitolo dei principi lasciasse
intendere.

Il vincolo dell'\textbf{assenza di autenticazione} non ha pagato. Il beneficio
dichiarato --- meno modi in cui un esperimento pu\`o fallire per ragioni non
attinenti --- \`e reale ma piccolo, e nessun risultato di questo documento
dipende da esso. \`E anche la sola rinuncia che impedisce alla piattaforma di
essere usata fuori da un laboratorio. Rivalutandola alla luce dei risultati,
\`e la prima da rimuovere.

\section{Sulla natura dei risultati negativi}

Tre dei risultati principali di questo lavoro sono negativi: il gradiente sul
tempo di servizio non ha ottimo interno, la ricerca del minimo non pu\`o
attribuire il proprio effetto, il ciclo chiuso non pu\`o porre la domanda.
Ciascuno ha richiesto un'implementazione completa e una campagna di misura per
essere stabilito, e ciascuno ha eliminato un'intera famiglia di soluzioni
apparentemente ragionevoli.

Il progetto li conserva nella propria documentazione insieme alle affermazioni
che hanno confutato. \`E una pratica che vale la pena nominare esplicitamente,
perch\'e la letteratura sulle piattaforme raramente la adotta: il costo di non
registrarli \`e che la stessa famiglia di soluzioni viene ritentata, e il segnale
che le esclude non \`e derivabile dalla lettura del codice finale --- che, per
costruzione, contiene solo ci\`o che ha funzionato.

\chapter{Lavori futuri}
\label{cap:28-lavori-futuri}

Le direzioni che seguono sono ordinate per rapporto fra valore atteso e costo,
e ciascuna \`e legata a un vincolo del \cref{cap:03-requisiti-principi} o a una
questione lasciata aperta dal \cref{cap:27-discussione}.

\section{Chiudere le questioni aperte}

Sono gli esperimenti che il progetto pu\`o eseguire con l'apparato che gi\`a
possiede, e che risolverebbero affermazioni oggi non stabilite.

\begin{description}[style=nextline, leftmargin=0pt]
  \item[\textbf{F1. Ripetere il confronto \code{SOJOURN}.}] Una sola esecuzione
  per modo non stabilisce una differenza del 10--15\%. Con cinque esecuzioni per
  modo la differenza diventerebbe difendibile o svanirebbe. \`E l'intervento a
  costo pi\`u basso e valore pi\`u alto dell'elenco.

  \item[\textbf{F2. Osservare \code{SOJOURN} fuori dal sovraccarico.}] Il modo
  \`e progettato per trovare un ginocchio, e ogni sua esecuzione \`e stata a
  capacit\`a o oltre, dove il ginocchio non esiste. Un carico a ciclo aperto con
  picco sotto la capacit\`a misurata direbbe se il modello sceglie limiti sensati
  quando ha margine.

  \item[\textbf{F3. Esercitare i pesi di \code{BUDGETED}.}] Due funzioni con SLO
  e pesi diversi, e la verifica che i limiti concessi seguano i pesi quando il
  budget \`e scarso e li ignorino quando \`e abbondante. \`E la validazione del
  contributo dichiarato che oggi manca del tutto.

  \item[\textbf{F4. Tarare la soglia di guardia al 70\%.}] \`E stata scelta, non
  misurata, e sul profilo a ciclo chiuso ha scavalcato completamente il modello.
  Una scansione su alcuni valori sotto arrivi aperti mostrerebbe se il valore
  conta.

  \item[\textbf{F5. Isolare il meccanismo del cross-talk.}] Il
  \cref{cap:25-valutazione-concorrenza} ha escluso la contesa di CPU e ha
  indicato il control plane condiviso. I candidati --- thread dello scheduler,
  event loop, strutture concorrenti delle code --- sono distinguibili per
  profilazione.
\end{description}

\section{Rimuovere il vincolo che non ha pagato}

Il \cref{cap:27-discussione} ha concluso che l'assenza di autenticazione \`e la
sola delle tre rinunce fondamentali il cui beneficio sperimentale sia
trascurabile, e la sola che impedisca l'uso fuori da un laboratorio.

Introdurre un'autenticazione a token sull'API, con verifica sul percorso prima
della risoluzione del nome, \`e un intervento contenuto. Il costo che va misurato
--- e che va misurato \emph{prima} di dichiarare l'intervento innocuo --- \`e il
suo contributo alla latenza sul percorso caldo. \`E esattamente il tipo di
domanda che questa piattaforma \`e attrezzata a rispondere.

\section{Estendere il controllo della concorrenza}

\subsection{Un ambito globale per \code{SOJOURN}}

Le due innovazioni della Parte III sono oggi disgiunte: \code{BUDGETED} decide
globalmente ma dal tempo di servizio; \code{SOJOURN} decide dal sojourn ma per
funzione. La combinazione naturale --- fabbisogno stimato dal sojourn, allocato
da un budget di piattaforma con equit\`a max-min --- non esiste, e sarebbe la
sintesi dei due risultati.

Il \cref{cap:12-concurrency-control} ha osservato che l'interfaccia dei
controllori non \`e uniforme proprio perch\'e un modo ad ambito globale non si
riconduce alla firma di uno per funzione. Quella asimmetria \`e il primo
ostacolo da rimuovere.

\subsection{Il tetto al ginocchio per \code{SOJOURN}}

Il commento nel codice registra la scelta e la sua condizione di revisione:
\emph{``no explicit knee cap --- the Little term already asks for no more than
the load needs, and the configured ceiling bounds the rest. Add one if a
function is seen climbing while its throughput stays flat.''} \`E un'osservazione
che F2 potrebbe produrre.

\subsection{Un minimo su finestra al posto delle derive}

Tre componenti mantengono estremi storici con un decadimento forfettario:
la baseline del modo adattivo (0{,}2\% per tick), e i due estremi dello
stimatore a budget (1\% per tick). Ciascuno \`e commentato come ``la versione
economica'' di un minimo su finestra con scadenza, che \`e la costruzione che
BBR e CoDel~\cite{nichols2012codel} usano.

\section{Copertura della combinatoria modulare}

Lo spazio delle configurazioni ha $2^9$ punti e la suite ne verifica due
(\cref{cap:16-moduli-diagnostici}). Una verifica di costruibilit\`a e
avviabilit\`a su tutte le combinazioni valide \`e realizzabile con gli strumenti
presenti, e coprirebbe la classe di difetti che il sistema modulare rende
possibile: una combinazione che compila e non si avvia.

Nella stessa direzione va la verifica di corrispondenza fra \code{openapi.yaml} e
le rotte effettivamente registrate (\cref{cap:05-contratti-dati}), e il cancello
end-to-end sugli scaffold generati (\cref{cap:19-cli-scaffolding}). I tre
interventi condividono la stessa forma: verificare artefatti dichiarati anzich\'e
soltanto codice eseguito.

\section{Offload: da segnale a modello}

Il modulo \code{offload} decide dal rifiuto della coda locale, senza alcun
modello del costo di trasferimento (\cref{cap:13-offload}). \`E una scelta
difendibile perch\'e non richiede calibrazione, e lascia sul tavolo tutto ci\`o
che la letteratura sull'offloading edge--cloud ha prodotto.

Due estensioni sono a portata:

\begin{itemize}[leftmargin=*]
  \item \textbf{Includere la latenza remota nella decisione.} Il gateway pu\`o
  misurare il round trip verso il target; confrontarlo con l'attesa stimata
  localmente trasformerebbe la decisione binaria in un confronto.
  \item \textbf{Offload preventivo.} Oggi si scarica solo dopo un rifiuto. Una
  frazione del traffico deviata \emph{prima} della congestione eviterebbe di
  raggiungerla.
\end{itemize}

Entrambe richiedono la campagna sperimentale che il
\cref{cap:13-offload} segnala come mancante: p95 end-to-end e tasso di rifiuto
con e senza offload, al variare del RTT verso il target.

\section{Alta disponibilit\`a: il vincolo da non rimuovere per primo}

\`E la direzione pi\`u ovvia e quella che raccomanderei per ultima.

Replicare il control plane richiede di replicare il registro, lo store delle
esecuzioni, le chiavi di idempotenza e le code, e di riconciliare le decisioni di
controllori che osservano stati parzialmente sovrapposti. Il
\cref{cap:03-requisiti-principi} ha argomentato che le oscillazioni cos\`i
introdotte sono un fenomeno \emph{diverso} da quelli studiati qui, e
indistinguibili da essi.

Il valore di una versione ad alta disponibilit\`a \`e in un'altra direzione: il
coordinamento fra controllori replicati \`e un problema interessante di per s\'e,
e questa piattaforma sarebbe un buon banco per studiarlo. Ma sarebbe una
piattaforma diversa, e il \cref{cap:29-conclusioni} sostiene che la sua
utilit\`a attuale \`e proporzionale a ci\`o che non fa.

\section{Sintesi}

\begin{table}[htbp]
\centering
\caption{Direzioni future, per rapporto valore/costo.}
\label{tab:futuro}
\small
\begin{tabularx}{\textwidth}{@{}lXll@{}}
\toprule
\textbf{Id} & \textbf{Intervento} & \textbf{Costo} & \textbf{Valore} \\
\midrule
F1 & Ripetere il confronto \code{SOJOURN} & basso & alto \\
F3 & Esercitare i pesi di \code{BUDGETED} & basso & alto \\
F2 & \code{SOJOURN} fuori dal sovraccarico & basso & medio \\
F5 & Isolare il meccanismo del cross-talk & medio & alto \\
--- & Autenticazione a token & basso & alto (uso) \\
--- & \code{SOJOURN} ad ambito globale & alto & alto \\
--- & Combinatoria modulare verificata & medio & medio \\
--- & Offload con modello di costo & medio & medio \\
--- & Alta disponibilit\`a & molto alto & basso per gli obiettivi attuali \\
\bottomrule
\end{tabularx}
\end{table}

Le prime quattro righe hanno una propriet\`a comune: non richiedono codice nuovo.
Sono esperimenti da eseguire su ci\`o che esiste gi\`a, e la loro presenza in cima
alla lista \`e coerente con la tesi del documento --- che il valore di questa
piattaforma stia in ci\`o che consente di misurare.

\chapter{Conclusioni}
\label{cap:29-conclusioni}

Il \cref{cap:01-introduzione} ha aperto con una domanda: qual \`e la piattaforma
FaaS pi\`u piccola che sia ancora una piattaforma FaaS, e che cosa diventa
misurabile una volta ridotta a quella dimensione? Questo capitolo risponde alle
quattro domande di ricerca che ne discendevano.

\section{Le risposte}

\subsection{RQ1 --- Modularit\`a a costo zero sul percorso critico}

\textbf{S\`i, ed \`e verificabile.} La selezione modulare a tempo di build,
combinata con implementazioni no-op per ciascun punto di estensione, produce un
percorso di invocazione privo di rami condizionali sulla configurazione: la
catena di ripiego \emph{coda sincrona $\to$ coda asincrona $\to$ dispatch
diretto} non testa la presenza di un modulo, esegue il default del punto di
estensione precedente. La regola ArchUnit R3 rende la propriet\`a accertata dalla
build anzich\'e affermata.

Il costo esiste e \`e stato quantificato: due invocazioni della suite di test per
ogni verifica completa, e una copertura che riguarda i soli estremi dello spazio
delle configurazioni. La modularit\`a a costo zero sul percorso critico non \`e a
costo zero sulla catena di integrazione.

Le proporzioni, misurate di nuovo, non sono pi\`u quelle che questo capitolo
riportava: i dieci moduli sommano 6\,299 righe non vuote di produzione, contro
le 17\,334 del control plane e delle sue tre librerie obbligatorie
(\cref{cap:01-introduzione}). La riducibilit\`a della piattaforma non \`e quindi
la maggioranza del codice che si pu\`o non compilare, ma l'insieme delle
politiche che si possono non compilare; \`e un'affermazione pi\`u stretta e
verificata, non una proporzione.

\subsection{RQ2 --- Da quale segnale decidere}

\textbf{Dal tempo di sojourn, non da quello di servizio.} La risposta \`e
stabilita su due piani.

Analiticamente: il tempo di servizio cresce monotonamente con il limite di
concorrenza su qualunque risorsa condivisa, quindi una regola che riduce il
limite quando la latenza cresce non ha ottimo interno e cammina fino al proprio
pavimento. Il tempo di sojourn --- attesa pi\`u servizio --- l'ottimo interno ce
l'ha, perch\'e l'attesa decresce dove il servizio cresce.

Sperimentalmente: sotto carico l'attesa \`e stata misurata a 37--43\,ms contro un
servizio di circa 5\,ms, e un controllore pu\`o quindi sedere comodamente dentro
un SLO di servizio di 10\,ms mentre i suoi chiamanti attendono ottanta.

La realizzazione ha richiesto tre tentativi, e i due falliti sono contributi
quanto quello riuscito. Il gradiente sulla latenza end-to-end \`e instabile per
retroazione positiva. La ricerca del minimo per passi successivi \`e ben posta in
teoria e fallisce perch\'e non pu\`o distinguere ``la mia mossa ha aiutato'' da
``il carico \`e calato'' --- 36\,181 richieste rifiutate contro 32. Il progetto
che funziona \emph{calcola anzich\'e confrontare}: guardia di ammissione, legge di
Little su un tasso predetto, orizzonte di drenaggio proporzionale allo
scostamento dalla promessa. Sotto arrivi aperti ha servito il 3{,}2\% di
richieste in pi\`u, ne ha rifiutate il 15{,}7\% in meno e ha ridotto il p95
end-to-end del 10{,}6\% --- con \emph{meno} concorrenza media.

Il risultato \`e di una sola esecuzione per modo, e il
\cref{cap:27-discussione} lo classifica come suggerito e non stabilito.

\subsection{RQ3 --- Concorrenza come risorsa condivisa}

\textbf{Separare la stima del fabbisogno dall'allocazione funziona; la sua
superiorit\`a non \`e stata dimostrata.}

Il modo \code{BUDGETED} spezza la domanda in due --- quanto serve a una funzione
per il proprio SLO, quanto la piattaforma pu\`o dare --- e risolve la seconda con
equit\`a max-min pesata. Misurato, mantiene il totale sotto il budget e sposta
capacit\`a fra le funzioni al muoversi del carico, cosa che due controllori
indipendenti non possono fare per costruzione. Ha concesso alla funzione che non
ne aveva bisogno met\`a della concorrenza a parit\`a di throughput.

Due riserve. Il confronto \`e stato eseguito su un impianto a ciclo chiuso che,
come RQ4 ha poi mostrato, non poteva distinguere i segnali dei due controllori.
E i pesi --- l'argomento principale per cui il meccanismo esiste --- non sono mai
stati esercitati con valori diversi.

La campagna di co-tenancy ha per\`o prodotto il risultato pi\`u solido del lavoro,
e lo ha prodotto in negativo: il cross-talk fra funzioni co-residenti \`e reale
($-28\%$ di throughput sulla funzione stabile) e \textbf{non} \`e contesa di CPU
(vincolandole agli stessi quattro core, lo 0{,}6\% di differenza). La premessa
del modo a budget --- che esista una risorsa condivisa da dividere --- \`e
confermata; la sua identit\`a non \`e la CPU.

\subsection{RQ4 --- Il costo reale di un control plane}

\textbf{Meno di un core e circa 200\,MB di heap per 300 richieste al secondo, e
la scelta della catena di build conta pi\`u della scelta del framework.}

Tre risultati.

Per una JVM Spring, lo heap --- i dati --- \`e il 12\% della memoria residente. Il
resto \`e il runtime che descrive se stesso. Ottimizzare le allocazioni agisce
sull'ottavo del problema.

Il collector seriale di GraalVM Community spende il 31{,}8\% del tempo di parete
a raccogliere sotto carico sostenuto, con pause fino a 1{,}77\,s, e dimezza il
throughput. G1 lo elimina: 98{,}5\% del throughput della JVM, coda a 123\,ms,
avvio in 0{,}07\,s. Compilato a \code{-O3}, il control plane nativo supera la
JVM --- 9\,251 contro 8\,205 richieste al secondo, con un avvio ventitr\'e volte
pi\`u rapido.

Il terzo risultato \`e metodologico e vale oltre questo progetto: \emph{misure di
memoria senza un limite del container descrivono appetito, non fabbisogno}. Con
un limite a 1\,GB e uno heap esplicito, la stessa build consegna 3\,089
richieste al secondo in un terzo della memoria che l'esecuzione senza vincoli
aveva occupato.

\section{La tesi, rivista}

La tesi del \cref{cap:01-introduzione} era che un insieme selezionato di rinunce
produce un sistema pi\`u utile come strumento di misura di quanto lo sarebbe la
sua versione completa. A valle dei risultati, il bilancio \`e disomogeneo e vale
la pena dichiararlo.

Il \textbf{processo singolo} ha pagato pi\`u di quanto promesso. Ha reso
attribuibile una differenza di pochi millisecondi di attesa a una politica di
ammissione --- che \`e l'intero contenuto del \cref{cap:25-valutazione-concorrenza}
--- e si \`e rivelato, per esclusione, anche il sospettato principale del
cross-talk fra funzioni co-residenti. Il vincolo ha prodotto lo strumento e il
fenomeno.

Lo \textbf{stato in memoria} ha pagato in parte, a un costo pi\`u alto del
previsto: tre orizzonti di ritenzione e un intero apparato di dimenticanza per
evitare che uno store che non persiste diventi uno store che non dimentica.

L'\textbf{assenza di autenticazione} non ha pagato. Nessun risultato di questo
documento vi dipende, ed \`e la sola rinuncia che impedisca l'uso fuori da un
laboratorio.

\section{Il contributo che non era previsto}

Il contributo di cui vado pi\`u convinto non compare fra i sei elencati nel
\cref{cap:01-introduzione}, ed \`e la pratica di conservare i risultati negativi
insieme alle affermazioni che hanno confutato.

Tre dei risultati principali sono negativi. Ciascuno ha richiesto
un'implementazione completa e una campagna di misura, e ciascuno ha eliminato una
famiglia intera di soluzioni apparentemente ragionevoli: il gradiente sul tempo
di servizio, la ricerca del minimo per confronto fra istanti, il carico a ciclo
chiuso per valutare politiche di ammissione. A questi si aggiungono un errore di
calibrazione che ha fabbricato 86\,302 rifiuti, un test invalidato da un flag
inefficace, e un confondente --- due funzioni omonime con runtime diversi ---
che ha invalidato analisi precedenti.

Nessuno di questi \`e derivabile dalla lettura del codice finale, che per
costruzione contiene solo ci\`o che ha funzionato. Il costo di non registrarli \`e
che la stessa famiglia di soluzioni viene ritentata, e il segnale che le esclude
va riprodotto da capo.

\section{Chiusura}

\nf{} \`e piccolo di proposito. Il suo percorso critico attraversa un solo
processo, e le sue politiche --- code, autoscaling, controllo della concorrenza,
offload --- si possono non compilare. Non compete con le piattaforme mature su completezza e
non aspira a farlo.

Ci\`o che ha prodotto \`e una risposta precisa a una domanda che le piattaforme
complete rendono difficile porre: da quale grandezza misurata debba decidere un
controllore di ammissione, e che cosa succede quando decide dalla grandezza
sbagliata. La risposta --- il tempo di sojourn, non quello di servizio --- \`e
banale una volta enunciata, ed \`e costata due progetti falliti e la scoperta che
il generatore di carico stava impedendo alla domanda di essere posta.

\`E il tipo di risultato che una piattaforma abbastanza piccola da essere letta
per intero consente di ottenere, e che una piattaforma abbastanza grande da
essere utile in produzione avrebbe nascosto sotto tre livelli di proxy.


\appendix
\chapter{Riferimento API}
\label{cap:A-api}

Sintesi dell'API HTTP del control plane, derivata da \code{openapi.yaml} nella
radice del repository e dai controller. Come osservato nel
\cref{cap:05-contratti-dati}, il file di specifica \`e mantenuto a mano e la sua
corrispondenza con le rotte registrate non \`e verificata automaticamente.

\section{Gestione delle funzioni}

\begin{center}
\small
\begin{adjustbox}{max width=\textwidth}
\begin{tabular}{@{}llp{6.4cm}@{}}
\toprule
\textbf{Op.} & \textbf{Percorso} & \textbf{Esiti} \\
\midrule
\code{GET} & \code{/v1/functions} & \code{200} elenco \\
\code{POST} & \code{/v1/functions} &
\code{201} creata; \code{409} nome gi\`a presente; \code{422} specifica non
valida; \code{424} validazione immagine fallita; \code{503} nessun provider
disponibile \\
\code{GET} & \code{/v1/functions/\{name\}} & \code{200}; \code{404} \\
\code{PATCH} & \code{/v1/functions/\{name\}} &
\code{200}; \code{400} campo non modificabile o valore invalido; \code{404} \\
\code{DELETE} & \code{/v1/functions/\{name\}} & \code{204}; \code{404} \\
\code{GET} & \code{/v1/functions/\{name\}/replicas} &
\code{200}; \code{400}; \code{404}; \code{503} \\
\code{PUT} & \code{/v1/functions/\{name\}/replicas} &
\code{200}; \code{400}; \code{404}; \code{503} \\
\bottomrule
\end{tabular}
\end{adjustbox}
\end{center}

Il corpo della registrazione \`e un \code{FunctionSpec}
(\cref{cap:05-contratti-dati}). Il corpo del \code{PATCH} \`e un
\code{FunctionUpdateRequest} e accetta soltanto quattro campi
(\cref{cap:06-nucleo}); le propriet\`a sconosciute sono rifiutate.

\section{Invocazione}

\begin{center}
\small
\begin{adjustbox}{max width=\textwidth}
\begin{tabular}{@{}llp{6.2cm}@{}}
\toprule
\textbf{Op.} & \textbf{Percorso} & \textbf{Esiti} \\
\midrule
\code{POST} & \code{/v1/functions/\{name\}:invoke} &
\code{200} (o il codice deciso dalla funzione); \code{404}; \code{408};
\code{429}; \code{500}; \code{502}; \code{504} \\
\code{POST} & \code{/v1/functions/\{name\}:enqueue} &
\code{202}; \code{404}; \code{429}; \code{500}; \code{501} \\
\code{GET} & \code{/v1/executions/\{executionId\}} & \code{200}; \code{404} \\
\code{POST} & \code{/v1/internal/executions/\{id\}:complete} & \code{204} \\
\bottomrule
\end{tabular}
\end{adjustbox}
\end{center}

Codici con semantica specifica:

\begin{center}
\small
\begin{adjustbox}{max width=\textwidth}
\begin{tabular}{@{}lp{9cm}@{}}
\toprule
\textbf{Codice} & \textbf{Condizione} \\
\midrule
\code{429} & Limite di frequenza globale, oppure rifiuto della coda sincrona.
Distinguibili dagli header (sotto). \\
\code{501} & Il percorso asincrono richiede il modulo \code{async-queue}, assente
in questo binario (\cref{cap:16-moduli-diagnostici}). \\
\code{502} & L'offload verso l'istanza remota \`e fallito
(\cref{cap:13-offload}). Nessun fallback locale. \\
\code{504} & L'istanza remota non ha risposto entro il budget. \\
\bottomrule
\end{tabular}
\end{adjustbox}
\end{center}

\section{Header}

\subsection{Richiesta}

\begin{center}
\small
\begin{adjustbox}{max width=\textwidth}
\begin{tabular}{@{}lp{8.4cm}@{}}
\toprule
\textbf{Header} & \textbf{Significato} \\
\midrule
\code{Idempotency-Key} & Deduplica l'invocazione sulla funzione
(\cref{cap:06-nucleo}). \\
\code{X-Trace-Id} & Identificatore di traccia propagato fino al runtime. \\
\code{X-Timeout-Ms} & Sovrascrive il timeout della specifica per questa
invocazione. \\
\code{X-NanoFaaS-Offload-Hop} & Marca una richiesta gi\`a inoltrata: vieta un
secondo salto. \\
\code{traceparent}, \code{tracestate} & Contesto di traccia W3C, propagato
sull'offload. \\
\bottomrule
\end{tabular}
\end{adjustbox}
\end{center}

Tutti gli altri header del chiamante, esclusi quelli hop-by-hop, sono resi
visibili all'handler nel campo \code{headers} della richiesta, che viene
\emph{sovrascritto} e mai fuso con quanto dichiarato nel corpo
(\cref{cap:04-architettura-control-plane}).

\subsection{Risposta}

\begin{center}
\small
\begin{adjustbox}{max width=\textwidth}
\begin{tabular}{@{}lp{8.2cm}@{}}
\toprule
\textbf{Header} & \textbf{Significato} \\
\midrule
\code{X-Execution-Id} & Identificatore dell'esecuzione. \\
\code{X-NanoFaaS-Function-Status} & \code{true} se il codice di stato \`e stato
deciso dall'handler (\cref{cap:17-runtime-funzione}). \\
\code{X-NanoFaaS-Encoding} & Codifica dell'output, per esempio \code{base64}. \\
\code{X-NanoFaaS-Offloaded} & URL dell'istanza remota che ha servito. \\
\code{Retry-After} & Su \code{429} dalla coda sincrona. \\
\code{X-Queue-Reject-Reason} & \code{depth}, \code{est\_wait} o \code{timeout}
(\cref{cap:10-sync-queue}). \\
\bottomrule
\end{tabular}
\end{adjustbox}
\end{center}

\subsection{Dal runtime al control plane}

\begin{center}
\small
\begin{adjustbox}{max width=\textwidth}
\begin{tabular}{@{}ll@{}}
\toprule
\textbf{Header} & \textbf{Significato} \\
\midrule
\code{X-Cold-Start} & L'invocazione ha pagato un'inizializzazione. \\
\code{X-Init-Duration-Ms} & Durata di quell'inizializzazione. \\
\code{X-NanoFaaS-Function-Status} & Il codice di stato \`e intenzionale. \\
\code{X-Dispatch-Attempt} & Numero di tentativo, sulla callback. \\
\bottomrule
\end{tabular}
\end{adjustbox}
\end{center}

\section{API di amministrazione}

Attiva soltanto con \code{nanofaas.admin.runtime-config.enabled=true}; altrimenti
i percorsi rispondono \code{404} (\cref{cap:14-runtime-config}).

\begin{center}
\small
\begin{adjustbox}{max width=\textwidth}
\begin{tabular}{@{}llp{6cm}@{}}
\toprule
\textbf{Op.} & \textbf{Percorso} & \textbf{Esiti} \\
\midrule
\code{GET} & \code{/v1/admin/runtime-config} & \code{200}; \code{404} disattivata \\
\code{POST} & \code{/v1/admin/runtime-config/validate} & \code{200};
\code{400}; \code{422} \\
\code{PATCH} & \code{/v1/admin/runtime-config} & \code{200}; \code{400}
(\code{expectedRevision} mancante o durata malformata); \code{404}; \code{409}
revisione non corrispondente; \code{422}; \code{503} applicazione fallita \\
\bottomrule
\end{tabular}
\end{adjustbox}
\end{center}

Le durate usano il formato ISO-8601 (\code{PT5S}).

\section{Management}

Sulla porta \code{8081}:

\begin{center}
\small
\begin{adjustbox}{max width=\textwidth}
\begin{tabular}{@{}ll@{}}
\toprule
\code{GET /actuator/health/liveness} & Sonda di liveness \\
\code{GET /actuator/health/readiness} & Sonda di readiness \\
\code{GET /actuator/prometheus} & Metriche (\cref{cap:C-metriche}) \\
\code{GET /modules/build-metadata} & Identit\`a della build, se il modulo \`e
compilato \\
\bottomrule
\end{tabular}
\end{adjustbox}
\end{center}

\chapter{Configurazione completa}
\label{cap:B-configurazione}

\section{Configurazione di piattaforma}

Da \code{platform/control-plane/src/main/resources/application.yml}. Sono omessi
l'elenco \code{spring.autoconfigure.exclude} (le auto-configurazioni di Spring
Boot di cui il control plane non usa alcun bean: client HTTP, TLS, esecutori di
task, multipart, sessioni, AOP, \emph{info} e spazio su disco) e il blocco
\code{management.endpoint}, che abilita solo \code{health} con le probe e
\code{prometheus}. Le chiavi sono raggruppate per argomento; i commenti sono
del paper, non del file.

\begin{lstlisting}[language=yamlnf]
server:
  port: 8080
management:
  server:
    port: 8081
  endpoints:
    web:
      exposure:
        include: health,prometheus

nanofaas:
  metrics:
    profile: basic              # basic | advanced | soak
  admin:
    runtime-config:
      enabled: false
  scheduler:
    strategy: ${NANOFAAS_SCHEDULER_STRATEGY:}   # vuoto = nessuna scelta esplicita
  deployment:
    default-backend: ""         # k8s | container-local; vuoto = risoluzione implicita
    wakeup:
      timeout: 30s
      poll-interval: 250ms
      ready-observation-max-age: 5s
  registry:
    path: "${NANOFAAS_REGISTRY_PATH:build/nanofaas/functions.json}"
  defaults:
    timeoutMs: 30000
    concurrency: 4
    queueSize: 100
    maxRetries: 3
  retry:
    initial-backoff: 100ms
    max-backoff: 2s
  rate:
    maxPerSecond: 1000000
  k8s:
    namespace: ""
    callbackUrl: "http://control-plane.default.svc.cluster.local:8080/v1/internal/executions"
  scaling:
    poll-interval-ms: 5000
    default-min-replicas: 1
    default-max-replicas: 10
  execution-store:
    ttl: 5m
    sync-ttl: 30s
    max-lifetime: 30m
    max-outcomes: 100000
    max-keys: 100000
    max-outcome-bytes: 0        # 0 = max-outcomes x 116 byte
  http-client:
    max-connections: 500
    pending-acquire-max-count: 0        # 0 = 2 x max-connections
    pending-acquire-timeout-ms: 45000
    max-idle-time-ms: 30000
    max-life-time-ms: 0                 # 0 = rotazione disabilitata
    eviction-interval-ms: 5000
    inactive-pool-dispose-interval-ms: 5000
    pool-inactivity-ms: 30000
  container-local:
    proxy:
      max-request-bytes: 2097152
      max-response-bytes: 4194304
      max-buffered-bytes: 33554432
      inbound-read-timeout: 5s
      response-write-timeout: 5s

sync-queue:
  enabled: true
  admission-enabled: true
  max-depth: 200
  max-estimated-wait: 2s
  max-queue-wait: 2s
  retry-after-seconds: 2
  throughput-window: 30s
  per-function-min-samples: 50
\end{lstlisting}

\paragraph{Che cosa governano i blocchi nuovi.}
\code{nanofaas.scheduler.strategy} sceglie la strategia su cui parte il motore
composto (\code{per-function} o \code{shared-queue}, \cref{sec:09-switch}); il
valore vuoto non \`e una terza strategia ma ``nessuna scelta esplicita'' e il
motore deriva la mappatura dai moduli di coda presenti nell'artefatto. Un
identificatore che nessuna strategia fornisce fa fallire l'avvio. Il cambio a
caldo passa dall'endpoint di amministrazione (\cref{cap:A-api}); l'ordine di
grandezza del budget di preparazione \`e una costante del motore (50\,ms), non
una chiave, e le soglie di verifica sono congelate in
\code{docs/experiments/scheduler-switching-2026-09/budgets.json}.
\code{nanofaas.retry.*} governa il backoff tra i tentativi
(\cref{sec:06-riprova}). \code{nanofaas.execution-store.*} governa la
ritenzione degli esiti: \code{max-outcome-bytes} \`e l'unico limite che la
cache degli esiti fa rispettare, mentre \code{max-outcomes} serve solo a
derivarne il default; \code{max-keys} limita le chiavi di idempotenza, e a
budget esaurito una \emph{nuova} ammissione con chiave \`e rifiutata con 429
prima del dispatch. \code{nanofaas.http-client.*} descrive il pool di
connessioni verso le funzioni, per destinazione concreta e non in aggregato.
\code{nanofaas.registry.path} \`e il file su cui persiste il catalogo delle
funzioni. \code{nanofaas.deployment.wakeup.*} regola l'attesa dell'avvio di una
funzione scalata a zero.

\subsection{Tetti di capacit\`a di invocazione}

Il blocco \code{nanofaas.invocation-capacity} fissa i tetti di ci\`o che pu\`o
essere vivo contemporaneamente: sono limiti di occupazione, non limiti di
frequenza, e non promettono un RSS. I default sono calibrati sul minimo
supportato dal chart, un container da 512\,MiB.

\begin{center}
\small
\begin{adjustbox}{max width=\textwidth}
\begin{tabular}{@{}lrr@{}}
\toprule
\textbf{Chiave} (\code{nanofaas.invocation-capacity.}) & \textbf{Globale} &
\textbf{Per funzione} \\
\midrule
\code{executions-*} & 4096 & 512 \\
\code{canonical-input-bytes-*} & 134\,217\,728 & 33\,554\,432 \\
\code{physical-input-copy-bytes-*} & 67\,108\,864 & 16\,777\,216 \\
\code{waiters-*} & 8192 & 1024 \\
\midrule
\code{ingress-body-bytes} & \multicolumn{2}{r}{1\,048\,576 (1\,MiB)} \\
\code{max-input-references} & \multicolumn{2}{r}{64} \\
\code{retained-input-max-depth} & \multicolumn{2}{r}{32} \\
\code{retained-input-max-container-entries} & \multicolumn{2}{r}{16\,384} \\
\code{retained-input-max-visited-nodes} & \multicolumn{2}{r}{65\,536} \\
\code{retained-input-max-bytes-per-execution} & \multicolumn{2}{r}{1\,048\,576} \\
\bottomrule
\end{tabular}
\end{adjustbox}
\end{center}

\section{Chiavi dei moduli non presenti nel file di default}

Ogni modulo dichiara in \code{module.properties} se \`e abilitato di default
(\code{defaultEnabled}). Lo \`e ogni modulo tranne i due provider di deployment
locali, \code{container-deployment-provider} e
\code{containerd-deployment-provider}: \code{sync-queue} \`e quindi attivo in
ogni build che lo include, senza bisogno di impostare una chiave. Il profilo di
ammissione (\code{nanofaas.admission.profile}: \emph{direct},
\emph{function-queue} o \emph{sync-queue}), se non \`e dichiarato, \`e derivato
dalle strategie presenti: \code{per-function} d\`a \emph{function-queue},
altrimenti \code{shared-queue} d\`a \emph{sync-queue}, altrimenti \emph{direct}.

\begin{center}
\small
\begin{adjustbox}{max width=\textwidth}
\begin{tabular}{@{}llp{5.4cm}@{}}
\toprule
\textbf{Chiave} & \textbf{Default} & \textbf{Modulo} \\
\midrule
\code{nanofaas.concurrency-control.poll-interval-ms} & 5000 &
\code{concurrency-control} \\
\code{nanofaas.concurrency-control.default-target-in-flight-per-pod} & 2 & \\
\code{nanofaas.concurrency-control.total-budget} & $\max(8, 4 \times$ core$)$ & \\
\midrule
\code{nanofaas.offload.enabled} & \code{true} & \code{offload} \\
\code{nanofaas.offload.target-url} & --- & \\
\code{nanofaas.offload.pressure-enabled} & \code{true} & \\
\midrule
\code{nanofaas.container-local.runtime-adapter} & \code{docker} &
\code{container-deployment-provider} \\
\code{nanofaas.container-local.bind-host} & \code{127.0.0.1} & \\
\code{nanofaas.container-local.readiness-timeout} & 20\,s & \\
\code{nanofaas.container-local.readiness-poll-interval} & 250\,ms & \\
\code{nanofaas.container-local.network-name} & --- & \\
\code{nanofaas.container-local.cpuset} & --- & vincola le funzioni agli stessi
core (\cref{cap:15-deployment-provider}) \\
\midrule
\code{nanofaas.containerd.socket-path} & \code{\$XDG\_RUNTIME\_DIR/containerd/containerd.sock} &
\code{containerd-deployment-provider} \\
\code{nanofaas.containerd.namespace} & \code{nanofaas} & \\
\code{nanofaas.containerd.runtime-binary} & \code{crun} & \\
\code{nanofaas.containerd.snapshotter} & \code{native} & \\
\code{nanofaas.containerd.network-name} & \code{nanofaas} & \\
\code{nanofaas.containerd.bind-host} & \code{127.0.0.1} & \\
\code{nanofaas.containerd.cgroups-path}, \code{cpuset}, \code{callback-url}, \code{state-directory}, \code{cni-*}, \code{readiness-*}, \code{*-timeout} & vedi codice & \\
\bottomrule
\end{tabular}
\end{adjustbox}
\end{center}

\section{Configurazione per funzione}

Tutti i campi di \code{FunctionSpec}. \code{null} significa ``non specificato'' e
attiva il default di piattaforma.

\begin{lstlisting}[language=yamlnf]
name: <obbligatorio, non vuoto>
image: <obbligatorio, non vuoto>
command: [...]                  # default []
env: { ... }                    # default {}
resources:
  requests: { cpu: <positivo>, memoryMiB: <positivo> }
  limits:   { cpu: <positivo>, memoryMiB: <positivo> }   # richiesta <= limite
timeoutMs: <>= 1>
concurrency: <>= 1>
queueSize: <>= 1>
maxRetries: <>= 0>
endpointUrl: <obbligatorio per EXTERNAL>
executionMode: LOCAL | EXTERNAL | DEPLOYMENT      # default DEPLOYMENT
runtimeMode: HTTP | STDIO | FILE
runtimeCommand: <comando del watchdog>
imagePullSecrets: [...]
scalingConfig:
  strategy: HPA | INTERNAL | NONE                 # default INTERNAL
  minReplicas: <>= 0>                             # default 1
  maxReplicas: <>= 1>                             # default 10
  metrics:
    - type: queue_depth | in_flight | rps
      target: "<numero>"
  concurrencyControl:
    mode: FIXED | STATIC_PER_POD | ADAPTIVE_PER_POD | BUDGETED | SOJOURN
    targetInFlightPerPod: 2
    minTargetInFlightPerPod: 1
    maxTargetInFlightPerPod: 8
    upscaleCooldownMs: 30000
    downscaleCooldownMs: 60000
    highLoadThreshold: 0.5
    lowLoadThreshold: 0.15
    targetLatencyMs: 250       # servizio per BUDGETED, end-to-end per SOJOURN
    weight: 1.0
offload:
  enabled: <true|false>
  targetUrl: <url>
  mode: pressure | always
\end{lstlisting}

\section{Costanti non configurabili}

Valori fissati nel codice, elencati perch\'e sono i candidati naturali per un
esperimento che voglia variarli.

\begin{center}
\small
\begin{adjustbox}{max width=\textwidth}
\begin{tabular}{@{}llp{5cm}@{}}
\toprule
\textbf{Costante} & \textbf{Valore} & \textbf{Dove} \\
\midrule
Batch massimo per funzione (strategia \code{per-function}) & 2 & \code{async-queue} \\
Finestra di scansione della coda condivisa (\code{shared-queue}) & 64 & \code{sync-queue} \\
Budget di preparazione di un cambio di scheduler & 50\,ms & motore (\code{SchedulerEngine}) \\
Cooldown di scale-up / scale-down & 30\,s / 60\,s & \code{autoscaler} \\
Finestra di inferenza del cold start & 30\,s & \code{autoscaler} \\
Deriva della baseline & $1{,}002$ per tick & \code{concurrency-control} \\
Decadimento degli estremi & $1{,}01$ per tick & \code{concurrency-control} \\
Pavimento del gradiente & $0{,}5$ & \code{concurrency-control} \\
Soglia di guardia della coda & $0{,}7$ & \code{concurrency-control} \\
Orizzonte di drenaggio & 5\,s (min 1\,s) & \code{concurrency-control} \\
Holt: $\alpha$, $\beta$, $\phi$ & $0{,}5$, $0{,}3$, $0{,}8$ &
\code{concurrency-control} \\
Margine di timeout dell'offload & 50\,ms & \code{offload} \\
\bottomrule
\end{tabular}
\end{adjustbox}
\end{center}

\section{Selettori di build}

\begin{center}
\small
\begin{adjustbox}{max width=\textwidth}
\begin{tabular}{@{}lp{7.6cm}@{}}
\toprule
\textbf{Selettore} & \textbf{Effetto} \\
\midrule
\code{-PcontrolPlaneModules} & \code{none}, \code{all} (default) o lista; anche
\code{NANOFAAS\_CONTROL\_PLANE\_MODULES} \\
\code{-PnativeOptimization} & \code{3} (default) o \code{s} \\
\code{-PnativeGc} & \code{G1} richiede Oracle GraalVM; anche
\code{NATIVE\_GC} \\
\code{-PnativeMonitoring} & per esempio \code{jfr,jvmstat} \\
\code{-PnativeBuildMemory} & heap del \emph{compilatore}, non dell'immagine \\
\code{-PnativeParallelism} & lavoratori del compilatore \\
\bottomrule
\end{tabular}
\end{adjustbox}
\end{center}

\chapter{Catalogo delle metriche}
\label{cap:C-metriche}

Tutte le metriche sono esposte su \code{/actuator/prometheus}, porta
\code{8081}. La colonna \emph{profilo} indica se la metrica \`e presente nel
profilo \code{basic} (default), soltanto in \code{advanced} (che \`e cumulativo
e vale anche per \code{soak}) o soltanto in \code{soak}
(\cref{cap:21-osservabilita}). Il profilo attivo \`e leggibile dal gauge
\code{nanofaas\_metrics\_profile\_info}.

\section{Nucleo, per funzione}

\begin{center}
\small
\begin{adjustbox}{max width=\textwidth}
\begin{tabular}{@{}lllp{4.4cm}@{}}
\toprule
\textbf{Metrica} & \textbf{Tipo} & \textbf{Profilo} & \textbf{Note} \\
\midrule
\code{function\_enqueue\_total} & counter & basic & \\
\code{function\_dispatch\_total} & counter & basic & cumulativo; la metrica
\code{rps} \`e la sua derivata \\
\code{function\_success\_total} & counter & basic & esclude le risposte non-2xx
decise dall'handler \\
\code{function\_error\_total} & counter & basic & \\
\code{function\_retry\_total} & counter & basic & \\
\code{function\_timeout\_total} & counter & basic & \\
\code{function\_queue\_rejected\_total} & counter & basic & \\
\code{function\_latency\_ms} & timer & basic & \textbf{tempo di servizio}:
dispatch $\to$ completamento \\
\code{function\_e2e\_latency\_ms} & timer & basic & \textbf{sojourn}: ammissione
$\to$ completamento \\
\code{function\_queue\_wait\_ms} & timer & advanced & attesa in coda, per
tentativo dispatchato \\
\code{function\_init\_duration\_ms} & timer & advanced & inizializzazione
riportata dal runtime \\
\code{function\_cold\_start\_total} & counter & basic & \\
\code{function\_warm\_start\_total} & counter & basic & \\
\code{function\_admitted\_total}, \code{function\_refused\_total},
\code{function\_replayed\_total} & counter & basic & per percorso di
invocazione \\
\bottomrule
\end{tabular}
\end{adjustbox}
\end{center}

I due timer che il governatore di concorrenza legge dal registro
(\code{function\_latency\_ms} e \code{function\_e2e\_latency\_ms}) sono
volutamente presenti anche in \code{basic}. Nei profili non \code{basic} i
quattro timer per funzione ricevono anche gli istogrammi necessari ai
percentili.

\section{Stato dello store e idempotenza}

\begin{center}
\small
\begin{adjustbox}{max width=\textwidth}
\begin{tabular}{@{}lllp{4.4cm}@{}}
\toprule
\textbf{Metrica} & \textbf{Tipo} & \textbf{Profilo} & \textbf{Note} \\
\midrule
\code{execution\_store\_size} & gauge & basic & esiti trattenuti \\
\code{execution\_in\_flight\_records} & gauge & basic & esecuzioni non
concluse \\
\code{idempotency\_keys\_held} & gauge & basic & chiavi (legami e lapidi) \\
\code{idempotency\_key\_budget\_rejections} & counter & basic & nuove
ammissioni con chiave rifiutate per budget esaurito \\
\bottomrule
\end{tabular}
\end{adjustbox}
\end{center}

\section{Code}

\begin{center}
\small
\begin{adjustbox}{max width=\textwidth}
\begin{tabular}{@{}lllp{4.4cm}@{}}
\toprule
\textbf{Metrica} & \textbf{Tipo} & \textbf{Modulo} & \textbf{Note} \\
\midrule
\code{function\_queue\_depth} & gauge & \code{async-queue} & \\
\code{function\_inFlight} & gauge & \code{async-queue} & \\
\code{function\_dispatchable\_backlog} & gauge & \code{async-queue} &
backlog realmente dispatchabile dato il limite corrente \\
\code{function\_effective\_concurrency} & gauge & \code{async-queue} &
advanced; ci\`o che il controllore ha deciso \\
\code{sync\_queue\_depth} & gauge & \code{sync-queue} & \code{function=""}
per la serie globale; nullo fuori dal profilo \code{SYNC\_QUEUE} \\
\code{sync\_queue\_admitted\_total} & counter & \code{sync-queue} & \\
\code{sync\_queue\_rejected\_total} & counter & \code{sync-queue} & anche il
timeout in coda (motivo \code{TIMEOUT}) \\
\bottomrule
\end{tabular}
\end{adjustbox}
\end{center}

\section{Scheduler}

\begin{center}
\small
\begin{adjustbox}{max width=\textwidth}
\begin{tabular}{@{}llp{5.6cm}@{}}
\toprule
\textbf{Metrica} & \textbf{Tipo} & \textbf{Note} \\
\midrule
\code{scheduler\_active} & gauge & una serie per strategia disponibile
(\code{strategy}): 1 se attiva \\
\code{scheduler\_switch\_total} & counter & cambi di strategia, per
\code{outcome} \\
\code{scheduler\_switch\_duration} & timer & durata del cambio; istogramma
fuori da \code{basic} \\
\bottomrule
\end{tabular}
\end{adjustbox}
\end{center}

\section{Politiche}

\begin{center}
\small
\begin{adjustbox}{max width=\textwidth}
\begin{tabular}{@{}lllp{4.2cm}@{}}
\toprule
\textbf{Metrica} & \textbf{Tipo} & \textbf{Modulo} & \textbf{Note} \\
\midrule
\code{function\_target\_inflight\_per\_pod} & gauge &
\code{concurrency-control} & advanced \\
\code{function\_concurrency\_controller\_mode} & gauge &
\code{concurrency-control} & advanced \\
\code{gateway\_service\_target\_load} & gauge & \code{autoscaler} & advanced \\
\code{function\_scaling\_recommended\_replicas},
\code{function\_scaling\_desired\_replicas},
\code{function\_scaling\_limited}, \code{function\_scaling\_ratio\_milli} &
gauge & \code{autoscaler} & la decisione dello scaler interno, prima e dopo il
clamp; il rapporto \`e scalato per 1000 \\
\code{nanofaas\_offload\_total} & counter & \code{offload} & etichette
\code{function}, \code{trigger} \\
\code{nanofaas\_offload\_failure\_total} & counter & \code{offload} & \\
\code{controlplane\_runtime\_config\_revision} & gauge &
\code{runtime-config} & advanced \\
\code{controlplane\_runtime\_config\_updates\_total} & counter &
\code{runtime-config} & etichetta \code{status} \\
\code{controlplane\_runtime\_config\_apply\_duration\_seconds} & timer &
\code{runtime-config} & \\
\bottomrule
\end{tabular}
\end{adjustbox}
\end{center}

\section{Sistema}

\begin{center}
\small
\begin{adjustbox}{max width=\textwidth}
\begin{tabular}{@{}llp{6.4cm}@{}}
\toprule
\textbf{Metrica} & \textbf{Tipo} & \textbf{Note} \\
\midrule
\code{function\_scheduler\_*}, \code{function\_queue\_offer\_duration*},
\code{function\_queue\_poll\_duration*}, \code{function\_dispatch\_slot\_hold\_*},
\code{scheduler\_visit\_duration*}, \code{scheduler\_idle\_duration*} & vari &
strumentazione diagnostica del percorso caldo: negata in \code{basic} \\
\code{removed\_function\_*} & vari & registro separato, non esportato
(\cref{cap:06-nucleo}) \\
\code{jvm\_*}, \code{process\_*} & vari & standard Micrometer; sulle build
native G1 i binder del collector non riportano nulla
(\cref{cap:24-footprint}) \\
\bottomrule
\end{tabular}
\end{adjustbox}
\end{center}

\section{Profilo \code{soak}}

Solo con \code{nanofaas.metrics.profile=soak}, nove gauge aggregati senza
etichette. Sei nel control plane:
\code{invocation\_execution\_reservations},
\code{invocation\_canonical\_input\_bytes},
\code{invocation\_physical\_input\_copy\_bytes},
\code{execution\_waiters\_retained}, \code{execution\_expiry\_queue\_depth},
\code{function\_capacity\_retired\_generations}. Tre nell'SDK Java:
\code{runtime\_active\_handlers}, \code{runtime\_pending\_callbacks},
\code{runtime\_pending\_callback\_bytes}.

\section{Interrogazioni di riferimento}

\begin{lstlisting}[language=bash]
# tasso di errore per funzione (percentuale)
rate(function_error_total{function="F"}[5m])
  / (rate(function_success_total{function="F"}[5m])
     + rate(function_error_total{function="F"}[5m])) * 100

# p95 della latenza di servizio (richiede il profilo advanced)
histogram_quantile(0.95, sum by (le) (rate(function_latency_ms_bucket{function="F"}[5m])))

# durata media dell'inizializzazione (richiede advanced)
rate(function_init_duration_ms_sum{function="F"}[1h])
  / rate(function_init_duration_ms_count{function="F"}[1h])

# arretrato asincrono
function_queue_depth{function!=""}

# pressione di ammissione sincrona
rate(sync_queue_rejected_total{function="F"}[5m])
  / rate(sync_queue_admitted_total{function="F"}[5m])

# carico per replica: il segnale dell'autoscaler
rate(function_dispatch_total{function="F"}[1m])
  / max by (function) (kube_deployment_status_replicas{deployment="fn-F"})

# attivita' di offload per causa
rate(nanofaas_offload_total[5m])

# quota della latenza dovuta all'attesa
rate(function_queue_wait_ms_sum{function="F"}[5m])
  / rate(function_e2e_latency_ms_sum{function="F"}[5m])
\end{lstlisting}

L'ultima interrogazione non compare nella documentazione operativa del progetto
ed \`e aggiunta qui: \`e la forma diretta del risultato centrale del
\cref{cap:12-concurrency-control}, e su un sistema sotto carico dovrebbe
avvicinarsi a $6/7$.

\chapter{Record sperimentali grezzi}
\label{cap:D-record-sperimentali}

\section{Serie di release}

Profilo \code{azure-d8s-v5+d2s-v5-amd64-native-loadtest-v1}, scenario
\code{autoscaling-cycle-k8s}, tre esecuzioni per versione, mediane.
Fonte: \code{docs/performance/history.md} e
\code{docs/performance/releases/*.json}.

\begin{center}
\small
\begin{adjustbox}{max width=\textwidth}
\begin{tabular}{@{}lrrrr@{}}
\toprule
\textbf{Versione} & \textbf{Throughput (req/s)} & \textbf{Tasso errore} &
\textbf{p95 (ms)} & \textbf{Repliche picco} \\
\midrule
0.17.0  & 241{,}904 & 0{,}000000 & 6{,}496  & 3 \\
0.17.1  & 269{,}069 & 0{,}000000 & 5{,}772  & 3 \\
v0.18.1 & 239{,}130 & 0{,}000000 & 6{,}607  & 4 \\
v0.18.2 & 249{,}341 & 0{,}000000 & 5{,}125  & 3 \\
v0.18.3 & 232{,}407 & 0{,}000000 & 5{,}131  & 4 \\
v0.19.0 & 300{,}483 & 0{,}000000 & 11{,}256 & 5 \\
\bottomrule
\end{tabular}
\end{adjustbox}
\end{center}

\subsection{Dettaglio della v0.19.0}

\begin{center}
\small
\begin{adjustbox}{max width=\textwidth}
\begin{tabular}{@{}lr@{}}
\toprule
\textbf{Campo} & \textbf{Valore} \\
\midrule
\code{throughputRps} & 300{,}4828762829388 \\
\code{errorRate} & 0{,}0 \\
\code{latencyP50Ms} & 4{,}432555 \\
\code{latencyP95Ms} & 11{,}256492 \\
\code{latencyP99Ms} & 21{,}048210 \\
\code{queueWaitSeconds} & 0{,}0008152916783986485 \\
\code{coldStarts} & 5{,}0 \\
\code{peakReplicas} & 5{,}0 \\
\code{controlPlaneCpuPeak} & 0{,}664 \\
\code{controlPlaneHeapPeakBytes} & 199\,495\,680 \\
\code{runCount} & 3 \\
\code{sourceCommit} & \code{7a5a669\ldots} \\
\bottomrule
\end{tabular}
\end{adjustbox}
\end{center}

Soglie del gate: \code{errorRateMax} $=0{,}3\%$,
\code{p95MaxIncreasePercent} $=15\%$, \code{throughputMaxLossPercent} $=10\%$.

\section{Confronto \code{BUDGETED} / \code{ADAPTIVE\_PER\_POD}}

Due funzioni, quattro core condivisi, coda 100 ciascuna, tetto 8, budget 12,
ciclo chiuso con picco 103 per funzione su tre finestre.

\begin{center}
\small
\begin{adjustbox}{max width=\textwidth}
\begin{tabular}{@{}lrr@{}}
\toprule
 & \textbf{\code{ADAPTIVE\_PER\_POD}} & \textbf{\code{BUDGETED}} \\
\midrule
richieste & 899\,684 & 893\,841 \\
rifiutate & 7 (0{,}001\%) & 33 (0{,}004\%) \\
p95 end-to-end & 80{,}2\,ms & 79{,}7\,ms \\
p99 end-to-end & 90{,}2\,ms & 87{,}2\,ms \\
\bottomrule
\end{tabular}
\end{adjustbox}
\end{center}

\begin{center}
\small
\begin{adjustbox}{max width=\textwidth}
\begin{tabular}{@{}lrrrrrr@{}}
\toprule
\textbf{Finestra} & \textbf{AD. java} & \textbf{AD. lite} & \textbf{somma} &
\textbf{BU. java} & \textbf{BU. lite} & \textbf{somma} \\
\midrule
una sola & 7{,}8 & 4{,}1 & --- & 8{,}0 & 1{,}2 & --- \\
antifase & 8{,}0 & 5{,}0 & 13{,}0 & 8{,}0 & 2{,}9 & 10{,}9 \\
in fase  & 7{,}9 & 5{,}5 & 13{,}4 & 8{,}0 & 3{,}7 & 11{,}7 \\
\bottomrule
\end{tabular}
\end{adjustbox}
\end{center}

\section{\code{SOJOURN}: ricerca contro modello}

\subsection{Versione a ricerca (fallita)}

\begin{center}
\small
\begin{adjustbox}{max width=\textwidth}
\begin{tabular}{@{}lrr@{}}
\toprule
 & \textbf{\code{ADAPTIVE\_PER\_POD}} & \textbf{\code{SOJOURN} a ricerca} \\
\midrule
rifiutate & 32 (0{,}004\%) & 36\,181 (4{,}275\%) \\
limite [min--max] & 7{,}8 [4--8] & 4{,}0 [1--8] \\
\bottomrule
\end{tabular}
\end{adjustbox}
\end{center}

\subsection{Versione a modello, arrivi aperti}

\begin{center}
\small
\begin{adjustbox}{max width=\textwidth}
\begin{tabular}{@{}lrrr@{}}
\toprule
 & \textbf{\code{ADAPTIVE\_PER\_POD}} & \textbf{\code{SOJOURN}} & \textbf{$\Delta$} \\
\midrule
offerte & 659\,973 & 659\,995 & --- \\
servite & 549\,856 & 567\,210 & $+3{,}2\%$ \\
rifiutate & 110\,117 (16{,}7\%) & 92\,785 (14{,}1\%) & $-15{,}7\%$ \\
p95 end-to-end & 120{,}48\,ms & 107{,}68\,ms & $-10{,}6\%$ \\
p99 end-to-end & 157{,}67\,ms & 133{,}82\,ms & $-15{,}1\%$ \\
attesa in coda & 35{,}61\,ms & 31{,}03\,ms & $-12{,}9\%$ \\
tempo di servizio & 1{,}69\,ms & 2{,}00\,ms & $+18\%$ \\
limite medio & 3{,}9 & 2{,}8 & \\
\code{dropped\_iterations} & 21 & 0 & \\
\bottomrule
\end{tabular}
\end{adjustbox}
\end{center}

Una sola esecuzione per modo.

\section{Co-tenancy}

\begin{center}
\small
\begin{adjustbox}{max width=\textwidth}
\begin{tabular}{@{}lrr@{}}
\toprule
 & \textbf{4 core propri} & \textbf{core 0--3 condivisi} \\
\midrule
java da sola & 1\,762 rps / 4{,}57\,ms & 1\,762 rps / 4{,}57\,ms \\
java con vicina & 1\,274 rps / 6{,}23\,ms & 1\,274 rps / 6{,}23\,ms \\
combinato & 2\,769 rps & 2\,751 rps \\
\bottomrule
\end{tabular}
\end{adjustbox}
\end{center}

\section{Runtime a confronto}

\begin{center}
\small
\begin{adjustbox}{max width=\textwidth}
\begin{tabular}{@{}lrr@{}}
\toprule
 & \textbf{\code{word-stats-java}} & \textbf{\code{word-stats-java-lite}} \\
\midrule
residente a riposo & 133\,MiB & 2{,}7\,MiB \\
residente a caldo & 160\,MiB & 51{,}6\,MiB \\
tempo di servizio & 5{,}57\,ms & 2{,}15\,ms \\
\bottomrule
\end{tabular}
\end{adjustbox}
\end{center}

\section{Memoria di una JVM Spring}

\code{word-stats-java} a 190\,MiB residenti.

\begin{center}
\small
\begin{adjustbox}{max width=\textwidth}
\begin{tabular}{@{}lr@{}}
\toprule
heap in uso & 22{,}4\,MB \\
Metaspace (9\,375 classi) & 53{,}0\,MB \\
Compressed class space & 6{,}6\,MB \\
CodeHeap & 13{,}5\,MB \\
23 stack di thread & $\sim$23\,MB \\
\bottomrule
\end{tabular}
\end{adjustbox}
\end{center}

\section{Control plane: JVM contro nativo}

\begin{center}
\small
\begin{adjustbox}{max width=\textwidth}
\begin{tabular}{@{}lrrr@{}}
\toprule
 & \textbf{JVM} & \textbf{nativo serial} & \textbf{nativo G1} \\
\midrule
throughput & 8\,205 rps & 4\,983 rps & 8\,085 rps \\
mediana & 5{,}98\,ms & 4{,}71\,ms & 6{,}08\,ms \\
p95 & 16{,}06\,ms & 27{,}95\,ms & 16{,}51\,ms \\
massimo & 92\,ms & 1{,}77\,s & 123\,ms \\
avvio & 1{,}585\,s & 0{,}192\,s & 0{,}072\,s \\
\bottomrule
\end{tabular}
\end{adjustbox}
\end{center}

\begin{center}
\small
\begin{adjustbox}{max width=\textwidth}
\begin{tabular}{@{}lrr@{}}
\toprule
 & \textbf{JVM} & \textbf{nativo serial} \\
\midrule
raccolte complete & 2 ($\sim$50\,ms) & 25 (883\,ms) \\
tempo totale di raccolta & 6{,}6\,s & 22{,}1\,s su 60 \\
quota del tempo di parete & 0{,}9\,\% & 31{,}8\,\% \\
\bottomrule
\end{tabular}
\end{adjustbox}
\end{center}

\subsection{Pause G1 da registrazione JFR (90\,s)}

\begin{center}
\small
\begin{adjustbox}{max width=\textwidth}
\begin{tabular}{@{}lrrrr@{}}
\toprule
\textbf{Operazione} & \textbf{N} & \textbf{Totale} & \textbf{Media} &
\textbf{Max} \\
\midrule
\code{G1 wrapper} & 43 & 835{,}1\,ms & 19{,}4\,ms & 97{,}6\,ms \\
\code{Collect for allocation} & 31 & 651{,}1\,ms & 21{,}0\,ms & 94{,}5\,ms \\
\code{Try init concurrent mark} & 2 & 37{,}3\,ms & 18{,}6\,ms & 22{,}5\,ms \\
\midrule
totale & 76 & 1\,524\,ms & & 1{,}69\,\% \\
\bottomrule
\end{tabular}
\end{adjustbox}
\end{center}

\section{Livello di ottimizzazione}

\begin{center}
\small
\begin{adjustbox}{max width=\textwidth}
\begin{tabular}{@{}lrr@{}}
\toprule
 & \textbf{\code{-Os}} & \textbf{\code{-O3}} \\
\midrule
throughput & 8\,398 rps & 9\,251 rps \\
p95 & 15{,}64\,ms & 14{,}79\,ms \\
massimo & 148{,}7\,ms & 127{,}2\,ms \\
avvio & 0{,}173\,s & 0{,}068\,s \\
residente a riposo & 120{,}6\,MiB & 39{,}3\,MiB \\
immagine & 195\,MB & 397\,MB \\
area di codice & 47{,}78\,MB & 166{,}59\,MB \\
unit\`a di compilazione & 105\,256 & 84\,694 \\
\bottomrule
\end{tabular}
\end{adjustbox}
\end{center}

\section{Limiti di memoria}

\begin{center}
\small
\begin{adjustbox}{max width=\textwidth}
\begin{tabular}{@{}llrrrr@{}}
\toprule
\textbf{Limite} & \textbf{Build} & \textbf{rps} & \textbf{p95} &
\textbf{Max} & \textbf{Memoria} \\
\midrule
500\,MB & JVM & 1\,515 & 324\,ms & 1{,}99\,s & 455/500 \\
        & nativo serial & 1\,370 & 32\,ms & 30{,}0\,s & 489/500 \\
        & nativo G1 default & 520 & 61\,ms & 30{,}0\,s & 145/500 \\
        & nativo G1 \code{-Xmx350m} & 1\,445 & 96\,ms & 30{,}0\,s & 378/500 \\
1\,GB & JVM & 2\,450 & 46\,ms & 1{,}34\,s & 805/1024 \\
      & nativo G1 \code{-Xmx700m} & 3\,089 & 44\,ms & 135\,ms & 736/1024 \\
2\,GB & JVM & 2\,712 & 47{,}5\,ms & 451\,ms & 727/2048 \\
      & nativo G1 \code{-Xmx1400m} & 3\,400 & 39{,}5\,ms & 151\,ms & 1\,201/2048 \\
nessuno & JVM & 8\,205 & 16{,}1\,ms & 92\,ms & 1{,}8\,GiB \\
        & nativo G1 & 8\,085 & 16{,}5\,ms & 123\,ms & 2{,}7\,GiB \\
\bottomrule
\end{tabular}
\end{adjustbox}
\end{center}

\section{Cambio a caldo dello scheduler}

Campagna di settembre 2026; fonte: \code{docs/experiments/scheduler-switching-2026-09/}
(\code{FINAL.md}, \code{budgets.json}, \code{raw/}). I budget sono congelati in
\code{budgets.json} prima della misura; la tabella riporta ogni criterio con il
valore osservato, compresi quelli non superati (analisi in
\cref{sec:25-switch}).

\begin{center}
\small
\begin{adjustbox}{max width=\textwidth}
\begin{tabular}{@{}lrrl@{}}
\toprule
\textbf{Criterio} & \textbf{Congelato} & \textbf{Osservato} & \textbf{Esito} \\
\midrule
pausa massima di uno switch & 250\,ms & 3{,}98\,ms & passa \\
p99 della pausa (in processo) & 100\,ms & 0{,}266\,ms & passa \\
p99 della pausa (timer della piattaforma) & 100\,ms & 0{,}00099\,s & passa,
ma cieco sotto 1\,ms \\
indici vivi & 2 & 2 & passa \\
throughput utile & $-5\%$ & peggiore $+2{,}09\%$ & passa \\
heap dopo GC & $+10\%$ & peggiore $+0{,}01\%$ & passa \\
p99 in regime & $+5\%$ & peggiore $+9{,}51\%$ & oltre \\
CPU per completamento & $+10\%$ & peggiore $+22{,}06\%$ & oltre \\
\bottomrule
\end{tabular}
\end{adjustbox}
\end{center}

\subsection{Soak del 22 settembre 2026}

\begin{center}
\small
\begin{adjustbox}{max width=\textwidth}
\begin{tabular}{@{}lr@{}}
\toprule
\textbf{Campo} & \textbf{Valore} \\
\midrule
switch committati & 1\,000 \\
switch rifiutati & 0 \\
durata & 5\,394{,}6\,s \\
indici vivi (massimo) & 2 \\
pausa massima (driver) & 0{,}087\,ms \\
criteri valutati / falliti & 21 / 3 \\
criteri falliti & \code{p24-rss-return} sui tre ruoli \\
RSS sopra la mediana di riferimento: control plane & $+1{,}6\%$ \\
RSS sopra la mediana di riferimento: funzione Java & $+12{,}9\%$ \\
RSS sopra la mediana di riferimento: funzione JavaScript & $+4{,}3\%$ \\
\bottomrule
\end{tabular}
\end{adjustbox}
\end{center}

\section{Backoff delle riprove}

Campagna del 24 settembre 2026; fonte:
\code{docs/experiments/retry-backoff-2026-09/} (\code{RESULTS.md}, \code{raw/},
con \code{SHA256SUMS}). Riassunto delle fasi su VM k3s in
\cref{sec:25-backoff}; qui il dettaglio dell'upstream controllato
(\code{429} con \code{Retry-After: 1} per un secondo, poi \code{200};
300 esecuzioni per revisione).

\begin{center}
\small
\begin{adjustbox}{max width=\textwidth}
\begin{tabular}{@{}lrr@{}}
\toprule
\textbf{Grandezza} & \textbf{Base} & \textbf{Con backoff} \\
\midrule
tentativi totali & 1\,195 & 600 \\
intervalli sotto il secondo & 892 & 0 \\
distanza minima fra tentativi & 0{,}0007\,s & 1{,}0007\,s \\
successi (VM k3s, tre fasi da 300) & 742/900 & 900/900 \\
riprove (VM k3s, somma delle tre fasi) & 1\,184 & 48 \\
\bottomrule
\end{tabular}
\end{adjustbox}
\end{center}

Gate prestazionale sul codice finale, p99 in regime: quattro bracci, limite
congelato $+5\%$; il ramo \code{shared-queue} invariato risulta $+8{,}06\%$
(mediana 2{,}571 $\to$ 2{,}779\,ms), gli altri tre $-7{,}42\%$, $-3{,}26\%$,
$-4{,}56\%$. Tutti i confronti di CPU passano.

\chapter{Glossario}
\label{cap:E-glossario}

\begin{description}[style=nextline, leftmargin=0pt, itemsep=2pt]

\item[\textbf{Ammissione}] La decisione se accettare una richiesta in coda. In
\nf{} \`e presa dal modulo \code{sync-queue} per profondit\`a e per attesa
stimata (\cref{cap:10-sync-queue}).

\item[\textbf{Attesa} (\emph{queue wait})] Intervallo fra l'ammissione della
richiesta e la sua consegna al runtime. Metrica
\code{function\_queue\_wait\_ms}.

\item[\textbf{Budget di concorrenza}] Numero massimo di invocazioni
contemporaneamente in volo su tutte le funzioni in modo \code{BUDGETED}. \`E una
dichiarazione di capacit\`a della piattaforma, non una grandezza inferita dal
carico (\cref{cap:12-concurrency-control}).

\item[\textbf{Ciclo aperto / ciclo chiuso}] Nel primo gli arrivi sono
programmati dall'orologio; nel secondo ogni utente virtuale attende la risposta
prima di inviare. Un carico a ciclo chiuso non pu\`o distinguere politiche di
ammissione (\cref{cap:22-metodo-sperimentale}).

\item[\textbf{Cold start}] Ritardo dovuto all'inizializzazione di un'istanza non
ancora pronta. Metrica \code{function\_cold\_start\_ms}.

\item[\textbf{Compilazione anticipata} (\emph{AOT})] Fase di build di Spring
che risolve i bean e le loro condizioni prima dell'esecuzione. Nell'immagine
nativa le condizioni sulle propriet\`a sono quindi congelate al valore che
avevano in build (\cref{cap:20-build-packaging}).

\item[\textbf{Contratto di saturazione}] Politica comune con cui ogni runtime di
funzione rifiuta il lavoro oltre limiti finiti di conteggio e di byte, con
codici \code{429}, \code{413}, \code{500} e \code{503} e un corpus JSON
condiviso che tutti gli SDK verificano (\cref{sec:18-saturazione}).

\item[\textbf{Control plane}] Il processo che registra le funzioni, orchestra le
invocazioni e governa le politiche. In \nf{} \`e uno solo.

\item[\textbf{Degradazione osservabile}] Riduzione di funzionalit\`a che il
sistema dichiara --- via log, metrica o codice di stato --- anzich\'e
nascondere. \`E il criterio che distingue le degradazioni accettabili da quelle
che fanno rifiutare l'avvio (\cref{cap:03-requisiti-principi}).

\item[\textbf{Deriva della baseline}] Meccanismo per cui un estremo storico
(minimo di latenza, massimo di throughput) viene lentamente riportato verso il
presente, cos\`i che un intervallo anomalo non lo fissi per sempre
(\cref{cap:12-concurrency-control}).

\item[\textbf{Dispatch}] La consegna dell'invocazione al runtime della funzione.

\item[\textbf{Equit\`a max-min pesata}] Criterio di ripartizione in cui si
massimizza l'allocazione minima, chi chiede meno della propria quota la ottiene
per intero, e l'eccedenza \`e ridistribuita~\cite{bertsekas1992datanetworks}.

\item[\textbf{Finestra di ammissione}] La quantit\`a $C + Q$ --- limite di
concorrenza pi\`u dimensione della coda --- oltre la quale una funzione rifiuta.
Il fatto che il limite ne faccia parte \`e il secondo lavoro del limite
(\cref{cap:12-concurrency-control}).

\item[\textbf{Ginocchio}] Il valore del limite di concorrenza a cui il servizio
satura e oltre il quale l'aggiunta produce attesa anzich\'e lavoro. Calcolato
come $\lambda_{\max} S_{\min}$ per la legge di Little.

\item[\textbf{Idempotenza}] Propriet\`a per cui due invocazioni con la stessa
chiave corrispondono a una sola esecuzione (\cref{cap:06-nucleo}).

\item[\textbf{Legge di Little}] $L = \lambda W$: il numero medio nel sistema
uguaglia il tasso di arrivo per il tempo medio di permanenza~\cite{little1961proof}.

\item[\textbf{Limite configurato / effettivo}] Il primo \`e ci\`o che l'utente ha
dichiarato ed \`e un tetto invalicabile; il secondo \`e ci\`o che un controllore
adattivo applica in questo momento (\cref{cap:09-async-queue}).

\item[\textbf{Modulo}] Unit\`a opzionale del control plane, selezionata a tempo
di build e caricata via SPI (\cref{cap:08-sistema-moduli}).

\item[\textbf{Offload}] Inoltro trasparente di un'invocazione sincrona a
un'istanza remota, per politica della funzione o per pressione locale
(\cref{cap:13-offload}).

\item[\textbf{Provider di deployment}] Backend che realizza l'intento
\code{DEPLOYMENT}: Kubernetes, runtime container locale o containerd. I tre
moduli sono mutuamente esclusivi (\cref{cap:15-deployment-provider}).

\item[\textbf{Ricetta}] File YAML versionato (schema v2) che dice quale control plane
costruire --- moduli e modo JVM o nativo --- e quali implementazioni di funzione
impacchettare con esso; non avvia n\'e registra nulla. Quattro comandi Gradle
elencano il catalogo, validano, assemblano e pubblicano
(\cref{sec:20-ricette}).

\item[\textbf{Replay}] Riproduzione immediata dell'esito di un'esecuzione gi\`a
terminata, per una richiesta con chiave di idempotenza corrispondente.

\item[\textbf{Servizio} (\emph{service time})] Intervallo fra dispatch e
completamento, comprensivo del round trip HTTP. Metrica
\code{function\_latency\_ms}. \emph{Non} include l'attesa in coda.

\item[\textbf{Sojourn}] Attesa pi\`u servizio: ci\`o che il chiamante osserva al
netto della rete. Metrica \code{function\_e2e\_latency\_ms}. \`E la grandezza da
cui il modo omonimo decide.

\item[\textbf{SPI} (\emph{Service Provider Interface})] Meccanismo Java di
scoperta di implementazioni via \code{ServiceLoader}, usato per caricare i
moduli.

\item[\textbf{Watchdog}] Binario Rust che fa da entrypoint di ogni container di
funzione gestito e adatta i tre protocolli di runtime a un unico contratto HTTP
(\cref{cap:17-runtime-funzione}).

\end{description}


\bibliographystyle{plain}
\bibliography{bibliografia}

\end{document}
